% Generated mechanically from the frozen morphisms.tex target.
% Do not edit this generated file; edit the bound locale source instead.

% OK, start here.
%

\setcounter{chapter}{28}
\chapter{概形的态射}



\phantomsection
\label{morphisms-section-phantom}


\section{引言}
\label{morphisms-section-introduction}

\noindent
本章介绍若干类概形态射。基本参考文献为 \cite{EGA}。




















\section{闭浸入}
\label{morphisms-section-closed-immersions}

\noindent
本节阐明此前所得的若干概形闭浸入结果。回顾，概形态射
$i : Z \to X$ 按定义是闭浸入，如果：(a) $i$ 诱导到 $X$
的某个闭子集上的同胚；(b)
$i^\sharp : \mathcal{O}_X \to i_*\mathcal{O}_Z$ 满射；并且 (c)
$i^\sharp$ 的核局部由截面生成；见《概形》定义
\ref{schemes-definition-immersion} 与
\ref{schemes-definition-closed-immersion-locally-ringed-spaces}。
事实表明，在已知 $Z$ 与 $X$ 都是概形时，有许多不同方式可以刻画
闭浸入。

\begin{lemma}
\label{morphisms-lemma-closed-immersion}
设 $i : Z \to X$ 为概形态射。下列条件等价：
\begin{enumerate}
\item 态射 $i$ 是闭浸入。
\item 对每个仿射开集 $\Spec(R) = U \subset X$，都存在理想
$I \subset R$，使得 $i^{-1}(U) = \Spec(R/I)$ 作为
$U = \Spec(R)$ 上的概形成立。
\item 存在仿射开覆盖 $X = \bigcup_{j \in J} U_j$，其中
$U_j = \Spec(R_j)$，并且对每个 $j \in J$ 都存在理想
$I_j \subset R_j$，使得 $i^{-1}(U_j) = \Spec(R_j/I_j)$ 作为
$U_j = \Spec(R_j)$ 上的概形成立。
\item 态射 $i$ 诱导 $Z$ 到 $X$ 的某个闭子集上的同胚，并且
$i^\sharp : \mathcal{O}_X \to i_*\mathcal{O}_Z$ 满射。
\item 态射 $i$ 诱导 $Z$ 到 $X$ 的某个闭子集上的同胚，映射
$i^\sharp : \mathcal{O}_X \to i_*\mathcal{O}_Z$ 满射，并且核
$\Ker(i^\sharp)\subset \mathcal{O}_X$ 是准凝聚理想层。
\item 态射 $i$ 诱导 $Z$ 到 $X$ 的某个闭子集上的同胚，映射
$i^\sharp : \mathcal{O}_X \to i_*\mathcal{O}_Z$ 满射，并且核
$\Ker(i^\sharp)\subset \mathcal{O}_X$ 是局部由截面生成的理想层。
\end{enumerate}
\end{lemma}

\begin{proof}
条件 (6) 就是闭浸入的定义；见《概形》定义
\ref{schemes-definition-closed-immersion-locally-ringed-spaces}
与 \ref{schemes-definition-immersion}。所以
(6) $\Leftrightarrow$ (1)。由《概形》引理
\ref{schemes-lemma-closed-subspace-scheme}，有
(1) $\Rightarrow$ (2)。显然，(2) $\Rightarrow$ (3)。

\medskip\noindent
假设 (3)。每个态射
$\Spec(R_j/I_j) \to \Spec(R_j)$ 都是闭浸入；见《概形》例
\ref{schemes-example-closed-immersion-affines}。因此，
$i^{-1}(U_j) \to U_j$ 是到其像上的同胚，并且
$i^\sharp|_{U_j}$ 满射。所以 $i$ 是到其像上的同胚，并且
$i^\sharp$ 满射，因为这些性质可以局部检验。由此断定
(3) $\Rightarrow$ (4)。

\medskip\noindent
蕴含 (4) $\Rightarrow$ (1) 是《概形》引理
\ref{schemes-lemma-characterize-closed-immersions}。
蕴含 (5) $\Rightarrow$ (6) 显然。蕴含
(6) $\Rightarrow$ (5) 则由《概形》引理
\ref{schemes-lemma-closed-subspace-scheme} 得出。
\end{proof}



\begin{lemma}
\label{morphisms-lemma-closed-immersion-ideals}
设 $X$ 为概形。设 $i : Z \to X$ 与 $i' : Z' \to X$ 为闭浸入，
并考虑理想层
$\mathcal{I} = \Ker(i^\sharp)$ 与
$\mathcal{I}' = \Ker((i')^\sharp)$，它们都是
$\mathcal{O}_X$ 的理想层。
\begin{enumerate}
\item 态射 $i : Z \to X$ 可分解为 $Z \to Z' \to X$，
其中相应态射为某个 $a : Z \to Z'$，当且仅当
$\mathcal{I}' \subset \mathcal{I}$。若这成立，则 $a$ 是闭浸入。
\item $Z \cong Z'$ 作为 $X$ 上的同构成立，当且仅当
$\mathcal{I} = \mathcal{I}'$。
\end{enumerate}
\end{lemma}

\begin{proof}
这由《概形》第 \ref{schemes-section-closed-immersion} 节中关于闭子空间的
讨论得出，特别是《概形》引理
\ref{schemes-lemma-closed-immersion} 与
\ref{schemes-lemma-characterize-closed-subspace}。
也可以直接由上面引理 \ref{morphisms-lemma-closed-immersion} 中的刻画 (3) 得出。
\end{proof}

\begin{lemma}
\label{morphisms-lemma-closed-immersion-bijection-ideals}
设 $X$ 为概形。设 $\mathcal{I} \subset \mathcal{O}_X$ 为理想层。
下列条件等价：
\begin{enumerate}
\item $\mathcal{I}$ 作为 $\mathcal{O}_X$-模层局部由截面生成；
\item $\mathcal{I}$ 作为 $\mathcal{O}_X$-模层准凝聚；并且
\item 存在概形的闭浸入 $i : Z \to X$，其相应理想层
$\Ker(i^\sharp)$ 等于 $\mathcal{I}$。
\end{enumerate}
\end{lemma}

\begin{proof}
(1) 与 (2) 的等价性直接来自《概形》引理
\ref{schemes-lemma-closed-subspace-scheme}。若 (1) 成立，则由《概形》定义
\ref{schemes-definition-closed-subspace} 与例
\ref{schemes-example-closed-subspace}，存在闭子空间
$i : Z \to X$，使得 $\mathcal{I} = \Ker(i^\sharp)$。
由《概形》引理 \ref{schemes-lemma-closed-subspace-scheme}，
这是概形的闭浸入，故 (3) 成立。反过来，若 (3) 成立，
则由《概形》引理 \ref{schemes-lemma-closed-subspace-scheme}，
(2) 成立（该引理适用，因为概形的闭浸入尤其是局部环层空间的
闭浸入）。
\end{proof}

\begin{lemma}
\label{morphisms-lemma-base-change-closed-immersion}
闭浸入的基变换仍是闭浸入。
\end{lemma}

\begin{proof}
见《概形》引理 \ref{schemes-lemma-base-change-immersion}。
\end{proof}

\begin{lemma}
\label{morphisms-lemma-composition-closed-immersion}
闭浸入的复合仍是闭浸入。
\end{lemma}

\begin{proof}
我们已在《概形》引理
\ref{schemes-lemma-composition-immersion} 中见过这一点，
不过这里再给出一个证明。确切地说，这由引理
\ref{morphisms-lemma-closed-immersion} 中闭浸入的刻画 (3) 得出。
因为若 $I \subset R$ 是理想，且
$\overline{J} \subset R/I$ 是理想，则有 $\overline{J} = J/I$，
其中 $J \subset R$ 是某个包含 $I$ 的理想，并且
$(R/I)/\overline{J} = R/J$。
\end{proof}

\begin{lemma}
\label{morphisms-lemma-closed-immersion-quasi-compact}
闭浸入拟紧。
\end{lemma}

\begin{proof}
本引理与《概形》引理
\ref{schemes-lemma-closed-immersion-quasi-compact} 重复。
\end{proof}

\begin{lemma}
\label{morphisms-lemma-closed-immersion-separated}
闭浸入分离。
\end{lemma}

\begin{proof}
本引理是《概形》引理
\ref{schemes-lemma-immersions-monomorphisms} 的特殊情形。
\end{proof}




\section{浸入}
\label{morphisms-section-immersions}

\noindent
本节汇集若干关于浸入的事实。

\begin{lemma}
\label{morphisms-lemma-immersion-permanence}
设 $Z \to Y \to X$ 为概形态射。
\begin{enumerate}
\item 若 $Z \to X$ 是浸入，则 $Z \to Y$ 是浸入。
\item 若 $Z \to X$ 是拟紧浸入且 $Y \to X$ 拟分离，
则 $Z \to Y$ 是拟紧浸入。
\item 若 $Z \to X$ 是闭浸入且 $Y \to X$ 分离，
则 $Z \to Y$ 是闭浸入。
\end{enumerate}
\end{lemma}

\begin{proof}
在每种情形中，证明都归结为考察交换图
$$
\xymatrix{
Z \ar[r] \ar[rd] & Y \times_X Z \ar[r] \ar[d] & Z \ar[d] \\
& Y \ar[r] & X
}
$$
其中上方两条水平箭头的复合为恒等态射。我们证明 (1)。
第一条水平箭头是 $Y \times_X Z \to Z$ 的截面，因而由《概形》引理
\ref{schemes-lemma-section-immersion}，它是浸入。箭头
$Y \times_X Z \to Y$ 是 $Z \to X$ 的基变换，故为浸入
（《概形》引理 \ref{schemes-lemma-base-change-immersion}）。
最后，浸入的复合仍是浸入
（《概形》引理 \ref{schemes-lemma-composition-immersion}）。
这证明了 (1)。另两个结论以完全相同的方式证明。
\end{proof}

\begin{lemma}
\label{morphisms-lemma-factor-quasi-compact-immersion}
设 $h : Z \to X$ 为浸入。若 $h$ 拟紧，则可将
$h = i \circ j$ 分解，其中 $j : Z \to \overline{Z}$ 为开浸入，
而 $i : \overline{Z} \to X$ 为闭浸入。
\end{lemma}

\begin{proof}
注意，$h$ 拟紧且拟分离（见《概形》引理
\ref{schemes-lemma-immersions-monomorphisms}）。因此，由《概形》引理
\ref{schemes-lemma-push-forward-quasi-coherent}，
$h_*\mathcal{O}_Z$ 是准凝聚 $\mathcal{O}_X$-模层。由此可知
$\mathcal{I} = \Ker(\mathcal{O}_X \to h_*\mathcal{O}_Z)$
是准凝聚理想层；见《概形》第
\ref{schemes-section-quasi-coherent} 节。设
$\overline{Z} \subset X$ 为与 $\mathcal{I}$ 对应的闭子概形；
见引理 \ref{morphisms-lemma-closed-immersion-bijection-ideals}。
由《概形》引理 \ref{schemes-lemma-characterize-closed-subspace}，
态射 $h$ 可分解为 $h = i \circ j$，其中
$i : \overline{Z} \to X$ 为包含态射。为说明 $j$ 是开浸入，
选取开子概形 $U \subset X$，使得 $h$ 诱导 $Z$ 到 $U$ 的闭浸入。
此时显然，$\mathcal{I}|_U$ 正是与闭浸入 $Z \to U$ 对应的理想层。
因此，$Z = \overline{Z} \cap U$。
\end{proof}

\begin{lemma}
\label{morphisms-lemma-factor-reduced-immersion}
设 $h : Z \to X$ 为浸入。若 $Z$ 约化，则可将
$h = i \circ j$ 分解，其中 $j : Z \to \overline{Z}$ 为开浸入，
而 $i : \overline{Z} \to X$ 为闭浸入。
\end{lemma}

\begin{proof}
设 $\overline{Z} \subset X$ 为 $h(Z)$ 的闭包，并赋予约化诱导闭子概形
结构；见《概形》定义
\ref{schemes-definition-reduced-induced-scheme}。
由《概形》引理 \ref{schemes-lemma-map-into-reduction}，
态射 $h$ 可分解为 $h = i \circ j$，其中
$i : \overline{Z} \to X$ 为包含态射，而
$j : Z \to \overline{Z}$。由浸入的定义可知，存在开子概形
$U \subset X$，使得 $h$ 可经到 $U$ 的闭浸入分解。因此，
$\overline{Z} \cap U$ 与 $h(Z)$ 都是 $U$ 的约化闭子概形，
并且具有相同的底闭集。于是由《概形》引理
\ref{schemes-lemma-reduced-closed-subscheme} 中的唯一性可知
$h(Z) \cong \overline{Z} \cap U$。所以 $j$ 诱导 $Z$ 到
$\overline{Z} \cap U$ 的同构。换言之，$j$ 是开浸入。
\end{proof}



\begin{example}
\label{morphisms-example-thibaut}
下面给出一个不能表示为先开浸入后闭浸入之复合的浸入例子。
设 $k$ 为域。设 $X = \Spec(k[x_1, x_2, x_3, \ldots])$。
设 $U = \bigcup_{n = 1}^{\infty} D(x_n)$。则
$U \to X$ 是开浸入。考虑理想
$$
I_n =
(x_1^n, x_2^n, \ldots, x_{n - 1}^n, x_n - 1, x_{n + 1}, x_{n + 2}, \ldots)
\subset
k[x_1, x_2, x_3, \ldots][1/x_n].
$$
注意，$I_n k[x_1, x_2, x_3, \ldots][1/x_nx_m] = (1)$
对任意 $m \not = n$ 都成立。因此，准凝聚理想
$\widetilde I_n$ 在 $D(x_n)$ 上给出，并且它们在
$D(x_nx_m)$ 上相容；确切地说，
$\widetilde I_n|_{D(x_nx_m)} = \mathcal{O}_{D(x_n x_m)}$
在 $n \not = m$ 时成立。
故这些理想粘合为准凝聚理想层
$\mathcal{I} \subset \mathcal{O}_U$。设 $Z \subset U$ 为与
$\mathcal{I}$ 对应的闭子概形。于是 $Z \to X$ 是浸入。

\medskip\noindent
我们断言，不能将 $Z \to X$ 分解为
$Z \to \overline{Z} \to X$，其中 $\overline{Z} \to X$ 为闭浸入，
而 $Z \to \overline{Z}$ 为开浸入。事实上，$\overline{Z}$ 必须由
某个理想 $I \subset k[x_1, x_2, x_3, \ldots]$ 定义，使得
$I_n = I k[x_1, x_2, x_3, \ldots][1/x_n]$。但是，唯一的元素
$f \in k[x_1, x_2, x_3, \ldots]$，其最终属于所有 $I_n$，
是 $0$！因此，$I$ 不存在。
\end{example}

\begin{lemma}
\label{morphisms-lemma-check-immersion}
设 $f : Y \to X$ 为概形态射。若对所有 $y \in Y$，都存在开子概形
$f(y) \in U \subset X$，使得
$f|_{f^{-1}(U)} : f^{-1}(U) \to U$ 是浸入，则 $f$ 是浸入。
\end{lemma}

\begin{proof}
这一陈述直接来自《概形》第 \ref{schemes-section-immersions} 节末尾
关于闭子概形的讨论，不过我们也给出详细证明。设
$Z \subset X$ 为 $f(Y)$ 的闭包。由于取闭包与限制到开集可交换，
由假设可知 $f(Y) \subset Z$ 是开集。故
$Z' = Z \setminus f(Y)$ 是闭集。因此，
$X' = X \setminus Z'$ 是 $X$ 的开子概形，并且 $f$ 可分解为
$f : Y \to X'$，随后接包含态射。若 $y \in Y$ 且
$U \subset X$ 如引理陈述中所述，则
$U' = X' \cap U$ 是 $f'(y)$ 的开邻域，使得
$(f')^{-1}(U') \to U'$ 是像为闭集的浸入
（引理 \ref{morphisms-lemma-immersion-permanence}）。因此它是闭浸入；
见《概形》引理 \ref{schemes-lemma-immersion-when-closed}。
由于闭浸入性在目标上是局部的
（例如由引理 \ref{morphisms-lemma-closed-immersion}），可断定
$f'$ 如所要求的是闭浸入。
\end{proof}






\section{闭浸入与准凝聚层}
\label{morphisms-section-closed-immersions-quasi-coherent}

\noindent
下列引理终于对概形上的准凝聚层完成了《模》引理
\ref{modules-lemma-i-star-exact} 对阿贝尔群层所完成的工作。
另见《模》第 \ref{modules-section-closed-immersion} 节中的讨论。

\begin{lemma}
\label{morphisms-lemma-i-star-equivalence}
设 $i : Z \to X$ 为概形的闭浸入。设
$\mathcal{I} \subset \mathcal{O}_X$ 为切出 $Z$ 的准凝聚理想层。
函子
$$
i_* :
\QCoh(\mathcal{O}_Z)
\longrightarrow
\QCoh(\mathcal{O}_X)
$$
正合且全忠实，其本质像由准凝聚 $\mathcal{O}_X$-模
$\mathcal{G}$ 中满足 $\mathcal{I}\mathcal{G} = 0$ 的那些模组成。
\end{lemma}

\begin{proof}
闭浸入拟紧且分离；见引理
\ref{morphisms-lemma-closed-immersion-quasi-compact} 与
\ref{morphisms-lemma-closed-immersion-separated}。因此《概形》引理
\ref{schemes-lemma-push-forward-quasi-coherent} 适用，并且
$Z$ 上准凝聚层的正像确实是 $X$ 上的准凝聚层。

\medskip\noindent
由《模》引理 \ref{modules-lemma-i-star-equivalence}，函子
$i_*$ 全忠实。

\medskip\noindent
现在来描述函子 $i_*$ 的本质像。我们有
$\mathcal{I}(i_*\mathcal{F}) = 0$ 对任意准凝聚
$\mathcal{O}_Z$-模成立；
例如可由《模》引理 \ref{modules-lemma-i-star-equivalence} 得出。
接着，设 $\mathcal{G}$ 为任意准凝聚 $\mathcal{O}_X$-模，且满足
$\mathcal{I}\mathcal{G} = 0$。只需证明典范映射
$$
\mathcal{G} \longrightarrow i_* i^*\mathcal{G}
$$
是同构\footnote{这已在更一般的情形中于《模》引理
\ref{modules-lemma-i-star-equivalence} 的证明中得到。}。
在概形与准凝聚模的情形中，在 $X$ 上仿射局部地工作，并使用引理
\ref{morphisms-lemma-closed-immersion} 与《概形》引理
\ref{schemes-lemma-widetilde-pullback}，即可归结为证明下列代数陈述：
给定环 $R$、理想 $I$ 以及 $R$-模 $N$，并假设 $IN = 0$，典范映射
$$
N \longrightarrow N \otimes_R R/I,\quad
n \longmapsto n \otimes 1
$$
是 $R$-模同构。这个简单代数事实的证明从略。
\end{proof}

\noindent
设 $i : Z \to X$ 为闭浸入。由于上引理，我们常滥用记号，用
$\mathcal{F}$ 表示层 $i_*\mathcal{F}$，并将其视为 $X$ 上的层。

\begin{lemma}
\label{morphisms-lemma-largest-quasi-coherent-subsheaf}
设 $X$ 为概形。设 $\mathcal{F}$ 为准凝聚
$\mathcal{O}_X$-模。设 $\mathcal{G} \subset \mathcal{F}$ 为
$\mathcal{O}_X$-子模。存在唯一的准凝聚
$\mathcal{O}_X$-子模 $\mathcal{G}' \subset \mathcal{G}$，具有下列性质：
对每个准凝聚 $\mathcal{O}_X$-模 $\mathcal{H}$，映射
$$
\Hom_{\mathcal{O}_X}(\mathcal{H}, \mathcal{G}')
\longrightarrow
\Hom_{\mathcal{O}_X}(\mathcal{H}, \mathcal{G})
$$
是双射。特别地，$\mathcal{G}'$ 是最大的准凝聚
$\mathcal{O}_X$-子模，它作为 $\mathcal{F}$ 的子模包含于
$\mathcal{G}$。
\end{lemma}

\begin{proof}
设 $\mathcal{G}_a$（$a \in A$）为所有准凝聚
$\mathcal{O}_X$-子模中包含于 $\mathcal{G}$ 的那些所成的集合。则映射
$$
\bigoplus\nolimits_{a \in A} \mathcal{G}_a \longrightarrow \mathcal{F}
$$
的像 $\mathcal{G}'$ 是准凝聚的，因为 $X$ 上准凝聚层之间映射的像
是准凝聚的，且准凝聚层的直和仍准凝聚；见《概形》第
\ref{schemes-section-quasi-coherent} 节。模 $\mathcal{G}'$ 包含于
$\mathcal{G}$。因此，这就是最大的准凝聚 $\mathcal{O}_X$-模，
并包含于 $\mathcal{G}$。

\medskip\noindent
为证明该公式，设 $\mathcal{H}$ 为准凝聚 $\mathcal{O}_X$-模，并设
$\alpha : \mathcal{H} \to \mathcal{G}$ 为 $\mathcal{O}_X$-模映射。
复合 $\mathcal{H} \to \mathcal{G} \to \mathcal{F}$ 的像是准凝聚的，
因为它是准凝聚层之间映射的像。因此它包含于 $\mathcal{G}'$。
故 $\alpha$ 如所要求地经 $\mathcal{G}'$ 分解。
\end{proof}

\begin{lemma}
\label{morphisms-lemma-i-upper-shriek}
设 $i : Z \to X$ 为概形的闭浸入。存在函子
\footnote{这很可能是非标准记号。}
$i^! : \QCoh(\mathcal{O}_X) \to \QCoh(\mathcal{O}_Z)$，
它是 $i_*$ 的右伴随。（比较《模》引理
\ref{modules-lemma-i-star-right-adjoint}。）
\end{lemma}

\begin{proof}
给定准凝聚 $\mathcal{O}_X$-模 $\mathcal{G}$，考察子层
$\mathcal{H}_Z(\mathcal{G})$；它是 $\mathcal{G}$ 中被
$\mathcal{I}$ 零化的局部截面所组成的子层。由引理
\ref{morphisms-lemma-largest-quasi-coherent-subsheaf}，存在典范的最大准凝聚
$\mathcal{O}_X$-子模 $\mathcal{H}_Z(\mathcal{G})'$。由构造，对任意
准凝聚 $\mathcal{O}_Z$-模 $\mathcal{F}$，都有
$$
\Hom_{\mathcal{O}_X}(i_*\mathcal{F}, \mathcal{H}_Z(\mathcal{G})')
=
\Hom_{\mathcal{O}_X}(i_*\mathcal{F}, \mathcal{G})
$$
。因此可以令 $i^!\mathcal{G} = i^*(\mathcal{H}_Z(\mathcal{G})')$。
细节从略。
\end{proof}

\noindent
利用准凝聚理想层与闭子概形之间的 $1$-to-$1$ 对应
（见引理 \ref{morphisms-lemma-closed-immersion-bijection-ideals}），
可以定义闭子概形的概形论交与概形论并。

\begin{definition}
\label{morphisms-definition-scheme-theoretic-intersection-union}
设 $X$ 为概形。设 $Z, Y \subset X$ 为分别对应于准凝聚理想层
$\mathcal{I}, \mathcal{J} \subset \mathcal{O}_X$ 的闭子概形。
$Z$ 与 $Y$ 的{\it 概形论交}是 $X$ 中由
$\mathcal{I} + \mathcal{J}$ 切出的闭子概形。
$Z$ 与 $Y$ 的{\it 概形论并}是 $X$ 中由
$\mathcal{I} \cap \mathcal{J}$ 切出的闭子概形。
\end{definition}

\noindent

\begin{lemma}
\label{morphisms-lemma-scheme-theoretic-intersection}
设 $X$ 为概形。设 $Z, Y \subset X$ 为闭子概形。设
$Z \cap Y$ 为 $Z$ 与 $Y$ 的概形论交。则
$Z \cap Y \to Z$ 与 $Z \cap Y \to Y$ 都是闭浸入，并且
$$
\xymatrix{
Z \cap Y \ar[r] \ar[d] & Z \ar[d] \\
Y \ar[r] & X
}
$$
是概形的笛卡尔图，即 $Z \cap Y = Z \times_X Y$。
\end{lemma}

\begin{proof}
态射 $Z \cap Y \to Z$ 与 $Z \cap Y \to Y$ 由引理
\ref{morphisms-lemma-closed-immersion-ideals} 可知都是闭浸入。
设 $U = \Spec(A)$ 为 $X$ 的仿射开集，并设 $Z \cap U$ 与
$Y \cap U$ 分别对应于理想 $I \subset A$ 与 $J \subset A$。
则 $Z \cap Y \cap U$ 对应于 $I + J \subset A$。由于
$A/I \otimes_A A/J = A/(I + J)$，由《概形》第
\ref{schemes-section-fibre-products} 节中对概形纤维积的描述可知，
该图是笛卡尔的。
\end{proof}

\begin{lemma}
\label{morphisms-lemma-scheme-theoretic-union}
设 $S$ 为概形。设 $X, Y \subset S$ 为闭子概形。设
$X \cup Y$ 为 $X$ 与 $Y$ 的概形论并。设
$X \cap Y$ 为 $X$ 与 $Y$ 的概形论交。则
$X \to X \cup Y$ 与 $Y \to X \cup Y$ 都是闭浸入，存在短正合列
$$
0 \to \mathcal{O}_{X \cup Y} \to \mathcal{O}_X \times \mathcal{O}_Y
\to \mathcal{O}_{X \cap Y} \to 0
$$
作为 $\mathcal{O}_S$-模的短正合列，并且图
$$
\xymatrix{
X \cap Y \ar[r] \ar[d] & X \ar[d] \\
Y \ar[r] & X \cup Y
}
$$
在概形范畴中是余笛卡尔的，即
$X \cup Y = X \amalg_{X \cap Y} Y$。
\end{lemma}

\begin{proof}
态射 $X \to X \cup Y$ 与 $Y \to X \cup Y$ 由引理
\ref{morphisms-lemma-closed-immersion-ideals} 可知都是闭浸入。
在短正合列中，我们利用引理 \ref{morphisms-lemma-i-star-equivalence} 的等价，
把 $S$ 的闭子概形上的准凝聚模看作 $S$ 上的准凝聚模。
对于该列的第一个映射，我们使用典范映射
$\mathcal{O}_{X \cup Y} \to \mathcal{O}_X$ 与
$\mathcal{O}_{X \cup Y} \to \mathcal{O}_Y$；
对于第二个映射，我们使用典范映射
$\mathcal{O}_X \to \mathcal{O}_{X \cap Y}$ 以及典范映射
$\mathcal{O}_Y \to \mathcal{O}_{X \cap Y}$ 的负映射。
于是可仿射局部地检验正合性。设 $U = \Spec(A)$ 为 $S$ 的仿射开集，
并设 $X \cap U$ 与 $Y \cap U$ 分别对应于理想
$I \subset A$ 与 $J \subset A$。则
$(X \cup Y) \cap U$ 对应于 $I \cap J \subset A$，而
$X \cap Y \cap U$ 对应于 $I + J \subset A$。
因此，正合性来自下列序列的正合性：
$$
0 \to A/I \cap J \to A/I \times A/J \to A/(I + J) \to 0
$$
为证明该图是余笛卡尔的，设给定概形 $T$ 以及概形态射
$f : X \to T$、$g : Y \to T$，它们作为态射
$X \cap Y \to T$ 相同。目标：证明存在唯一态射
$h : X \cup Y \to T$ 与 $f$ 和 $g$ 相同。为构造 $h$，可以在
$X \cup Y$ 上仿射局部地工作；见《概形》第
\ref{schemes-section-glueing-schemes} 节。若 $s \in X$ 且
$s \not \in Y$，则 $X \to X \cup Y$ 在 $s$ 的某邻域上是同构，
此时显然可构造 $h$。对 $s \in Y$ 且 $s \not \in X$ 也同样。
若 $s \in X \cap Y$，可选取仿射开集
$V = \Spec(B) \subset T$，它包含 $f(s) = g(s)$。
随后可选取仿射开集 $U = \Spec(A) \subset S$，它包含 $s$，
并使得 $f(X \cap U)$ 与 $g(Y \cap U)$ 都包含于 $V$。态射
$f|_{X \cap U}$ 与 $g|_{Y \cap V}$ 到 $V$，并对应于环映射
$$
B \to A/I
\quad\text{且}\quad
B \to A/J
$$
，它们作为到 $A/(I + J)$ 的映射相同。由上面显示的短正合列，
这些环同态存在唯一的提升，得到所需的环映射
$B \to A/I \cap J$。
\end{proof}









\section{模的支集}
\label{morphisms-section-support}

\noindent
本节汇集概形上准凝聚模支集的若干初等结果。回顾，模层的支集已在
《模》第 \ref{modules-section-support} 节中定义。另一方面，模的支集
已在《代数》第 \ref{algebra-section-support} 节中定义。二者相符。

\begin{lemma}
\label{morphisms-lemma-support-affine-open}
设 $X$ 为概形。设 $\mathcal{F}$ 为 $X$ 上的准凝聚层。
设 $\Spec(A) = U \subset X$ 为仿射开集，并令
$M = \Gamma(U, \mathcal{F})$。设 $x \in U$，并设
$\mathfrak p \subset A$ 为对应素理想。下列条件等价：
\begin{enumerate}
\item $\mathfrak p$ 属于 $M$ 的支集；以及
\item $x$ 属于 $\mathcal{F}$ 的支集。
\end{enumerate}
\end{lemma}

\begin{proof}
这来自等式 $\mathcal{F}_x = M_{\mathfrak p}$；见《概形》引理
\ref{schemes-lemma-spec-sheaves} 以及相应定义。
\end{proof}

\begin{lemma}
\label{morphisms-lemma-support-closed-specialization}
设 $X$ 为概形。设 $\mathcal{F}$ 为 $X$ 上的准凝聚层。
$\mathcal{F}$ 的支集在特殊化下封闭。
\end{lemma}

\begin{proof}
若 $x' \leadsto x$ 是特殊化且 $\mathcal{F}_x = 0$，则
$\mathcal{F}_{x'}$ 为零，因为 $\mathcal{F}_{x'}$ 是模
$\mathcal{F}_x$ 的局部化。因此，$\text{Supp}(\mathcal{F})$ 的补集
在一般化下封闭。
\end{proof}

\noindent
对有限型准凝聚模而言，支集是闭的，可以在纤维上检验，并与基变换
可交换。

\begin{lemma}
\label{morphisms-lemma-support-finite-type}
设 $\mathcal{F}$ 为概形 $X$ 上的有限型准凝聚模。则：
\begin{enumerate}
\item $\mathcal{F}$ 的支集是闭的。
\item 对 $x \in X$，有
$$
x \in \text{Supp}(\mathcal{F})
\Leftrightarrow
\mathcal{F}_x \not = 0
\Leftrightarrow
\mathcal{F}_x \otimes_{\mathcal{O}_{X, x}} \kappa(x) \not = 0.
$$
\item 对任意概形态射 $f : Y \to X$，拉回
$f^*\mathcal{F}$ 也为有限型，并且
$\text{Supp}(f^*\mathcal{F}) = f^{-1}(\text{Supp}(\mathcal{F}))$。
\end{enumerate}
\end{lemma}

\begin{proof}
部分 (1) 是《模》引理
\ref{modules-lemma-support-finite-type-closed} 的重新表述。
也可结合引理 \ref{morphisms-lemma-support-affine-open}、《性质》引理
\ref{properties-lemma-finite-type-module} 以及《代数》引理
\ref{algebra-lemma-support-closed} 看出这一点。(2) 中第一个等价是
支集的定义，第二个等价来自 Nakayama 引理；见《代数》引理
\ref{algebra-lemma-NAK}。设 $f : Y \to X$ 为概形态射。注意，
$f^*\mathcal{F}$ 由《模》引理
\ref{modules-lemma-pullback-finite-type} 可知是有限型。
对于最后一个断言，设 $y \in Y$ 的像为 $x \in X$。回顾
$$
(f^*\mathcal{F})_y =
\mathcal{F}_x \otimes_{\mathcal{O}_{X, x}} \mathcal{O}_{Y, y},
$$
见《层》引理 \ref{sheaves-lemma-stalk-pullback-modules}。
因此，$(f^*\mathcal{F})_y \otimes \kappa(y)$ 非零，当且仅当
$\mathcal{F}_x \otimes \kappa(x)$ 非零。由 (2)，这推出
$x \in \text{Supp}(\mathcal{F})$ 当且仅当
$y \in \text{Supp}(f^*\mathcal{F})$，这正是断言 (3) 的内容。
\end{proof}

\begin{lemma}
\label{morphisms-lemma-scheme-theoretic-support}
设 $\mathcal{F}$ 为概形 $X$ 上的有限型准凝聚模。存在最小闭子概形
$i : Z \to X$，使得存在准凝聚 $\mathcal{O}_Z$-模
$\mathcal{G}$，满足 $i_*\mathcal{G} \cong \mathcal{F}$。此外：
\begin{enumerate}
\item 若 $\Spec(A) \subset X$ 是任意仿射开集，并且
$\mathcal{F}|_{\Spec(A)} = \widetilde{M}$，则
$Z \cap \Spec(A) = \Spec(A/I)$，其中 $I = \text{Ann}_A(M)$。
\item 准凝聚层 $\mathcal{G}$ 在唯一同构意义下唯一。
\item 准凝聚层 $\mathcal{G}$ 是有限型的。
\item $\mathcal{G}$ 与 $\mathcal{F}$ 的支集都是 $Z$。
\end{enumerate}
\end{lemma}

\begin{proof}
假设 $i' : Z' \to X$ 是满足引理中仿射开集上描述的闭子概形。
那么由引理 \ref{morphisms-lemma-i-star-equivalence} 可知，有
$\mathcal{F} \cong i'_*\mathcal{G}'$，其中 $\mathcal{G}'$ 是
$Z'$ 上某个唯一的准凝聚层。此外，显然
$Z'$ 是具有此性质的最小闭子概形（仍由同一引理）。
最后，使用《性质》引理
\ref{properties-lemma-finite-type-module} 与《代数》引理
\ref{algebra-lemma-finite-over-subring} 可知 $\mathcal{G}'$ 是有限型的。
由《代数》引理 \ref{algebra-lemma-support-closed}，有
$\text{Supp}(\mathcal{G}') = Z$。因此，为证明本引理，只需说明
(1) 中的刻画确实定义了闭子概形。为此，又只需证明给定规则产生
准凝聚理想层；见引理
\ref{morphisms-lemma-closed-immersion-bijection-ideals}。
这归结为下列代数事实：若 $A$ 为环，$f \in A$，且 $M$ 为有限
$A$-模，则 $\text{Ann}_A(M)_f = \text{Ann}_{A_f}(M_f)$。
证明从略。
\end{proof}

\begin{definition}
\label{morphisms-definition-scheme-theoretic-support}
设 $X$ 为概形。设 $\mathcal{F}$ 为有限型准凝聚
$\mathcal{O}_X$-模。$\mathcal{F}$ 的{\it 概形论支集}是引理
\ref{morphisms-lemma-scheme-theoretic-support} 中构造的闭子概形
$Z \subset X$。
\end{definition}

\noindent
在此情形中，我们常把 $\mathcal{F}$ 看作 $Z$ 上的有限型准凝聚层
（通过引理 \ref{morphisms-lemma-i-star-equivalence} 的范畴等价）。










\section{概形论像}
\label{morphisms-section-scheme-theoretic-image}

\noindent
注意：本节中的一些材料处于极其一般的层次，其表现可能不同于预期。

\begin{lemma}
\label{morphisms-lemma-scheme-theoretic-image}
设 $f : X \to Y$ 为概形态射。存在闭子概形 $Z \subset Y$，
使得 $f$ 经 $Z$ 分解，并且对任意其他闭子概形 $Z' \subset Y$，
若 $f$ 经 $Z'$ 分解，则有 $Z \subset Z'$。
\end{lemma}

\begin{proof}
设 $\mathcal{I} = \Ker(\mathcal{O}_Y \to f_*\mathcal{O}_X)$。
若 $\mathcal{I}$ 准凝聚，就取 $Z$ 为由 $\mathcal{I}$ 决定的闭子概形；
见引理 \ref{morphisms-lemma-closed-immersion-bijection-ideals}。
由引理 \ref{morphisms-lemma-closed-immersion-ideals}，这样确实可行。
一般而言，同一引理要求证明存在最大准凝聚理想层
$\mathcal{I}'$，它包含于 $\mathcal{I}$。这由引理
\ref{morphisms-lemma-largest-quasi-coherent-subsheaf} 得出。
\end{proof}

\begin{definition}
\label{morphisms-definition-scheme-theoretic-image}
设 $f : X \to Y$ 为概形态射。$f$ 的{\it 概形论像}是最小闭子概形
$Z \subset Y$，且 $f$ 经它分解；见上面的引理
\ref{morphisms-lemma-scheme-theoretic-image}。
\end{definition}

\noindent
对概形态射 $f : X \to Y$ 及其概形论像 $Z$，我们常以
$f : X \to Z$ 表示 $f$ 经其概形论像的分解。若态射 $f$ 不拟紧，
则一般而言：
\begin{enumerate}
\item 集合论包含 $\overline{f(X)} \subset Z$ 不是等式，即
$f(X) \subset Z$ 不是稠密子集；并且
\item 概形论像的构造不与到 $Y$ 的开子概形的限制可交换。
\end{enumerate}
在《例》第 \ref{examples-section-scheme-theoretic-image} 节中，
读者可找到同时表现这两个现象的例子。这些现象甚至可对浸入出现；
见《例》第 \ref{examples-section-strange-immersion} 节。
不过，下一个引理表明，态射拟紧时两个灾难都会避免。

\begin{lemma}
\label{morphisms-lemma-quasi-compact-scheme-theoretic-image}
设 $f : X \to Y$ 为概形态射。设 $Z \subset Y$ 为 $f$ 的概形论像。
若 $f$ 拟紧，则：
\begin{enumerate}
\item 理想层
$\mathcal{I} = \Ker(\mathcal{O}_Y \to f_*\mathcal{O}_X)$
是准凝聚的；
\item 概形论像 $Z$ 是由 $\mathcal{I}$ 决定的闭子概形；
\item 对任意开集 $U \subset Y$，态射
$f|_{f^{-1}(U)} : f^{-1}(U) \to U$ 的概形论像等于
$Z \cap U$；并且
\item 像 $f(X) \subset Z$ 是 $Z$ 的稠密子集，换言之，态射
$X \to Z$ 是支配的（见定义 \ref{morphisms-definition-dominant}）。
\end{enumerate}
\end{lemma}

\begin{proof}
部分 (4) 由部分 (3) 得出。为证明 (3)，只需证明 (1)，因为
$\mathcal{I}$ 的形成与限制到 $Y$ 的开子概形可交换。
而若 (1) 成立，则我们已在引理
\ref{morphisms-lemma-scheme-theoretic-image} 的证明中证明 (2)。
因此只需证明 $\mathcal{I}$ 准凝聚。由于准凝聚性是局部性质，
可以假设 $Y$ 仿射。由于 $f$ 拟紧，可取有限仿射开覆盖
$X = \bigcup_{i = 1, \ldots, n} U_i$。以 $f'$ 表示复合
$$
X' = \coprod U_i \longrightarrow X \longrightarrow Y.
$$
于是 $f_*\mathcal{O}_X$ 是 $f'_*\mathcal{O}_{X'}$ 的子层，因而
$\mathcal{I} = \Ker(\mathcal{O}_Y \to f'_*\mathcal{O}_{X'})$。
由《概形》引理 \ref{schemes-lemma-push-forward-quasi-coherent}，
层 $f'_*\mathcal{O}_{X'}$ 在 $Y$ 上准凝聚。因此得证。
\end{proof}

\begin{example}
\label{morphisms-example-scheme-theoretic-image}
若 $A \to B$ 是核为 $I$ 的环映射，则
$\Spec(B) \to \Spec(A)$ 的概形论像是闭子概形
$\Spec(A/I)$，它位于 $\Spec(A)$ 中。这由引理
\ref{morphisms-lemma-quasi-compact-scheme-theoretic-image} 得出。
\end{example}

\noindent
若态射拟紧，则概形论像只添加像中各点的特殊化点。

\begin{lemma}
\label{morphisms-lemma-reach-points-scheme-theoretic-image}
设 $f : X \to Y$ 为拟紧态射。设 $Z$ 为 $f$ 的概形论像。
设 $z \in Z$\footnote{由引理
\ref{morphisms-lemma-quasi-compact-scheme-theoretic-image}，从集合论看，
$Z$ 与 $f(X)$ 在 $Y$ 中的闭包相同。}。存在赋值环 $A$，其分式域为
$K$，并且存在交换图
$$
\xymatrix{
\Spec(K) \ar[rr] \ar[d] & & X \ar[d] \ar[ld] \\
\Spec(A) \ar[r] & Z \ar[r] & Y
}
$$
使得 $\Spec(A)$ 的闭点映到 $z$。特别地，$Z$ 中任意一点都是
$f(X)$ 中某一点的特殊化。
\end{lemma}

\begin{proof}
设 $z \in \Spec(R) = V \subset Y$ 为 $z$ 的仿射开邻域。
由引理 \ref{morphisms-lemma-quasi-compact-scheme-theoretic-image}，
交 $Z \cap V$ 是 $f^{-1}(V) \to V$ 的概形论像。因此可将 $Y$
替换为 $V$，并假设 $Y = \Spec(R)$ 仿射。在此情形中，$X$ 拟紧，
因为 $f$ 拟紧。设 $X = U_1 \cup \ldots \cup U_n$ 为有限仿射开覆盖。
写成 $U_i = \Spec(A_i)$。令
$I = \Ker(R \to A_1 \times \ldots \times A_n)$。
再次由引理 \ref{morphisms-lemma-quasi-compact-scheme-theoretic-image} 可知，
$Z$ 对应于闭子概形 $\Spec(R/I)$，它位于 $Y$ 中。若
$\mathfrak p \subset R$ 为与 $z$ 对应的素理想，则
$I_{\mathfrak p} \subset R_{\mathfrak p}$ 不是等式。因此
（由于局部化正合，见《代数》命题
\ref{algebra-proposition-localization-exact}）可知
$R_{\mathfrak p} \to
(A_1)_{\mathfrak p} \times \ldots \times (A_n)_{\mathfrak p}$
不是零映射。所以某个环 $(A_i)_{\mathfrak p}$ 非零。
因此存在 $i$ 以及素理想 $\mathfrak q_i \subset A_i$，它位于某个
素理想 $\mathfrak p_i \subset \mathfrak p$ 之上。由《代数》引理
\ref{algebra-lemma-dominate}，可选择赋值环
$A \subset K = \kappa(\mathfrak q_i)$，它支配局部环
$R_{\mathfrak p}/\mathfrak p_iR_{\mathfrak p} \subset \kappa(\mathfrak q_i)$。
这给出所需交换图。若干细节从略。
\end{proof}

\begin{lemma}
\label{morphisms-lemma-factor-factor}
设
$$
\xymatrix{
X_1 \ar[d] \ar[r]_{f_1} & Y_1 \ar[d] \\
X_2 \ar[r]^{f_2} & Y_2
}
$$
为概形的交换图。设 $Z_i \subset Y_i$（$i = 1, 2$）为 $f_i$
的概形论像。则态射 $Y_1 \to Y_2$ 诱导态射 $Z_1 \to Z_2$ 以及交换图
$$
\xymatrix{
X_1 \ar[r] \ar[d] & Z_1 \ar[d] \ar[r] & Y_1 \ar[d] \\
X_2 \ar[r] & Z_2 \ar[r] & Y_2
}
$$
\end{lemma}

\begin{proof}
$Z_2$ 在 $Y_1$ 中的概形论逆像是 $Y_1$ 的闭子概形，并且 $f_1$
经它分解。因此 $Z_1$ 包含于其中。引理得证。
\end{proof}

\begin{lemma}
\label{morphisms-lemma-scheme-theoretic-image-reduced}
设 $f : X \to Y$ 为概形态射。若 $X$ 约化，则 $f$ 的概形论像是
$\overline{f(X)}$ 上的约化诱导概形结构。
\end{lemma}

\begin{proof}
这是因为 $\overline{f(X)}$ 上的约化诱导概形结构显然是 $Y$ 的
最小闭子概形，并且 $f$ 经它分解；见《概形》引理
\ref{schemes-lemma-map-into-reduction}。
\end{proof}

\begin{lemma}
\label{morphisms-lemma-scheme-theoretic-image-of-partial-section}
设 $f : X \to Y$ 为概形的分离态射。设 $V \subset Y$ 为逆拟紧开集。
设 $s : V \to X$ 为满足 $f \circ s = \text{id}_V$ 的态射。
设 $Y'$ 为 $s$ 的概形论像。则 $Y' \to Y$ 在 $V$ 上为同构。
\end{lemma}

\begin{proof}
假设 $V$ 在 $Y$ 中逆拟紧（《拓扑》定义
\ref{topology-definition-quasi-compact}），这意味着
$V \to Y$ 是拟紧态射。由《概形》引理
\ref{schemes-lemma-quasi-compact-permanence}，态射
$s : V \to X$ 拟紧。因此，由引理
\ref{morphisms-lemma-quasi-compact-scheme-theoretic-image}，概形论像
$Y'$（它是 $s$ 的概形论像）的构造与到开集的限制可交换。特别地，
$Y' \cap f^{-1}(V)$ 是分离态射 $f^{-1}(V) \to V$ 的某个截面的
概形论像。由于分离态射的截面是闭浸入
（《概形》引理 \ref{schemes-lemma-section-immersion}），
可断定 $Y' \cap f^{-1}(V) \to V$ 如所要求地为同构。
\end{proof}






\section{概形论闭包与稠密性}
\label{morphisms-section-scheme-theoretic-closure}

\noindent
我们采用 \cite[IV, Definition 11.10.2]{EGA} 中的下列定义。

\begin{definition}
\label{morphisms-definition-scheme-theoretically-dense}
设 $X$ 为概形。设 $U \subset X$ 为开子概形。
\begin{enumerate}
\item 态射 $U \to X$ 的概形论像称为{\it $U$ 在 $X$ 中的概形论闭包}。
\item 称 $U$ 在 $X$ 中{\it 概形论稠密}，如果对每个开集
$V \subset X$，$U \cap V$ 在 $V$ 中的概形论闭包等于 $V$。
\end{enumerate}
\end{definition}

\noindent
采用此定义，$U$ 在 $X$ 中概形论稠密{\bf 并不}等价于 $U$ 的概形论
闭包为 $X$；见例
\ref{morphisms-example-scheme-theretically-dense-not-dense}。这略显不雅；
不过可参见下面的引理
\ref{morphisms-lemma-scheme-theoretically-dense-quasi-compact} 与
\ref{morphisms-lemma-reduced-scheme-theoretically-dense}。另一方面，采用此定义，
$U$ 在 $X$ 中概形论稠密，当且仅当对每个开集 $V \subset X$，环映射
$\mathcal{O}_X(V) \to \mathcal{O}_X(U \cap V)$ 都是单射；见下面的引理
\ref{morphisms-lemma-characterize-scheme-theoretically-dense}。特别地，概形论稠密
蕴含稠密，这是令人满意的。

\begin{example}
\label{morphisms-example-scheme-theretically-dense-not-dense}
下面给出一个概形论闭包为 $X$，但其底拓扑空间中并不稠密的例子。
设 $k$ 为域。令
$A = k[x, z_1, z_2, \ldots]/(x^n z_n)$，并令
$I = (z_1, z_2, \ldots) \subset A$。考虑仿射概形
$X = \Spec(A)$ 与开子概形 $U = X \setminus V(I)$。
由于 $A \to \prod_n A_{z_n}$ 单射，可知 $U$ 的概形论闭包是 $X$。
考虑态射 $X \to \Spec(k[x])$。此态射满射
（令所有 $z_n = 0$ 即可看出）。但此态射到 $U$ 的限制不满射，
因为它映到点 $x = 0$。因此，$U$ 不可能在 $X$ 中拓扑稠密。
\end{example}

\begin{lemma}
\label{morphisms-lemma-scheme-theoretically-dense-quasi-compact}
设 $X$ 为概形。设 $U \subset X$ 为开子概形。
若包含态射 $U \to X$ 拟紧，则 $U$ 在 $X$ 中概形论稠密，当且仅当
$U$ 在 $X$ 中的概形论闭包是 $X$。
\end{lemma}

\begin{proof}
这由引理 \ref{morphisms-lemma-quasi-compact-scheme-theoretic-image} 的部分 (3)
得出。
\end{proof}

\begin{example}
\label{morphisms-example-scheme-theoretic-closure}
设 $A$ 为环且 $X = \Spec(A)$。设 $f_1, \ldots, f_n \in A$，并设
$U = D(f_1) \cup \ldots \cup D(f_n)$。令
$I = \Ker(A \to \prod A_{f_i})$。则 $U$ 在 $X$ 中的概形论闭包是
闭子概形 $\Spec(A/I)$，它位于 $X$ 中。注意，$U \to X$ 拟紧。因此由引理
\ref{morphisms-lemma-scheme-theoretically-dense-quasi-compact} 可知，$U$ 在
$X$ 中概形论稠密，当且仅当 $I = 0$。
\end{example}

\begin{lemma}
\label{morphisms-lemma-characterize-scheme-theoretically-dense}
设 $j : U \to X$ 为概形的开浸入。则 $U$ 在 $X$ 中概形论稠密，
当且仅当 $\mathcal{O}_X \to j_*\mathcal{O}_U$ 单射。
\end{lemma}

\begin{proof}
若 $\mathcal{O}_X \to j_*\mathcal{O}_U$ 单射，则限制到任意开集
$V$ 后仍然如此，其中该开集是 $X$ 的开集。因此，$U \cap V$ 在 $V$
中的概形论闭包等于 $V$；
见引理 \ref{morphisms-lemma-scheme-theoretic-image} 的证明。反过来，假设
$U \cap V$ 的概形论闭包等于 $V$，并且对所有开集 $V$ 都如此。假设
$\mathcal{O}_X \to j_*\mathcal{O}_U$ 不单射。则可找到仿射开集
$\Spec(A) = V \subset X$ 与非零元素 $f \in A$，使得 $f$ 在
$\Gamma(V \cap U, \mathcal{O}_X)$ 中的像为零。在此情形中，
$V \cap U$ 在 $V$ 中的概形论闭包显然包含于
$\Spec(A/(f))$，矛盾。
\end{proof}

\begin{lemma}
\label{morphisms-lemma-intersection-scheme-theoretically-dense}
设 $X$ 为概形。若 $U$、$V$ 是 $X$ 的概形论稠密开子概形，
则 $U \cap V$ 也是。
\end{lemma}

\begin{proof}
设 $W \subset X$ 为任意开集。考虑映射
$\mathcal{O}_X(W) \to \mathcal{O}_X(W \cap V)
\to \mathcal{O}_X(W \cap V \cap U)$。
由引理 \ref{morphisms-lemma-characterize-scheme-theoretically-dense}，
两个映射都单射。因此复合单射。故再次由引理
\ref{morphisms-lemma-characterize-scheme-theoretically-dense}，
$U \cap V$ 在 $X$ 中概形论稠密。
\end{proof}

\begin{lemma}
\label{morphisms-lemma-quasi-compact-immersion}
设 $h : Z \to X$ 为浸入。假设 $h$ 拟紧，或 $Z$ 约化。
设 $\overline{Z} \subset X$ 为 $h$ 的概形论像。则态射
$Z \to \overline{Z}$ 是开浸入，它将 $Z$ 认同为 $\overline{Z}$ 的
概形论稠密开子概形。此外，$Z$ 在 $\overline{Z}$ 中拓扑稠密。
\end{lemma}

\begin{proof}
由引理 \ref{morphisms-lemma-factor-quasi-compact-immersion} 或引理
\ref{morphisms-lemma-factor-reduced-immersion}，可将 $Z \to X$ 分解为
$Z \to \overline{Z}_1 \to X$，其中 $Z \to \overline{Z}_1$ 为开浸入，
而 $\overline{Z}_1 \to X$ 为闭浸入。另一方面，设
$Z \to \overline{Z} \subset X$ 为 $Z \to X$ 的概形论像。
可断定 $\overline{Z} \subset \overline{Z}_1$。由于 $Z$ 是
$\overline{Z}_1$ 的开子概形，故 $Z$ 也是 $\overline{Z}$ 的
开子概形。在 $Z$ 约化的情形中，已知
$Z \subset \overline{Z}_1$ 拓扑稠密；这是由引理
\ref{morphisms-lemma-factor-reduced-immersion} 的证明中对
$\overline{Z}_1$ 的构造所得。
因此，$\overline{Z}_1$ 与 $\overline{Z}$ 具有相同的底拓扑空间。
于是 $\overline{Z} \subset \overline{Z}_1$ 是到约化概形的闭浸入，
并在底拓扑空间上诱导双射，故它是同构。
在 $Z \to X$ 拟紧的情形中，论证如下：$Z$ 在 $\overline{Z}$ 中
概形论稠密这一断言由引理
\ref{morphisms-lemma-quasi-compact-scheme-theoretic-image} 的部分 (3) 得出。
最后一个断言由引理
\ref{morphisms-lemma-quasi-compact-scheme-theoretic-image} 的部分 (4) 得出。
\end{proof}

\begin{lemma}
\label{morphisms-lemma-reduced-scheme-theoretically-dense}
设 $X$ 为约化概形，并设 $U \subset X$ 为开子概形。下列条件等价：
\begin{enumerate}
\item $U$ 在 $X$ 中拓扑稠密；
\item $U$ 在 $X$ 中的概形论闭包是 $X$；并且
\item $U$ 在 $X$ 中概形论稠密。
\end{enumerate}
\end{lemma}

\begin{proof}
这由引理 \ref{morphisms-lemma-quasi-compact-immersion} 以及下述事实得出：
若闭子概形 $Z$ 位于 $X$ 中，且其底拓扑空间等于 $X$，则作为概形
它必等于 $X$。
\end{proof}

\begin{lemma}
\label{morphisms-lemma-reduced-subscheme-closure}
设 $X$ 为概形，并设 $U \subset X$ 为约化开子概形。下列条件等价：
\begin{enumerate}
\item $U$ 在 $X$ 中的概形论闭包是 $X$；以及
\item $U$ 在 $X$ 中概形论稠密。
\end{enumerate}
若这些条件成立，则 $X$ 是约化概形。
\end{lemma}

\begin{proof}
这由引理 \ref{morphisms-lemma-quasi-compact-immersion} 以及下述事实得出：
由引理 \ref{morphisms-lemma-scheme-theoretic-image-reduced}，$U$ 在 $X$ 中的
概形论闭包约化。
\end{proof}


\begin{lemma}
\label{morphisms-lemma-equality-of-morphisms}
设 $S$ 为概形。设 $X$、$Y$ 为 $S$ 上的概形。
设 $f, g : X \to Y$ 为 $S$ 上的概形态射。设 $U \subset X$
为满足 $f|_U = g|_U$ 的开子概形。若 $U$ 在 $X$ 中的概形论闭包
是 $X$，且 $Y \to S$ 分离，则 $f = g$。
\end{lemma}

\begin{proof}
这由定义以及《概形》引理
\ref{schemes-lemma-where-are-they-equal} 得出。
\end{proof}




\section{支配态射}
\label{morphisms-section-dominant}

\noindent
如果习惯了簇的支配态射概念，可能会发现概形态射为支配态射的定义
与预期略有不同。

\begin{definition}
\label{morphisms-definition-dominant}
概形态射 $f : X \to S$ 称为{\it 支配的}，如果 $f$ 的像是 $S$
的稠密子集。
\end{definition}

\noindent
例如，若 $k$ 是无限域，且 $\lambda_1, \lambda_2, \ldots$ 是 $k$
中两两不同元素的可数族，则态射
$$
\coprod\nolimits_{i = 1, 2, \ldots } \Spec(k)
\longrightarrow
\Spec(k[x])
$$
是支配的，其中第 $i$ 个因子映到点 $x = \lambda_i$。

\begin{lemma}
\label{morphisms-lemma-generic-points-in-image-dominant}
设 $f : X \to S$ 为概形态射。若 $S$ 的每个不可约分支的每个泛点
都属于 $f$ 的像，则 $f$ 支配。
\end{lemma}

\begin{proof}
这是一个拓扑事实，直接来自以下两点：概形的底拓扑空间清醒，
见《概形》引理 \ref{schemes-lemma-scheme-sober}；并且 $S$ 的每个点
都包含在 $S$ 的某个不可约分支中，见《拓扑》引理
\ref{topology-lemma-irreducible}。
\end{proof}

\noindent
若态射拟紧，则只有在目标泛点属于像时态射才支配这一预期确实成立。

\begin{lemma}
\label{morphisms-lemma-quasi-compact-dominant}
\begin{slogan}
像包含泛点的态射是支配的
\end{slogan}
设 $f : X \to S$ 为概形的拟紧态射。则 $f$ 支配，当且仅当对每个
不可约分支 $Z \subset S$，$Z$ 的泛点属于 $f$ 的像。
\end{lemma}

\begin{proof}
设 $V \subset S$ 为仿射开集。由于 $f$ 拟紧，可以选择有限多个仿射
开集 $U_i \subset f^{-1}(V)$（$i = 1, \ldots, n$）覆盖
$f^{-1}(V)$。考虑仿射概形的态射
$$
f' :
\coprod\nolimits_{i = 1, \ldots, n} U_i
\longrightarrow
V.
$$
仿射概形的不交并仍仿射；见《概形》引理
\ref{schemes-lemma-disjoint-union-affines}。$V$ 的不可约分支的泛点，
恰好是 $S$ 中与 $V$ 相交的那些不可约分支的泛点。此外，$f$ 支配
当且仅当 $f'$ 支配，不论上面如何选择 $V, n, U_i$。
因此已将本引理归结为仿射概形态射的情形。仿射情形即《代数》引理
\ref{algebra-lemma-image-dense-generic-points}。
\end{proof}

\begin{lemma}
\label{morphisms-lemma-base-change-dominant}
设 $f : X \to S$ 为概形的拟紧支配态射。设 $g : S' \to S$
为概形态射，并设 $f' : X' \to S'$ 为 $f$ 沿 $g$ 的基变换。
若一般化可沿 $g$ 提升，则 $f'$ 支配。
\end{lemma}

\noindent
我们将看到，若 $g$ 开（引理 \ref{morphisms-lemma-open-generizing}）或平坦
（引理 \ref{morphisms-lemma-generalizations-lift-flat}），则一般化可沿 $g$ 提升。

\begin{proof}
注意，由《概形》引理
\ref{schemes-lemma-quasi-compact-preserved-base-change}，$f'$ 拟紧。
设 $\eta' \in S'$ 为 $S'$ 某个不可约分支的泛点。
若一般化可沿 $g$ 提升，则 $\eta = g(\eta')$ 是 $S$ 某个不可约分支
的泛点。由引理 \ref{morphisms-lemma-quasi-compact-dominant} 可知，$\eta$
属于 $f$ 的像。于是由《概形》引理
\ref{schemes-lemma-points-fibre-product}，$\eta'$ 属于 $f'$ 的像。
再由引理 \ref{morphisms-lemma-quasi-compact-dominant}，$f'$ 支配。
\end{proof}

\noindent
其中一部分对任意支配态射（不必拟紧）仍成立。

\begin{lemma}
\label{morphisms-lemma-open-base-change-dominant}
设 $f : X \to S$ 为概形的支配态射。设 $g : S' \to S$
为概形的开态射，并设 $f' : X' \to S'$ 为 $f$ 沿 $g$ 的基变换。
则 $f'$ 支配。
\end{lemma}

\begin{proof}
若 $f'$ 不支配，则 $f'(X')$ 包含于某个闭子集 $B$ 中，且该子集
不等于 $S'$。于是 $S'-B$ 是 $S'$ 中非空的开集。由于 $g$ 开，
$g(S'-B)$ 是 $S$ 中非空的开集。由于 $f(X)$ 在 $S$ 中稠密，
存在点 $x\in X$，使得 $f(x) \in g(S'-B)$。令
$W = \Spec(\kappa (f(x)))$，并用 $X_W$、$S_W$、$S'_W$ 表示这些概形
从 $S$ 到 $W$ 的拉回。则 $X_W$ 与 $(S'-B)_W$ 都非空，因而
$X_W \times_W (S'-B)_W$ 非空，因为 $W$ 是域的谱。
这与 $f'(X \times_S S')$ 包含于 $B \subset S'$ 矛盾。
所以 $f'$ 确实支配。
\end{proof}

\begin{lemma}
\label{morphisms-lemma-quasi-compact-generic-point-not-in-image}
设 $f : X \to S$ 为概形的拟紧态射。设 $\eta \in S$ 为 $S$
某个不可约分支的泛点。若 $\eta \not \in f(X)$，则存在开邻域
$V \subset S$，它是 $\eta$ 的邻域，并使得
$f^{-1}(V) = \emptyset$。
\end{lemma}

\begin{proof}
设 $Z \subset S$ 为 $f$ 的概形论像。需要证明
$\eta \not \in Z$。这由引理
\ref{morphisms-lemma-reach-points-scheme-theoretic-image} 得出，也可如下看出。
由引理 \ref{morphisms-lemma-quasi-compact-scheme-theoretic-image}，态射
$X \to Z$ 支配；由引理 \ref{morphisms-lemma-quasi-compact-dominant}，
这意味着 $Z$ 的所有不可约分支的所有泛点都属于 $X \to Z$ 的像。
由假设可知 $\eta \not \in Z$，因为若 $\eta$ 属于 $Z$，它将是
$Z$ 某个不可约分支的泛点。
\end{proof}

\noindent
还有一种情形，其中支配等同于所有不可约分支的泛点都属于像。

\begin{lemma}
\label{morphisms-lemma-dominant-finite-number-irreducible-components}
设 $f : X \to S$ 为概形态射。假设 $X$ 只有有限多个不可约分支。
则 $f$ 支配（当且）仅当对每个不可约分支 $Z \subset S$，$Z$ 的泛点
属于 $f$ 的像。若如此，则 $S$ 也只有有限多个不可约分支。
\end{lemma}

\begin{proof}
假设 $f$ 支配。设 $X = Z_1 \cup Z_2 \cup \ldots \cup Z_n$ 为 $X$
分解成不可约分支的分解。设 $\xi_i \in Z_i$ 为其泛点，因而
$Z_i = \overline{\{\xi_i\}}$。注意，$f(Z_i)$ 是 $S$ 的不可约子集。
因此
$$
S = \overline{f(X)} = \bigcup \overline{f(Z_i)} =
\bigcup \overline{\{f(\xi_i)\}}
$$
是有限多个不可约子集的并，而这些子集的泛点属于 $f$ 的像。
引理得证。
\end{proof}

\begin{lemma}
\label{morphisms-lemma-dominant-between-integral}
设 $f : X \to Y$ 为整概形之间的态射。下列条件等价：
\begin{enumerate}
\item $f$ 支配；
\item $f$ 将 $X$ 的泛点映到 $Y$ 的泛点；
\item 对某些非空仿射开集 $U \subset X$ 与 $V \subset Y$，若
$f(U) \subset V$，则环映射
$\mathcal{O}_Y(V) \to \mathcal{O}_X(U)$ 单射；
\item 对所有非空仿射开集 $U \subset X$ 与 $V \subset Y$，若
$f(U) \subset V$，则环映射
$\mathcal{O}_Y(V) \to \mathcal{O}_X(U)$ 单射；
\item 对某个 $x \in X$，其像为 $y = f(x) \in Y$，局部环映射
$\mathcal{O}_{Y, y} \to \mathcal{O}_{X, x}$ 单射；并且
\item 对所有 $x \in X$，其像为 $y = f(x) \in Y$，局部环映射
$\mathcal{O}_{Y, y} \to \mathcal{O}_{X, x}$ 单射。
\end{enumerate}
\end{lemma}

\begin{proof}
(1) 与 (2) 的等价性由引理
\ref{morphisms-lemma-dominant-finite-number-irreducible-components} 得出。
设 $U \subset X$ 与 $V \subset Y$ 为满足 $f(U) \subset V$ 的非空
仿射开集。回顾，环 $A = \mathcal{O}_X(U)$ 与
$B = \mathcal{O}_Y(V)$ 都是整环。态射 $f|_U : U \to V$ 对应于环映射
$\varphi : B \to A$。$X$ 与 $Y$ 的泛点分别对应于素理想
$(0) \subset A$ 与 $(0) \subset B$。因此，(2) 等价于条件
$(0) = \varphi^{-1}((0))$，即等价于 $\varphi$ 单射。
这样便看出 (2)、(3) 与 (4) 等价。类似地，给定 (5) 中的 $x$ 与
$y$，局部环 $\mathcal{O}_{X, x}$ 与 $\mathcal{O}_{Y, y}$ 都是整环，
而素理想 $(0) \subset \mathcal{O}_{X, x}$ 与
$(0) \subset \mathcal{O}_{Y, y}$ 对应于 $X$ 与 $Y$ 的泛点
（通过将 $x$ 处局部环的谱认同为特殊化到 $x$ 的点集；见《概形》引理
\ref{schemes-lemma-specialize-points}）。因此可以完全按照上面的方式
论证，得到 (2)、(5) 与 (6) 等价。
\end{proof}






\section{满射态射}
\label{morphisms-section-surjective}

\begin{definition}
\label{morphisms-definition-surjective}
若概形态射在底拓扑空间上满射，则称它{\it 满射}。
\end{definition}

\begin{lemma}
\label{morphisms-lemma-composition-surjective}
满射态射的复合满射。
\end{lemma}

\begin{proof}
从略。
\end{proof}

\begin{lemma}
\label{morphisms-lemma-when-point-maps-to-pair}
设 $X$ 与 $Y$ 为基概形 $S$ 上的概形。给定点 $x \in X$ 与
$y \in Y$，存在 $X \times_S Y$ 中的点，在投影下分别映到 $x$ 与
$y$，当且仅当 $x$ 与 $y$ 位于 $S$ 的同一点之上。
\end{lemma}

\begin{proof}
该条件显然必要，反向则来自《概形》引理
\ref{schemes-lemma-points-fibre-product} 的证明。
\end{proof}

\begin{lemma}
\label{morphisms-lemma-base-change-surjective}
满射态射的基变换满射。
\end{lemma}

\begin{proof}
设 $f: X \to Y$ 为基概形 $S$ 上的概形态射。
若 $S' \to S$ 为概形态射，设 $p: X_{S'} \to X$ 与
$q: Y_{S'} \to Y$ 为典范投影。交换方块
$$
\xymatrix{
X_{S'} \ar[d]_{f_{S'}} \ar[r]_p & X \ar[d]^{f} \\
Y_{S'} \ar[r]^{q} & Y.
}
$$
将 $X_{S'}$ 认同为 $X \to Y$ 与 $Y_{S'} \to Y$ 的纤维积。
设 $Z$ 为 $X$ 的底拓扑空间的子集。则
$q^{-1}(f(Z)) = f_{S'}(p^{-1}(Z))$，因为
$y' \in q^{-1}(f(Z))$ 当且仅当 $q(y') = f(x)$ 对某个
$x \in Z$ 成立；由引理 \ref{morphisms-lemma-when-point-maps-to-pair}，
这又当且仅当存在 $x' \in X_{S'}$，使得
$f_{S'}(x') = y'$ 且 $p(x') = x$。特别地，取 $Z = X$，可知若
$f$ 满射，则其基变换 $f_{S'}: X_{S'} \to Y_{S'}$ 也满射。
\end{proof}

\begin{example}
\label{morphisms-example-injective-not-preserved-base-change}
双射性在基变换下不稳定，因而单射性也不稳定。例如考虑双射
$\Spec(\mathbf{C}) \to \Spec(\mathbf{R})$。基变换
$\Spec(\mathbf{C} \otimes_{\mathbf{R}} \mathbf{C}) \to
\Spec(\mathbf{C})$
不单射，因为存在同构
$\mathbf{C} \otimes_{\mathbf{R}} \mathbf{C} \cong \mathbf{C} \times \mathbf{C}$
（该分解来自幂等元
$\frac{1 \otimes 1 + i \otimes i}{2}$），因而
$\Spec(\mathbf{C} \otimes_{\mathbf{R}} \mathbf{C})$ 有两个点。
\end{example}

\begin{lemma}
\label{morphisms-lemma-surjection-from-quasi-compact}
设
$$
\xymatrix{
X \ar[rr]_f \ar[rd]_p & &
Y \ar[dl]^q \\
& Z
}
$$
为概形态射的交换图。若 $f$ 满射且 $p$ 拟紧，则 $q$ 拟紧。
\end{lemma}

\begin{proof}
设 $W \subset Z$ 为拟紧开集。由假设，$p^{-1}(W)$ 拟紧。
因此由《拓扑》引理 \ref{topology-lemma-image-quasi-compact}，
逆像 $q^{-1}(W) = f(p^{-1}(W))$ 也拟紧。引理得证。
\end{proof}

\section{根态射与泛单射态射}
\label{morphisms-section-radicial}

\noindent
本节定义概形态射何时为{\it 根态射}，以及何时为
{\it 泛单射}，并证明这两个概念等价。之所以同时引入二者，
是因为对代数空间而言也有相应概念，而它们未必总是等价。

\begin{definition}
\label{morphisms-definition-universally-injective}
设 $f : X \to S$ 为一个态射。
\begin{enumerate}
\item 称 $f$ 为{\it 泛单射}，如果且仅如果对概形的任意态射
$S' \to S$，基变换 $f' : X_{S'} \to S'$ 在底层拓扑空间上为单射。
\item 称 $f$ 为{\it 根态射}，如果 $f$ 作为拓扑空间之间的映射
为单射，且对每个 $x \in X$，域扩张
$\kappa(x)/\kappa(f(x))$ 都纯不可分。
\end{enumerate}
\end{definition}

\begin{lemma}
\label{morphisms-lemma-universally-injective}
\begin{reference}
\cite[Proposition 3.7.1]{EGA1-second}
\end{reference}
设 $f : X \to S$ 为概形态射。下列条件等价：
\begin{enumerate}
\item 对每个域 $K$，诱导映射
$\Mor(\Spec(K), X) \to \Mor(\Spec(K), S)$
都是单射。
\item 态射 $f$ 泛单射。
\item 态射 $f$ 是根态射。
\item 对角态射 $\Delta_{X/S} : X \longrightarrow X \times_S X$
满射。
\end{enumerate}
\end{lemma}

\begin{proof}
设 $K$ 为一个域，$s : \Spec(K) \to S$ 为一个态射。
给出一个态射 $x : \Spec(K) \to X$ 使得 $f \circ x = s$，
等价于给出投影 $X_K = \Spec(K) \times_S X \to \Spec(K)$
的一个截面，而这又等价于给出一个点 $x \in X_K$，
其剩余域为 $K$。因此 (2) 蕴含 (1)。

\medskip\noindent
反之，设 (1) 成立。假设 $x, x' \in X_{S'}$ 映到同一点
$s' \in S'$。选取域的交换图
$$
\xymatrix{
K & \kappa(x) \ar[l] \\
\kappa(x') \ar[u] & \kappa(s') \ar[l] \ar[u]
}
$$
由《概形》引理 \ref{schemes-lemma-characterize-points}，
得到两个态射 $a, a' : \Spec(K) \to X_{S'}$。其中一个对应于
点 $x$ 及嵌入 $\kappa(x) \subset K$，另一个对应于点 $x'$ 及嵌入
$\kappa(x') \subset K$。此外，$f' \circ a = f' \circ a'$。
条件 (1) 于是推出 $a$ 与 $a'$ 同 $X_{S'} \to X$ 的复合相等。
由于 $X_{S'}$ 是 $S'$ 与 $X$ 在 $S$ 上的纤维积，故
$a = a'$，从而 $x = x'$。因此 (1) 蕴含 (2)。

\medskip\noindent
若有两个不同的点 $x, x' \in X$ 映到同一个点 $s$，
则 (2) 不成立。若对某个 $s = f(x)$、$x \in X$，域扩张
$\kappa(x)/\kappa(s)$ 不是纯不可分的，则可以找到域扩张
$K/\kappa(s)$，使得 $\kappa(x)$ 有两个 $\kappa(s)$-同态到
$K$。由《概形》引理
\ref{schemes-lemma-characterize-points}，这意味着映射
$\Mor(\Spec(K), X) \to \Mor(\Spec(K), S)$
不是单射，因而 (1) 不成立。因此，等价条件 (1) 与 (2)
蕴含 $f$ 是根态射，即蕴含 (3)。

\medskip\noindent
假设 (3)。由《概形》引理
\ref{schemes-lemma-characterize-points}，一个态射
$\Spec(K) \to X$ 由一对 $(x, \kappa(x) \to K)$ 给出。
性质 (3) 恰好说明，把这一对 $(x, \kappa(x) \to K)$
对应到 $(s, \kappa(s) \to \kappa(x) \to K)$ 的映射是单射。
换言之，(1) 成立。至此已知 (1)、(2)、(3) 彼此等价。

\medskip\noindent
最后证明 (4) 与 (1)、(2)、(3) 等价。$X \times_S X$ 的一个点
由四元组 $(x_1, x_2, s, \mathfrak p)$ 给出，其中
$x_1, x_2 \in X$，$f(x_1) = f(x_2) = s$，并且
$\mathfrak p \subset \kappa(x_1) \otimes_{\kappa(s)} \kappa(x_2)$
是一个素理想；见《概形》引理
\ref{schemes-lemma-points-fibre-product}。若 $f$ 泛单射，
则在泛单射的定义中取 $S'=X$。由于 $\Delta_{X/S}$ 是单射态射
$X \times_S X  \longrightarrow X$ 的一个截面，它必为满射。
反之，若 $\Delta_{X/S}$ 满射，则总有 $x_1 = x_2 = x$，
且恰有一个这样的素理想 $\mathfrak p$；这说明
$\kappa(s) \subset \kappa(x)$ 纯不可分。因此 $f$ 是根态射。
或者，若 $\Delta_{X/S}$ 满射，则对任意 $S' \to S$，
其基变换 $\Delta_{X_{S'}/S'}$ 都满射，这推出 $f$ 泛单射。
引理证毕。
\end{proof}

\begin{lemma}
\label{morphisms-lemma-universally-injective-separated}
泛单射态射是分离的。
\end{lemma}

\begin{proof}
将引理 \ref{morphisms-lemma-universally-injective} 与以下事实结合即可：
$X \to S$ 分离，当且仅当 $\Delta_{X/S}$ 的像在
$X \times_S X$ 中闭；见《概形》定义
\ref{schemes-definition-separated} 及其后的讨论。
\end{proof}

\begin{lemma}
\label{morphisms-lemma-base-change-universally-injective}
泛单射态射的基变换仍泛单射。
\end{lemma}

\begin{proof}
这是形式的。
\end{proof}

\begin{lemma}
\label{morphisms-lemma-composition-universally-injective}
根态射的复合仍为根态射，因而等价的泛单射条件在复合下也保持。
\end{lemma}

\begin{proof}
略。
\end{proof}










\section{仿射态射}
\label{morphisms-section-affine}

\begin{definition}
\label{morphisms-definition-affine}
若概形态射 $f : X \to S$ 满足：$S$ 的每个仿射开集的逆像
都是 $X$ 的仿射开集，则称它为{\it 仿射的}。
\end{definition}

\begin{lemma}
\label{morphisms-lemma-affine-separated}
仿射态射是分离且拟紧的。
\end{lemma}

\begin{proof}
设 $f : X \to S$ 仿射。由《概形》引理
\ref{schemes-lemma-quasi-compact-affine} 立即得到拟紧性。
我们用《概形》引理 \ref{schemes-lemma-characterize-separated}
证明 $f$ 分离。设 $x_1, x_2 \in X$ 为 $X$ 中映到同一点
$s \in S$ 的两个点。任取仿射开集 $W \subset S$，使其包含 $s$。
依假设，$f^{-1}(W)$ 仿射。在上述引理中取
$U = V = f^{-1}(W)$ 即可。
\end{proof}

\begin{lemma}
\label{morphisms-lemma-characterize-affine}
\begin{reference}
\cite[II, Corollary 1.3.2]{EGA}
\end{reference}
设 $f : X \to S$ 为概形态射。下列条件等价：
\begin{enumerate}
\item 态射 $f$ 仿射。
\item 存在仿射开覆盖 $S = \bigcup W_j$，使每个
$f^{-1}(W_j)$ 都仿射。
\item 存在一个准凝聚 $\mathcal{O}_S$-代数层
$\mathcal{A}$，以及同构
$X \cong \underline{\Spec}_S(\mathcal{A})$；它是 $S$ 上概形的同构。记号见《构造》
第 \ref{constructions-section-spec} 节。
\end{enumerate}
此外，在这种情形下，$X = \underline{\Spec}_S(f_*\mathcal{O}_X)$。
\end{lemma}

\begin{proof}
(1) 蕴含 (2) 是显然的。

\medskip\noindent
设 $S = \bigcup_{j \in J} W_j$ 是仿射开覆盖，且每个
$f^{-1}(W_j)$ 都仿射。由《概形》引理
\ref{schemes-lemma-quasi-compact-affine}，$f$ 拟紧。由《概形》引理
\ref{schemes-lemma-characterize-quasi-separated}，态射 $f$ 拟分离。
因此由《概形》引理 \ref{schemes-lemma-push-forward-quasi-coherent}，
层 $\mathcal{A} = f_*\mathcal{O}_X$ 是准凝聚
$\mathcal{O}_S$-代数层。于是有概形
$g : Y = \underline{\Spec}_S(\mathcal{A}) \to S$，它在 $S$ 上。
恒等映射
$\text{id} : \mathcal{A} = f_*\mathcal{O}_X \to f_*\mathcal{O}_X$
由相对谱的定义给出态射 $can : X \to Y$，它在 $S$ 上；见《构造》
引理 \ref{constructions-lemma-canonical-morphism}。由假设和刚引用的引理，
限制 $can|_{f^{-1}(W_j)} : f^{-1}(W_j) \to g^{-1}(W_j)$
是同构。因此 $can$ 是同构。我们已经证明 (2) 蕴含 (3)。

\medskip\noindent
假设 (3)。由《构造》引理
\ref{constructions-lemma-spec-properties}，每个仿射开集的逆像都仿射，
所以按定义该态射仿射。
\end{proof}

\begin{remark}
\label{morphisms-remark-direct-argument}
也可如下直接证明上述引理 \ref{morphisms-lemma-characterize-affine} 中
(2) 蕴含 (1)。设 $S = \bigcup W_j$ 是仿射开覆盖，且每个
$f^{-1}(W_j)$ 都仿射。先如上述证明那样论证
$\mathcal{A} = f_*\mathcal{O}_X$ 准凝聚。
设 $\Spec(R) = V \subset S$ 为仿射开集。需要证明
$f^{-1}(V)$ 仿射。令
$A = \mathcal{A}(V) = f_*\mathcal{O}_X(V) = \mathcal{O}_X(f^{-1}(V))$。
由《概形》引理 \ref{schemes-lemma-morphism-into-affine}，存在
典范态射 $\psi : f^{-1}(V) \to \Spec(A)$，它在
$\Spec(R) = V$ 上。
由《概形》引理 \ref{schemes-lemma-good-subcover}，存在整数
$n \geq 0$、标准开覆盖
$V = \bigcup_{i = 1, \ldots, n} D(h_i)$（$h_i \in R$），
以及映射 $a : \{1, \ldots, n\} \to J$，使得每个 $D(h_i)$
也是仿射概形 $W_{a(i)}$ 的标准开集。仿射概形态射下标准开集的逆像
仍为标准开集，见《代数》引理 \ref{algebra-lemma-spec-functorial}。
因此 $f^{-1}(D(h_i))$ 是 $f^{-1}(W_{a(i)})$ 的标准开集，
特别地，$f^{-1}(D(h_i))$ 仿射。由于 $\mathcal{A}$ 准凝聚，
$A_{h_i} = \mathcal{A}(D(h_i)) = \mathcal{O}_X(f^{-1}(D(h_i)))$，
所以 $f^{-1}(D(h_i))$ 是 $A_{h_i}$ 的谱。由此，态射 $\psi$
把开集 $f^{-1}(D(h_i))$ 与开集 $\Spec(A_{h_i})$ 同构起来，
后者是 $\Spec(A)$ 的开集。由于 $f^{-1}(V) = \bigcup f^{-1}(D(h_i))$
且 $\Spec(A) = \bigcup \Spec(A_{h_i})$，结论成立。
\end{remark}


\begin{lemma}
\label{morphisms-lemma-affine-equivalence-algebras}
设 $S$ 为概形。存在范畴的反等价
$$
\begin{matrix}
\text{Schemes affine} \\
\text{相对于 }S
\end{matrix}
\longleftrightarrow
\begin{matrix}
\text{quasi-coherent sheaves} \\
\text{属于 }\mathcal{O}_S\text{-algebras}
\end{matrix}
$$
它把 $f : X \to S$ 对应到层 $f_*\mathcal{O}_X$。
此外，这一等价与任意基变换相容。
\end{lemma}

\begin{proof}
从右到左的函子由 $\underline{\Spec}_S$ 给出。
由引理 \ref{morphisms-lemma-characterize-affine} 和《构造》引理
\ref{constructions-lemma-spec-properties} 的 (3)，这两个函子互逆。
最后一个断言即《构造》引理
\ref{constructions-lemma-spec-properties} 的 (2)。
\end{proof}

\begin{lemma}
\label{morphisms-lemma-quasi-coherence-scalar-restriction}
\begin{reference}
\cite[I, Proposition 9.6.1]{EGA}
\end{reference}
设 $S$ 为概形，$\mathcal{A}$ 为准凝聚 $\mathcal{O}_S$-代数。
一个 $\mathcal{A}$-模作为 $\mathcal{O}_S$-模准凝聚，当且仅当
它作为 $\mathcal{A}$-模准凝聚。
\end{lemma}

\begin{proof}
设 $\mathcal{F}$ 为一个 $\mathcal{A}$-模。若 $\mathcal{F}$ 作为
$\mathcal{A}$-模准凝聚，则对每个 $s \in S$，存在开邻域
$U$（它包含 $s$）以及 $\mathcal{A}|_U$-模的正合列
$$
\bigoplus\nolimits_J\mathcal{A}|_U\to
\bigoplus\nolimits_I\mathcal{A}|_U\to
\mathcal{F}|_U
$$
它同时也是 $\mathcal{O}_U$-模的正合列。因此，$\mathcal{F}|_U$
作为概形上两个准凝聚 $\mathcal{O}_U$-模之间态射的余核而准凝聚。
所以 $\mathcal{F}$ 作为 $\mathcal{O}_X$-模准凝聚。

\medskip\noindent
反之，假设 $\mathcal{F}$ 作为 $\mathcal{O}_X$-模准凝聚。
取仿射开集 $\Spec(R) = U \subset S$。有 $\mathcal{O}_U$-模同构
$\mathcal{A}|_U \cong \widetilde{A}$ 与
$\mathcal{F}|_U \cong \widetilde{M}$，其中 $R$-代数 $A$ 和
$R$-模 $M$ 适当选取。$\mathcal{A}$-模在 $\mathcal{F}$ 上的结构
转化为 $A$-模在 $M$ 上的结构，并与给定的 $R$-模结构相容
（细节略）。选取 $A$-模的正合列
$$
\bigoplus\nolimits_J A \to
\bigoplus\nolimits_I A \to
M \to 0
$$
由于函子 $\widetilde{\ }$ 正合，这给出 $\widetilde{A}$-模的正合列
$$
\bigoplus\nolimits_J \widetilde{A}\to
\bigoplus\nolimits_I \widetilde{A}\to
\widetilde{M} \to 0
$$
这说明 $\mathcal{F}$ 作为 $\mathcal{A}$-模准凝聚。
\end{proof}

\begin{lemma}
\label{morphisms-lemma-affine-equivalence-modules}
设 $f : X \to S$ 为概形的仿射态射，并令
$\mathcal{A} = f_*\mathcal{O}_X$。函子
$\mathcal{F} \mapsto f_*\mathcal{F}$ 诱导范畴等价
$$
\left\{
\begin{matrix}
\text{category of quasi-coherent}\\
\mathcal{O}_X\text{-模}
\end{matrix}
\right\}
\longrightarrow
\left\{
\begin{matrix}
\text{category of quasi-coherent}\\
\mathcal{A}\text{-模}
\end{matrix}
\right\}
$$
此外，一个 $\mathcal{A}$-模作为 $\mathcal{O}_S$-模准凝聚，
当且仅当它作为 $\mathcal{A}$-模准凝聚。
\end{lemma}

\begin{proof}
最后一个断言即引理 \ref{morphisms-lemma-quasi-coherence-scalar-restriction}。
由引理 \ref{morphisms-lemma-affine-separated} 和《概形》引理
\ref{schemes-lemma-push-forward-quasi-coherent}，正向像
$f_*\mathcal{F}$（其中 $\mathcal{O}_X$-模 $\mathcal{F}$ 准凝聚）
是准凝聚 $\mathcal{O}_S$-模。因此得到引理所述的函子。

\medskip\noindent
我们为这一函子构造拟逆 $h$。设 $\mathcal{G}$ 为准凝聚
$\mathcal{A}$-模。定义
$$
h(\mathcal{G}) = f^*\mathcal{G} \otimes_{f^*\mathcal{A}} \mathcal{O}_X =
f^*\mathcal{G} \otimes_{f^*f_*\mathcal{O}_X} \mathcal{O}_X
$$
说明如下：拉回 $f^*\mathcal{A} = f^*f_*\mathcal{O}_X$ 是
$\mathcal{O}_X$-代数，伴随映射
$f^*f_*\mathcal{O}_X \to \mathcal{O}_X$ 是代数同态，而拉回
$f^*\mathcal{G}$ 是 $f^*\mathcal{A}$-模。由于拉回和换环保持
拟凝聚性，$h(\mathcal{G})$ 准凝聚。存在函子性映射
$$
h(f_*\mathcal{F}) =
f^*f_*\mathcal{F} \otimes_{f^*f_*\mathcal{O}_X} \mathcal{O}_X \to \mathcal{F}
$$
它来自伴随映射 $f^*f_*\mathcal{F} \to \mathcal{F}$；还存在函子性映射
$$
\mathcal{G} \to f_*h(\mathcal{G}) \quad\text{adjoint to the map}\quad
f^*\mathcal{G} \to f^*\mathcal{G} \otimes_{f^*f_*\mathcal{O}_X} \mathcal{O}_X
$$
它把局部截面 $s$（它属于 $f^*\mathcal{F}$）送到 $s \otimes 1$。
要完成证明，只需证明这些映射对上述 $\mathcal{F}$ 和
$\mathcal{G}$ 都是同构。这可以在一个仿射覆盖的各成员上检验，
也就是在 $X$ 和 $S$ 都仿射时检验。

\medskip\noindent
使上述论证成立的关键代数观察如下。设 $R \to A$ 为环同态，
$N$ 为 $A$-模，则
$$
(N \otimes_R A) \otimes_{(A \otimes_R A)} A = N
$$
即等式左端是把 $h$ 作用于准凝聚 $\widetilde{A}$-模
$\widetilde{N}$ 所得的结果；该模在 $\Spec(R)$ 上。细节略。
\end{proof}

\begin{lemma}
\label{morphisms-lemma-composition-affine}
仿射态射的复合仍为仿射态射。
\end{lemma}

\begin{proof}
设 $f : X \to Y$ 与 $g : Y \to Z$ 为仿射态射。
设 $U \subset Z$ 为仿射开集。$g^{-1}(U)$ 由关于 $g$ 的假设仿射；
继而 $f^{-1}(g^{-1}(U))$ 由关于 $f$ 的假设仿射。
因此 $(g \circ f)^{-1}(U)$ 仿射。
\end{proof}

\begin{lemma}
\label{morphisms-lemma-base-change-affine}
仿射态射的基变换仍为仿射态射。
\end{lemma}

\begin{proof}
设 $f : X \to S$ 为仿射态射，$S' \to S$ 为任意态射。
以 $f' : X_{S'} = S' \times_S X \to S'$ 表示 $f$ 的基变换。
对每个 $s' \in S'$，都存在仿射开邻域 $s' \in V \subset S'$，
它映入某个仿射开集 $U \subset S$。由假设，$f^{-1}(U)$ 仿射。
由《概形》第 \ref{schemes-section-fibre-products} 节的内容，
$f^{-1}(U)_V = V \times_U f^{-1}(U)$ 仿射且等于
$(f')^{-1}(V)$。这证明 $S'$ 有一个仿射开覆盖，其在 $f'$ 下的逆像
均仿射。最后应用上述引理 \ref{morphisms-lemma-characterize-affine}。
\end{proof}

\begin{lemma}
\label{morphisms-lemma-closed-immersion-affine}
闭浸入是仿射态射。
\end{lemma}

\begin{proof}
这一事实最初见于《概形》引理
\ref{schemes-lemma-closed-immersion-affine-case}；完整陈述见《概形》引理
\ref{schemes-lemma-closed-subspace-scheme}。
\end{proof}

\begin{lemma}
\label{morphisms-lemma-affine-s-open}
设 $X$ 为概形，$\mathcal{L}$ 为可逆 $\mathcal{O}_X$-模，且
$s \in \Gamma(X, \mathcal{L})$。包含态射 $j : X_s \to X$ 仿射。
\end{lemma}

\begin{proof}
这由《性质》引理 \ref{properties-lemma-affine-cap-s-open} 及定义推出。
\end{proof}

\begin{lemma}
\label{morphisms-lemma-affine-permanence}
设 $g : X \to Y$ 为 $S$ 上的概形态射。
\begin{enumerate}
\item 若 $X$ 在 $S$ 上仿射，且 $\Delta : Y \to Y \times_S Y$ 仿射，
则 $g$ 仿射。
\item 若 $X$ 在 $S$ 上仿射，且 $Y$ 在 $S$ 上分离，则 $g$ 仿射。
\item 从仿射概形到具有仿射对角态射的概形的态射是仿射的。
\item 从仿射概形到分离概形的态射是仿射的。
\end{enumerate}
\end{lemma}

\begin{proof}
证明 (1)。基变换 $X \times_S Y \to Y$ 由引理
\ref{morphisms-lemma-base-change-affine} 仿射。态射
$(1, g) : X \to X \times_S Y$ 是 $Y \to Y \times_S Y$
沿态射 $X \times_S Y \to Y \times_S Y$ 的基变换。因此由引理
\ref{morphisms-lemma-base-change-affine}，它仿射。仿射态射的复合仿射
（见引理 \ref{morphisms-lemma-composition-affine}），故 (1) 成立。
由 (1) 得 (2)，因为闭浸入仿射（见引理
\ref{morphisms-lemma-closed-immersion-affine}），而 $Y/S$ 分离意味着
$\Delta$ 是闭浸入。(3) 与 (4) 分别是 (1) 与 (2) 的特殊情形。
\end{proof}

\begin{lemma}
\label{morphisms-lemma-morphism-affines-affine}
仿射概形之间的态射是仿射的。
\end{lemma}

\begin{proof}
在引理 \ref{morphisms-lemma-affine-permanence} 中取
$S = \Spec(\mathbf{Z})$ 即得。也可由引理
\ref{morphisms-lemma-characterize-affine} 中 (1) 与 (2) 的等价直接得到。
\end{proof}

\begin{lemma}
\label{morphisms-lemma-Artinian-affine}
设 $S$ 为概形，$A$ 为阿廷环。任意态射 $\Spec(A) \to S$ 都仿射。
\end{lemma}

\begin{proof}
略。
\end{proof}

\begin{lemma}
\label{morphisms-lemma-get-affine}
设 $j : Y \to X$ 为概形的浸入。假设存在开集 $U \subset X$，
其补集为 $Z = X \setminus U$，且满足
\begin{enumerate}
\item $U \to X$ 仿射；
\item $j^{-1}(U) \to U$ 仿射；
\item $j(Y) \cap Z$ 闭。
\end{enumerate}
则 $j$ 仿射。特别地，若 $X$ 仿射，则 $Y$ 也仿射。
\end{lemma}

\begin{proof}
由《概形》定义 \ref{schemes-definition-immersion}，存在开子概形
$W \subset X$，使 $j$ 分解为闭浸入 $i : Y \to W$，随后接包含态射
$W \to X$。由于闭浸入仿射（引理
\ref{morphisms-lemma-closed-immersion-affine}），对每个 $x \in W$，
$x$ 在 $X$ 中都有一个仿射开邻域，其在 $j$ 下的逆像仿射。
若 $x \in U$，则由假设 (2) 同样成立。最后，设 $x \in Z$
且 $x \not \in W$。于是 $x \not \in j(Y) \cap Z$。
由假设 (3)，可找到仿射开邻域 $V \subset X$，它包含 $x$
且与 $j(Y) \cap Z$ 不相交。此时
$j^{-1}(V) = j^{-1}(V \cap U)$，由假设 (1)、(2) 它仿射。
因此由引理 \ref{morphisms-lemma-characterize-affine}，$j$ 仿射。
\end{proof}

\section{丰沛可逆模族}
\label{morphisms-section-families-ample-invertible-modules}

\noindent
本节简述丰沛可逆模族的概念。

\begin{definition}
\label{morphisms-definition-family-ample-invertible-modules}
\begin{reference}
\cite[II Definition 2.2.4]{SGA6}
\end{reference}
设 $X$ 为概形，$\{\mathcal{L}_i\}_{i \in I}$ 为一族可逆
$\mathcal{O}_X$-模。若满足下列条件，则称
$\{\mathcal{L}_i\}_{i \in I}$ 为{\it $X$ 上的丰沛可逆模族}：
\begin{enumerate}
\item $X$ 拟紧；
\item 对每个 $x \in X$，存在 $i \in I$、$n \geq 1$ 以及
$s \in \Gamma(X, \mathcal{L}_i^{\otimes n})$，使得
$x \in X_s$ 且 $X_s$ 仿射。
\end{enumerate}
\end{definition}

\noindent
若 $\{\mathcal{L}_i\}_{i \in I}$ 是概形 $X$ 上的丰沛可逆模族，
则存在有限子集 $I' \subset I$，使得
$\{\mathcal{L}_i\}_{i \in I'}$ 仍是 $X$ 上的丰沛可逆模族
（这由拟紧性立即推出）。由下一个引理，具有丰沛可逆模族的概形
具有仿射对角态射，因而尤其拟分离。

\begin{lemma}
\label{morphisms-lemma-affine-diagonal}
设 $X$ 为概形，并且对每个点 $x \in X$，都存在可逆
$\mathcal{O}_X$-模 $\mathcal{L}$ 和全局截面
$s \in \Gamma(X, \mathcal{L})$，使得 $x \in X_s$ 且
$X_s$ 仿射。那么 $X$ 的对角态射仿射。
\end{lemma}

\begin{proof}
给定可逆 $\mathcal{O}_X$-模 $\mathcal{L}$、$\mathcal{M}$ 以及全局截面
$s \in \Gamma(X, \mathcal{L})$、
$t \in \Gamma(X, \mathcal{M})$，使得 $X_s$ 与 $X_t$ 仿射，
我们只需证明 $X_s \cap X_t$ 仿射。事实上，把引理
\ref{morphisms-lemma-characterize-affine} 应用于
$\Delta : X \to X \times X$，并利用
$\Delta^{-1}(X_s \times X_t) = X_s \cap X_t$，即可看出
$\Delta$ 仿射。而 $X_s \cap X_t$ 仿射则由《性质》引理
\ref{properties-lemma-affine-cap-s-open} 得到。
\end{proof}

\begin{remark}
\label{morphisms-remark-affine-s-opens-cover-family}
《性质》引理 \ref{properties-lemma-affine-s-opens-cover-quasi-separated}
表明，具有一个丰沛可逆模的概形是分离的。对于具有丰沛可逆模族的概形，
这一断言并不成立。具体地，取《概形》例
\ref{schemes-example-affine-space-zero-doubled} 中的 $X$，令
$n = 1$，即零点加倍的仿射直线。沿用该例的记号，但以 $x$
表示 $x_1$，以 $y$ 表示 $y_1$。对每个整数 $n$，都有一个可逆层
$\mathcal{L}_n$ 在 $X$ 上；它在 $X_1$ 与 $X_2$ 上平凡，且其转移函数
$U_{12} \to U_{21}$ 为 $f(x) \mapsto y^n f(y)$。
$\mathcal{L}_n$ 的全局截面是二元组
$(f(x), g(y)) \in k[x] \oplus
k[y]$，满足 $y^n f(y) = g(y)$。截面 $s = (1,
y)$ 属于 $\mathcal{L}_1$，截面 $t = (x, 1)$ 属于
$\mathcal{L}_{-1}$；二者确定一个仿射开覆盖，
因为 $X_s = X_1$ 且 $X_t = X_2$。因此 $X$ 具有丰沛可逆模族，
但它不分离。
\end{remark}

\section{拟仿射态射}
\label{morphisms-section-quasi-affine}

\noindent
回忆：若概形 $X$ 拟紧且同构于某个仿射概形的开子概形，
则称它为{\it 拟仿射的}；见《性质》定义
\ref{properties-definition-quasi-affine}。

\begin{definition}
\label{morphisms-definition-quasi-affine}
若概形态射 $f : X \to S$ 满足：$S$ 的每个仿射开集的逆像
都是拟仿射概形，则称它为{\it 拟仿射的}。
\end{definition}

\begin{lemma}
\label{morphisms-lemma-quasi-affine-separated}
拟仿射态射是分离且拟紧的。
\end{lemma}

\begin{proof}
设 $f : X \to S$ 拟仿射。由《概形》引理
\ref{schemes-lemma-quasi-compact-affine} 立即得到拟紧性。
设 $U \subset S$ 为仿射开集。若能证明 $f^{-1}(U)$ 是分离概形，
则 $f$ 分离（《概形》引理
\ref{schemes-lemma-characterize-separated} 表明分离性在基底上是局部的）。
依假设，$f^{-1}(U)$ 同构于某个仿射概形的开子概形。
仿射概形分离，故其每个开子概形也分离，如所需。
\end{proof}

\begin{lemma}
\label{morphisms-lemma-characterize-quasi-affine}
设 $f : X \to S$ 为概形态射。下列条件等价：
\begin{enumerate}
\item 态射 $f$ 拟仿射。
\item 存在仿射开覆盖 $S = \bigcup W_j$，使每个
$f^{-1}(W_j)$ 都拟仿射。
\item 存在准凝聚 $\mathcal{O}_S$-代数层 $\mathcal{A}$ 以及拟紧开浸入
$$
\xymatrix{
X \ar[rr] \ar[rd] & & \underline{\Spec}_S(\mathcal{A}) \ar[dl] \\
& S &
}
$$
在 $S$ 上。
\item 与 (3) 相同，但取 $\mathcal{A} = f_*\mathcal{O}_X$，
并令水平箭头为《构造》引理
\ref{constructions-lemma-canonical-morphism} 的典范态射。
\end{enumerate}
\end{lemma}

\begin{proof}
(1) 蕴含 (2) 以及 (4) 蕴含 (3) 都是显然的。

\medskip\noindent
设 $S = \bigcup_{j \in J} W_j$ 是仿射开覆盖，且每个
$f^{-1}(W_j)$ 都拟仿射。由《概形》引理
\ref{schemes-lemma-quasi-compact-affine}，$f$ 拟紧。由《概形》引理
\ref{schemes-lemma-characterize-quasi-separated}，态射 $f$ 拟分离。
因此由《概形》引理 \ref{schemes-lemma-push-forward-quasi-coherent}，
层 $\mathcal{A} = f_*\mathcal{O}_X$ 是准凝聚
$\mathcal{O}_X$-代数层。于是有概形
$g : Y = \underline{\Spec}_S(\mathcal{A}) \to S$，它在 $S$ 上。
恒等映射
$\text{id} : \mathcal{A} = f_*\mathcal{O}_X \to f_*\mathcal{O}_X$
由相对谱的定义给出态射 $can : X \to Y$，它在 $S$ 上；见《构造》
引理 \ref{constructions-lemma-canonical-morphism}。由假设、刚引用的引理
以及《性质》引理 \ref{properties-lemma-quasi-affine}，限制
$can|_{f^{-1}(W_j)} : f^{-1}(W_j) \to g^{-1}(W_j)$ 是拟紧开浸入。
因此 $can$ 是拟紧开浸入。我们已经证明 (2) 蕴含 (4)。

\medskip\noindent
假设 (3)。任取仿射开集 $U \subset S$。由《构造》引理
\ref{constructions-lemma-spec-properties}，$U$ 在相对谱中的逆像仿射。
因此 $f^{-1}(U)$ 拟仿射（注意拟紧性也已编码在 (3) 中）。
故 (3) 蕴含 (1)。
\end{proof}

\begin{lemma}
\label{morphisms-lemma-composition-quasi-affine}
拟仿射态射的复合仍为拟仿射态射。
\end{lemma}

\begin{proof}
设 $f : X \to Y$ 与 $g : Y \to Z$ 为拟仿射态射。
设 $U \subset Z$ 为仿射开集。$g^{-1}(U)$ 由关于 $g$ 的假设拟仿射。
设 $j : g^{-1}(U) \to V$ 是到仿射概形 $V$ 的拟紧开浸入。
由上述引理 \ref{morphisms-lemma-characterize-quasi-affine}，
$f^{-1}(g^{-1}(U))$ 是相对谱
$\underline{\Spec}_{g^{-1}(U)}(\mathcal{A})$ 的拟紧开子概形，其中有某个
准凝聚 $\mathcal{O}_{g^{-1}(U)}$-代数层 $\mathcal{A}$。
由《概形》引理 \ref{schemes-lemma-push-forward-quasi-coherent}，层
$\mathcal{A}' = j_*\mathcal{A}$ 是准凝聚 $\mathcal{O}_V$-代数层，
且满足 $j^*\mathcal{A}' = \mathcal{A}$。因此得到交换图
$$
\xymatrix{
f^{-1}(g^{-1}(U)) \ar[r] &
\underline{\Spec}_{g^{-1}(U)}(\mathcal{A})
\ar[r] \ar[d] &
\underline{\Spec}_V(\mathcal{A}') \ar[d] \\
& g^{-1}(U) \ar[r]^j & V
}
$$
其中方块为纤维方块；见《构造》引理
\ref{constructions-lemma-spec-properties}。注意右上角是仿射概形。
因此 $(g \circ f)^{-1}(U)$ 拟仿射。
\end{proof}

\begin{lemma}
\label{morphisms-lemma-base-change-quasi-affine}
拟仿射态射的基变换仍为拟仿射态射。
\end{lemma}

\begin{proof}
设 $f : X \to S$ 为拟仿射态射。由上述引理
\ref{morphisms-lemma-characterize-quasi-affine}，可以找到准凝聚
$\mathcal{O}_S$-代数层 $\mathcal{A}$ 和拟紧开浸入
$X \to \underline{\Spec}_S(\mathcal{A})$，后者在 $S$ 上。
设 $g : S' \to S$ 为任意态射。以
$f' : X_{S'} = S' \times_S X \to S'$ 表示 $f$ 的基变换。
拟紧开浸入的基变换仍为拟紧开浸入，故
$X_{S'} \to \underline{\Spec}_{S'}(g^*\mathcal{A})$
是拟紧开浸入（这里使用了《概形》引理
\ref{schemes-lemma-quasi-compact-preserved-base-change}、
\ref{schemes-lemma-base-change-immersion} 以及《构造》引理
\ref{constructions-lemma-spec-properties}）。再次由引理
\ref{morphisms-lemma-characterize-quasi-affine}，$X_{S'} \to S'$ 拟仿射。
\end{proof}

\begin{lemma}
\label{morphisms-lemma-quasi-compact-immersion-quasi-affine}
拟紧浸入是拟仿射态射。
\end{lemma}

\begin{proof}
设 $X \to S$ 为拟紧浸入。需要证明每个仿射开集的逆像拟仿射。
因此，可假设 $S$ 是仿射概形，并证明 $X$ 拟仿射。
由引理 \ref{morphisms-lemma-quasi-compact-immersion}，态射 $X \to S$ 分解为
$X \to Z \to S$，其中 $Z$ 是 $S$ 的闭子概形，且 $X \subset Z$
是拟紧开集。由于 $S$ 仿射，引理 \ref{morphisms-lemma-closed-immersion}
推出 $Z$ 仿射。结论得证。
\end{proof}

\begin{lemma}
\label{morphisms-lemma-affine-quasi-affine}
设 $S$ 为概形，$X$ 为仿射概形。态射 $f : X \to S$ 拟仿射，
当且仅当它拟紧。特别地，从仿射概形到拟分离概形的任意态射都拟仿射。
\end{lemma}

\begin{proof}
设 $V \subset S$ 为仿射开集。则 $f^{-1}(V)$ 是仿射概形 $X$ 的
开子概形，所以它拟仿射当且仅当它拟紧。这证明第一个断言。
若 $f : X \to S$ 中 $X$ 仿射而 $S$ 拟分离，则该态射的拟紧性
由《概形》引理 \ref{schemes-lemma-quasi-compact-permanence} 应用于
$X \to S \to \Spec(\mathbf{Z})$ 得到。
\end{proof}

\begin{lemma}
\label{morphisms-lemma-quasi-affine-permanence}
设 $g : X \to Y$ 为 $S$ 上的概形态射。若 $X$ 在 $S$ 上拟仿射，
且 $Y$ 在 $S$ 上拟分离，则 $g$ 拟仿射。特别地，从拟仿射概形
到拟分离概形的任意态射都拟仿射。
\end{lemma}

\begin{proof}
基变换 $X \times_S Y \to Y$ 由引理
\ref{morphisms-lemma-base-change-quasi-affine} 拟仿射。态射
$X \to X \times_S Y$ 是拟紧浸入，因为 $Y \to S$ 拟分离；见《概形》引理
\ref{schemes-lemma-section-immersion}。由引理
\ref{morphisms-lemma-quasi-compact-immersion-quasi-affine}，拟紧浸入拟仿射；
而拟仿射态射的复合拟仿射（见引理
\ref{morphisms-lemma-composition-quasi-affine}）。结论得证。
\end{proof}

\section{由环同态性质定义的态射类型}
\label{morphisms-section-properties-ring-maps}

\noindent
本节研究环同态的哪些性质可用来定义概形态射的局部性质。

\begin{definition}
\label{morphisms-definition-property-local}
设 $P$ 为环同态的一个性质。
\begin{enumerate}
\item 若下列条件成立，则称 $P$ 是{\it 局部的}：
\begin{enumerate}
\item 对任意环同态 $R \to A$ 以及任意 $f \in R$，有
$P(R \to A) \Rightarrow P(R_f \to A_f)$。
\item 对任意环 $R$、$A$，任意 $f \in R$、$a\in A$，以及任意环同态
$R_f \to A$，有 $P(R_f \to A) \Rightarrow P(R \to A_a)$。
\item 对任意环同态 $R \to A$，若 $a_i \in A$ 且
$(a_1, \ldots, a_n) = A$，则
$\forall i, P(R \to A_{a_i}) \Rightarrow P(R \to A)$。
\end{enumerate}
\item 称 $P$ 在{\it 基变换下稳定}，如果对任意环同态
$R \to A$、$R \to R'$ 都有
$P(R \to A) \Rightarrow P(R' \to R' \otimes_R A)$。
\item 称 $P$ 在{\it 复合下稳定}，如果对任意环同态
$A \to B$、$B \to C$ 都有
$P(A \to B) \wedge P(B \to C) \Rightarrow P(A \to C)$。
\end{enumerate}
\end{definition}

\begin{definition}
\label{morphisms-definition-locally-P}
设 $P$ 为环同态的一个性质，$f : X \to S$ 为概形态射。
称 $f$ {\it 局部为 $P$ 型}，如果对任意 $x \in X$，
存在仿射开邻域 $U$，它是 $x$ 在 $X$ 中的邻域，并映入某个仿射开集
$V \subset S$，而诱导环同态
$\mathcal{O}_S(V) \to \mathcal{O}_X(U)$ 具有性质 $P$。
\end{definition}

\noindent
除非性质 $P$ 是局部性质，否则这不是一个“好”定义。
即使 $P$ 是局部性质，我们也不会自动用这个定义称一个态射
“局部为 $P$ 型”，除非还在别处明确给出相应定义。

\begin{lemma}
\label{morphisms-lemma-locally-P}
设 $f : X \to S$ 为概形态射，$P$ 为环同态的一个性质。
设 $U$ 为 $X$ 的仿射开集，$V$ 为 $S$ 的仿射开集，且
$f(U) \subset V$。若 $f$ 局部为 $P$ 型且 $P$ 是局部的，
则 $P(\mathcal{O}_S(V) \to \mathcal{O}_X(U))$ 成立。
\end{lemma}

\begin{proof}
因为 $f$ 局部为 $P$ 型，所以对每个 $u \in U$，存在仿射开集
$U_u \subset X$，它映入仿射开集 $V_u \subset S$，且
$P(\mathcal{O}_S(V_u) \to \mathcal{O}_X(U_u))$ 成立。
选取开邻域 $U'_u \subset U \cap U_u$，使它包含 $u$ 且同时是
$U$ 与 $U_u$ 中的标准仿射开集；见《概形》引理
\ref{schemes-lemma-standard-open-two-affines}。由定义
\ref{morphisms-definition-property-local} (1)(b)，
$P(\mathcal{O}_S(V_u) \to \mathcal{O}_X(U'_u))$ 成立。
因此可以假设 $U_u \subset U$ 是标准仿射开集。
选取开邻域 $V'_u \subset V \cap V_u$，使它包含 $f(u)$ 且同时是
$V$ 与 $V_u$ 中的标准仿射开集；见《概形》引理
\ref{schemes-lemma-standard-open-two-affines}。于是
$U'_u = f^{-1}(V'_u) \cap U_u$ 是 $U_u$（因而也是 $U$）中的
标准仿射开集，并由定义 \ref{morphisms-definition-property-local} (1)(a) 有
$P(\mathcal{O}_S(V'_u) \to \mathcal{O}_X(U'_u))$。
因此可以假设 $U_u \subset U$ 与 $V_u \subset V$ 都是标准仿射开集。
再应用一次定义 \ref{morphisms-definition-property-local} (1)(b)，得到
$P(\mathcal{O}_S(V) \to \mathcal{O}_X(U_u))$。因为 $U$ 拟紧，
可选取有限多个点 $u_1, \ldots, u_n \in U$，使得
$$
U = U_{u_1} \cup \ldots \cup U_{u_n}.
$$
由定义 \ref{morphisms-definition-property-local} (1)(c)，得
$P(\mathcal{O}_S(V) \to \mathcal{O}_X(U))$。
\end{proof}

\begin{lemma}
\label{morphisms-lemma-locally-P-characterize}
设 $P$ 为环同态的局部性质，$f : X \to S$ 为概形态射。
下列条件等价：
\begin{enumerate}
\item 态射 $f$ 局部为 $P$ 型。
\item 对每一对仿射开集 $U \subset X$、$V \subset S$，若
$f(U) \subset V$，则 $P(\mathcal{O}_S(V) \to \mathcal{O}_X(U))$ 成立。
\item 存在开覆盖 $S = \bigcup_{j \in J} V_j$ 以及开覆盖
$f^{-1}(V_j) = \bigcup_{i \in I_j} U_i$，使每个态射
$U_i \to V_j$（$j\in J, i\in I_j$）都局部为 $P$ 型。
\item 存在仿射开覆盖 $S = \bigcup_{j \in J} V_j$ 以及仿射开覆盖
$f^{-1}(V_j) = \bigcup_{i \in I_j} U_i$，使得
$P(\mathcal{O}_S(V_j) \to \mathcal{O}_X(U_i))$ 对所有
$j\in J, i\in I_j$ 都成立。
\end{enumerate}
此外，若 $f$ 局部为 $P$ 型，则对任意开子概形
$U \subset X$、$V \subset S$，只要 $f(U) \subset V$，限制
$f|_U : U \to V$ 就局部为 $P$ 型。
\end{lemma}

\begin{proof}
这由上述引理 \ref{morphisms-lemma-locally-P} 推出。
\end{proof}

\begin{lemma}
\label{morphisms-lemma-composition-type-P}
设 $P$ 为环同态的一个性质，并假设 $P$ 是局部的且在复合下稳定。
局部为 $P$ 型的态射之复合仍局部为 $P$ 型。
\end{lemma}

\begin{proof}
设 $f : X \to Y$ 与 $g : Y \to Z$ 为局部 $P$ 型的态射。
设 $x \in X$。选取仿射开邻域 $W \subset Z$，使它包含
$g(f(x))$；再选取仿射开邻域 $V \subset g^{-1}(W)$，使它包含
$f(x)$；最后选取仿射开邻域 $U \subset f^{-1}(V)$，使它包含
$x$。由引理
\ref{morphisms-lemma-locally-P-characterize}，环同态
$\mathcal{O}_Z(W) \to \mathcal{O}_Y(V)$ 与
$\mathcal{O}_Y(V) \to \mathcal{O}_X(U)$ 满足 $P$。
故 $\mathcal{O}_Z(W) \to \mathcal{O}_X(U)$ 满足 $P$，
因为假设 $P$ 在复合下稳定。
\end{proof}

\begin{lemma}
\label{morphisms-lemma-base-change-type-P}
设 $P$ 为环同态的一个性质，并假设 $P$ 是局部的且在基变换下稳定。
局部为 $P$ 型的态射之基变换仍局部为 $P$ 型。
\end{lemma}

\begin{proof}
设 $f : X \to S$ 为局部 $P$ 型的态射，$S' \to S$ 为任意态射。
以 $f' : X_{S'} = S' \times_S X \to S'$ 表示 $f$ 的基变换。
对每个 $s' \in S'$，都存在仿射开邻域 $s' \in V' \subset S'$，
它映入某个仿射开集 $V \subset S$。由引理
\ref{morphisms-lemma-locally-P-characterize}，开集 $f^{-1}(V)$ 是仿射开集
$U_i$ 的并，并且环同态
$\mathcal{O}_S(V) \to \mathcal{O}_X(U_i)$ 都满足 $P$。
由《概形》第 \ref{schemes-section-fibre-products} 节的内容，
$f^{-1}(U)_{V'} = V' \times_V f^{-1}(V)$ 是仿射开集
$V' \times_V U_i$ 的并。由于
$\mathcal{O}_{X_{S'}}(V' \times_V U_i) =
\mathcal{O}_{S'}(V') \otimes_{\mathcal{O}_S(V)} \mathcal{O}_X(U_i)$，
环同态
$\mathcal{O}_{S'}(V') \to \mathcal{O}_{X_{S'}}(V' \times_V U_i)$
满足 $P$，因为假设 $P$ 在基变换下稳定。
\end{proof}

\begin{lemma}
\label{morphisms-lemma-properties-local}
环同态 $R \to A$ 的下列性质是局部的：
\begin{enumerate}
\item （局部环上的同构。）
对每个素理想 $\mathfrak q$（它属于 $A$），若它位于
$\mathfrak p \subset R$ 之上，则环同态 $R \to A$ 诱导同构
$R_{\mathfrak p} \to A_{\mathfrak q}$。
\item （开浸入。）
对每个素理想 $\mathfrak q$（它属于 $A$），存在 $f \in R$，满足
$\varphi(f) \not \in \mathfrak q$，且环同态 $\varphi : R \to A$
诱导同构 $R_f \to A_f$。
\item （既约纤维。）
对每个素理想 $\mathfrak p$（它属于 $R$），纤维环
$A \otimes_R \kappa(\mathfrak p)$ 是既约的。
\item （维数至多为 $n$ 的纤维。）
对每个素理想 $\mathfrak p$（它属于 $R$），纤维环
$A \otimes_R \kappa(\mathfrak p)$ 的 Krull 维数至多为 $n$。
\item （靶上局部诺特。）
环同态 $R \to A$ 具有性质：$A$ 是诺特环。
\item 按需在此添加更多性质\footnote{但仅添加那些尚未在别处单独处理的性质。}。
\end{enumerate}
\end{lemma}

\begin{proof}
略。
\end{proof}

\begin{lemma}
\label{morphisms-lemma-properties-base-change}
环同态的下列性质在基变换下稳定：
\begin{enumerate}
\item （局部环上的同构。）
对每个素理想 $\mathfrak q$（它属于 $A$），若它位于
$\mathfrak p \subset R$ 之上，则环同态 $R \to A$ 诱导同构
$R_{\mathfrak p} \to A_{\mathfrak q}$。
\item （开浸入。）
对每个素理想 $\mathfrak q$（它属于 $A$），存在 $f \in R$，满足
$\varphi(f) \not \in \mathfrak q$，且环同态 $\varphi : R \to A$
诱导同构 $R_f \to A_f$。
\item 按需在此添加更多性质\footnote{但仅添加那些尚未在别处单独处理的性质。}。
\end{enumerate}
\end{lemma}

\begin{proof}
略。
\end{proof}

\begin{lemma}
\label{morphisms-lemma-properties-composition}
环同态的下列性质在复合下稳定：
\begin{enumerate}
\item （局部环上的同构。）
对每个素理想 $\mathfrak q$（它属于 $A$），若它位于
$\mathfrak p \subset R$ 之上，则环同态 $R \to A$ 诱导同构
$R_{\mathfrak p} \to A_{\mathfrak q}$。
\item （开浸入。）
对每个素理想 $\mathfrak q$（它属于 $A$），存在 $f \in R$，满足
$\varphi(f) \not \in \mathfrak q$，且环同态 $\varphi : R \to A$
诱导同构 $R_f \to A_f$。
\item （靶上局部诺特。）
环同态 $R \to A$ 具有性质：$A$ 是诺特环。
\item 按需在此添加更多性质\footnote{但仅添加那些尚未在别处单独处理的性质。}。
\end{enumerate}
\end{lemma}

\begin{proof}
略。
\end{proof}

\section{有限型态射}
\label{morphisms-section-finite-type}

\noindent
回忆：若环同态 $R \to A$ 使得 $A$ 同构于
$R[x_1, \ldots, x_n]$ 的一个商（作为 $R$-代数），则称它为有限型；见《代数》定义
\ref{algebra-definition-finite-type}。

\begin{definition}
\label{morphisms-definition-finite-type}
设 $f : X \to S$ 为概形态射。
\begin{enumerate}
\item 称 $f$ {\it 在 $x \in X$ 处为有限型}，如果存在仿射开邻域
$\Spec(A) = U \subset X$，它包含 $x$，以及仿射开集
$\Spec(R) = V \subset S$，满足 $f(U) \subset V$，且诱导环同态
$R \to A$ 为有限型。
\item 若 $f$ 在 $X$ 的每一点处都为有限型，则称它{\it 局部为有限型}。
\item 若 $f$ 局部为有限型且拟紧，则称它为{\it 有限型}。
\end{enumerate}
\end{definition}

\begin{lemma}
\label{morphisms-lemma-locally-finite-type-characterize}
设 $f : X \to S$ 为概形态射。下列条件等价：
\begin{enumerate}
\item 态射 $f$ 局部为有限型。
\item 对所有仿射开集 $U \subset X$、$V \subset S$，若
$f(U) \subset V$，则环同态
$\mathcal{O}_S(V) \to \mathcal{O}_X(U)$ 为有限型。
\item 存在开覆盖 $S = \bigcup_{j \in J} V_j$ 以及开覆盖
$f^{-1}(V_j) = \bigcup_{i \in I_j} U_i$，使每个态射
$U_i \to V_j$（$j\in J, i\in I_j$）都局部为有限型。
\item 存在仿射开覆盖 $S = \bigcup_{j \in J} V_j$ 以及仿射开覆盖
$f^{-1}(V_j) = \bigcup_{i \in I_j} U_i$，使环同态
$\mathcal{O}_S(V_j) \to \mathcal{O}_X(U_i)$ 对所有
$j\in J, i\in I_j$ 都为有限型。
\end{enumerate}
此外，若 $f$ 局部为有限型，则对任意开子概形
$U \subset X$、$V \subset S$，只要 $f(U) \subset V$，限制
$f|_U : U \to V$ 就局部为有限型。
\end{lemma}

\begin{proof}
若证明“$R \to A$ 为有限型”是局部性质，结论即由上述引理
\ref{morphisms-lemma-locally-P} 得到。我们检验定义
\ref{morphisms-definition-property-local} 的条件 (a)、(b)、(c)。
由《代数》引理 \ref{algebra-lemma-base-change-finiteness}，有限型性
在基变换下稳定，故 (a) 成立。由《代数》引理
\ref{algebra-lemma-compose-finite-type}，有限型性在复合下稳定；
而对任意环 $R$，环同态 $R \to R_f$ 显然为有限型。
故 (b) 成立。最后，(c) 由《代数》引理
\ref{algebra-lemma-cover-upstairs} 得到。
\end{proof}

\begin{lemma}
\label{morphisms-lemma-composition-finite-type}
两个局部有限型态射的复合局部为有限型。两个有限型态射的复合也为有限型。
\end{lemma}

\begin{proof}
在引理 \ref{morphisms-lemma-locally-finite-type-characterize} 的证明中，
我们已经看到有限型性是环同态的局部性质。因此，引理的第一个断言
由引理 \ref{morphisms-lemma-composition-type-P} 以及有限型性在环同态复合下稳定
这一事实推出；见《代数》引理
\ref{algebra-lemma-compose-finite-type}。再结合上述结论以及拟紧态射的复合
仍拟紧这一事实（见《概形》引理
\ref{schemes-lemma-composition-quasi-compact}），可知有限型态射的复合
为有限型。
\end{proof}

\begin{lemma}
\label{morphisms-lemma-base-change-finite-type}
局部有限型态射的基变换局部为有限型。有限型态射的基变换也为有限型。
\end{lemma}

\begin{proof}
在引理 \ref{morphisms-lemma-locally-finite-type-characterize} 的证明中，
我们已经看到有限型性是环同态的局部性质。因此，引理的第一个断言
由引理 \ref{morphisms-lemma-base-change-type-P} 以及有限型性在环同态基变换下稳定
这一事实推出；见《代数》引理
\ref{algebra-lemma-base-change-finiteness}。再结合上述结论以及拟紧态射的基变换
仍拟紧这一事实（见《概形》引理
\ref{schemes-lemma-quasi-compact-preserved-base-change}），可知有限型态射的基变换
为有限型态射。
\end{proof}

\begin{lemma}
\label{morphisms-lemma-immersion-locally-finite-type}
闭浸入为有限型。浸入局部为有限型。
\end{lemma}

\begin{proof}
这是因为开浸入是局部同构，而闭浸入显然为有限型。
\end{proof}

\begin{lemma}
\label{morphisms-lemma-finite-type-noetherian}
设 $f : X \to S$ 为态射。若 $S$（局部）诺特且 $f$（局部）为有限型，
则 $X$（局部）诺特。
\end{lemma}

\begin{proof}
有限型诺特环代数仍为诺特环，故结论立即成立；见《代数》引理
\ref{algebra-lemma-Noetherian-permanence}。（也可使用：拟紧态射的靶若拟紧，
则源也拟紧。）
\end{proof}

\begin{lemma}
\label{morphisms-lemma-finite-type-Noetherian-quasi-separated}
设 $f : X \to S$ 局部为有限型，且 $S$ 局部诺特。则 $f$ 拟分离。
\end{lemma}

\begin{proof}
事实上，$X$ 拟分离；见《性质》引理
\ref{properties-lemma-locally-Noetherian-quasi-separated} 及上述引理
\ref{morphisms-lemma-finite-type-noetherian}。再应用《概形》引理
\ref{schemes-lemma-compose-after-separated}，得到 $f$ 拟分离。
\end{proof}

\begin{lemma}
\label{morphisms-lemma-permanence-finite-type}
设 $X \to Y$ 为基概形 $S$ 上的概形态射。若 $X$ 在 $S$ 上
局部为有限型，则 $X \to Y$ 局部为有限型。
\end{lemma}

\begin{proof}
通过引理 \ref{morphisms-lemma-locally-finite-type-characterize}，这转化为下列
代数事实：给定环同态 $A \to B \to C$，若 $A \to C$ 为有限型，
则 $B \to C$ 为有限型。（见《代数》引理
\ref{algebra-lemma-compose-finite-type}。）
\end{proof}

\section{有限型点与 Jacobson 概形}
\label{morphisms-section-points-finite-type}

\noindent
设 $S$ 为概形。有限型点 $s$（属于 $S$）是使态射
$\Spec(\kappa(s)) \to S$ 为有限型的点。研究它的原因是：
在某种意义和某些情形下，有限型点可以代替闭点。例如，它们总是足够多。
此外，一个概形是 Jacobson 概形，当且仅当所有有限型点都是闭点。

\begin{lemma}
\label{morphisms-lemma-point-finite-type}
设 $S$ 为概形，$k$ 为域，$f : \Spec(k) \to S$ 为态射。
下列条件等价：
\begin{enumerate}
\item 态射 $f$ 为有限型。
\item 态射 $f$ 局部为有限型。
\item 存在仿射开集 $U = \Spec(R)$（它属于 $S$），使 $f$ 对应于有限环同态
$R \to k$。
\item 存在仿射开集 $U = \Spec(R)$（它属于 $S$），使 $f$ 的像由
一个闭点 $u$ 组成，该点在 $U$ 中，且域扩张 $k/\kappa(u)$ 有限。
\end{enumerate}
\end{lemma}

\begin{proof}
由于 $\Spec(k)$ 仅含一点，故从它出发的任意态射都拟紧，
所以 (1) 与 (2) 的等价是显然的。

\medskip\noindent
设 $f$ 局部为有限型。任取仿射开集 $\Spec(R) = U \subset S$，
使 $f$ 的像包含在 $U$ 中且环同态 $R \to k$ 为有限型。
设 $\mathfrak p \subset R$ 为其核。于是 $R/\mathfrak p \subset k$
为有限型。由《代数》引理
\ref{algebra-lemma-field-finite-type-over-domain}，存在
$\overline{f} \in R/\mathfrak p$，使 $(R/\mathfrak p)_{\overline{f}}$
为域，且 $(R/\mathfrak p)_{\overline{f}} \to k$ 为有限域扩张。
若 $f \in R$ 是 $\overline{f}$ 的一个提升，则 $k$ 是有限
$R_f$-模。因此 (2) $\Rightarrow$ (3)。

\medskip\noindent
设 $\Spec(R) = U \subset S$ 为仿射开集，并且 $f$ 对应于有限环同态
$R \to k$。由引理 \ref{morphisms-lemma-locally-finite-type-characterize}，
$f$ 局部为有限型。因此 (3) $\Rightarrow$ (2)。

\medskip\noindent
设 $R \to k$ 有限。$R \to k$ 的像是一个域，并且由《代数》引理
\ref{algebra-lemma-integral-under-field}，$k$ 在该域上有限。
因此 $R \to k$ 的核是极大理想。故 (3) $\Rightarrow$ (4)。

\medskip\noindent
(4) $\Rightarrow$ (3) 是立即的。
\end{proof}

\begin{lemma}
\label{morphisms-lemma-artinian-finite-type}
设 $S$ 为概形，$A$ 为剩余域是 $\kappa$ 的阿廷局部环，
$f : \Spec(A) \to S$ 为概形态射。则 $f$ 为有限型，当且仅当复合
$\Spec(\kappa) \to \Spec(A) \to S$ 为有限型。
\end{lemma}

\begin{proof}
由于态射 $\Spec(\kappa) \to \Spec(A)$ 为有限型，若 $f$ 为有限型，
则复合 $\Spec(\kappa) \to S$ 也为有限型，这是清楚的
（见引理 \ref{morphisms-lemma-composition-finite-type}）。反之，注意
$\Spec(A) \to S$ 的像包含在某个仿射开集
$U = \Spec(B)$ 中，该开集属于 $S$，因为 $\Spec(A)$ 只有一个点。最后对
$B \to A$ 应用《代数》引理
\ref{algebra-lemma-essentially-of-finite-type-into-artinian-local}。
\end{proof}

\noindent
回忆：给定一点 $s$，它属于概形 $S$，则有典范态射
$\Spec(\kappa(s)) \to S$；见《概形》第
\ref{schemes-section-points} 节。

\begin{definition}
\label{morphisms-definition-finite-type-point}
设 $S$ 为概形。若一点 $s$（属于 $S$）的典范态射
$\Spec(\kappa(s)) \to S$ 为有限型，则称该点为{\it 有限型点}。
用 $S_{\text{ft-pts}}$ 表示 $S$ 的有限型点集。
\end{definition}

\noindent
有限型点集可以刻画如下。

\begin{lemma}
\label{morphisms-lemma-identify-finite-type-points}
设 $S$ 为概形。我们有
$$
S_{\text{ft-pts}} = \bigcup\nolimits_{U \subset S\text{ 开}} U_0
$$
其中 $U_0$ 是 $U$ 的闭点集。这里可以让 $U$ 遍历 $S$ 的所有开集，
也可以只遍历所有仿射开集。
\end{lemma}

\begin{proof}
这由引理 \ref{morphisms-lemma-point-finite-type} 立即得到。
\end{proof}

\begin{lemma}
\label{morphisms-lemma-finite-type-points-morphism}
设 $f : T \to S$ 为概形态射。若 $f$ 局部为有限型，则
$f(T_{\text{ft-pts}}) \subset S_{\text{ft-pts}}$。
\end{lemma}

\begin{proof}
若 $T$ 是域的谱，这就是引理 \ref{morphisms-lemma-point-finite-type}。
一般情形则由局部有限型态射的复合仍局部为有限型得到
（引理 \ref{morphisms-lemma-composition-finite-type}）。
\end{proof}

\begin{lemma}
\label{morphisms-lemma-finite-type-points-surjective-morphism}
设 $f : T \to S$ 为概形态射。若 $f$ 局部为有限型且满射，则
$f(T_{\text{ft-pts}}) = S_{\text{ft-pts}}$。
\end{lemma}

\begin{proof}
由引理 \ref{morphisms-lemma-finite-type-points-morphism}，有
$f(T_{\text{ft-pts}}) \subset S_{\text{ft-pts}}$。设 $s \in S$
为有限型点。由于 $f$ 满射，概形
$T_s = \Spec(\kappa(s)) \times_S T$ 非空；因此由引理
\ref{morphisms-lemma-identify-finite-type-points}，它有有限型点 $t \in T_s$。
现在 $T_s \to T$ 是有限型态射，因为它是 $s \to S$ 的基变换
（引理 \ref{morphisms-lemma-base-change-finite-type}）。因此由引理
\ref{morphisms-lemma-finite-type-points-morphism}，$t$ 在 $T$ 中的像是有限型点，
并且按构造映到 $s$。
\end{proof}

\begin{lemma}
\label{morphisms-lemma-enough-finite-type-points}
设 $S$ 为概形。对任意局部闭子集 $T \subset S$，有
$$
T \not = \emptyset
\Rightarrow
T \cap S_{\text{ft-pts}} \not = \emptyset.
$$
特别地，对任意闭子集 $T \subset S$，$T \cap S_{\text{ft-pts}}$
在 $T$ 中稠密。
\end{lemma}

\begin{proof}
注意 $T$ 带有概形结构（见《概形》引理
\ref{schemes-lemma-reduced-closed-subscheme}），使 $T \to S$ 是局部闭浸入。
任意局部闭浸入都局部为有限型，见引理
\ref{morphisms-lemma-immersion-locally-finite-type}。因此由引理
\ref{morphisms-lemma-finite-type-points-morphism}，有
$T_{\text{ft-pts}} \subset S_{\text{ft-pts}}$。最后，$T$ 的任意非空
仿射开集都至少有一个闭点；由引理
\ref{morphisms-lemma-identify-finite-type-points}，该点是 $T$ 的有限型点。
\end{proof}

\noindent
因此，《拓扑》第 \ref{topology-section-space-jacobson} 节的大部分内容
在把闭点集替换为有限型点集后仍成立。事实上，若 $S$ 是 Jacobson
概形，则有限型点正好恢复为闭点。

\begin{lemma}
\label{morphisms-lemma-jacobson-finite-type-points}
设 $S$ 为概形。下列条件等价：
\begin{enumerate}
\item 概形 $S$ 是 Jacobson 的；
\item $S_{\text{ft-pts}}$ 是 $S$ 的闭点集；
\item 对所有局部有限型态射 $T \to S$，闭点映到闭点；
\item 对所有局部有限型态射 $T \to S$，闭点 $t \in T$ 映到闭点
$s \in S$，且 $\kappa(s) \subset \kappa(t)$ 有限。
\end{enumerate}
\end{lemma}

\begin{proof}
显然有 (4) $\Rightarrow$ (3) $\Rightarrow$ (2)。引理
\ref{morphisms-lemma-enough-finite-type-points} 表明 (2) 蕴含 (1)。
故只需证明 (1) 蕴含 (4)。设 $T \to S$ 局部为有限型。
取闭点 $t \in T$，并令 $s \in S$ 为其像。选取仿射开邻域
$\Spec(R) = U \subset S$，它包含 $s$；再选取仿射开邻域
$\Spec(A) = V \subset T$，它包含 $t$，并使 $V$ 映入 $U$。诱导环同态
$R \to A$ 为有限型（见引理
\ref{morphisms-lemma-locally-finite-type-characterize}），且由《性质》引理
\ref{properties-lemma-locally-jacobson}，$R$ 是 Jacobson 环。
于是结论由《代数》命题
\ref{algebra-proposition-Jacobson-permanence} 得到。
\end{proof}

\begin{lemma}
\label{morphisms-lemma-Jacobson-universally-Jacobson}
设 $S$ 为 Jacobson 概形。任意在 $S$ 上局部为有限型的概形都是
Jacobson 概形。
\end{lemma}

\begin{proof}
这由《代数》命题 \ref{algebra-proposition-Jacobson-permanence}
直接得到（同时使用《性质》引理
\ref{properties-lemma-locally-jacobson} 和引理
\ref{morphisms-lemma-locally-finite-type-characterize}）。
\end{proof}

\begin{lemma}
\label{morphisms-lemma-ubiquity-Jacobson-schemes}
下列各类概形是 Jacobson 概形：
\begin{enumerate}
\item 任意在域上局部有限型的概形。
\item 任意在 $\mathbf{Z}$ 上局部有限型的概形。
\item 任意在具有无穷多个素理想的 $1$ 维诺特整环上局部有限型的概形。
\item 形如 $\Spec(R) \setminus \{\mathfrak m\}$ 的概形，其中
$(R, \mathfrak m)$ 是诺特局部环；以及任意在它上面局部有限型的概形。
\end{enumerate}
\end{lemma}

\begin{proof}
以下不再明言地使用引理 \ref{morphisms-lemma-Jacobson-universally-Jacobson}。
域的谱显然是 Jacobson 概形。$\mathbf{Z}$ 的谱是 Jacobson 概形，
见《代数》引理 \ref{algebra-lemma-pid-jacobson}。
(3) 见《代数》引理 \ref{algebra-lemma-noetherian-dim-1-Jacobson}。
(4) 见《性质》引理
\ref{properties-lemma-complement-closed-point-Jacobson}。
\end{proof}

\section{普遍链状概形}
\label{morphisms-section-universally-catenary}

\noindent
回顾：若对拓扑空间 $X$ 的每一对不可约闭子集 $T \subset T'$，
都存在一条不可约闭子集的极大链
$$
T = T_0 \subset T_1 \subset \ldots \subset T_e = T'
$$
且每条这样的链都具有相同长度，则称该空间为{\it 链状的}。
见《拓扑学》定义 \ref{topology-definition-catenary}。
回顾：若一个概形的底拓扑空间是链状的，则称该概形为链状概形。
见《性质》定义 \ref{properties-definition-catenary}。

\begin{definition}
\label{morphisms-definition-universally-catenary}
令 $S$ 为一个概形，并假设 $S$ 局部诺特。我们称 $S$
为{\it 普遍链状概形}，如果对每个局部有限型态射
$X \to S$，概形 $X$ 都是链状概形。
\end{definition}

\noindent
这是一个比链状性“更好”的概念，因为存在链状但并非普遍链状的诺特概形。
见《例》节 \ref{examples-section-non-catenary-Noetherian-local}。
许多概形都是普遍链状的，见下文引理 \ref{morphisms-lemma-ubiquity-uc}。

\medskip\noindent
回顾：称环 $A$ 为{\it 链环}，如果对任意一对素理想
$\mathfrak p \subset \mathfrak q$，都存在一条素理想的极大链
$$
\mathfrak p =
\mathfrak p_0 \subset \ldots \subset \mathfrak p_e
= \mathfrak q
$$
且所有这样的链都具有相同长度。
见《代数》定义 \ref{algebra-definition-catenary}。
我们已在《性质》节 \ref{properties-section-catenary}
中讨论了链状概形与链环之间的关系。回顾：称环 $A$
为{\it 普遍链环}，如果 $A$ 是诺特环，并且对每个有限型环同态
$A \to B$，环 $B$ 都是链环。见《代数》定义
\ref{algebra-definition-universally-catenary}。
代数几何中出现的许多重要环都具有这一性质。

\begin{lemma}
\label{morphisms-lemma-universally-catenary-local}
令 $S$ 为局部诺特概形。下列条件等价：
\begin{enumerate}
\item $S$ 是普遍链状概形；
\item 存在 $S$ 的一个开覆盖，其每个成员都是普遍链状概形；
\item 对每个仿射开集 $\Spec(R) = U \subset S$，环 $R$
都是普遍链环；以及
\item 存在仿射开覆盖 $S = \bigcup U_i$，使得每个 $U_i$
都是一个普遍链环的谱。
\end{enumerate}
此外，在这种情形下，任何在 $S$ 上局部有限型的概形也都是普遍链状概形。
\end{lemma}

\begin{proof}
由引理 \ref{morphisms-lemma-immersion-locally-finite-type}，开浸入是局部有限型的。
局部有限型态射的复合仍然局部有限型（引理
\ref{morphisms-lemma-composition-finite-type}）。因此显然，若 $S$ 是普遍链状概形，
则其任意开子概形以及任意在 $S$ 上局部有限型的概形也都是普遍链状概形。
这就证明了引理的最后一个断言，并证明了 (1) 蕴含 (2)。

\medskip\noindent
若 $\Spec(R)$ 是普遍链状概形，则每个概形 $\Spec(A)$，
其中 $A$ 是有限型 $R$-代数，都为链状概形。因此由《代数》引理
\ref{algebra-lemma-catenary}，所有这些环 $A$ 都是链环。
故 $R$ 是普遍链环。结合以上论述，我们得到 (1) 蕴含 (3)，
而 (2) 蕴含 (4)。当然，(3) 显然蕴含 (4)。

\medskip\noindent
为完成证明，我们来证明 (4) 蕴含 (1)。假设 (4)，并令
$X \to S$ 为局部有限型态射。可以找到仿射开覆盖
$X = \bigcup V_j$，使每个 $V_j \to S$ 的像都包含于某个 $U_i$ 中。
由引理 \ref{morphisms-lemma-locally-finite-type-characterize}，诱导环同态
$\mathcal{O}(U_i) \to \mathcal{O}(V_j)$ 是有限型的。
因此 $\mathcal{O}(V_j)$ 是链环。于是由《性质》引理
\ref{properties-lemma-catenary-local}，$X$ 是链状概形。
\end{proof}

\begin{lemma}
\label{morphisms-lemma-universally-catenary-local-rings-universally-catenary}
令 $S$ 为局部诺特概形。下列条件等价：
\begin{enumerate}
\item $S$ 是普遍链状概形；以及
\item 所有局部环 $\mathcal{O}_{S, s}$（这里所指的概形是 $S$）
都是普遍链环。
\end{enumerate}
\end{lemma}

\begin{proof}
假设 $S$ 的所有局部环都是普遍链环。令 $f : X \to S$ 为局部有限型态射。
我们知道，$X$ 是链状概形，当且仅当 $\mathcal{O}_{X, x}$
对所有 $x \in X$ 都是链环。若 $f(x) = s$，则
$\mathcal{O}_{X, x}$ 在 $\mathcal{O}_{S, s}$ 上本质有限型。
因此 $\mathcal{O}_{X, x}$ 是链环，因为我们假设
$\mathcal{O}_{S, s}$ 是普遍链环。

\medskip\noindent
反过来，假设 $S$ 是普遍链状概形。令 $s \in S$。
由引理 \ref{morphisms-lemma-universally-catenary-local}，可以把 $S$
替换为 $s$ 的一个仿射开邻域。设 $S = \Spec(R)$，并设 $s$
对应素理想 $\mathfrak p$。任意有限型 $R_{\mathfrak p}$-代数
$A'$ 都具有形式 $A_{\mathfrak p}$，其中某个有限型 $R$-代数记为 $A$。
由假设（如果愿意，也可再用引理
\ref{morphisms-lemma-universally-catenary-local}），环 $A$ 是链环，
从而 $A'$（$A$ 的一个局部化）也是链环。因此
$R_{\mathfrak p}$ 是普遍链环。
\end{proof}

\begin{lemma}
\label{morphisms-lemma-catenary-check-irreducible}
令 $S$ 为局部诺特概形。则 $S$ 是普遍链状概形，当且仅当
$S$ 的各不可约分支都是普遍链状概形。
\end{lemma}

\begin{proof}
略。仿射情形见《代数》引理
\ref{algebra-lemma-catenary-check-irreducible}。
\end{proof}

\begin{lemma}
\label{morphisms-lemma-ubiquity-uc}
下列各类概形都是普遍链状概形：
\begin{enumerate}
\item 任意在域上局部有限型的概形。
\item 任意在 Cohen--Macaulay 概形上局部有限型的概形。
\item 任意在 $\mathbf{Z}$ 上局部有限型的概形。
\item 任意在 $1$ 维诺特整环上局部有限型的概形。
\item 等等。
\end{enumerate}
\end{lemma}

\begin{proof}
这些结论都源于 Cohen--Macaulay 环是普遍链环这一事实，
见《代数》引理 \ref{algebra-lemma-CM-ring-catenary}。
还要使用引理 \ref{morphisms-lemma-universally-catenary-local} 的最后一个断言。
略去一些细节。
\end{proof}

\section{再论 Nagata 概形}
\label{morphisms-section-nagata}

\noindent
关于 Nagata 概形与泛 Japanese 概形的定义和基本性质，
见《性质》节 \ref{properties-section-nagata}。

\begin{lemma}
\label{morphisms-lemma-finite-type-nagata}
令 $f : X \to S$ 为一个态射。若 $S$ 是 Nagata 概形，且 $f$
局部有限型，则 $X$ 是 Nagata 概形。若 $S$ 是泛 Japanese 概形，
且 $f$ 局部有限型，则 $X$ 是泛 Japanese 概形。
\end{lemma}

\begin{proof}
对于“泛 Japanese”，这来自《代数》引理
\ref{algebra-lemma-universally-japanese}。
对于“Nagata”，这来自《代数》命题
\ref{algebra-proposition-nagata-universally-japanese}。
\end{proof}

\begin{lemma}
\label{morphisms-lemma-ubiquity-nagata}
下列各类概形都是 Nagata 概形：
\begin{enumerate}
\item 任意在域上局部有限型的概形。
\item 任意在诺特完备局部环上局部有限型的概形。
\item 任意在 $\mathbf{Z}$ 上局部有限型的概形。
\item 任意在特征为零的 Dedekind 环上局部有限型的概形。
\item 等等。
\end{enumerate}
\end{lemma}

\begin{proof}
由引理 \ref{morphisms-lemma-finite-type-nagata}，只需证明以上提到的环都是
Nagata 环。为此见《代数》命题
\ref{algebra-proposition-ubiquity-nagata}。
\end{proof}

\section{再论奇异轨迹}
\label{morphisms-section-singular-locus}

\noindent
我们要寻找一个判据，使得对任意在基底上局部有限型的概形，其正则轨迹都是开的。
定义如下。

\begin{definition}
\label{morphisms-definition-J}
令 $X$ 为局部诺特概形。称 $X$ 是{\it J-2 的}，如果对每个局部有限型态射
$Y \to X$，正则轨迹 $\text{Reg}(Y)$ 在 $Y$ 中都是开的。
\end{definition}

\noindent
这是诺特环的相应概念的类比，见《代数续篇》定义
\ref{more-algebra-definition-J}。

\begin{lemma}
\label{morphisms-lemma-J}
令 $X$ 为局部诺特概形。下列条件等价：
\begin{enumerate}
\item $X$ 是 J-2 的；
\item 存在 $X$ 的一个开覆盖，其每个成员都是 J-2 概形；
\item 对每个仿射开集 $\Spec(R) = U \subset X$，环 $R$ 都是 J-2 的；以及
\item 存在仿射开覆盖 $S = \bigcup U_i$，使得每个
$\mathcal{O}(U_i)$ 对所有 $i$ 都是 J-2 的。
\end{enumerate}
此外，在这种情形下，任何在 $X$ 上局部有限型的概形也都是 J-2 的。
\end{lemma}

\begin{proof}
由引理 \ref{morphisms-lemma-immersion-locally-finite-type}，开浸入是局部有限型的。
局部有限型态射的复合仍然局部有限型（引理
\ref{morphisms-lemma-composition-finite-type}）。因此显然，若 $X$ 是 J-2 的，
则任意开子概形以及任意在 $X$ 上局部有限型的概形也都是 J-2 的。
这证明了引理的最后一个断言。

\medskip\noindent
若 $\Spec(R)$ 是 J-2 的，则对每个有限型 $R$-代数 $A$，
概形 $\Spec(A)$ 的正则轨迹都是开的。因此按照定义，$R$ 是 J-2 的
（见《代数续篇》定义 \ref{more-algebra-definition-J}）。
结合以上论述，我们得到 (1) 蕴含 (3)，而 (2) 蕴含 (4)。
当然，(1) $\Rightarrow$ (2) 和 (3) $\Rightarrow$ (4) 都是显然的。

\medskip\noindent
为完成证明，我们来证明 (4) 蕴含 (1)。假设 (4)，并令
$Y \to X$ 为局部有限型态射。可以找到仿射开覆盖
$Y = \bigcup V_j$，使每个 $V_j \to X$ 的像都包含于某个 $U_i$ 中。
由引理 \ref{morphisms-lemma-locally-finite-type-characterize}，诱导环同态
$\mathcal{O}(U_i) \to \mathcal{O}(V_j)$ 是有限型的。因此
$V_j = \Spec(\mathcal{O}(V_j))$ 的正则轨迹是开的。由于
$\text{Reg}(Y) \cap V_j = \text{Reg}(V_j)$，可知
$\text{Reg}(Y)$ 是开的，这正是所求。
\end{proof}

\begin{lemma}
\label{morphisms-lemma-ubiquity-J-2}
下列各类概形都是 J-2 的：
\begin{enumerate}
\item 任意在域上局部有限型的概形。
\item 任意在诺特完备局部环上局部有限型的概形。
\item 任意在 $\mathbf{Z}$ 上局部有限型的概形。
\item 任意在维数为 $1$ 的诺特局部环上局部有限型的概形。
\item 任意在维数为 $1$ 的 Nagata 环上局部有限型的概形。
\item 任意在特征为零的 Dedekind 环上局部有限型的概形。
\item 等等。
\end{enumerate}
\end{lemma}

\begin{proof}
由引理 \ref{morphisms-lemma-J}，只需证明以上提到的环都是 J-2 的。
为此见《代数续篇》命题
\ref{more-algebra-proposition-ubiquity-J-2}。
\end{proof}

\section{优秀概形}
\label{morphisms-section-excellent}

\noindent
回顾：若一个环是 G-环并且是 J-2 的，则称它为拟优秀环。
若一个环是拟优秀环并且是普遍链环，则称它为优秀环。

\begin{definition}
\label{morphisms-definition-excellent}
令 $X$ 为一个概形。
\begin{enumerate}
\item 称 $X$ 是{\it 拟优秀概形}，如果对每个 $x \in X$，都存在仿射开邻域
$x \in U \subset X$，使环 $\mathcal{O}_X(U)$ 是拟优秀环
（见《代数续篇》定义 \ref{more-algebra-definition-excellent}）。
\item 称 $X$ 是{\it 优秀概形}，如果对每个 $x \in X$，都存在仿射开邻域
$x \in U \subset X$，使环 $\mathcal{O}_X(U)$ 是优秀环
（见《代数续篇》定义 \ref{more-algebra-definition-excellent}）。
\end{enumerate}
\end{definition}

\noindent
拟优秀概形是局部诺特的（G-环是诺特环）。

\begin{lemma}
\label{morphisms-lemma-characterize-quasi-excellent}
令 $X$ 为一个概形。下列条件等价：
\begin{enumerate}
\item $X$ 是拟优秀概形；以及
\item $X$ 是 G-概形并且是 J-2 的。
\end{enumerate}
\end{lemma}

\begin{proof}
将拟优秀概形的定义（定义 \ref{morphisms-definition-excellent}）、
拟优秀环的定义（《代数续篇》定义
\ref{more-algebra-definition-excellent}）、G-概形的定义
（《性质》定义 \ref{properties-definition-G-scheme}）以及引理
\ref{morphisms-lemma-J} 结合起来，即得 (1) $\Rightarrow$ (2)。
反过来，若 $X$ 是 G-概形并且是 J-2 的，则对任意 $x \in X$，
都可以找到仿射开邻域 $x \in U \subset X$，使
$\mathcal{O}_X(U)$ 是 G-环。由引理 \ref{morphisms-lemma-J}，环
$\mathcal{O}_X(U)$ 也是 J-2 的，因而是拟优秀环。
\end{proof}

\begin{lemma}
\label{morphisms-lemma-quasi-excellent-nagata}
拟优秀概形是 Nagata 概形。
\end{lemma}

\begin{proof}
见《代数续篇》引理
\ref{more-algebra-lemma-quasi-excellent-nagata}。
\end{proof}

\begin{lemma}
\label{morphisms-lemma-characterize-excellent}
令 $X$ 为一个概形。下列条件等价：
\begin{enumerate}
\item $X$ 是优秀概形；以及
\item $X$ 是拟优秀概形并且是普遍链状概形。
\end{enumerate}
\end{lemma}

\begin{proof}
假设 (1)。由于优秀环是拟优秀环，显然 $X$ 是拟优秀概形。
由于优秀环是普遍链环，由引理
\ref{morphisms-lemma-universally-catenary-local} 还可推出 $X$ 是普遍链状概形。
假设 (2)。令 $x \in X$。由于 $X$ 是拟优秀概形，存在仿射开邻域
$x \in U \subset X$，使 $\mathcal{O}_X(U)$ 是拟优秀环。
由于 $X$ 是普遍范畴的，由引理
\ref{morphisms-lemma-universally-catenary-local}，环 $\mathcal{O}_X(U)$
也是普遍链环，因而是优秀环。
\end{proof}

\begin{lemma}
\label{morphisms-lemma-locally-excellent}
令 $X$ 为一个概形。下列条件等价：
\begin{enumerate}
\item 概形 $X$ 是（拟）优秀概形。
\item 对每个仿射开集 $U \subset X$，环 $\mathcal{O}_X(U)$
都是（拟）优秀环。
\item 存在仿射开覆盖 $X = \bigcup U_i$，使每个
$\mathcal{O}_X(U_i)$ 都是（拟）优秀环。
\item 存在开覆盖 $X = \bigcup X_j$，使每个开子概形
$X_j$ 都是（拟）优秀概形。
\end{enumerate}
此外，若 $X$ 是（拟）优秀概形，则每个开子概形也都是（拟）优秀概形。
\end{lemma}

\begin{proof}
由引理 \ref{morphisms-lemma-characterize-quasi-excellent}，拟优秀等价于既是
G-概形又是 J-2 的。因此在拟优秀情形，该引理由引理
\ref{morphisms-lemma-J} 和《性质》引理
\ref{properties-lemma-locally-G-scheme} 得出。
在优秀情形，使用引理 \ref{morphisms-lemma-characterize-excellent}、
拟优秀情形以及引理 \ref{morphisms-lemma-universally-catenary-local}。
\end{proof}

\begin{lemma}
\label{morphisms-lemma-finite-type-excellent}
令 $f : X \to S$ 为一个态射。若 $S$ 是（拟）优秀概形且 $f$
局部有限型，则 $X$ 是（拟）优秀概形。
\end{lemma}

\begin{proof}
见《代数续篇》引理
\ref{more-algebra-lemma-finite-type-over-excellent}。
\end{proof}

\begin{lemma}
\label{morphisms-lemma-ubiquity-excellent}
下列各类概形都是优秀概形：
\begin{enumerate}
\item 任意在域上局部有限型的概形。
\item 任意在诺特完备局部环上局部有限型的概形。
\item 任意在 $\mathbf{Z}$ 上局部有限型的概形。
\item 任意在特征为零的 Dedekind 环上局部有限型的概形。
\end{enumerate}
\end{lemma}

\begin{proof}
由引理 \ref{morphisms-lemma-locally-excellent} 和
\ref{morphisms-lemma-finite-type-excellent}，只需证明以上提到的环都是优秀环。
为此见《代数续篇》命题
\ref{more-algebra-proposition-ubiquity-excellent}。
\end{proof}

\section{拟有限态射}
\label{morphisms-section-quasi-finite}

\noindent
对拟有限态射进行扎实处理，是后续许多发展的基础。这将导向 Zariski
主定理的各种版本、纤维维数的行为、\'etale 态射的下降等诸多内容。
阅读本节之前，最好先看一下《代数》节
\ref{algebra-section-quasi-finite} 中的代数结果。

\medskip\noindent
回顾：若有限型环同态 $R \to A$ 的纤维中由素理想 $\mathfrak q$
确定的点是孤立点，则称该环同态在 $\mathfrak q$ 处拟有限；见《代数》定义
\ref{algebra-definition-quasi-finite}。

\begin{definition}
\label{morphisms-definition-quasi-finite}
\begin{reference}
\cite[II Definition 6.2.3]{EGA}
\end{reference}
令 $f : X \to S$ 为概形态射。
\begin{enumerate}
\item 称 $f$ 在点 $x \in X$ 处{\it 拟有限}，如果存在仿射邻域
$\Spec(A) = U \subset X$（它是 $x$ 的邻域）以及仿射开集
$\Spec(R) = V \subset S$，使得 $f(U) \subset V$，环同态
$R \to A$ 是有限型的，并且 $R \to A$ 在 $A$ 中对应于 $x$
的素理想处拟有限（见上文）。
\item 称 $f$ {\it 局部拟有限}，如果 $f$ 在每一点 $x$
（它属于 $X$）处都拟有限。
\item 称 $f$ {\it 拟有限}，如果 $f$ 是有限型的，并且每一点 $x$
都是其所在纤维的孤立点。
\end{enumerate}
\end{definition}

\noindent
显然，局部拟有限态射是局部有限型的。下文将看到，对一个局部有限型态射
$f$，它在 $x$ 处拟有限，当且仅当 $x$ 是其所在纤维的孤立点。
此外，一个态射拟有限的点所成的集合是开的；我们将在关于 Zariski
主定理的节 \ref{morphisms-section-Zariski} 中看到这一点。

\begin{lemma}
\label{morphisms-lemma-algebraic-residue-field-extension-closed-point-fibre}
令 $f : X \to S$ 为概形态射。令 $x \in X$ 为一点，并置
$s = f(x)$。若 $\kappa(x)/\kappa(s)$ 是代数域扩张，则：
\begin{enumerate}
\item $x$ 是其所在纤维的闭点；以及
\item 若此外 $s$ 是 $S$ 的闭点，则 $x$ 是 $X$ 的闭点。
\end{enumerate}
\end{lemma}

\begin{proof}
第二个断言由第一个断言和初等拓扑立即得到。根据《概形》引理
\ref{schemes-lemma-fibre-topological}，为了证明第一个断言，
可以把 $X$ 替换为 $X_s$，把 $S$ 替换为 $\Spec(\kappa(s))$。
因此可以假设 $S = \Spec(k)$ 是一个域的谱。在这种情形下，令
$\Spec(A) = U \subset X$ 为任意包含 $x$ 的仿射开集。点 $x$
对应于素理想 $\mathfrak q \subset A$，且
$\kappa(\mathfrak q)/k$ 是代数域扩张。由《代数》引理
\ref{algebra-lemma-finite-residue-extension-closed}，可知
$\mathfrak q$ 是极大理想，即 $x \in U$ 是闭点。由于仿射开集构成
$X$ 的拓扑基，可知 $\{x\}$ 是闭集。
\end{proof}

\noindent
下一个引理是 Hilbert 零点定理的一个版本。

\begin{lemma}
\label{morphisms-lemma-closed-point-fibre-locally-finite-type}
令 $f : X \to S$ 为概形态射。令 $x \in X$ 为一点，并置
$s = f(x)$。假设 $f$ 局部有限型。则 $x$ 是其所在纤维的闭点，
当且仅当 $\kappa(x)/\kappa(s)$ 是有限域扩张。
\end{lemma}

\begin{proof}
若该扩张有限，则由上面的引理
\ref{morphisms-lemma-algebraic-residue-field-extension-closed-point-fibre}，
$x$ 是其所在纤维的闭点。反过来，假设 $x$ 是其所在纤维的闭点。
选取仿射开集 $\Spec(A) = U \subset X$ 和
$\Spec(R) = V \subset S$，使得 $f(U) \subset V$。
由引理 \ref{morphisms-lemma-locally-finite-type-characterize}，环同态
$R \to A$ 是有限型的。令 $\mathfrak q \subset A$ 和
$\mathfrak p \subset R$ 分别为对应于 $x$ 和 $s$ 的素理想。
考虑纤维环 $\overline{A} = A \otimes_R \kappa(\mathfrak p)$。
令 $\overline{\mathfrak q}$ 为 $\overline{A}$ 中对应于
$\mathfrak q$ 的素理想。$x$ 是其所在纤维闭点这一假设蕴含
$\overline{\mathfrak q}$ 是 $\overline{A}$ 的极大理想。由于
$\overline{A}$ 是域 $\kappa(\mathfrak p)$ 上的有限型代数，
由 Hilbert 零点定理（见《代数》定理
\ref{algebra-theorem-nullstellensatz}），可知
$\kappa(\overline{\mathfrak q})$ 是 $\kappa(\mathfrak p)$
的有限扩张。由于 $\kappa(s) = \kappa(\mathfrak p)$ 且
$\kappa(x) = \kappa(\mathfrak q) = \kappa(\overline{\mathfrak q})$，
结论成立。
\end{proof}

\begin{lemma}
\label{morphisms-lemma-base-change-closed-point-fibre-locally-finite-type}
令 $f : X \to S$ 为局部有限型的概形态射。令 $g : S' \to S$
为任意态射，并记其基变换为 $f' : X' \to S'$。若
$x' \in X'$ 映到的点 $x \in X$ 在 $X_{f(x)}$ 中是闭点，
则 $x'$ 在 $X'_{f'(x')}$ 中是闭点。
\end{lemma}

\begin{proof}
剩余域 $\kappa(x')$ 是
$\kappa(f'(x')) \otimes_{\kappa(f(x))} \kappa(x)$ 的一个商，
见《概形》引理 \ref{schemes-lemma-points-fibre-product}。
因此它是 $\kappa(f'(x'))$ 的有限扩张，因为由引理
\ref{morphisms-lemma-closed-point-fibre-locally-finite-type}，
$\kappa(x)$ 是 $\kappa(f(x))$ 的有限扩张。
再应用该引理一次，便可知 $x'$ 在其所在纤维中是闭点。
\end{proof}

\begin{lemma}
\label{morphisms-lemma-residue-field-quasi-finite}
令 $f : X \to S$ 为概形态射。令 $x \in X$ 为一点，并置
$s = f(x)$。若 $f$ 在 $x$ 处拟有限，则剩余域扩张
$\kappa(x)/\kappa(s)$ 是有限的。
\end{lemma}

\begin{proof}
这直接来自《代数》定义 \ref{algebra-definition-quasi-finite}。
\end{proof}

\begin{lemma}
\label{morphisms-lemma-quasi-finite-at-point-characterize}
令 $f : X \to S$ 为概形态射。令 $x \in X$ 为一点，并置
$s = f(x)$。令 $X_s$ 为 $f$ 在 $s$ 处的纤维。假设 $f$
局部有限型。下列条件等价：
\begin{enumerate}
\item 态射 $f$ 在 $x$ 处拟有限。
\item 点 $x$ 在 $X_s$ 中是孤立点。
\item 点 $x$ 在 $X_s$ 中是闭点，并且不存在点
$x' \in X_s$ 满足 $x' \not = x$ 且特化到 $x$。
\item 对任意一对仿射开集
$\Spec(A) = U \subset X$、$\Spec(R) = V \subset S$，若
$f(U) \subset V$，且 $x \in U$ 对应于 $\mathfrak q \subset A$，
则环同态 $R \to A$ 在 $\mathfrak q$ 处拟有限。
\end{enumerate}
\end{lemma}

\begin{proof}
假设 $f$ 在 $x$ 处拟有限。由假设，存在开集
$U \subset X$、$V \subset S$，使得 $f(U) \subset V$、$x \in U$，
并且 $x$ 是 $U_s$ 的孤立点。因此 $\{x\} \subset U_s$ 是开子集。
由于 $U_s = U \cap X_s \subset X_s$ 也是开的，可知
$\{x\} \subset X_s$ 也是开子集。所以 $x$ 是 $X_s$ 的孤立点。

\medskip\noindent
注意，由引理 \ref{morphisms-lemma-ubiquity-Jacobson-schemes}
（以及引理 \ref{morphisms-lemma-base-change-finite-type}），$X_s$
是 Jacobson 概形。若 $x$ 在 $X_s$ 中是孤立点，即
$\{x\} \subset X_s$ 是开的，则 $\{x\}$ 包含一个闭点
（由 Jacobson 性质），因此 $x$ 在 $X_s$ 中是闭点。
显然，不存在点 $x' \in X_s$，它不同于 $x$ 且特化到 $x$。

\medskip\noindent
假设 $x$ 在 $X_s$ 中是闭点，并且不存在点 $x' \in X_s$，
它不同于 $x$ 且特化到 $x$。考虑一对仿射开集
$\Spec(A) = U \subset X$、$\Spec(R) = V \subset S$，满足
$f(U) \subset V$ 且 $x \in U$。令 $\mathfrak q \subset A$
对应于 $x$，并令 $\mathfrak p \subset R$ 对应于 $s$。
由引理 \ref{morphisms-lemma-locally-finite-type-characterize}，环同态
$R \to A$ 是有限型的。考虑纤维环
$\overline{A} = A \otimes_R \kappa(\mathfrak p)$。
令 $\overline{\mathfrak q}$ 为 $\overline{A}$ 中对应于
$\mathfrak q$ 的素理想。由于 $\Spec(\overline{A})$ 是纤维
$X_s$ 的开子概形，可知 $\overline{q}$ 是 $\overline{A}$ 的极大理想，
并且 $\Spec(\overline{A})$ 中没有点特化到
$\overline{\mathfrak q}$。这蕴含
$\dim(\overline{A}_{\overline{q}}) = 0$。因此由《代数》定义
\ref{algebra-definition-quasi-finite}，可知 $R \to A$ 在
$\mathfrak q$ 处拟有限；按照定义，这就是说 $X \to S$ 在 $x$
处拟有限。

\medskip\noindent
至此，我们已经证明条件 (1)--(3) 等价。显然 (4) 蕴含 (1)。
(2) 也显然蕴含 (4)，因为若 $x$ 是 $X_s$ 的孤立点，则它在
$U_s$ 中也为孤立点，其中 $U$ 是任意包含它的开集。
\end{proof}

\begin{lemma}
\label{morphisms-lemma-finite-fibre}
令 $f : X \to S$ 为概形态射，并令 $s \in S$。假设：
\begin{enumerate}
\item $f$ 局部有限型；以及
\item $f^{-1}(\{s\})$ 是有限集。
\end{enumerate}
则 $X_s$ 是有限离散拓扑空间，并且 $f$ 在 $X$ 中位于 $s$
上方的每一点处都拟有限。
\end{lemma}

\begin{proof}
设 $T$ 为一个概形，它 (a) 在域 $k$ 上局部有限型，并且 (b) 只有有限多个点。
则引理 \ref{morphisms-lemma-ubiquity-Jacobson-schemes} 表明 $T$ 是 Jacobson 概形。
有限 Jacobson 空间是离散的，见《拓扑学》引理
\ref{topology-lemma-finite-jacobson}。把这一论述应用于纤维 $X_s$；
它在 $\Spec(\kappa(s))$ 上局部有限型，由此得到第一个断言。
最后，应用引理 \ref{morphisms-lemma-quasi-finite-at-point-characterize}
即得第二个断言。
\end{proof}

\begin{lemma}
\label{morphisms-lemma-locally-quasi-finite-fibres}
\begin{slogan}
具有离散纤维的有限型态射是拟有限的。
\end{slogan}
令 $f : X \to S$ 为概形态射。假设 $f$ 局部有限型。
则下列条件等价：
\begin{enumerate}
\item $f$ 局部拟有限；
\item 对每个 $s \in S$，纤维 $X_s$ 都是离散拓扑空间；以及
\item 对每个态射 $\Spec(k) \to S$（其中 $k$ 是域），基变换
$X_k$ 的底拓扑空间都是离散的。
\end{enumerate}
\end{lemma}

\begin{proof}
(3) 蕴含 (2) 是立即的。引理
\ref{morphisms-lemma-quasi-finite-at-point-characterize} 表明 (2) 与 (1) 等价。
假设 (2)，并令 $\Spec(k) \to S$ 如 (3) 所述。以 $s \in S$
记 $\Spec(k) \to S$ 的像。则 $X_k$ 是 $X_s$ 经
$\Spec(k) \to \Spec(\kappa(s))$ 的基变换。因此由引理
\ref{morphisms-lemma-base-change-closed-point-fibre-locally-finite-type}，
$X_k$ 的每一点都是闭点。由于 $X_k \to \Spec(k)$ 局部有限型
（由引理 \ref{morphisms-lemma-base-change-finite-type}），可以应用引理
\ref{morphisms-lemma-quasi-finite-at-point-characterize}，得到 $X_k$ 的每一点
都是孤立点，即 $X_k$ 的底拓扑空间是离散的。
\end{proof}

\begin{lemma}
\label{morphisms-lemma-quasi-finite-locally-quasi-compact}
令 $f : X \to S$ 为概形态射。则 $f$ 拟有限，当且仅当 $f$
局部拟有限且拟紧。
\end{lemma}

\begin{proof}
假设 $f$ 拟有限。由定义 \ref{morphisms-definition-finite-type}，它是拟紧的。
令 $x \in X$。由引理 \ref{morphisms-lemma-quasi-finite-at-point-characterize}，
可知 $f$ 在 $x$ 处拟有限。因此 $f$ 既拟紧又局部拟有限。

\medskip\noindent
假设 $f$ 拟紧且局部拟有限。则 $f$ 是有限型的。令 $x \in X$
为一点。由引理 \ref{morphisms-lemma-quasi-finite-at-point-characterize}，
可知 $x$ 是其所在纤维的孤立点。引理得证。
\end{proof}

\begin{lemma}
\label{morphisms-lemma-quasi-finite}
令 $f : X \to S$ 为概形态射。下列条件等价：
\begin{enumerate}
\item $f$ 拟有限；以及
\item $f$ 局部有限型、拟紧，并且具有有限纤维。
\end{enumerate}
\end{lemma}

\begin{proof}
假设 $f$ 拟有限。特别地，$f$ 局部有限型且拟紧（因为它是有限型的）。
令 $s \in S$。由于每个 $x \in X_s$ 都是 $X_s$ 的孤立点，可见
$X_s = \bigcup_{x \in X_s} \{x\}$ 是一个开覆盖。由于 $f$ 拟紧，
纤维 $X_s$ 是拟紧的。因此 $X_s$ 是有限的。

\medskip\noindent
反过来，假设 $f$ 局部有限型、拟紧且具有有限纤维。
由引理 \ref{morphisms-lemma-finite-fibre}，它局部拟有限。因此由引理
\ref{morphisms-lemma-quasi-finite-locally-quasi-compact}，它拟有限。
\end{proof}

\noindent
回顾：若环同态 $R \to A$ 是有限型的，并且在 $A$ 的{\it 所有}素理想处
都拟有限，则称该环同态拟有限；见《代数》定义
\ref{algebra-definition-quasi-finite}。

\begin{lemma}
\label{morphisms-lemma-locally-quasi-finite-characterize}
令 $f : X \to S$ 为概形态射。下列条件等价：
\begin{enumerate}
\item 态射 $f$ 局部拟有限。
\item 对每一对仿射开集 $U \subset X$、$V \subset S$，若
$f(U) \subset V$，则环同态
$\mathcal{O}_S(V) \to \mathcal{O}_X(U)$ 都拟有限。
\item 存在开覆盖 $S = \bigcup_{j \in J} V_j$，以及开覆盖
$f^{-1}(V_j) = \bigcup_{i \in I_j} U_i$，使每个态射
$U_i \to V_j$（$j\in J, i\in I_j$）都局部拟有限。
\item 存在仿射开覆盖 $S = \bigcup_{j \in J} V_j$，以及仿射开覆盖
$f^{-1}(V_j) = \bigcup_{i \in I_j} U_i$，使环同态
$\mathcal{O}_S(V_j) \to \mathcal{O}_X(U_i)$ 对所有
$j\in J, i\in I_j$ 都拟有限。
\end{enumerate}
此外，若 $f$ 局部拟有限，则对任何开子概形
$U \subset X$、$V \subset S$，若 $f(U) \subset V$，则限制
$f|_U : U \to V$ 都局部拟有限。
\end{lemma}

\begin{proof}
对环同态 $R \to A$，定义 $P(R \to A)$ 表示“$R \to A$ 拟有限”
（见引理之前的评注）。我们断言 $P$ 是环同态的局部性质。
来验证定义 \ref{morphisms-definition-property-local} 的条件 (a)、(b) 和 (c)。
在引理 \ref{morphisms-lemma-locally-finite-type-characterize} 的证明中，
我们已经看到性质“有限型”满足 (a)、(b) 和 (c)。注意，对于有限型环同态
$R \to A$，$R \to A$ 在 $\mathfrak q$ 处拟有限这一性质，只依赖于局部环
$A_{\mathfrak q}$ 作为 $R_{\mathfrak p}$-代数的结构，其中
$\mathfrak p = R \cap \mathfrak q$（通常的记号滥用）。由这些评注，
定义 \ref{morphisms-definition-property-local} 的 (a)、(b) 和 (c) 立即成立。
例如，设 $R \to A$ 为环同态，并且所有环同态
$R \to A_{a_i}$ 都拟有限，其中 $a_1, \ldots, a_n \in A$
生成单位理想。
可知 $R \to A$ 是有限型的。此外，对任意素理想
$\mathfrak q \subset A$，局部环 $A_{\mathfrak q}$ 作为 $R$-代数，
同构于局部环 $(A_{a_i})_{\mathfrak q_i}$，其中 $i$ 是某个指标，
且 $\mathfrak q_i \subset A_{a_i}$ 是某个素理想。因此
$R \to A$ 在 $\mathfrak q$ 处拟有限。

\medskip\noindent
可知，引理 \ref{morphisms-lemma-locally-P} 可应用于上一段中的 $P$。
因此只需证明：$f$ 局部拟有限，等价于 $f$ 局部为
$P$ 型。由于
$P(R \to A)$ 表示“$R \to A$ 拟有限”，也就是 $R \to A$
在 $A$ 的每个素理想处拟有限，这来自引理
\ref{morphisms-lemma-quasi-finite-at-point-characterize}。
\end{proof}

\begin{lemma}
\label{morphisms-lemma-composition-quasi-finite}
两个局部拟有限态射的复合仍然局部拟有限。对拟有限态射也有相同结论。
\end{lemma}

\begin{proof}
在引理 \ref{morphisms-lemma-locally-quasi-finite-characterize} 的证明中，
我们看到 $P = $“拟有限”是环同态的局部性质，并且概形态射局部拟有限，
当且仅当它局部为 $P$ 型（如定义 \ref{morphisms-definition-locally-P}）。
因此，引理的第一个断言来自引理 \ref{morphisms-lemma-composition-type-P}，
以及拟有限是一个在复合下稳定的环同态性质这一事实；见《代数》引理
\ref{algebra-lemma-quasi-finite-composition}。
结合上述结论、引理 \ref{morphisms-lemma-quasi-finite-locally-quasi-compact}，
以及拟紧态射的复合拟紧这一事实（见《概形》引理
\ref{schemes-lemma-composition-quasi-compact}），可知拟有限态射的复合
仍然拟有限。
\end{proof}

\noindent
稍后将看到（引理 \ref{morphisms-lemma-quasi-finite-points-open}），
下一个引理中的集合 $U$ 是开的。

\begin{lemma}
\label{morphisms-lemma-base-change-quasi-finite}
\begin{slogan}
（局部）拟有限态射在基变换下稳定。
\end{slogan}
令 $f : X \to S$ 为概形态射。令 $g : S' \to S$ 为概形态射。
以 $f' : X' \to S'$ 记 $f$ 经 $g$ 的基变换，
并以 $g' : X' \to X$ 记投影。假设 $X$ 在 $S$ 上局部有限型。
\begin{enumerate}
\item 令 $U \subset X$（相应地，$U' \subset X'$）为 $f$
（相应地，$f'$）拟有限的点所成的集合。则
$U' = U \times_S S' = (g')^{-1}(U)$。
\item 局部拟有限态射的基变换局部拟有限。
\item 拟有限态射的基变换拟有限。
\end{enumerate}
\end{lemma}

\begin{proof}
第一和第二个断言来自相应的代数结果，见《代数》引理
\ref{algebra-lemma-quasi-finite-base-change}（并结合引理
\ref{morphisms-lemma-base-change-finite-type} 所给出的 $f'$ 也局部有限型这一事实）。
结合上述结论、引理 \ref{morphisms-lemma-quasi-finite-locally-quasi-compact}，
以及拟紧态射的基变换拟紧这一事实（见《概形》引理
\ref{schemes-lemma-quasi-compact-preserved-base-change}），可知拟有限态射的
基变换拟有限。
\end{proof}

\begin{lemma}
\label{morphisms-lemma-quasi-finite-at-a-finite-number-of-points}
令 $f : X \to S$ 为有限型概形态射，并令 $s \in S$。
$X$ 中位于 $s$ 上方且 $f$ 在其处拟有限的点至多只有有限多个。
\end{lemma}

\begin{proof}
纤维 $X_s$ 是域上的有限型概形，因而是诺特概形
（引理 \ref{morphisms-lemma-finite-type-noetherian}）。所以 $X_s$
上的拓扑是诺特的（《性质》引理
\ref{properties-lemma-Noetherian-topology}），由初等拓扑可知它至多有
有限多个孤立点。因此本引理由引理
\ref{morphisms-lemma-quasi-finite-at-point-characterize} 得出。
\end{proof}

\begin{lemma}
\label{morphisms-lemma-monomorphism-loc-finite-type-loc-quasi-finite}
令 $f : X \to Y$ 为概形态射。若 $f$ 局部有限型且为单态射，
则 $f$ 是分离且局部拟有限的。
\end{lemma}

\begin{proof}
由《概形》引理 \ref{schemes-lemma-monomorphism-separated}，
单态射是分离的。单态射是单射，因此例如由引理
\ref{morphisms-lemma-quasi-finite-at-point-characterize}，可知 $f$
在每个 $x \in X$ 处拟有限。
\end{proof}

\begin{lemma}
\label{morphisms-lemma-immersion-locally-quasi-finite}
任意浸入都局部拟有限。
\end{lemma}

\begin{proof}
这是因为开浸入是局部同构，而闭浸入显然拟有限。
\end{proof}

\begin{lemma}
\label{morphisms-lemma-permanence-quasi-finite}
令 $X \to Y$ 为基概形 $S$ 上的概形态射。令 $x \in X$。
若 $X \to S$ 在 $x$ 处拟有限，则 $X \to Y$ 在 $x$ 处拟有限。
若 $X$ 在 $S$ 上局部拟有限，则 $X \to Y$ 局部拟有限。
\end{lemma}

\begin{proof}
通过引理 \ref{morphisms-lemma-locally-quasi-finite-characterize}，这转化为如下代数事实：
给定环同态 $A \to B \to C$，若 $A \to C$ 拟有限，则
$B \to C$ 拟有限。这来自《代数》引理
\ref{algebra-lemma-four-rings}，其中取 $R = A$、$S = S' = C$
和 $R' = B$。
\end{proof}

\begin{lemma}
\label{morphisms-lemma-quasi-finite-local-source}
令 $f : X \to Y$ 和 $g : Y \to S$ 为概形态射。若 $f$ 满射、
$g \circ f$ 局部拟有限且 $g$ 局部有限型，则 $g : Y \to S$
局部拟有限。
\end{lemma}

\begin{proof}
令 $x \in X$ 的像依次为 $y \in Y$ 和 $s \in S$。
由于 $g \circ f$ 局部拟有限，由引理
\ref{morphisms-lemma-residue-field-quasi-finite}，扩张
$\kappa(x)/\kappa(s)$ 是有限的。因此
$\kappa(y)/\kappa(s)$ 是有限的。由引理
\ref{morphisms-lemma-algebraic-residue-field-extension-closed-point-fibre}，
$y$ 是 $Y_s$ 的闭点。由于 $f$ 满射，可知 $Y$ 的每一点在其
$S$ 上的纤维中都是闭点。因此由引理
\ref{morphisms-lemma-quasi-finite-at-point-characterize}，可知 $g$
在每一点处拟有限。
\end{proof}

\section{有限表示态射}
\label{morphisms-section-finite-presentation}

\noindent
回顾：若环同态 $R \to A$ 使得 $A$ 同构于
$R[x_1, \ldots, x_n]/(f_1, \ldots, f_m)$，且这是一个 $R$-代数同构，
其中 $n, m$
为某些整数且 $f_j$ 为某些多项式，则称该环同态为有限表示的；
见《代数》定义 \ref{algebra-definition-finite-type}。

\begin{definition}
\label{morphisms-definition-finite-presentation}
令 $f : X \to S$ 为概形态射。
\begin{enumerate}
\item 称 $f$ 在 $x \in X$ 处{\it 有限表示}，如果存在仿射开邻域
$\Spec(A) = U \subset X$（它是 $x$ 的邻域）以及仿射开集
$\Spec(R) = V \subset S$，满足 $f(U) \subset V$，使得诱导环同态
$R \to A$ 是有限表示的。
\item 称 $f$ {\it 局部有限表示}，如果它在 $X$ 的每一点处都有限表示。
\item 称 $f$ {\it 有限表示}，如果它局部有限表示、拟紧且拟分离。
\end{enumerate}
\end{definition}

\noindent
注意，有限表示态射{\bf 不}只是局部有限表示的拟紧态射。
稍后我们将把局部有限表示态射刻画为如下态射：对于仿射概形的
任意有向系 $T_i$（它们都在 $S$ 上），都有
$$
\colim \Mor_S(T_i, X) = \Mor_S(\lim T_i, X)
$$
成立。见《极限》命题
\ref{limits-proposition-characterize-locally-finite-presentation}。
在《极限》节 \ref{limits-section-descending-relative} 中，我们证明：
若 $S = \lim_i S_i$ 是仿射概形的极限，则任意概形 $X$，若它在
$S$ 上有限表示，都能下降为概形 $X_i$；它在 $S_i$ 上，其中 $i$
为某个指标。

\begin{lemma}
\label{morphisms-lemma-locally-finite-presentation-characterize}
令 $f : X \to S$ 为概形态射。下列条件等价：
\begin{enumerate}
\item 态射 $f$ 局部有限表示。
\item 对每一对仿射开集 $U \subset X$、$V \subset S$，若
$f(U) \subset V$，则环同态
$\mathcal{O}_S(V) \to \mathcal{O}_X(U)$ 是有限表示的。
\item 存在开覆盖 $S = \bigcup_{j \in J} V_j$，以及开覆盖
$f^{-1}(V_j) = \bigcup_{i \in I_j} U_i$，使每个态射
$U_i \to V_j$（$j\in J, i\in I_j$）都局部有限表示。
\item 存在仿射开覆盖 $S = \bigcup_{j \in J} V_j$，以及仿射开覆盖
$f^{-1}(V_j) = \bigcup_{i \in I_j} U_i$，使环同态
$\mathcal{O}_S(V_j) \to \mathcal{O}_X(U_i)$ 对所有
$j\in J, i\in I_j$ 都是有限表示的。
\end{enumerate}
此外，若 $f$ 局部有限表示，则对任何开子概形 $U \subset X$、
$V \subset S$，若 $f(U) \subset V$，则限制 $f|_U : U \to V$
局部有限表示。
\end{lemma}

\begin{proof}
只要证明性质“$R \to A$ 是有限表示的”为局部性质，结论便来自引理
\ref{morphisms-lemma-locally-P-characterize}。来验证定义
\ref{morphisms-definition-property-local} 的条件 (a)、(b) 和 (c)。
由《代数》引理 \ref{algebra-lemma-base-change-finiteness}，
有限表示性在基变换下稳定，故 (a) 成立。由《代数》引理
\ref{algebra-lemma-compose-finite-type}，有限表示性在复合下稳定；
并且对任意环 $R$，环同态 $R \to R_f$ 显然是有限表示的。
故 (b) 成立。最后，由《代数》引理
\ref{algebra-lemma-cover-upstairs}，性质 (c) 成立。
\end{proof}

\begin{lemma}
\label{morphisms-lemma-composition-finite-presentation}
两个局部有限表示态射的复合仍然局部有限表示。
有限表示态射也有同样的结论。
\end{lemma}

\begin{proof}
在引理 \ref{morphisms-lemma-locally-finite-presentation-characterize} 的证明中，
我们看到有限表示性是环同态的局部性质。因此引理的第一个断言来自引理
\ref{morphisms-lemma-composition-type-P}，以及有限表示性是一个在复合下稳定的
环同态性质这一事实；见《代数》引理
\ref{algebra-lemma-compose-finite-type}。
结合上述结论以及拟紧拟分离态射的复合仍然拟紧拟分离这一事实
（见《概形》引理 \ref{schemes-lemma-composition-quasi-compact} 和
\ref{schemes-lemma-separated-permanence}），可知有限表示态射的复合
仍然有限表示。
\end{proof}

\begin{lemma}
\label{morphisms-lemma-base-change-finite-presentation}
局部有限表示态射的基变换仍然局部有限表示。
有限表示态射也有同样的结论。
\end{lemma}

\begin{proof}
在引理 \ref{morphisms-lemma-locally-finite-presentation-characterize} 的证明中，
我们看到有限表示性是环同态的局部性质。因此引理的第一个断言来自引理
\ref{morphisms-lemma-composition-type-P}，以及有限表示性是一个在基变换下稳定的
环同态性质这一事实；见《代数》引理
\ref{algebra-lemma-base-change-finiteness}。
结合上述结论以及拟紧拟分离态射的基变换仍然拟紧拟分离这一事实
（见《概形》引理
\ref{schemes-lemma-quasi-compact-preserved-base-change} 和
\ref{schemes-lemma-separated-permanence}），可知有限表示态射的基变换
仍然是有限表示态射。
\end{proof}

\begin{lemma}
\label{morphisms-lemma-open-immersion-locally-finite-presentation}
任意开浸入都局部有限表示。
\end{lemma}

\begin{proof}
这是因为开浸入是局部同构。
\end{proof}

\begin{lemma}
\label{morphisms-lemma-quasi-compact-open-immersion-finite-presentation}
开浸入有限表示，当且仅当它拟紧。
\end{lemma}

\begin{proof}
我们已经看到（引理
\ref{morphisms-lemma-open-immersion-locally-finite-presentation}），开浸入局部有限表示。
我们也已经看到（《概形》引理
\ref{schemes-lemma-immersions-monomorphisms}），浸入是分离的，因而拟分离。
结合这一点和定义 \ref{morphisms-definition-finite-presentation}，即得结论。
\end{proof}

\begin{lemma}
\label{morphisms-lemma-closed-immersion-finite-presentation}
\begin{slogan}
有限表示的闭浸入对应于有限型拟凝聚理想层。
\end{slogan}
闭浸入 $i : Z \to X$ 有限表示，当且仅当相伴的拟凝聚理想层
$\mathcal{I} = \Ker(\mathcal{O}_X \to i_*\mathcal{O}_Z)$
作为 $\mathcal{O}_X$-模是有限型的。
\end{lemma}

\begin{proof}
在任意仿射开集 $\Spec(R) \subset X$ 上，我们有
$i^{-1}(\Spec(R)) = \Spec(R/I)$ 和
$\mathcal{I} = \widetilde{I}$。此外，$\mathcal{I}$ 是有限型的，
当且仅当对每个这样的仿射开集，$I$ 都是有限 $R$-模
（见《性质》引理 \ref{properties-lemma-finite-type-module}）。
而 $R/I$ 在 $R$ 上有限表示，当且仅当 $I$ 是有限 $R$-模。
因此结论成立。
\end{proof}

\begin{lemma}
\label{morphisms-lemma-finite-presentation-finite-type}
局部有限表示态射局部有限型。有限表示态射有限型。
\end{lemma}

\begin{proof}
略。
\end{proof}

\begin{lemma}
\label{morphisms-lemma-noetherian-finite-type-finite-presentation}
\begin{slogan}
在局部诺特基底上，有限型就是有限表示。
\end{slogan}
令 $f : X \to S$ 为一个态射。
\begin{enumerate}
\item 若 $S$ 局部诺特且 $f$ 局部有限型，则 $f$ 局部有限表示。
\item 若 $S$ 局部诺特且 $f$ 有限型，则 $f$ 有限表示。
\end{enumerate}
\end{lemma}

\begin{proof}
第一个断言来自以下事实：诺特环上的有限型环是有限表示的；见《代数》引理
\ref{algebra-lemma-Noetherian-finite-type-is-finite-presentation}。
假设 $f$ 有限型且 $S$ 局部诺特。则由 (1)，$f$ 拟紧且局部有限表示。
所以只需证明 $f$ 拟分离。这来自引理
\ref{morphisms-lemma-finite-type-Noetherian-quasi-separated}
（以及引理 \ref{morphisms-lemma-finite-presentation-finite-type}）。
\end{proof}

\begin{lemma}
\label{morphisms-lemma-finite-presentation-quasi-compact-quasi-separated}
令 $S$ 为拟紧拟分离概形。若 $X$ 在 $S$ 上有限表示，则 $X$
拟紧且拟分离。
\end{lemma}

\begin{proof}
略。
\end{proof}

\begin{lemma}
\label{morphisms-lemma-finite-presentation-permanence}
令 $f : X \to Y$ 为 $S$ 上的概形态射。
\begin{enumerate}
\item 若 $X$ 在 $S$ 上局部有限表示，并且 $Y$ 在 $S$ 上局部有限型，
则 $f$ 局部有限表示。
\item 若 $X$ 在 $S$ 上有限表示，并且 $Y$ 拟分离且在 $S$ 上局部有限型，
则 $f$ 有限表示。
\end{enumerate}
\end{lemma}

\begin{proof}
证明 (1)。通过引理 \ref{morphisms-lemma-locally-finite-presentation-characterize}，
这转化为如下代数事实：给定环同态 $A \to B \to C$，若
$A \to C$ 是有限表示的，且 $A \to B$ 是有限型的，则
$B \to C$ 是有限表示的。见《代数》引理
\ref{algebra-lemma-compose-finite-type}。

\medskip\noindent
(2) 来自 (1) 以及《概形》引理
\ref{schemes-lemma-compose-after-separated} 和
\ref{schemes-lemma-quasi-compact-permanence}。
\end{proof}

\begin{lemma}
\label{morphisms-lemma-diagonal-morphism-finite-type}
令 $f : X \to Y$ 为概形态射，其对角态射为
$\Delta : X \to X \times_Y X$。若 $f$ 局部有限型，则
$\Delta$ 局部有限表示。若 $f$ 拟分离且局部有限型，则
$\Delta$ 有限表示。
\end{lemma}

\begin{proof}
注意，$\Delta$ 是 $X$ 上的概形态射（经由第二投影
$X \times_Y X \to X$）。假设 $f$ 局部有限型。
注意，$X$ 在 $X$ 上有限表示，而 $X \times_Y X$ 在 $X$ 上局部有限型
（由引理 \ref{morphisms-lemma-base-change-finite-type}）。因此第一个断言来自引理
\ref{morphisms-lemma-finite-presentation-permanence}。
第二个断言来自第一个断言、各定义，以及对角态射是单态射因而分离这一事实
（《概形》引理 \ref{schemes-lemma-monomorphism-separated}）。
\end{proof}

\section{可构造集}
\label{morphisms-section-constructible}

\noindent
我们已在《性质》节 \ref{properties-section-constructible}
中讨论了概形的可构造集与局部可构造集。本节证明一些关于（局部）可构造集
的像与逆像的结果。主要结果是 Chevalley 定理：局部可构造集在有限表示态射
下的像仍然局部可构造。

\begin{lemma}
\label{morphisms-lemma-inverse-image-constructible}
令 $f : X \to Y$ 为概形态射，并令 $E \subset Y$ 为子集。
若 $E$ 在 $Y$ 中（局部）可构造，则 $f^{-1}(E)$ 在 $X$
中（局部）可构造。
\end{lemma}

\begin{proof}
为了证明每个可构造子集的逆像都是可构造的，只需证明每个逆拟紧开集
$V$（它是 $Y$ 的开集）的逆像在 $X$ 中逆拟紧；见《拓扑学》引理
\ref{topology-lemma-inverse-images-constructibles}。$V$ 在 $Y$
中逆拟紧的含义正是开浸入 $V \to Y$ 拟紧。因此基变换
$f^{-1}(V) = X \times_Y V \to X$ 也是拟紧的；见《概形》引理
\ref{schemes-lemma-quasi-compact-preserved-base-change}。
所以 $f^{-1}(V)$ 在 $X$ 中逆拟紧。
假设 $E$ 在 $Y$ 中局部可构造。选取 $x \in X$。选取仿射邻域
$V$（它是 $f(x)$ 的邻域），并选取仿射邻域 $U \subset X$
（它是 $x$ 的邻域），使得
$f(U) \subset V$。于是把 $f|_U : U \to V$ 看成以 $V$ 为靶的态射。
由《性质》引理 \ref{properties-lemma-locally-constructible}，
$E \cap V$ 在 $V$ 中可构造。由可构造情形，
$(f|_U)^{-1}(E \cap V)$ 在 $U$ 中可构造。由于
$(f|_U)^{-1}(E \cap V) = f^{-1}(E) \cap U$，结论成立。
\end{proof}

\begin{lemma}
\label{morphisms-lemma-chevalley}
令 $f : X \to Y$ 为概形态射。假设：
\begin{enumerate}
\item $f$ 拟紧且局部有限表示；以及
\item $Y$ 拟紧且拟分离。
\end{enumerate}
则 $X$ 的每个可构造子集的像在 $Y$ 中可构造。
\end{lemma}

\begin{proof}
由《性质》引理
\ref{properties-lemma-constructible-quasi-compact-quasi-separated}，
只需在 $Y$ 为仿射概形时证明本引理。在这种情形下，$X$ 拟紧。
因此可以写成 $X = U_1 \cup \ldots \cup U_n$，其中每个 $U_i$
都是 $X$ 的仿射开集。若 $E \subset X$ 可构造，则每个
$E \cap U_i$ 也可构造；见《拓扑学》引理
\ref{topology-lemma-open-immersion-constructible-inverse-image}。
由于 $f(E) = \bigcup f(E \cap U_i)$，且可构造集的有限并仍可构造，
问题归结为 $X$ 为仿射概形的情形。在这种情形下，结论就是《代数》定理
\ref{algebra-theorem-chevalley}。
\end{proof}

\begin{theorem}[Chevalley 定理]
\label{morphisms-theorem-chevalley}
\begin{reference}
\cite[IV, Theorem 1.8.4]{EGA}
\end{reference}
令 $f : X \to Y$ 为概形态射。假设 $f$ 拟紧且局部有限表示。
则每个局部可构造子集的像仍然局部可构造。
\end{theorem}

\begin{proof}
令 $E \subset X$ 局部可构造。必须证明 $f(E)$ 也局部可构造。
我们将证明 $f(E) \cap V$ 是可构造的，其中 $V \subset Y$
为任意仿射开集。由此可归结到 $Y$ 为仿射概形的情形。在这种情形下，
$X$ 拟紧。因此可以写成 $X = U_1 \cup \ldots \cup U_n$，
其中每个 $U_i$ 都是 $X$ 的仿射开集。若 $E \subset X$
局部可构造，则每个 $E \cap U_i$ 都可构造；见《性质》引理
\ref{properties-lemma-locally-constructible}。由于
$f(E) = \bigcup f(E \cap U_i)$，且可构造集的有限并仍可构造，
问题归结为 $X$ 为仿射概形的情形。在这种情形下，结论就是《代数》定理
\ref{algebra-theorem-chevalley}。
\end{proof}

\begin{lemma}
\label{morphisms-lemma-constructible-containing-open}
令 $X$ 为概形，令 $x \in X$，并令 $E \subset X$ 为局部可构造子集。
若 $\{x' \mid x' \leadsto x\} \subset E$，则 $E$ 包含 $x$
的一个开邻域。
\end{lemma}

\begin{proof}
假设 $\{x' \mid x' \leadsto x\} \subset E$。可以假设 $X$ 为仿射概形。
在这种情形下，由《性质》引理
\ref{properties-lemma-locally-constructible}，$E$ 可构造。
特别地，补集 $E^c$ 也可构造。由《代数》引理
\ref{algebra-lemma-constructible-is-image}，可以找到仿射概形态射
$f : Y \to X$，使得 $E^c = f(Y)$。令 $Z \subset X$ 为 $f$
的概形论像。由引理 \ref{morphisms-lemma-reach-points-scheme-theoretic-image}
和假设 $\{x' \mid x' \leadsto x\} \subset E$，可知
$x \not \in Z$。因此 $X \setminus Z \subset E$ 是 $x$
的开邻域，并且包含于 $E$。
\end{proof}





\section{开态射}
\label{morphisms-section-open}

\begin{definition}
\label{morphisms-definition-open}
令 $f : X \to S$ 为态射。
\begin{enumerate}
\item 称 $f$ 为{\it 开态射}，如果它在底层拓扑空间上诱导的映射为开映射。
\item 称 $f$ 为{\it 泛开态射}，如果对任意概形态射 $S' \to S$，
基变换 $f' : X_{S'} \to S'$ 都是开态射。
\end{enumerate}
\end{definition}

\noindent
根据《拓扑学》引理
\ref{topology-lemma-closed-open-map-specialization}，
泛化可沿某些类型的拓扑空间开映射提升。事实上，泛化可沿任意概形的开态射提升
（见引理 \ref{morphisms-lemma-open-generizing}）。我们还将看到，泛化可沿概形的平坦态射
提升（引理 \ref{morphisms-lemma-generalizations-lift-flat}）。这有时反过来推出该态射是开态射。

\begin{lemma}
\label{morphisms-lemma-locally-finite-presentation-universally-open}
令 $f : X \to S$ 为态射。
\begin{enumerate}
\item 若 $f$ 局部有限表示，且泛化可沿 $f$ 提升，则 $f$ 是开态射。
\item 若 $f$ 局部有限表示，且泛化可沿 $f$ 的每个基变换提升，
则 $f$ 是泛开态射。
\end{enumerate}
\end{lemma}

\begin{proof}
只需证明第一个断言。问题可归结到 $X$ 和 $S$ 都是仿射概形的情形。
在这种情形下，结论来自《代数》引理
\ref{algebra-lemma-going-up-down-specialization} 和命题
\ref{algebra-proposition-fppf-open}。
\end{proof}

\noindent
关于有限表示平坦态射的情形，另见引理 \ref{morphisms-lemma-fppf-open}。

\begin{lemma}
\label{morphisms-lemma-composition-open}
（泛）开态射的复合仍为（泛）开态射。
\end{lemma}

\begin{proof}
略。
\end{proof}

\begin{lemma}
\label{morphisms-lemma-scheme-over-field-universally-open}
令 $k$ 为域，令 $X$ 为 $k$ 上的概形。结构态射
$X \to \Spec(k)$ 是泛开态射。
\end{lemma}

\begin{proof}
令 $S \to \Spec(k)$ 为态射。必须证明基变换
$X_S \to S$ 是开态射。这个问题在 $S$ 和 $X$ 上都是局部的，
故可假设 $S$ 和 $X$ 都是仿射概形。在这种情形下，结论就是《代数》引理
\ref{algebra-lemma-map-into-tensor-algebra-open}。
\end{proof}

\begin{lemma}
\label{morphisms-lemma-open-generizing}
\begin{reference}
由 \cite[IV, Corollary 1.10.4]{EGA} 中 (a) $\Rightarrow$ (b) 得出。
\end{reference}
令 $\varphi : X \to Y$ 为概形态射。若 $\varphi$ 是开态射，
则 $\varphi$ 是保泛化的（即泛化可沿 $\varphi$ 提升）。
若 $\varphi$ 是泛开态射，则 $\varphi$ 是泛保泛化的。
\end{lemma}

\begin{proof}
假设 $\varphi$ 是开态射。令 $y' \leadsto y$ 为 $Y$ 中点的一次特化。
令 $x \in X$ 且 $\varphi(x) = y$。选取仿射开集
$U \subset X$ 和 $V \subset Y$，使得 $\varphi(U) \subset V$ 且
$x \in U$。于是 $y' \in V$。因此可将 $X$ 替换为 $U$，将 $Y$
替换为 $V$，并假设 $X$、$Y$ 都是仿射概形。仿射情形就是《代数》引理
\ref{algebra-lemma-open-going-down}（结合《代数》引理
\ref{algebra-lemma-going-up-down-specialization}）。
\end{proof}

\begin{lemma}
\label{morphisms-lemma-descent-quasi-compact}
令 $f : X \to Y$ 为概形态射。令 $g : Y' \to Y$ 为开满射，
并假设基变换 $f' : X' \to Y'$ 拟紧。则 $f$ 拟紧。
\end{lemma}

\begin{proof}
令 $V \subset Y$ 为拟紧开集。由于 $g$ 开且满，可以找到拟紧开集
$W' \subset W$，使得 $g(W') = V$。依假设，$(f')^{-1}(W')$
拟紧。$(f')^{-1}(W')$ 在 $X$ 中的像等于 $f^{-1}(V)$；见引理
\ref{morphisms-lemma-when-point-maps-to-pair}。因此 $f^{-1}(V)$ 作为拟紧空间的像
是拟紧的；见《拓扑学》引理 \ref{topology-lemma-image-quasi-compact}。
所以 $f$ 拟紧。
\end{proof}





\section{商态射}
\label{morphisms-section-submersive}

\begin{definition}
\label{morphisms-definition-submersive}
令 $f : X \to Y$ 为概形态射。
\begin{enumerate}
\item 称 $f$ 为{\it 商态射}\footnote{这与微分流形中的浸没概念大不相同。}，
如果底层拓扑空间之间的连续映射是商映射；见《拓扑学》定义
\ref{topology-definition-submersive}。
\item 称 $f$ 为{\it 泛商态射}，如果对每个概形态射 $Y' \to Y$，
基变换 $Y' \times_Y X \to Y'$ 都是商态射。
\end{enumerate}
\end{definition}

\noindent
注意，商态射特别是满射。

\begin{lemma}
\label{morphisms-lemma-base-change-universally-submersive}
概形的泛商态射沿任意概形态射的基变换仍为泛商态射。
\end{lemma}

\begin{proof}
这由定义立即得到。
\end{proof}

\begin{lemma}
\label{morphisms-lemma-composition-universally-submersive}
两个概形的（泛）商态射的复合仍为（泛）商态射。
\end{lemma}

\begin{proof}
略。
\end{proof}


\section{平坦态射}
\label{morphisms-section-flat}

\noindent
平坦性是代数几何中最重要的技术工具之一。本节引入这一概念。除引理
\ref{morphisms-lemma-fppf-open} 外，我们有意只讨论直接的观察。非常重要的一类结果，
即平坦性判别准则，见《代数》各节
\ref{algebra-section-criteria-flatness}、
\ref{algebra-section-flatness-artinian}、
\ref{algebra-section-more-flatness-criteria}，以及《态射进阶》节
\ref{more-morphisms-section-criterion-flat-fibres}。另有一章专门讨论概形平坦态射
的进阶材料，即《平坦性进阶》节 \ref{flat-section-introduction}。

\medskip\noindent
回忆：模 $M$ 在环 $R$ 上称为{\it 平坦的}，如果函子
$-\otimes_R M : \text{Mod}_R \to \text{Mod}_R$ 正合。环同态
$R \to A$ 称为{\it 平坦的}，如果 $A$ 作为 $R$-模平坦。见《代数》定义
\ref{algebra-definition-flat}。

\begin{definition}
\label{morphisms-definition-flat}
令 $f : X \to S$ 为概形态射。令 $\mathcal{F}$ 为拟凝聚
$\mathcal{O}_X$-模层。
\begin{enumerate}
\item 称 $f$ 在点 $x \in X$ 处{\it 平坦}，如果局部环
$\mathcal{O}_{X, x}$ 在局部环 $\mathcal{O}_{S, f(x)}$ 上平坦。
\item 称 $\mathcal{F}$ {\it 在 $S$ 上于点 $x \in X$ 处平坦}，
如果茎 $\mathcal{F}_x$ 是平坦的 $\mathcal{O}_{S, f(x)}$-模。
\item 称 $f$ {\it 平坦}，如果 $f$ 在 $X$ 的每一点处平坦。
\item 称 $\mathcal{F}$ {\it 在 $S$ 上平坦}，如果 $\mathcal{F}$
在 $S$ 上于每一点 $x$ 处平坦，而这些点取自 $X$。
\end{enumerate}
\end{definition}

\noindent
因此，$f$ 平坦当且仅当结构层 $\mathcal{O}_X$ 在 $S$ 上平坦。

\begin{lemma}
\label{morphisms-lemma-flat-module-characterize}
令 $f : X \to S$ 为概形态射。令 $\mathcal{F}$ 为拟凝聚
$\mathcal{O}_X$-模层。下列条件等价：
\begin{enumerate}
\item 层 $\mathcal{F}$ 在 $S$ 上平坦。
\item 对每对仿射开集 $U \subset X$、$V \subset S$，若
$f(U) \subset V$，则 $\mathcal{O}_S(V)$-模 $\mathcal{F}(U)$ 平坦。
\item 存在开覆盖 $S = \bigcup_{j \in J} V_j$ 以及开覆盖
$f^{-1}(V_j) = \bigcup_{i \in I_j} U_i$，使得每个模层
$\mathcal{F}|_{U_i}$ 在 $V_j$ 上平坦，其中所有
$j\in J, i\in I_j$ 均如此。
\item 存在仿射开覆盖 $S = \bigcup_{j \in J} V_j$ 以及仿射开覆盖
$f^{-1}(V_j) = \bigcup_{i \in I_j} U_i$，使得
$\mathcal{F}(U_i)$ 是平坦的 $\mathcal{O}_S(V_j)$-模，其中所有
$j\in J, i\in I_j$ 均如此。
\end{enumerate}
此外，若 $\mathcal{F}$ 在 $S$ 上平坦，则对任意开子概形
$U \subset X$、$V \subset S$，若 $f(U) \subset V$，限制
$\mathcal{F}|_U$ 在 $V$ 上平坦。
\end{lemma}

\begin{proof}
令 $R \to A$ 为环同态，令 $M$ 为 $A$-模。若 $M$ 在 $R$ 上平坦，
则对所有素理想 $\mathfrak q$，模 $M_{\mathfrak q}$ 在
$R_{\mathfrak p}$ 上平坦，其中 $\mathfrak p$ 是 $R$ 中位于
$\mathfrak q$ 之下的素理想。反之，若 $M_{\mathfrak q}$ 在
$R_{\mathfrak p}$ 上平坦，且这对每个素理想 $\mathfrak q$
（它属于 $A$）均成立，
则 $M$ 在 $R$ 上平坦。见《代数》引理
\ref{algebra-lemma-flat-localization}。这一等价性容易推出本引理的各个断言。
\end{proof}

\begin{lemma}
\label{morphisms-lemma-flat-characterize}
令 $f : X \to S$ 为概形态射。下列条件等价：
\begin{enumerate}
\item 态射 $f$ 平坦。
\item 对每对仿射开集 $U \subset X$、$V \subset S$，若
$f(U) \subset V$，则环同态
$\mathcal{O}_S(V) \to \mathcal{O}_X(U)$ 平坦。
\item 存在开覆盖 $S = \bigcup_{j \in J} V_j$ 以及开覆盖
$f^{-1}(V_j) = \bigcup_{i \in I_j} U_i$，使得每个态射
$U_i \to V_j$ 都平坦，其中 $j\in J, i\in I_j$。
\item 存在仿射开覆盖 $S = \bigcup_{j \in J} V_j$ 以及仿射开覆盖
$f^{-1}(V_j) = \bigcup_{i \in I_j} U_i$，使得
$\mathcal{O}_S(V_j) \to \mathcal{O}_X(U_i)$ 平坦，其中所有
$j\in J, i\in I_j$ 均如此。
\end{enumerate}
此外，若 $f$ 平坦，则对任意开子概形 $U \subset X$、$V \subset S$，
若 $f(U) \subset V$，限制 $f|_U : U \to V$ 平坦。
\end{lemma}

\begin{proof}
这是上面引理 \ref{morphisms-lemma-flat-module-characterize} 的特殊情形。
\end{proof}

\begin{lemma}
\label{morphisms-lemma-pushforward-flat-affine}
令 $f : X \to Y$ 为基概形 $S$ 上的仿射概形态射。令
$\mathcal{F}$ 为拟凝聚 $\mathcal{O}_X$-模。则 $\mathcal{F}$ 在
$S$ 上平坦，当且仅当 $f_*\mathcal{F}$ 在 $S$ 上平坦。
\end{lemma}

\begin{proof}
由引理 \ref{morphisms-lemma-flat-module-characterize} 以及 $f$ 是仿射态射这一事实，
问题归结到仿射情形。设 $X \to Y \to S$ 对应于环同态
$C \leftarrow B \leftarrow A$。令 $N$ 为一个 $C$-模，它对应于
$\mathcal{F}$。回忆，$f_*\mathcal{F}$ 对应于把 $N$ 看成 $B$-模；见《概形》引理
\ref{schemes-lemma-widetilde-pullback}。因此结论显然。
\end{proof}

\begin{lemma}
\label{morphisms-lemma-composition-module-flat}
令 $X \to Y \to Z$ 为概形态射。令 $\mathcal{F}$ 为拟凝聚
$\mathcal{O}_X$-模。令 $x \in X$，其像为 $y$，位于 $Y$ 中。
若 $\mathcal{F}$ 在 $Y$ 上于点 $x$ 处平坦，并且 $Y$ 在 $Z$
上于点 $y$ 处平坦，则 $\mathcal{F}$ 在 $Z$ 上于点 $x$ 处平坦。
\end{lemma}

\begin{proof}
见《代数》引理 \ref{algebra-lemma-composition-flat}。
\end{proof}

\begin{lemma}
\label{morphisms-lemma-composition-flat}
平坦态射的复合仍平坦。
\end{lemma}

\begin{proof}
这是引理 \ref{morphisms-lemma-composition-module-flat} 的特殊情形。
\end{proof}

\begin{lemma}
\label{morphisms-lemma-base-change-module-flat}
令 $f : X \to S$ 为概形态射。令 $\mathcal{F}$ 为拟凝聚
$\mathcal{O}_X$-模层。令 $g : S' \to S$ 为概形态射。
记投影为 $g' : X' = X_{S'} \to X$。令 $x' \in X'$ 为一点，
其像为 $x = g'(x') \in X$。若 $\mathcal{F}$ 在 $S$ 上于点
$x$ 处平坦，则 $(g')^*\mathcal{F}$ 在 $S'$ 上于点 $x'$ 处平坦。
特别地，若 $\mathcal{F}$ 在 $S$ 上平坦，则
$(g')^*\mathcal{F}$ 在 $S'$ 上平坦。
\end{lemma}

\begin{proof}
见《代数》引理 \ref{algebra-lemma-flat-base-change}。
\end{proof}

\begin{lemma}
\label{morphisms-lemma-base-change-flat}
平坦态射的基变换仍平坦。
\end{lemma}

\begin{proof}
这是引理 \ref{morphisms-lemma-base-change-module-flat} 的特殊情形。
\end{proof}

\begin{lemma}
\label{morphisms-lemma-generalizations-lift-flat}
令 $f : X \to S$ 为概形的平坦态射。则泛化可沿 $f$ 提升；见《拓扑学》定义
\ref{topology-definition-lift-specializations}。
\end{lemma}

\begin{proof}
见《代数》节 \ref{algebra-section-going-up}。
\end{proof}

\begin{lemma}
\label{morphisms-lemma-fppf-open}
局部有限表示的平坦态射是泛开态射。
\end{lemma}

\begin{proof}
这来自引理 \ref{morphisms-lemma-generalizations-lift-flat} 以及上面的引理
\ref{morphisms-lemma-locally-finite-presentation-universally-open}。
也可以如下直接论证。

\medskip\noindent
令 $f : X \to S$ 平坦且局部有限表示。由引理
\ref{morphisms-lemma-base-change-flat} 和
\ref{morphisms-lemma-base-change-finite-presentation}，$f$ 的任意基变换都平坦且
局部有限表示。因此只需证明 $f$ 是开态射。为证明 $f$ 是开态射，
只需证明可以用仿射开集覆盖 $X$，即 $X = \bigcup U_i$，使得
$U_i \to S$ 是开态射。可以覆盖 $X$，所用仿射开集为 $U_i \subset X$，
使得每个 $U_i$ 都映入仿射开集 $V_i \subset S$，并且所诱导的环同态
$\mathcal{O}_S(V_i) \to \mathcal{O}_X(U_i)$ 平坦且有限表示（引理
\ref{morphisms-lemma-flat-characterize} 和
\ref{morphisms-lemma-locally-finite-presentation-characterize}）。于是由《代数》命题
\ref{algebra-proposition-fppf-open}，$U_i \to V_i$ 是开态射，证明完成。
\end{proof}

\begin{lemma}
\label{morphisms-lemma-pf-flat-module-open}
令 $f : X \to Y$ 为概形态射。令 $\mathcal{F}$ 为拟凝聚
$\mathcal{O}_X$-模。假设 $f$ 局部有限表示，$\mathcal{F}$ 有限型，
$X = \text{Supp}(\mathcal{F})$，并且 $\mathcal{F}$ 在 $Y$ 上平坦。
则 $f$ 是泛开态射。
\end{lemma}

\begin{proof}
由引理 \ref{morphisms-lemma-base-change-module-flat}、
\ref{morphisms-lemma-base-change-finite-presentation} 和
\ref{morphisms-lemma-support-finite-type}，这些假设在基变换下保持。
由引理 \ref{morphisms-lemma-locally-finite-presentation-universally-open}，
只需证明泛化可沿 $f$ 提升。这来自《代数》引理
\ref{algebra-lemma-going-down-flat-module}。
\end{proof}

\begin{lemma}
\label{morphisms-lemma-fpqc-quotient-topology}
\begin{reference}
\cite[Expose VIII, Corollaire 4.3]{SGA1} 和
\cite[IV, Corollaire 2.3.12]{EGA}
\end{reference}
令 $f : X \to Y$ 为拟紧、满且平坦的态射。子集 $T \subset Y$ 为开集
（相应地，闭集），当且仅当 $f^{-1}(T)$ 为开集（相应地，闭集）。换言之，
$f$ 是商态射。
\end{lemma}

\begin{proof}
这个问题在 $Y$ 上是局部的，故可假设 $Y$ 为仿射概形。在这种情形下，
$X$ 也拟紧，因为 $f$ 拟紧。把 $X = X_1 \cup \ldots \cup X_n$
写成仿射开集的有限并。于是
$f' : X' = X_1 \amalg \ldots \amalg X_n \to Y$ 是仿射概形之间的
满平坦态射。注意，对 $T \subset Y$，有
$(f')^{-1}(T) = f^{-1}(T) \cap X_1 \amalg \ldots \amalg f^{-1}(T) \cap X_n$。
因此 $f^{-1}(T)$ 为开集，当且仅当 $(f')^{-1}(T)$ 为开集。
所以可假设 $X$ 和 $Y$ 都是仿射概形。

\medskip\noindent
令 $f : \Spec(B) \to \Spec(A)$ 为仿射概形的满态射，对应于平坦环同态
$A \to B$。假设 $f^{-1}(T)$ 为闭集，比如
$f^{-1}(T) = V(J)$，其中 $J \subset B$ 为理想。于是
$T = f(f^{-1}(T)) = f(V(J))$ 是 $\Spec(B/J) \to \Spec(A)$ 的像
（这里用到了 $f$ 的满性）。另一方面，泛化可沿 $f$ 提升（引理
\ref{morphisms-lemma-generalizations-lift-flat}）。因此，由《拓扑学》引理
\ref{topology-lemma-lift-specializations-images}，
$Y \setminus T = f(X \setminus f^{-1}(T))$ 在泛化下稳定。所以 $T$
在特化下稳定（《拓扑学》引理
\ref{topology-lemma-open-closed-specialization}）。由《代数》引理
\ref{algebra-lemma-image-stable-specialization-closed}，$T$ 是闭集。
\end{proof}

\begin{lemma}
\label{morphisms-lemma-flat-permanence}
令 $h : X \to Y$ 为 $S$ 上的概形态射。令 $\mathcal{G}$ 为 $Y$
上的拟凝聚层。令 $x \in X$，且 $y = h(x) \in Y$。若 $h$ 在 $x$
处平坦，则
$$
\mathcal{G}\text{ flat over }S\text{ 在 }y
\Leftrightarrow
h^*\mathcal{G}\text{ flat over }S\text{ 在 }x.
$$
特别地：若 $h$ 满且平坦，则 $\mathcal{G}$ 在 $S$ 上平坦，当且仅当
$h^*\mathcal{G}$ 在 $S$ 上平坦。若 $h$ 满且平坦，并且 $X$ 在 $S$
上平坦，则 $Y$ 在 $S$ 上平坦。
\end{lemma}

\begin{proof}
可以应用《代数》引理
\ref{algebra-lemma-flatness-descends-more-general} 来证明。下面给出直接证明。
令 $s \in S$ 为 $y$ 的像。考虑局部环同态
$\mathcal{O}_{S, s} \to \mathcal{O}_{Y, y} \to \mathcal{O}_{X, x}$。
由假设，环同态 $\mathcal{O}_{Y, y} \to \mathcal{O}_{X, x}$ 忠实平坦；
见《代数》引理 \ref{algebra-lemma-local-flat-ff}。令
$N = \mathcal{G}_y$。注意，
$h^*\mathcal{G}_x = N \otimes_{\mathcal{O}_{Y, y}} \mathcal{O}_{X, x}$；
见《层》引理 \ref{sheaves-lemma-stalk-pullback-modules}。令
$M' \to M$ 为 $\mathcal{O}_{S, s}$-模的单射。由上述忠实平坦性，有
\begin{align*}
\Ker(
M' \otimes_{\mathcal{O}_{S, s}} N \to M \otimes_{\mathcal{O}_{S, s}} N)
\otimes_{\mathcal{O}_{Y, y}} \mathcal{O}_{X, x} \\
=
\Ker(
M' \otimes_{\mathcal{O}_{S, s}} N
\otimes_{\mathcal{O}_{Y, y}} \mathcal{O}_{X, x}
\to
M \otimes_{\mathcal{O}_{S, s}} N
\otimes_{\mathcal{O}_{Y, y}} \mathcal{O}_{X, x})
\end{align*}
所以本引理的等价性来自《代数》引理 \ref{algebra-lemma-flat}
中平坦性的第二种刻画。
\end{proof}

\begin{lemma}
\label{morphisms-lemma-flat-pullback-support}
令 $f : Y \to X$ 为概形态射。令 $\mathcal{F}$ 为有限型拟凝聚
$\mathcal{O}_X$-模，其概形论支集为 $Z \subset X$。若 $f$ 平坦，
则 $f^{-1}(Z)$ 是 $f^*\mathcal{F}$ 的概形论支集。
\end{lemma}

\begin{proof}
使用引理 \ref{morphisms-lemma-scheme-theoretic-support} 给出的仿射概形上概形论支集的刻画，
问题归结到《代数》引理
\ref{algebra-lemma-annihilator-flat-base-change}。
\end{proof}

\begin{lemma}
\label{morphisms-lemma-flat-morphism-scheme-theoretically-dense-open}
令 $f : X \to Y$ 为概形的平坦态射。令 $V \subset Y$ 为逆拟紧开集，
且概形论稠密。则 $f^{-1}V$ 在 $X$ 中概形论稠密。
\end{lemma}

\begin{proof}
我们将使用引理 \ref{morphisms-lemma-characterize-scheme-theoretically-dense} 的刻画。
必须证明，对任意开集 $U \subset X$，映射
$\mathcal{O}_X(U) \to \mathcal{O}_X(U \cap f^{-1}V)$ 都是单射。
只需在 $U$ 是映入某个仿射开集 $W \subset Y$ 的仿射开集时证明这一点。
设 $W = \Spec(A)$ 且 $U = \Spec(B)$。由《代数》引理
\ref{algebra-lemma-qc-open}，有
$V \cap W = D(f_1) \cup \ldots \cup D(f_n)$，其中某些
$f_i \in A$。
因此必须证明
$B \to B_{f_1} \times \ldots \times B_{f_n}$ 是单射。已知
$A \to A_{f_1} \times \ldots \times A_{f_n}$ 是单射，并且
$A \to B$ 平坦。由于 $B_{f_i} = A_{f_i} \otimes_A B$，结论成立。
\end{proof}

\begin{lemma}
\label{morphisms-lemma-flat-base-change-scheme-theoretic-image}
\begin{slogan}
在拟紧情形下，取概形论像与平坦基变换可交换
\end{slogan}
令 $f : X \to Y$ 为概形的平坦态射。令 $g : V \to Y$ 为概形的拟紧态射。
令 $Z \subset Y$ 为 $g$ 的概形论像，并令 $Z' \subset X$ 为基变换
$V \times_Y X \to X$ 的概形论像。则 $Z' = f^{-1}Z$。
\end{lemma}

\begin{proof}
回忆，$Z$ 由下式所截：
$\mathcal{I} = \Ker(\mathcal{O}_Y \to g_*\mathcal{O}_V)$，
而 $Z'$ 由下式所截：
$\mathcal{I}' = \Ker(\mathcal{O}_X \to
(V \times_Y X \to X)_*\mathcal{O}_{V \times_Y X})$；见引理
\ref{morphisms-lemma-quasi-compact-scheme-theoretic-image}。因此这个问题在 $X$ 和
$Y$ 上是局部的，可以假设 $X$ 和 $Y$ 都是仿射概形。注意，可以将
$V$ 替换为 $\coprod V_i$，其中 $V = V_1 \cup \ldots \cup V_n$
是有限仿射开覆盖。因此可假设 $g$ 为仿射态射。在这种情形下，
$(V \times_Y X \to X)_*\mathcal{O}_{V \times_Y X}$ 是
$g_*\mathcal{O}_V$ 沿 $f$ 的拉回。由于 $f$ 平坦，可得
$f^*\mathcal{I} = \mathcal{I}'$，本引理成立。
\end{proof}




\section{平坦闭浸入}
\label{morphisms-section-flat-closed-immersions}

\noindent
概形的连通分支不总是开集，但它们总有典范概形结构。本节说明这一点。

\begin{lemma}
\label{morphisms-lemma-characterize-flat-closed-immersions}
令 $X$ 为概形。把 $X$ 的每个闭子概形对应到底层闭子集的规则定义了双射
$$
\left\{
\begin{matrix}
\text{closed subschemes }Z \subset X \\
\text{使得 }Z \to X\text{ 平坦}
\end{matrix}
\right\}
\leftrightarrow
\left\{
\begin{matrix}
\text{closed subsets }Z \subset X \\
\text{closed under generalizations}
\end{matrix}
\right\}
$$
若 $Z \subset X$ 是这样的闭子概形，则每个概形态射
$g : Y \to X$，若在集合论意义下 $g(Y) \subset Z$，都在概形论意义下
唯一地通过 $Z$ 分解。
\end{lemma}

\begin{proof}
双射的仿射情形就是《代数》引理
\ref{algebra-lemma-pure-open-closed-specializations}。对一般概形 $X$，
用仿射开集覆盖 $X$ 并黏合即可得到双射。细节略去。关于最后一个断言，
注意投影 $Z \times_{X, g} Y \to Y$ 是平坦的（引理
\ref{morphisms-lemma-base-change-flat}）闭浸入，并且在底层拓扑空间上是双射，
所以由证明第一部分所建立的双射，它必为同构。
\end{proof}

\begin{lemma}
\label{morphisms-lemma-flat-closed-immersions-finite-presentation}
有限表示的平坦闭浸入是某个开闭子概形的开浸入。
\end{lemma}

\begin{proof}
仿射情形就是《代数》引理
\ref{algebra-lemma-finitely-generated-pure-ideal}。一般情形下，
用仿射开集覆盖 $X$ 即得本引理。细节略去。
\end{proof}

\noindent
注意，连通分支 $T$（属于概形 $X$）是在泛化下稳定的闭子集。
因此下述定义是有意义的。

\begin{definition}
\label{morphisms-definition-scheme-structure-connected-component}
令 $X$ 为概形。令 $T \subset X$ 为连通分支。$T$ 上的
{\it 典范概形结构}是 $T$ 上唯一的概形结构，使得闭浸入
$T \to X$ 平坦；见引理
\ref{morphisms-lemma-characterize-flat-closed-immersions}。
\end{definition}

\noindent
利用前一引理，可以判定每个有限平坦 $\mathcal{O}_X$-模何时有限局部自由。

\begin{lemma}
\label{morphisms-lemma-finite-flat-is-finite-locally-free}
令 $X$ 为概形。下列条件等价：
\begin{enumerate}
\item 每个有限平坦拟凝聚 $\mathcal{O}_X$-模都是有限局部自由的；以及
\item 每个在泛化下闭合的闭子集 $Z \subset X$ 都是开集。
\end{enumerate}
\end{lemma}

\begin{proof}
在仿射情形下，这就是《代数》引理
\ref{algebra-lemma-finite-flat-module-finitely-presented}。
概形情形不能直接从仿射情形推出，所以我们直接重复论证。

\medskip\noindent
假设 (1)。考虑闭浸入 $i : Z \to X$，并假设 $i$ 平坦。于是
$i_*\mathcal{O}_Z$ 拟凝聚且平坦，故由 (1) 有限局部自由。因此
$Z = \text{Supp}(i_*\mathcal{O}_Z)$ 也是开集，从而 (2) 成立。
所以由引理 \ref{morphisms-lemma-characterize-flat-closed-immersions}
中对平坦闭浸入的刻画，可得 (1) $\Rightarrow$ (2)。

\medskip\noindent
反之，假设 $X$ 满足 (2)。令 $\mathcal{F}$ 为有限平坦拟凝聚
$\mathcal{O}_X$-模。令 $Z = \text{Supp}(\mathcal{F})$；这个集合是
$\mathcal{F}$ 的支集，并且为闭集；见《模》引理
\ref{modules-lemma-support-finite-type-closed}。另一方面，若
$x \leadsto x'$ 是一次特化，则由《代数》引理
\ref{algebra-lemma-finite-flat-local}，模 $\mathcal{F}_{x'}$ 在
$\mathcal{O}_{X, x'}$ 上自由，并且
$$
\mathcal{F}_x =
\mathcal{F}_{x'} \otimes_{\mathcal{O}_{X, x'}} \mathcal{O}_{X, x}.
$$
因此
$x' \in \text{Supp}(\mathcal{F}) \Rightarrow x \in \text{Supp}(\mathcal{F})$，
换言之，该支集在泛化下闭合。由于 $X$ 满足 (2)，可知
$\mathcal{F}$ 的支集是开闭集。模
$\wedge^i(\mathcal{F})$，$i = 1, 2, 3, \ldots$，也都是有限平坦拟凝聚
$\mathcal{O}_X$-模；见《模》节
\ref{modules-section-symmetric-exterior}。注意
$\text{Supp}(\wedge^{i + 1}(\mathcal{F})) \subset
\text{Supp}(\wedge^i(\mathcal{F}))$。
因此存在开闭子集的分解
$$
X = U_0 \amalg U_1 \amalg U_2 \amalg \ldots
$$
使得 $\wedge^i(\mathcal{F})$ 的支集都是
$U_i \cup U_{i + 1} \cup \ldots$，对所有 $i$ 均如此。令 $x$ 为 $X$ 的一点，设
$x \in U_r$。注意
$\wedge^i(\mathcal{F})_x \otimes \kappa(x) =
\wedge^i(\mathcal{F}_x \otimes \kappa(x))$。
因此，$x \in U_r$ 推出 $\mathcal{F}_x \otimes \kappa(x)$
是维数为 $r$ 的向量空间。由 Nakayama 引理（见《代数》引理
\ref{algebra-lemma-NAK}），可以选取仿射开邻域
$U \subset U_r \subset X$（它是 $x$ 的邻域），以及截面
$s_1, \ldots, s_r \in \mathcal{F}(U)$，使得所诱导的映射
$$
\mathcal{O}_U^{\oplus r} \longrightarrow \mathcal{F}|_U, \quad
(f_1, \ldots, f_r) \longmapsto \sum f_i s_i
$$
满射。这意味着 $\wedge^r(\mathcal{F}|_U)$ 是有限平坦拟凝聚
$\mathcal{O}_U$-模，其支集是整个 $U$。由上面的论证，它由单个元素
$s_1 \wedge \ldots \wedge s_r$ 生成。因此
$\wedge^r(\mathcal{F}|_U) \cong \mathcal{O}_U/\mathcal{I}$，
其中 $\mathcal{I}$ 是某个拟凝聚理想层，使得
$\mathcal{O}_U/\mathcal{I}$ 在 $\mathcal{O}_U$ 上平坦，并且
$V(\mathcal{I}) = U$。应用引理
\ref{morphisms-lemma-characterize-flat-closed-immersions}，可得
$\mathcal{I} = 0$。所以 $s_1 \wedge \ldots \wedge s_r$ 是
$\wedge^r(\mathcal{F}|_U)$ 的一组基，进而上述映射既单又满。
这证明了 $\mathcal{F}$ 如所要求的是有限局部自由的。
\end{proof}






\section{一般平坦性}
\label{morphisms-section-generic-flatness}

\noindent
整基底上的有限型概形在基底的某个稠密开集上平坦。在《代数》节
\ref{algebra-section-generic-flatness} 中，我们证明了 Noether 情形、
有限表示态射的情形以及一般情形。本节只陈述并证明一般情形。不过，
命题 \ref{morphisms-proposition-generic-flatness-reduced} 将取代这一结果；
它表明只假设基底约化，结论仍然成立。

\begin{proposition}[一般平坦性]
\label{morphisms-proposition-generic-flatness}
令 $f : X \to S$ 为概形态射。令 $\mathcal{F}$ 为拟凝聚
$\mathcal{O}_X$-模层。假设：
\begin{enumerate}
\item $S$ 是整概形；
\item $f$ 有限型；并且
\item $\mathcal{F}$ 是有限型 $\mathcal{O}_X$-模。
\end{enumerate}
则存在稠密开子概形 $U \subset S$，使得 $X_U \to U$ 平坦且有限表示，
并且 $\mathcal{F}|_{X_U}$ 在 $U$ 上平坦且在
$\mathcal{O}_{X_U}$ 上有限表示。
\end{proposition}

\begin{proof}
由于 $S$ 是整概形，它不可约（见《性质》引理
\ref{properties-lemma-characterize-integral}），且任意非空开集都稠密。
因此可以用 $S$ 的一个仿射开集替换 $S$，并假设
$S = \Spec(A)$ 为仿射概形。由于 $S$ 是整概形，可知 $A$ 是整环。
由于 $f$ 有限型，它拟紧，所以 $X$ 拟紧。因此可以找到有限仿射开覆盖
$X = \bigcup_{i = 1, \ldots, n} X_i$。写成 $X_i = \Spec(B_i)$。
由引理 \ref{morphisms-lemma-locally-finite-type-characterize}，$B_i$ 是有限型
$A$-代数。此外，存在有限型 $B_i$-模 $M_i$，使得
$\mathcal{F}|_{X_i}$ 是对应于 $B_i$-模 $M_i$ 的拟凝聚层；见《性质》引理
\ref{properties-lemma-finite-type-module}。接着，对每对指标 $i, j$，
选取理想 $I_{ij} \subset B_i$，使得
$X_i \setminus X_i \cap X_j = V(I_{ij})$ 在
$X_i = \Spec(B_i)$ 中成立。令
$M_{ij} = B_i/I_{ij}$，并把它看成 $B_i$-模。于是
$V(I_{ij}) = \text{Supp}(M_{ij})$，且 $M_{ij}$ 是有限 $B_i$-模。

\medskip\noindent
现在对各对 $(A \to B_i, M_{ij})$ 以及各对
$(A \to B_i, M_i)$ 应用《代数》引理
\ref{algebra-lemma-generic-flatness}。因此得到非零元
$f_{ij}, f_i \in A$，使得：(a)
$A_{f_{ij}} \to B_{i, f_{ij}}$ 平坦且有限表示，
$M_{ij, f_{ij}}$ 在 $A_{f_{ij}}$ 上平坦且在 $B_{i, f_{ij}}$
上有限表示；并且 (b) $B_{i, f_i}$ 在 $A_f$ 上平坦且有限表示，
$M_{i, f_i}$ 在 $B_{i, f_i}$ 上平坦且有限表示。令
$f = (\prod f_i) (\prod f_{ij})$。我们断言取 $U = D(f)$ 即可。

\medskip\noindent
为证明这一断言，可以将 $A$ 替换为 $A_f$，也就是施行基变换
$U = \Spec(A_f) \to S$。基变换后，每个 $A \to B_i$ 都平坦且
有限表示，并且 $M_i$、$M_{ij}$ 在 $A$ 上平坦且在 $B_i$ 上
有限表示。这已经证明 $X \to S$ 拟紧、局部有限表示且平坦，并证明
$\mathcal{F}$ 在 $S$ 上平坦且在 $\mathcal{O}_X$ 上有限表示；见引理
\ref{morphisms-lemma-locally-finite-presentation-characterize} 和《性质》引理
\ref{properties-lemma-finite-presentation-module}。由于 $M_{ij}$ 在
$B_i$ 上有限表示，可知
$X_i \cap X_j = X_i \setminus \text{Supp}(M_{ij})$ 是 $X_i$
的拟紧开集；见《代数》引理
\ref{algebra-lemma-support-finite-presentation-constructible}。
所以由《概形》引理 \ref{schemes-lemma-characterize-quasi-separated}，
$X \to S$ 拟分离。这证明了本命题。
\end{proof}

\noindent
事实上，在任意约化基底上也有一般平坦性版本。陈述如下。

\begin{proposition}[一般平坦性，约化情形]
\label{morphisms-proposition-generic-flatness-reduced}
令 $f : X \to S$ 为概形态射。令 $\mathcal{F}$ 为拟凝聚
$\mathcal{O}_X$-模层。假设：
\begin{enumerate}
\item $S$ 约化；
\item $f$ 有限型；并且
\item $\mathcal{F}$ 是有限型 $\mathcal{O}_X$-模。
\end{enumerate}
则存在稠密开子概形 $U \subset S$，使得 $X_U \to U$ 平坦且有限表示，
并且 $\mathcal{F}|_{X_U}$ 在 $U$ 上平坦且在
$\mathcal{O}_{X_U}$ 上有限表示。
\end{proposition}

\begin{proof}
对急于得知结果的读者：本证明重复命题
\ref{morphisms-proposition-generic-flatness} 的证明，只是用《代数》引理
\ref{algebra-lemma-generic-flatness-reduced} 代替《代数》引理
\ref{algebra-lemma-generic-flatness}。

\medskip\noindent
由于平坦和有限表示都在基底上是局部性质（见引理
\ref{morphisms-lemma-flat-module-characterize} 和
\ref{morphisms-lemma-locally-finite-presentation-characterize}），可以在 $S$
上仿射局部地工作。因此可假设 $S = \Spec(A)$，其中 $A$ 为约化环
（见《性质》引理 \ref{properties-lemma-characterize-reduced}）。
由于 $f$ 有限型，它拟紧，所以 $X$ 拟紧。因此可以找到有限仿射开覆盖
$X = \bigcup_{i = 1, \ldots, n} X_i$。写成 $X_i = \Spec(B_i)$。
由引理 \ref{morphisms-lemma-locally-finite-type-characterize}，$B_i$ 是有限型
$A$-代数。此外，存在有限型 $B_i$-模 $M_i$，使得
$\mathcal{F}|_{X_i}$ 是对应于 $B_i$-模 $M_i$ 的拟凝聚层；见《性质》引理
\ref{properties-lemma-finite-type-module}。接着，对每对指标 $i, j$，
选取理想 $I_{ij} \subset B_i$，使得
$X_i \setminus X_i \cap X_j = V(I_{ij})$ 在
$X_i = \Spec(B_i)$ 中成立。令
$M_{ij} = B_i/I_{ij}$，并把它看成 $B_i$-模。于是
$V(I_{ij}) = \text{Supp}(M_{ij})$，且 $M_{ij}$ 是有限 $B_i$-模。

\medskip\noindent
现在对各对 $(A \to B_i, M_{ij})$ 以及各对
$(A \to B_i, M_i)$ 应用《代数》引理
\ref{algebra-lemma-generic-flatness-reduced}。因此得到稠密开集
$U(A \to B_i, M_{ij}) \subset S$ 以及稠密开集
$U(A \to B_i, M_i) \subset S$，记号如《代数》方程
(\ref{algebra-equation-good-locus})。由于稠密开集的有限交仍为稠密开集，
可知
$$
U =
\bigcap\nolimits_{i, j} U(A \to B_i, M_{ij})
\quad\cap\quad
\bigcap\nolimits_i U(A \to B_i, M_i)
$$
在 $S$ 中开且稠密。我们断言 $U$ 就是所需开集。

\medskip\noindent
取 $u \in U$。由轨迹 $U(A \to B_i, M_{ij})$ 和
$U(A \to B, M_i)$ 的定义，存在 $f_{ij}, f_i \in A$，使得：
(a) $u \in D(f_i)$ 且 $u \in D(f_{ij})$；
(b) $A_{f_{ij}} \to B_{i, f_{ij}}$ 平坦且有限表示，
$M_{ij, f_{ij}}$ 在 $A_{f_{ij}}$ 上平坦且在 $B_{i, f_{ij}}$
上有限表示；并且
(c) $B_{i, f_i}$ 在 $A_{f_i}$ 上平坦且有限表示，
$M_{i, f_i}$ 在 $B_{i, f_i}$ 上有限表示且在 $A_{f_i}$ 上平坦。
令 $f = (\prod f_i) (\prod f_{ij})$。现在只需证明
$X \to S$ 在 $D(f)$ 上平坦且有限表示，并且 $\mathcal{F}$ 限制到
$X_{D(f)}$ 后在 $D(f)$ 上平坦且在 $X_{D(f)}$ 的结构层上有限表示。

\medskip\noindent
因此可以将 $A$ 替换为 $A_f$，也就是施行基变换
$\Spec(A_f) \to S$。基变换后，每个 $A \to B_i$ 都平坦且有限表示，
并且 $M_i$、$M_{ij}$ 在 $A$ 上平坦且在 $B_i$ 上有限表示。
这已经证明 $X \to S$ 拟紧、局部有限表示且平坦，并证明
$\mathcal{F}$ 在 $S$ 上平坦且在 $\mathcal{O}_X$ 上有限表示；见引理
\ref{morphisms-lemma-locally-finite-presentation-characterize} 和《性质》引理
\ref{properties-lemma-finite-presentation-module}。由于 $M_{ij}$ 在
$B_i$ 上有限表示，可知
$X_i \cap X_j = X_i \setminus \text{Supp}(M_{ij})$ 是 $X_i$
的拟紧开集；见《代数》引理
\ref{algebra-lemma-support-finite-presentation-constructible}。
所以由《概形》引理 \ref{schemes-lemma-characterize-quasi-separated}，
$X \to S$ 拟分离。这证明了本命题。
\end{proof}









\begin{remark}
\label{morphisms-remark-flattening}
上述结果是概形态射的更精细平坦化技术的第一步。Raynaud 和 Gruson
的文章 \cite{GruRay} 包含许多这方面的精彩结果。
\end{remark}







\section{态射与纤维的维数}
\label{morphisms-section-dimension-fibres}

\noindent
令 $X$ 为拓扑空间，令 $x \in X$。回忆，$\dim_x(X)$ 已定义为
$x$ 的开邻域（在 $X$ 中）的维数的最小值。见《拓扑学》定义
\ref{topology-definition-Krull}。

\begin{lemma}
\label{morphisms-lemma-dimension-fibre-at-a-point}
令 $f : X \to S$ 为概形态射。令 $x \in X$，并置 $s = f(x)$。
假设 $f$ 局部有限型。则
$$
\dim_x(X_s) =
\dim(\mathcal{O}_{X_s, x}) + \text{trdeg}_{\kappa(s)}(\kappa(x)).
$$
\end{lemma}

\begin{proof}
问题立即归结到 $S = s$ 且 $X$ 为仿射概形的情形。在这种情形下，
结论来自《代数》引理
\ref{algebra-lemma-dimension-at-a-point-finite-type-field}。
\end{proof}

\begin{lemma}
\label{morphisms-lemma-dimension-fibre-at-a-point-additive}
令 $f : X \to Y$ 和 $g : Y \to S$ 为概形态射。令 $x \in X$，
并置 $y = f(x)$、$s = g(y)$。假设 $f$ 和 $g$ 局部有限型。则
$$
\dim_x(X_s) \leq \dim_x(X_y) + \dim_y(Y_s).
$$
此外，若 $\mathcal{O}_{X_s, x}$ 在 $\mathcal{O}_{Y_s, y}$ 上平坦，
则等号成立；例如，若 $\mathcal{O}_{X, x}$ 在 $\mathcal{O}_{Y, y}$
上平坦，就满足这一条件。
\end{lemma}

\begin{proof}
注意
$\text{trdeg}_{\kappa(s)}(\kappa(x)) =
\text{trdeg}_{\kappa(y)}(\kappa(x)) + \text{trdeg}_{\kappa(s)}(\kappa(y))$。
因此由引理 \ref{morphisms-lemma-dimension-fibre-at-a-point}，该断言等价于
$$
\dim(\mathcal{O}_{X_s, x})
\leq
\dim(\mathcal{O}_{X_y, x}) + \dim(\mathcal{O}_{Y_s, y}).
$$
对此见《代数》引理
\ref{algebra-lemma-dimension-base-fibre-total}。平坦情形见《代数》引理
\ref{algebra-lemma-dimension-base-fibre-equals-total}。
\end{proof}

\begin{lemma}
\label{morphisms-lemma-dimension-fibre-after-base-change}
令
$$
\xymatrix{
X' \ar[r]_{g'} \ar[d]_{f'} & X \ar[d]^f \\
S' \ar[r]^g & S
}
$$
为概形的纤维积图。假设 $f$ 局部有限型。设 $x' \in X'$、
$x = g'(x')$、$s' = f'(x')$，并且 $s = g(s') = f(x)$。则：
\begin{enumerate}
\item $\dim_x(X_s) = \dim_{x'}(X'_{s'})$；
\item 若 $F$ 是态射 $X'_{s'} \to X_s$ 在 $x$ 上的纤维，则
$$
\dim(\mathcal{O}_{F, x'}) =
\dim(\mathcal{O}_{X'_{s'}, x'}) - \dim(\mathcal{O}_{X_s, x}) =
\text{trdeg}_{\kappa(s)}(\kappa(x)) -
\text{trdeg}_{\kappa(s')}(\kappa(x'))
$$
有 $\dim(\mathcal{O}_{X'_{s'}, x'}) \geq \dim(\mathcal{O}_{X_s, x})$，
且 $\text{trdeg}_{\kappa(s)}(\kappa(x)) \geq
\text{trdeg}_{\kappa(s')}(\kappa(x'))$。
\item 给定 $s', s, x$，存在一种对 $x'$ 的选择，使得
$\dim(\mathcal{O}_{X'_{s'}, x'}) = \dim(\mathcal{O}_{X_s, x})$，并且
$\text{trdeg}_{\kappa(s)}(\kappa(x)) = \text{trdeg}_{\kappa(s')}(\kappa(x'))$。
\end{enumerate}
\end{lemma}

\begin{proof}
(1) 立即来自《代数》引理
\ref{algebra-lemma-dimension-at-a-point-preserved-field-extension}。
(2) 和 (3) 来自《代数》引理
\ref{algebra-lemma-inequalities-under-field-extension}。
\end{proof}

\noindent
下一个引理来自一个非平凡的代数结果，即 Zariski 主定理的代数版本。

\begin{lemma}
\label{morphisms-lemma-openness-bounded-dimension-fibres}
\begin{reference}
\cite[IV Theorem 13.1.3]{EGA}
\end{reference}
令 $f : X \to S$ 为概形态射。令 $n \geq 0$。假设 $f$ 局部有限型。
集合
$$
U_n = \{x \in X \mid \dim_x X_{f(x)} \leq n\}
$$
在 $X$ 中是开集。
\end{lemma}

\begin{proof}
这立即来自《代数》引理
\ref{algebra-lemma-dimension-fibres-bounded-open-upstairs}
\end{proof}

\begin{lemma}
\label{morphisms-lemma-morphism-finite-type-bounded-dimension}
令 $f : X \to Y$ 为有限型态射，并假设 $Y$ 拟紧。则 $f$ 的纤维维数有界。
\end{lemma}

\begin{proof}
由引理 \ref{morphisms-lemma-openness-bounded-dimension-fibres}，集合
$U_n \subset X$ 中各点的纤维维数均为 $\leq n$，且该集合开。由于 $f$ 有限型，
每一点都包含在某个 $U_n$ 中（因为域上有限型代数的维数有限）。
由于 $Y$ 拟紧且 $f$ 有限型，可知 $X$ 拟紧。因此
$X = U_n$ 对某个 $n$ 成立。
\end{proof}

\begin{lemma}
\label{morphisms-lemma-openness-bounded-dimension-fibres-finite-presentation}
令 $f : X \to S$ 为概形态射。令 $n \geq 0$。假设 $f$ 局部有限表示。
引理 \ref{morphisms-lemma-openness-bounded-dimension-fibres} 中的开集
$$
U_n = \{x \in X \mid \dim_x X_{f(x)} \leq n\}
$$
在 $X$ 中逆拟紧。（见《拓扑学》定义
\ref{topology-definition-quasi-compact}。）
\end{lemma}

\begin{proof}
拓扑空间 $X$ 有一个由如下仿射开集组成的拓扑基：$U \subset X$，
且所诱导的态射 $f|_U : U \to S$ 通过某个仿射开集
$V \subset S$ 分解。因此只需证明 $U \cap U_n$ 拟紧，其中 $U$
是这样的开集。注意，$U_n \cap U$ 与开集
$\{x \in U \mid \dim_x U_{f(x)} \leq n\}$ 相同。
这把问题归结到 $X$ 和 $S$ 都是仿射概形的情形。在这种情形下，
本引理来自《代数》引理
\ref{algebra-lemma-dimension-fibres-bounded-quasi-compact-open-upstairs}
（以及引理 \ref{morphisms-lemma-locally-finite-presentation-characterize}）。
\end{proof}

\begin{lemma}
\label{morphisms-lemma-dimension-fibre-specialization}
令 $f : X \to S$ 为概形态射。令 $x \leadsto x'$ 为 $X$ 中点的一次
非平凡特化，并假设二者位于同一点 $s \in S$ 上。假设 $f$ 局部有限型。
则：
\begin{enumerate}
\item $\dim_x(X_s) \leq \dim_{x'}(X_s)$；
\item $\dim(\mathcal{O}_{X_s, x}) < \dim(\mathcal{O}_{X_s, x'})$；并且
\item $\text{trdeg}_{\kappa(s)}(\kappa(x)) >
\text{trdeg}_{\kappa(s)}(\kappa(x'))$。
\end{enumerate}
\end{lemma}

\begin{proof}
(1) 来自如下事实：$X_s$ 的任何含 $x'$ 的开集也含 $x$。
(2) 来自如下事实：$\mathcal{O}_{X_s, x}$ 是
$\mathcal{O}_{X_s, x'}$ 在一个素理想处的局部化，因此
$\mathcal{O}_{X_s, x}$ 中的任何素理想链都是
$\mathcal{O}_{X_s, x'}$ 中一条严格更长的素理想链的一部分。
最后一个不等式来自《代数》引理
\ref{algebra-lemma-tr-deg-specialization}。
\end{proof}






\section{指定相对维数的态射}
\label{morphisms-section-relative-dimension}

\noindent
为了方便地谈论指定相对维数的态射，引入下述概念。

\begin{definition}
\label{morphisms-definition-relative-dimension-d}
令 $f : X \to S$ 为概形态射。假设 $f$ 局部有限型。
\begin{enumerate}
\item 称 $f$ 具有{\it 相对维数 $\leq d$}（在 $x$ 处），如果
$\dim_x(X_{f(x)}) \leq d$。
\item 称 $f$ 具有{\it 相对维数 $\leq d$}，如果
$\dim_x(X_{f(x)}) \leq d$ 对所有 $x \in X$ 成立。
\item 称 $f$ 具有{\it 相对维数 $d$}，如果所有非空纤维 $X_s$
都等维且维数为 $d$。
\end{enumerate}
\end{definition}

\noindent
这个概念并非特别良态，但在许多情形下很好用。

\begin{lemma}
\label{morphisms-lemma-base-change-relative-dimension-d}
令 $f : X \to S$ 为局部有限型概形态射。若 $f$ 相对维数为 $d$，
则 $f$ 的任意基变换也如此。相对维数 $\leq d$ 的情形同样成立。
\end{lemma}

\begin{proof}
这立即来自引理 \ref{morphisms-lemma-dimension-fibre-after-base-change}。
\end{proof}

\begin{lemma}
\label{morphisms-lemma-composition-relative-dimension-d}
令 $f : X \to Y$、$g : Y \to Z$ 局部有限型。若 $f$ 相对维数
$\leq d$，且 $g$ 相对维数 $\leq e$，则 $g \circ f$ 相对维数
$\leq d + e$。若
\begin{enumerate}
\item $f$ 相对维数为 $d$；
\item $g$ 相对维数为 $e$；并且
\item $f$ 平坦，
\end{enumerate}
则 $g \circ f$ 相对维数为 $d + e$。
\end{lemma}

\begin{proof}
这立即来自引理 \ref{morphisms-lemma-dimension-fibre-at-a-point-additive}。
\end{proof}

\noindent
一般而言，不能把一个态射分解成相对维数分别为指定数值的各部分。
不过，当该态射的纤维是 Cohen--Macaulay 的并且该态射平坦且有限表示时，
可以作这种分解。

\begin{lemma}
\label{morphisms-lemma-flat-finite-presentation-CM-fibres-relative-dimension}
\begin{slogan}
Cohen--Macaulay 态射分解为纯相对维数的开闭部分
\end{slogan}
令 $f : X \to S$ 为概形态射。假设：
\begin{enumerate}
\item $f$ 平坦；
\item $f$ 局部有限表示；并且
\item 对所有 $s \in S$，纤维 $X_s$ 都是 Cohen--Macaulay 的
（《性质》定义 \ref{properties-definition-Cohen-Macaulay}）。
\end{enumerate}
则存在开闭子概形 $X_d \subset X$，使得
$X = \coprod_{d \geq 0} X_d$，并且 $f|_{X_d} : X_d \to S$
相对维数为 $d$。
\end{lemma}

\begin{proof}
这立即来自《代数》引理
\ref{algebra-lemma-relative-dimension-CM}。
\end{proof}

\begin{lemma}
\label{morphisms-lemma-locally-quasi-finite-rel-dimension-0}
令 $f : X \to S$ 为概形态射。假设 $f$ 局部有限型。令 $x \in X$，
且 $s = f(x)$。则 $f$ 在 $x$ 处拟有限，当且仅当
$\dim_x(X_s) = 0$。特别地，$f$ 局部拟有限，当且仅当 $f$ 相对维数为
$0$。
\end{lemma}

\begin{proof}
第一种证明。若 $f$ 在 $x$ 处拟有限，则 $\kappa(x)$ 是
$\kappa(s)$ 的有限扩张（由引理
\ref{morphisms-lemma-residue-field-quasi-finite}），且 $x$ 在 $X_s$ 中孤立
（由引理 \ref{morphisms-lemma-quasi-finite-at-point-characterize}），故由引理
\ref{morphisms-lemma-dimension-fibre-at-a-point} 有 $\dim_x(X_s) = 0$。
反之，若 $\dim_x(X_s) = 0$，则由引理
\ref{morphisms-lemma-dimension-fibre-at-a-point} 可知 $\kappa(s) \subset \kappa(x)$
是代数扩张，并且 $X_s$ 中没有其他点特化到 $x$。因此由引理
\ref{morphisms-lemma-algebraic-residue-field-extension-closed-point-fibre}，$x$
在其纤维中闭；再由引理
\ref{morphisms-lemma-quasi-finite-at-point-characterize} (3)，可得 $f$ 在 $x$
处拟有限。

\medskip\noindent
第二种证明。纤维 $X_s$ 是域上的局部有限型概形，因而局部 Noether
（引理 \ref{morphisms-lemma-finite-type-noetherian}）。结论现在来自引理
\ref{morphisms-lemma-quasi-finite-at-point-characterize} 和《性质》引理
\ref{properties-lemma-isolated-dim-0-locally-noetherian}。
\end{proof}

\begin{lemma}
\label{morphisms-lemma-rel-dimension-dimension}
令 $f : X \to Y$ 为局部 Noether 概形之间的态射，并假设它平坦、
局部有限型且相对维数为 $d$。对每一点 $x$（它属于 $X$），若其像为
$y$（位于 $Y$ 中），则有 $\dim_x(X) = \dim_y(Y) + d$。
\end{lemma}

\begin{proof}
把 $X$ 和 $Y$ 分别缩小到 $x$ 和 $y$ 的开邻域后，由概形在一点处的
维数的定义（《性质》定义 \ref{properties-definition-dimension}），
可以假设 $\dim(X) = \dim_x(X)$ 且 $\dim(Y) = \dim_y(Y)$。
由引理 \ref{morphisms-lemma-noetherian-finite-type-finite-presentation} 和
\ref{morphisms-lemma-fppf-open}，态射 $f$ 是开态射。因此可以缩小 $Y$，使得
$f$ 满。还需证明 $\dim(X) = \dim(Y) + d$。

\medskip\noindent
令 $a$ 为 $X$ 中一点，其像为 $b$，位于 $Y$ 中。由《代数》引理
\ref{algebra-lemma-dimension-base-fibre-equals-total}，
$$
\dim(\mathcal{O}_{X,a}) = \dim(\mathcal{O}_{Y,b}) + \dim(\mathcal{O}_{X_b, a}).
$$
对所有点 $a$（属于 $X$）取上确界，可得
$\dim(X) = \dim(Y) + d$，正如所需；见《性质》引理
\ref{properties-lemma-dimension}。
\end{proof}









\section{Syntomic 态射}
\label{morphisms-section-syntomic}

\noindent
$A$ 是域 $k$ 上的代数；称它为 $k$ 上的{\it 全局完全交}，如果
$A \cong k[x_1, \ldots, x_n]/(f_1, \ldots, f_c)$ 且
$\dim(A) = n - c$。$A$ 是域 $k$ 上的代数；称它为{\it 局部完全交}，
如果 $\Spec(A)$ 可以由标准开集覆盖，且每个标准开集都是 $k$
上的全局完全交。见《代数》节 \ref{algebra-section-lci}。
回忆，环同态 $R \to A$ 称为 {\it syntomic}，如果它有限表示、
平坦且其纤维为局部完全交环；见《代数》定义
\ref{algebra-definition-lci}。

\begin{definition}
\label{morphisms-definition-syntomic}
令 $f : X \to S$ 为概形态射。
\begin{enumerate}
\item 称 $f$ {\it 在 $x \in X$ 处 syntomic}，如果存在仿射开邻域
$\Spec(A) = U \subset X$（它是 $x$ 的邻域）以及仿射开集
$\Spec(R) = V \subset S$，使得 $f(U) \subset V$，且所诱导的环同态
$R \to A$ 是 syntomic 的。
\item 称 $f$ {\it syntomic}，如果它在 $X$ 的每一点处 syntomic。
\item 若 $S = \Spec(k)$ 且 $f$ syntomic，则称 $X$ 是
{\it $k$ 上的局部完全交}。
\item 仿射概形态射 $f : X \to S$ 称为{\it 标准 syntomic}，
如果存在全局相对完全交
$R \to R[x_1, \ldots, x_n]/(f_1, \ldots, f_c)$（见《代数》定义
\ref{algebra-definition-relative-global-complete-intersection}），使得
$X \to S$ 同构于
$$
\Spec(R[x_1, \ldots, x_n]/(f_1, \ldots, f_c)) \to \Spec(R).
$$
\end{enumerate}
\end{definition}

\noindent
文献中有时把 syntomic 态射称为{\it 平坦局部完全交态射}。
这是一类方便的态射。例如，可以用它们定义 syntomic 拓扑；
该拓扑比光滑拓扑和 \'etale 拓扑更细，但具有许多相同的形式性质。

\medskip\noindent
全局相对完全交（我们曾用它定义标准 syntomic 环同态）特别是平坦的。
在《态射进阶》节 \ref{more-morphisms-section-lci} 中，我们将考虑局部形如
$X \to S$ 的态射
$$
\Spec(R[x_1, \ldots, x_n]/(f_1, \ldots, f_c)) \to \Spec(R).
$$
其中 $f_1, \ldots, f_r$ 是 $R[x_1, \ldots, x_n]$ 中某个
Koszul--正则序列。这样的态射将称为{\it 局部完全交态射}。
一旦给出这个定义，态射 syntomic 当且仅当它是平坦的局部完全交态射。

\medskip\noindent
注意，syntomic 态射的定义中没有分离性或拟紧性假设。因此，
syntomic 是源上的局部性质。精确结果如下。

\begin{lemma}
\label{morphisms-lemma-syntomic-characterize}
令 $f : X \to S$ 为概形态射。下列条件等价：
\begin{enumerate}
\item 态射 $f$ syntomic。
\item 对每对仿射开集 $U \subset X$、$V \subset S$，若
$f(U) \subset V$，则环同态
$\mathcal{O}_S(V) \to \mathcal{O}_X(U)$ syntomic。
\item 存在开覆盖 $S = \bigcup_{j \in J} V_j$ 以及开覆盖
$f^{-1}(V_j) = \bigcup_{i \in I_j} U_i$，使得每个态射
$U_i \to V_j$ 都 syntomic，其中 $j\in J, i\in I_j$。
\item 存在仿射开覆盖 $S = \bigcup_{j \in J} V_j$ 以及仿射开覆盖
$f^{-1}(V_j) = \bigcup_{i \in I_j} U_i$，使得环同态
$\mathcal{O}_S(V_j) \to \mathcal{O}_X(U_i)$ syntomic，其中所有
$j\in J, i\in I_j$ 均如此。
\end{enumerate}
此外，若 $f$ syntomic，则对任意开子概形 $U \subset X$、$V \subset S$，
若 $f(U) \subset V$，限制 $f|_U : U \to V$ syntomic。
\end{lemma}

\begin{proof}
若能证明性质“$R \to A$ 是 syntomic 的”是局部的，结论就来自引理
\ref{morphisms-lemma-locally-P}。我们核验定义
\ref{morphisms-definition-property-local} 的条件 (a)、(b) 和 (c)。
由《代数》引理 \ref{algebra-lemma-base-change-syntomic}，
syntomic 在基变换下稳定，故条件 (a) 成立。由《代数》引理
\ref{algebra-lemma-composition-syntomic}，syntomic 在复合下稳定；
并且显然地，对任意环 $R$，环同态 $R \to R_f$ syntomic。
所以条件 (b) 成立。最后，由《代数》引理
\ref{algebra-lemma-local-syntomic}，条件 (c) 成立。
\end{proof}

\begin{lemma}
\label{morphisms-lemma-composition-syntomic}
两个 syntomic 态射的复合仍 syntomic。
\end{lemma}

\begin{proof}
在引理 \ref{morphisms-lemma-syntomic-characterize} 的证明中，我们看到
syntomic 是环同态的局部性质。因此，本引理的第一个断言来自引理
\ref{morphisms-lemma-composition-type-P}，并结合 syntomic 是在复合下稳定的
环同态性质这一事实；见《代数》引理
\ref{algebra-lemma-composition-syntomic}。
\end{proof}

\begin{lemma}
\label{morphisms-lemma-base-change-syntomic}
syntomic 态射的基变换仍 syntomic。
\end{lemma}

\begin{proof}
在引理 \ref{morphisms-lemma-syntomic-characterize} 的证明中，我们看到
syntomic 是环同态的局部性质。因此，本引理来自引理
\ref{morphisms-lemma-composition-type-P}，并结合 syntomic 是在基变换下稳定的
环同态性质这一事实；见《代数》引理
\ref{algebra-lemma-base-change-syntomic}。
\end{proof}

\begin{lemma}
\label{morphisms-lemma-open-immersion-syntomic}
任意开浸入都 syntomic。
\end{lemma}

\begin{proof}
这是因为开浸入是局部同构。
\end{proof}

\begin{lemma}
\label{morphisms-lemma-syntomic-locally-finite-presentation}
syntomic 态射局部有限表示。
\end{lemma}

\begin{proof}
这是因为按照定义，syntomic 环同态有限表示。
\end{proof}

\begin{lemma}
\label{morphisms-lemma-syntomic-flat}
syntomic 态射平坦。
\end{lemma}

\begin{proof}
这是因为按照定义，syntomic 环同态平坦。
\end{proof}

\begin{lemma}
\label{morphisms-lemma-syntomic-open}
syntomic 态射是泛开态射。
\end{lemma}

\begin{proof}
结合引理 \ref{morphisms-lemma-syntomic-locally-finite-presentation}、
\ref{morphisms-lemma-syntomic-flat} 和
\ref{morphisms-lemma-fppf-open}。
\end{proof}

\noindent
令 $k$ 为域。令 $A$ 为本质有限型的局部 $k$-代数（在 $k$ 上）。
回忆，$A$ 称为 $k$ 上的{\it 完全交}，如果可写成
$A \cong R/(f_1, \ldots, f_c)$，其中 $R$ 是本质有限型的正则局部环
（在 $k$ 上），且 $f_1, \ldots, f_c$ 是 $R$ 中的正则序列；
见《代数》定义 \ref{algebra-definition-lci-local-ring}。

\begin{lemma}
\label{morphisms-lemma-local-complete-intersection}
令 $k$ 为域。令 $X$ 为 $k$ 上的局部有限型概形。下列条件等价：
\begin{enumerate}
\item $X$ 是 $k$ 上的局部完全交；
\item 对每个 $x \in X$，存在仿射开集
$U = \Spec(R) \subset X$，它是 $x$ 的邻域，并且
$R \cong k[x_1, \ldots, x_n]/(f_1, \ldots, f_c)$ 是 $k$
上的全局完全交；并且
\item 对每个 $x \in X$，局部环 $\mathcal{O}_{X, x}$ 是 $k$
上的完全交。
\end{enumerate}
\end{lemma}

\begin{proof}
相应的代数结果见《代数》引理
\ref{algebra-lemma-lci-at-prime} 和
\ref{algebra-lemma-lci-global}。
\end{proof}

\noindent
下一个引理表明，每个 syntomic 态射在局部都是标准 syntomic 的。
因此可以用标准 syntomic 态射作为 syntomic 态射的{\it 局部模型}。
此外，它还表明，有限表示平坦态射 syntomic，当且仅当其纤维为局部完全交概形。

\begin{lemma}
\label{morphisms-lemma-syntomic-locally-standard-syntomic}
令 $f : X  \to S$ 为概形态射。令 $x \in X$ 为一点，其像为
$s = f(x)$。令 $V \subset S$ 为 $s$ 的仿射开邻域。下列条件等价：
\begin{enumerate}
\item 态射 $f$ 在 $x$ 处 syntomic。
\item 存在仿射开集 $U \subset X$，使得 $x \in U$、
$f(U) \subset V$，并且 $f|_U : U \to V$ 标准 syntomic。
\item 态射 $f$ 在 $x$ 处有限表示，局部环同态
$\mathcal{O}_{S, s} \to \mathcal{O}_{X, x}$ 平坦，并且
$\mathcal{O}_{X, x}/\mathfrak m_s \mathcal{O}_{X, x}$ 是
$\kappa(s)$ 上的完全交（见《代数》定义
\ref{algebra-definition-lci-local-ring}）。
\end{enumerate}
\end{lemma}

\begin{proof}
这来自各定义以及《代数》引理
\ref{algebra-lemma-syntomic}。
\end{proof}

\begin{lemma}
\label{morphisms-lemma-syntomic-flat-fibres}
令 $f : X \to S$ 为概形态射。若 $f$ 平坦、局部有限表示，且所有纤维
$X_s$ 都是局部完全交，则 $f$ syntomic。
\end{lemma}

\begin{proof}
由引理 \ref{morphisms-lemma-local-complete-intersection} 和
\ref{morphisms-lemma-syntomic-locally-standard-syntomic}，以及局部环同构
$
\mathcal{O}_{X, x}/\mathfrak m_s \mathcal{O}_{X, x}
\cong
\mathcal{O}_{X_s, x}
$，结论显然。
\end{proof}

\begin{lemma}
\label{morphisms-lemma-set-points-where-fibres-lci}
令 $f : X \to S$ 为概形态射。假设 $f$ 局部有限型。集合
$$
T = \{x \in X \mid \mathcal{O}_{X_{f(x)}, x}
\text{ is a complete intersection over }\kappa(f(x))\}
$$
的形成与任意基变换可交换：对任意态射 $g : S' \to S$，考虑基变换
$f' : X' \to S'$（它来自 $f$）以及投影 $g' : X' \to X$。
则集合 $T'$ 对应于态射 $f'$，且等于 $T' = (g')^{-1}(T)$。
特别地，若假设 $f$ 平坦且局部有限表示，则 $f$ 为 syntomic 的点所成
开集也有同一性质。
\end{lemma}

\begin{proof}
令 $s' \in S'$ 为一点，并令 $s = g(s')$。则有
$$
X'_{s'} =
\Spec(\kappa(s')) \times_{\Spec(\kappa(s))} X_s
$$
换言之，基变换的纤维就是纤维的基变换。因此第一部分等价于《代数》引理
\ref{algebra-lemma-lci-field-change-local}。第二部分来自第一部分，
因为在这种情形下，依引理
\ref{morphisms-lemma-syntomic-locally-standard-syntomic}，$T$ 正是 $f$
为 syntomic 的点所成集合。
\end{proof}

\begin{lemma}
\label{morphisms-lemma-standard-syntomic-relative-dimension}
令 $R$ 为环。令
$R \to A = R[x_1, \ldots, x_n]/(f_1, \ldots, f_c)$ 为全局相对完全交。
置 $S = \Spec(R)$ 且 $X = \Spec(A)$。考虑态射
$f : X \to S$，它与环同态 $R \to A$ 相伴。函数
$x \mapsto \dim_x(X_{f(x)})$ 是取值 $n - c$ 的常值函数。
\end{lemma}

\begin{proof}
由《代数》定义
\ref{algebra-definition-relative-global-complete-intersection}，
$R \to A$ 为全局相对完全交意味着所有非零纤维环的维数都是
$n - c$。因此，对素理想 $\mathfrak p$（属于 $R$），纤维环
$\kappa(\mathfrak p)[x_1, \ldots, x_n]/(\overline{f}_1, \ldots, \overline{f}_c)$
或者为零，或者是维数为 $n - c$ 的全局完全交环。由《代数》定义
\ref{algebra-definition-lci-field} 后面的讨论，这推出它等维且维数为
$n - c$。本引理得证。
\end{proof}

\begin{lemma}
\label{morphisms-lemma-syntomic-relative-dimension}
令 $f : X \to S$ 为 syntomic 态射。函数
$x \mapsto \dim_x(X_{f(x)})$ 在 $X$ 上局部常值。
\end{lemma}

\begin{proof}
由引理 \ref{morphisms-lemma-syntomic-locally-standard-syntomic}，态射 $f$
局部看起来像仿射概形之间的标准 syntomic 态射。因此结论来自引理
\ref{morphisms-lemma-standard-syntomic-relative-dimension}。
\end{proof}

\noindent
引理 \ref{morphisms-lemma-syntomic-relative-dimension} 表明下述定义是有意义的。

\begin{definition}
\label{morphisms-definition-syntomic-relative-dimension}
令 $d \geq 0$ 为整数。概形态射 $f : X \to S$ 被称为
{\it 相对维数 $d$ 的 syntomic 态射}，是指 $f$ syntomic，且函数满足
$\dim_x(X_{f(x)}) = d$，其中该等式对所有 $x \in X$ 成立。
\end{definition}

\noindent
换言之，$f$ syntomic，且非空纤维等维且维数为 $d$。

\begin{lemma}
\label{morphisms-lemma-syntomic-permanence}
令
$$
\xymatrix{
X \ar[rr]_f \ar[rd]_p & &
Y \ar[dl]^q \\
& S
}
$$
为概形态射的交换图。假设：
\begin{enumerate}
\item $f$ 满且 syntomic；
\item $p$ syntomic；并且
\item $q$ 局部有限表示\footnote{事实上，若 $f$ 满、平坦且局部有限表示，
并且 $p$ syntomic，则 $q$ 和 $f$ 都 syntomic；见《下降》引理
\ref{descent-lemma-syntomic-permanence}。}。
\end{enumerate}
则 $q$ syntomic。
\end{lemma}

\begin{proof}
由引理 \ref{morphisms-lemma-flat-permanence}，可知 $q$ 平坦。因此只需证明
$Y \to S$ 的纤维都是局部完全交；见引理
\ref{morphisms-lemma-syntomic-flat-fibres}。令 $s \in S$。考虑态射
$X_s \to Y_s$。这是态射 $X \to Y$ 的基变换，因而满且 syntomic
（引理 \ref{morphisms-lemma-base-change-syntomic}）。出于同一理由，$X_s$
在 $\kappa(s)$ 上 syntomic。此外，$Y_s$ 在 $\kappa(s)$ 上局部有限型
（引理 \ref{morphisms-lemma-base-change-finite-type}）。这样，问题归结到 $S$
是一个域的谱的情形。

\medskip\noindent
假设 $S = \Spec(k)$。令 $y \in Y$。选取仿射开集
$\Spec(A) \subset Y$，作为 $y$ 的邻域。令 $\Spec(B) \subset X$ 为仿射开集，使得
$f(\Spec(B)) \subset \Spec(A)$，并且包含某点 $x \in X$，满足
$f(x) = y$。选取满射 $k[x_1, \ldots, x_n] \to A$，其核为 $I$。
选取满射 $A[y_1, \ldots, y_m] \to B$；它又诱导满射
$k[x_i, y_j] \to B$，其核为 $J$。令
$\mathfrak q \subset k[x_i, y_j]$ 为对应于 $y \in \Spec(B)$ 的素理想，
并令 $\mathfrak p \subset k[x_i]$ 为对应于 $x \in \Spec(A)$ 的素理想。
由于 $x$ 映到 $y$，有 $\mathfrak p = \mathfrak q \cap k[x_i]$。
考虑下列局部环交换图：
$$
\xymatrix{
\mathcal{O}_{X, x} \ar@{=}[r] &
B_{\mathfrak q} &
k[x_1, \ldots, x_n, y_1, \ldots, y_m]_{\mathfrak q} \ar[l] \\
\mathcal{O}_{Y, y} \ar@{=}[r] \ar[u] & A_{\mathfrak p} \ar[u] &
k[x_1, \ldots, x_n]_{\mathfrak p} \ar[l] \ar[u]
}
$$
我们断言《代数》引理
\ref{algebra-lemma-lci-permanence-initial} 的假设均满足。条件 (1) 和 (2)
显然。由于 $X \to Y$ 平坦，条件 (4) 成立。条件 (3) 成立，
因为环 $\mathcal{O}_{Y, y}$ 和
$\mathcal{O}_{X_y, x} = \mathcal{O}_{X, x}/\mathfrak m_y\mathcal{O}_{X, x}$
由 $f$ 和 $p$ syntomic 这一假设都是完全交环；见引理
\ref{morphisms-lemma-syntomic-locally-standard-syntomic}。《代数》引理
\ref{algebra-lemma-lci-permanence-initial} 的结论恰好是
$\mathcal{O}_{Y, y}$ 为完全交环！所以再次由引理
\ref{morphisms-lemma-syntomic-locally-standard-syntomic}，可知 $Y$ 在 $k$
上于 $y$ 处 syntomic，正如所需。
\end{proof}

\section{浸入的余法层}
\label{morphisms-section-conormal-sheaf}

\noindent
令 $i : Z \to X$ 为闭浸入。令
$\mathcal{I} \subset \mathcal{O}_X$ 为相应的拟凝聚理想层。考虑 $X$
上拟凝聚层的短正合列
$$
0 \to \mathcal{I}^2 \to \mathcal{I} \to \mathcal{I}/\mathcal{I}^2 \to 0
$$
。由于层 $\mathcal{I}/\mathcal{I}^2$ 被 $\mathcal{I}$ 零化，故由引理
\ref{morphisms-lemma-i-star-equivalence}，它对应于 $Z$ 上的一个层。这个拟凝聚
$\mathcal{O}_Z$-模称为 $Z$ 在 $X$ 中的{\it 余法层}；按照
第 \ref{morphisms-section-closed-immersions-quasi-coherent} 节所述的记号滥用，
常将其简记为 $\mathcal{I}/\mathcal{I}^2$。

\medskip\noindent
若 $i : Z \to X$ 是（局部闭）浸入，则把 $i$ 的余法层定义为闭浸入
$i : Z \to X \setminus \partial Z$ 的余法层，其中
$\partial Z = \overline{Z} \setminus Z$。它常记作
$\mathcal{I}/\mathcal{I}^2$，这里 $\mathcal{I}$ 是闭浸入
$i : Z \to X \setminus \partial Z$ 的理想层。

\begin{definition}
\label{morphisms-definition-conormal-sheaf}
令 $i : Z \to X$ 为浸入。称 $\mathcal{C}_{Z/X}$ 为 $Z$ 在 $X$
中的{\it 余法层}，也称为 $i$ 的{\it 余法层}；它是上述拟凝聚
$\mathcal{O}_Z$-模 $\mathcal{I}/\mathcal{I}^2$。
\end{definition}

\noindent
在 \cite[IV Definition 16.1.2]{EGA} 中，这个层记作
$\mathcal{N}_{Z/X}$。我们不采用这一约定，因为希望保留记号
$\mathcal{N}_{Z/X}$ 表示{\it 该浸入的法层}。当余法层有限表示时，
它定义为
$$
\mathcal{N}_{Z/X} =
\SheafHom_{\mathcal{O}_Z}(\mathcal{C}_{Z/X}, \mathcal{O}_Z) =
\SheafHom_{\mathcal{O}_Z}(\mathcal{I}/\mathcal{I}^2, \mathcal{O}_Z)
$$
（否则法层甚至未必拟凝聚）。稍后我们会回到法层（此处待插未来引用）。

\begin{lemma}
\label{morphisms-lemma-affine-conormal}
令 $i : Z \to X$ 为浸入。$i$ 的余法层具有下列性质：
\begin{enumerate}
\item 令 $U \subset X$ 为任意开子概形，使 $i$ 可分解为
$Z \xrightarrow{i'} U \to X$，其中 $i'$ 是闭浸入。令
$\mathcal{I} = \Ker((i')^\sharp) \subset \mathcal{O}_U$。则
$$
\mathcal{C}_{Z/X} = (i')^*\mathcal{I}\quad\text{且}\quad
i'_*\mathcal{C}_{Z/X} = \mathcal{I}/\mathcal{I}^2
$$
\item 对任意仿射开集 $\Spec(R) = U \subset X$，若
$Z \cap U = \Spec(R/I)$，则有典范同构
$\Gamma(Z \cap U, \mathcal{C}_{Z/X}) = I/I^2$。
\end{enumerate}
\end{lemma}

\begin{proof}
大体上由定义即得。注意，对环 $R$ 及其理想 $I$（属于 $R$），有
$I/I^2 = I \otimes_R R/I$。细节从略。
\end{proof}

\begin{lemma}
\label{morphisms-lemma-conormal-functorial}
令
$$
\xymatrix{
Z \ar[r]_i \ar[d]_f & X \ar[d]^g \\
Z' \ar[r]^{i'} & X'
}
$$
为概形范畴中的交换图。假设 $i$、$i'$ 均为浸入。存在典范的
$\mathcal{O}_Z$-模同态
$$
f^*\mathcal{C}_{Z'/X'}
\longrightarrow
\mathcal{C}_{Z/X}
$$
，它由下述性质刻画：对每一对仿射开集
$(\Spec(R) = U \subset X, \Spec(R') = U' \subset X')$，若
$g(U) \subset U'$，且
$Z \cap U = \Spec(R/I)$ 和 $Z' \cap U' = \Spec(R'/I')$，则诱导同态
$$
\Gamma(Z' \cap U', \mathcal{C}_{Z'/X'}) = I'/I'^2
\longrightarrow
I/I^2 = \Gamma(Z \cap U, \mathcal{C}_{Z/X})
$$
由环同态 $g^\sharp : R' \to R$ 诱导；该环同态满足
$g^\sharp(I') \subset I$。
\end{lemma}

\begin{proof}
令 $\partial Z' = \overline{Z'} \setminus Z'$，并令
$\partial Z = \overline{Z} \setminus Z$。它们分别是 $X'$ 和 $X$
的闭子集。把 $X'$ 替换为 $X' \setminus \partial Z'$，并把 $X$
替换为 $X \setminus \big(g^{-1}(\partial Z') \cup \partial Z\big)$，
便可假设 $i$ 和 $i'$ 都是闭浸入。

\medskip\noindent
$g \circ i$ 分解通过 $i'$ 这一事实蕴含：$g^*\mathcal{I}'$ 映入
$\mathcal{I}$；这里所用的典范同态为 $g^*\mathcal{I}' \to \mathcal{O}_X$。
见《概形》引理
\ref{schemes-lemma-characterize-closed-subspace} 和
\ref{schemes-lemma-restrict-map-to-closed}。因此得到拟凝聚层的诱导同态
$g^*(\mathcal{I}'/(\mathcal{I}')^2) \to \mathcal{I}/\mathcal{I}^2$。
再由 $i$ 拉回，得到
$i^*g^*(\mathcal{I}'/(\mathcal{I}')^2) \to i^*(\mathcal{I}/\mathcal{I}^2)$。
注意 $i^*(\mathcal{I}/\mathcal{I}^2) = \mathcal{C}_{Z/X}$。另一方面，
$i^*g^*(\mathcal{I}'/(\mathcal{I}')^2) =
f^*(i')^*(\mathcal{I}'/(\mathcal{I}')^2) = f^*\mathcal{C}_{Z'/X'}$。
这就给出了所需同态。

\medskip\noindent
展开定义即可验证该同态在局部上由给定同态
$I'/(I')^2 \to I/I^2$ 描述，细节从略。还要注意，对任意
$x \in i(Z)$，确实存在仿射开邻域 $U$、$U'$，使得
$f(U) \subset U'$，并使 $Z \cap U$ 和 $U' \cap Z'$ 都闭，且
$x \in U$。证明从略。因此，本引理的要求确实刻画了该同态
（也可以用来定义它）。
\end{proof}

\begin{lemma}
\label{morphisms-lemma-conormal-functorial-flat}
令
$$
\xymatrix{
Z \ar[r]_i \ar[d]_f & X \ar[d]^g \\
Z' \ar[r]^{i'} & X'
}
$$
为概形范畴中的纤维积图，其中 $i$、$i'$ 为浸入。则引理
\ref{morphisms-lemma-conormal-functorial} 的典范同态
$f^*\mathcal{C}_{Z'/X'} \to \mathcal{C}_{Z/X}$ 是满射。若 $g$ 平坦，
则它是同构。
\end{lemma}

\begin{proof}
令 $R' \to R$ 为环同态，$I' \subset R'$ 为理想。置 $I = I'R$。
则 $I'/(I')^2 \otimes_{R'} R \to I/I^2$ 是满射。若
$R' \to R$ 平坦，则 $I = I' \otimes_{R'} R$ 且
$I^2 = (I')^2 \otimes_{R'} R$，从而可见该同态为同构。
\end{proof}

\begin{lemma}
\label{morphisms-lemma-transitivity-conormal}
令 $Z \to Y \to X$ 为概形的浸入。则有典范正合列
$$
i^*\mathcal{C}_{Y/X} \to
\mathcal{C}_{Z/X} \to
\mathcal{C}_{Z/Y} \to 0
$$
，其中同态来自引理 \ref{morphisms-lemma-conormal-functorial}，而
$i : Z \to Y$ 是第一个态射。
\end{lemma}

\begin{proof}
通过引理 \ref{morphisms-lemma-conormal-functorial}，这转化为下述代数事实。
设 $C \to B \to A$ 为满环同态。令 $I = \Ker(B \to A)$、
$J = \Ker(C \to A)$，并令 $K = \Ker(C \to B)$。则有正合列
$$
K/K^2 \otimes_B A \to J/J^2 \to I/I^2 \to 0.
$$
这立即由 $I = J/K$ 得出。
\end{proof}










\section{态射的微分层}
\label{morphisms-section-sheaf-differentials}

\noindent
我们建议读者参阅交换代数章中的相应一节
（《代数》第 \ref{algebra-section-differentials} 节），以及模层章中的相应一节
（《模》第 \ref{modules-section-differentials} 节）。

\begin{definition}
\label{morphisms-definition-sheaf-differentials}
令 $f : X \to S$ 为概形态射。{\it 微分层 $\Omega_{X/S}$}，即 $X$ 在
$S$ 上的微分层，是把 $f$ 视为环层空间态射时的微分层
（《模》定义 \ref{modules-definition-differentials}），并带有其{\it 泛
$S$-导子}
$$
\text{d}_{X/S} : \mathcal{O}_X \longrightarrow \Omega_{X/S}.
$$
\end{definition}

\noindent
$\Omega_{X/S}$ 是拟凝聚 $\mathcal{O}_X$-模。例如，它同构于对角态射
$\Delta : X \to X \times_S X$ 的余法层（引理
\ref{morphisms-lemma-differentials-diagonal}）。我们已通过泛性质定义 $X$ 在 $S$
上的微分模，即把它定义为泛导子的取值对象。若对相对微分层另有任何
满足这一泛性质的构造，则由 Yoneda 引理，它与上述构造典范同构。
为方便起见，我们在此重述这一泛性质。

\begin{lemma}
\label{morphisms-lemma-universal-derivation-universal}
令 $f : X \to S$ 为概形态射。映射
$$
\Hom_{\mathcal{O}_X}(\Omega_{X/S}, \mathcal{F})
\longrightarrow
\text{Der}_S(\mathcal{O}_X, \mathcal{F}),\quad
\alpha \longmapsto \alpha \circ \text{d}_{X/S}
$$
是函子 $\textit{Mod}(\mathcal{O}_X) \to \textit{Sets}$ 的同构。
\end{lemma}

\begin{proof}
这只是定义的重述。
\end{proof}

\begin{lemma}
\label{morphisms-lemma-differentials-restrict-open}
令 $f : X \to S$ 为概形态射。令 $U \subset X$、$V \subset S$ 为开子概形，
且 $f(U) \subset V$。则存在唯一同构
$\Omega_{X/S}|_U = \Omega_{U/V}$（作为 $\mathcal{O}_U$-模），使得
$\text{d}_{X/S}|_U = \text{d}_{U/V}$。
\end{lemma}

\begin{proof}
利用典范恒等式
$f^{-1}\mathcal{O}_S|_U = (f|_U)^{-1}\mathcal{O}_V$，这就是《模》引理
\ref{modules-lemma-localize-differentials} 的特殊情形。
\end{proof}

\noindent
以下我们将采用这些典范恒等式，直接用 $\Omega_{U/S}$ 或
$\Omega_{U/V}$ 表示 $\Omega_{X/S}$ 在 $U$ 上的限制。

\begin{lemma}
\label{morphisms-lemma-affine-case-derivation}
令 $R \to A$ 为环同态。令 $\mathcal{F}$ 为 $\mathcal{O}_X$-模层，
定义在 $X = \Spec(A)$ 上。置 $S = \Spec(R)$。把 $S$-导子作用于全局截面，
由此得到 $\mathcal{F}$ 的 $S$-导子之集（作用于 $\mathcal{F}$）与
$R$-导子之集（定义在 $M = \Gamma(X, \mathcal{F})$ 上）之间的双射。
\end{lemma}

\begin{proof}
令 $D : A \to M$ 为 $R$-导子。我们须证明，存在唯一的 $S$-导子，
作用于 $\mathcal{F}$，并且它在全局截面上给出 $D$。令
$U = D(f) \subset X$ 为标准仿射开集。$\Gamma(U, \mathcal{O}_X)$
的任意元素形如 $a/f^n$，其中 $a \in A$ 且 $n \geq 0$。由 Leibniz 法则，
在 $\Gamma(U, \mathcal{F})$ 中有
$$
D(a)|_U = a/f^n D(f^n)|_U + f^n D(a/f^n)
$$
。由于 $f$ 在 $\Gamma(U, \mathcal{F})$ 上可逆，这完全确定了
$D(a/f^n) \in \Gamma(U, \mathcal{F})$ 的值，从而证明唯一性。
存在性只需定义
$$
D(a/f^n) := (1/f^n) D(a)|_U - a/f^{2n} D(f^n)|_U
$$
，再证明它具有一切所需性质（在 $X$ 的标准开集基上）即可。细节从略。
\end{proof}

\begin{lemma}
\label{morphisms-lemma-differentials-affine}
令 $f : X \to S$ 为概形态射。对任意一对仿射开集
$\Spec(A) = U \subset X$、$\Spec(R) = V \subset S$，若 $f(U) \subset V$，
则存在唯一同构
$$
\Gamma(U, \Omega_{X/S}) = \Omega_{A/R}.
$$
它与 $\text{d}_{X/S}$ 以及 $\text{d} : A \to \Omega_{A/R}$ 相容。
\end{lemma}

\begin{proof}
由引理 \ref{morphisms-lemma-differentials-restrict-open}，可将 $X$、$S$ 分别替换为
$U$、$V$。因而可假设 $X = \Spec(A)$ 且 $S = \Spec(R)$，并须在
$U = X$、$V = S$ 时证明本引理。考虑 $A$-模
$M = \Gamma(X, \Omega_{X/S})$ 及其 $R$-导子
$\text{d}_{X/S} : A \to M$。令 $N$ 为另一个 $A$-模，并以
$\widetilde{N}$ 表示拟凝聚 $\mathcal{O}_X$-模，它与 $N$ 相伴；见《概形》
第 \ref{schemes-section-quasi-coherent-affine} 节。与
$\text{d}_{X/S} : A \to M$ 预复合，得到箭头
$$
\alpha : \Hom_A(M, N) \longrightarrow \text{Der}_R(A, N)
$$
。利用引理 \ref{morphisms-lemma-universal-derivation-universal} 和
\ref{morphisms-lemma-affine-case-derivation}，得到恒等式
$$
\Hom_{\mathcal{O}_X}(\Omega_{X/S}, \widetilde{N}) =
\text{Der}_S(\mathcal{O}_X, \widetilde{N}) =
\text{Der}_R(A, N)
$$
。取全局截面给出箭头
$\Hom_{\mathcal{O}_X}(\Omega_{X/S}, \widetilde{N}) \to \Hom_R(M, N)$。
把这个箭头与上述恒等式合并，得到箭头
$$
\beta : \text{Der}_R(A, N) \longrightarrow \Hom_R(M, N)
$$
。考察全局截面可知，$\alpha$ 与 $\beta$ 互为逆。因此
$\text{d}_{X/S} : A \to M$ 与 $\text{d} : A \to \Omega_{A/R}$
满足同一泛性质；见《代数》引理 \ref{algebra-lemma-universal-omega}。
于是 Yoneda 引理（《范畴》引理 \ref{categories-lemma-yoneda}）蕴含存在唯一的
$A$-模同构 $M \cong \Omega_{A/R}$，并且它与导子相容。
\end{proof}

\begin{remark}
\label{morphisms-remark-differentials-glue}
上述引理给出了构造微分模的第二种方法。具体地，令
$f : X \to S$ 为概形态射。考虑所有仿射开集 $U \subset X$，它们映入
$S$ 的某个仿射开集。它们构成 $X$ 的拓扑的一个基。因此，只须定义
$\Gamma(U, \Omega_{X/S})$（对这样的 $U$）。对这样的开集，我们直接置
$\Gamma(U, \Omega_{X/S}) = \Omega_{A/R}$；这里 $A$、$R$ 如上述引理
\ref{morphisms-lemma-differentials-affine} 中所述。这一办法可行，但要构造限制映射
并验证该方案满足层条件，还需要略多一些代数准备。
\end{remark}

\noindent
下面的引理给出了定义微分层的又一种方法，并特别表明 $\Omega_{X/S}$
在 $X$、$S$ 为概形时拟凝聚。

\begin{lemma}
\label{morphisms-lemma-differentials-diagonal}
令 $f : X \to S$ 为概形态射。$\Omega_{X/S}$ 与对角态射
$\Delta_{X/S} : X \longrightarrow X \times_S X$ 的余法层之间存在典范同构。
\end{lemma}

\begin{proof}
我们先构造一对“全局”层及层的全局同态，再在后面描述某些仿射开集上的构造。

\medskip\noindent
回忆 $\Delta = \Delta_{X/S} : X \to X \times_S X$ 是浸入；见《概形》引理
\ref{schemes-lemma-diagonal-immersion}。令 $\mathcal{J}$ 为该浸入的理想层；
它定义在某个开子概形 $W$ 上，该开子概形属于 $X \times_S X$，使得
$\Delta(X) \subset W$ 闭。取《概形》引理
\ref{schemes-lemma-diagonal-immersion} 的证明中找到的那个开集。注意，环层
$\mathcal{O}_W/\mathcal{J}^2$ 的支集包含于 $\Delta(X)$。此外，它位于层的短正合列
$$
0 \to \mathcal{J}/\mathcal{J}^2
\to \mathcal{O}_W/\mathcal{J}^2
\to \Delta_*\mathcal{O}_X
\to 0.
$$
中。利用 $\Delta^{-1}$，可把它看成满的 $f^{-1}\mathcal{O}_S$-代数层同态，
其核为 $\Delta$ 的余法层（见定义 \ref{morphisms-definition-conormal-sheaf} 和引理
\ref{morphisms-lemma-affine-conormal}）：
$$
0 \to \mathcal{C}_{X/X \times_S X}
\to \Delta^{-1}(\mathcal{O}_W/\mathcal{J}^2)
\to \mathcal{O}_X
\to 0
$$
。这把我们置于《模》引理
\ref{modules-lemma-double-structure-gives-derivation} 的情形中。投影态射
$p_i : X \times_S X \to X$（$i = 1, 2$）诱导环层同态
$(p_i)^\sharp : (p_i)^{-1}\mathcal{O}_X \to \mathcal{O}_{X \times_S X}$。
将它限制到 $W$ 并模去 $\mathcal{J}^2$，得到
$(p_i)^{-1}\mathcal{O}_X \to \mathcal{O}_W/\mathcal{J}^2$。由于
$\Delta^{-1}p_i^{-1}\mathcal{O}_X = \mathcal{O}_X$，得到同态
$$
s_i : \mathcal{O}_X \to \Delta^{-1}(\mathcal{O}_W/\mathcal{J}^2).
$$
$s_1$ 与 $s_2$ 都是同态
$\Delta^{-1}(\mathcal{O}_W/\mathcal{J}^2) \to \mathcal{O}_X$ 的截面，
正如《模》引理 \ref{modules-lemma-double-structure-gives-derivation} 中那样。
因此得到 $S$-导子
$\text{d} = s_2 - s_1 : \mathcal{O}_X \to \mathcal{C}_{X/X \times_S X}$。
由微分模的泛性质，得到唯一的 $\mathcal{O}_X$-线性同态
$$
\Omega_{X/S} \longrightarrow \mathcal{C}_{X/X \times_S X},\quad
f\text{d}g \longmapsto fs_2(g) - fs_1(g)
$$
。为证明该同态是同构，我们在适当的仿射开集上将它具体写出。可以覆盖
$X$，所用仿射开集为 $\Spec(A) = U \subset X$，使其像包含于仿射开集
$\Spec(R) = V \subset S$ 中。根据《概形》引理
\ref{schemes-lemma-diagonal-immersion} 的证明，
$U \times_V U \subset X \times_S X$ 是包含于上述 $W$ 中的仿射开集。此外，
$U \times_V U = \Spec(A \otimes_R A)$。层 $\mathcal{J}$ 对应于理想
$J = \Ker(A \otimes_R A \to A)$。该短正合列对应于
$A \otimes_R A$-模的短正合列
$$
0 \to J/J^2 \to (A \otimes_R A)/J^2 \to A \to 0
$$
。截面 $s_i$ 对应于环同态
$$
A \longrightarrow (A \otimes_R A)/J^2,\quad
s_1 : a \mapsto a \otimes 1,\quad
s_2 : a \mapsto 1 \otimes a.
$$
由引理 \ref{morphisms-lemma-affine-conormal}，有
$\Gamma(U, \mathcal{C}_{X/X \times_S X}) = J/J^2$；由引理
\ref{morphisms-lemma-differentials-affine}，有
$\Gamma(U, \Omega_{X/S}) = \Omega_{A/R}$。上述同态就是
$a \text{d}b \mapsto a \otimes b - ab \otimes 1$，它在《代数》引理
\ref{algebra-lemma-differentials-diagonal} 中已证明为同构。
\end{proof}

\begin{lemma}
\label{morphisms-lemma-functoriality-differentials}
令
$$
\xymatrix{
X' \ar[d] \ar[r]_f & X \ar[d] \\
S' \ar[r] & S
}
$$
为概形的交换图。将典范同态
$\mathcal{O}_X \to f_*\mathcal{O}_{X'}$ 与同态
$f_*\text{d}_{X'/S'} : f_*\mathcal{O}_{X'} \to f_*\Omega_{X'/S'}$ 复合，
得到一个 $S$-导子。因此得到典范的 $\mathcal{O}_X$-模同态
$\Omega_{X/S} \to f_*\Omega_{X'/S'}$；再由 $f_*$ 与 $f^*$ 的伴随性，
得到典范的 $\mathcal{O}_{X'}$-模同态
$$
c_f : f^*\Omega_{X/S} \longrightarrow \Omega_{X'/S'}.
$$
它由下述性质唯一刻画：$f^*\text{d}_{X/S}(h)$ 映到
$\text{d}_{X'/S'}(f^* h)$，其中 $h$ 是 $\mathcal{O}_X$ 的任意局部截面。
\end{lemma}

\begin{proof}
这是《模》引理
\ref{modules-lemma-functoriality-differentials-ringed-spaces} 的特殊情形。
对于概形，也可利用余法层的函子性（见引理
\ref{morphisms-lemma-conormal-functorial}）以及引理
\ref{morphisms-lemma-differentials-diagonal} 来定义 $c_f$。还可以利用本引理最后一行的
刻画，把仿射片上定义的同态粘合起来（见《代数》方程
(\ref{algebra-equation-functorial-omega})）。
\end{proof}

\begin{lemma}
\label{morphisms-lemma-triangle-differentials}
令 $f : X \to Y$、$g : Y \to S$ 为概形态射。则有典范正合列
$$
f^*\Omega_{Y/S} \to \Omega_{X/S} \to \Omega_{X/Y} \to 0
$$
，其中同态来自引理 \ref{morphisms-lemma-functoriality-differentials} 的应用。
\end{lemma}

\begin{proof}
这是《代数》引理 \ref{algebra-lemma-exact-sequence-differentials} 的层化版本。
也可以使用环层空间态射的一般版本；见《模》引理
\ref{modules-lemma-triangle-differentials}。
\end{proof}

\begin{lemma}
\label{morphisms-lemma-base-change-differentials}
令 $X \to S$ 为概形态射。令 $g : S' \to S$ 为概形态射。
令 $X' = X_{S'}$ 为 $X$ 的基变换。以 $g' : X' \to X$ 表示投影。
则引理 \ref{morphisms-lemma-functoriality-differentials} 的同态
$$
(g')^*\Omega_{X/S} \to \Omega_{X'/S'}
$$
为同构。
\end{lemma}

\begin{proof}
这是《代数》引理 \ref{algebra-lemma-differentials-base-change} 的层化版本。
\end{proof}

\begin{lemma}
\label{morphisms-lemma-differential-product}
令 $f : X \to S$ 与 $g : Y \to S$ 为具有同一目标的概形态射。令
$p : X \times_S Y \to X$ 与 $q : X \times_S Y \to Y$ 为投影态射。
引理 \ref{morphisms-lemma-functoriality-differentials} 给出的同态
$$
p^*\Omega_{X/S} \oplus q^*\Omega_{Y/S}
\longrightarrow
\Omega_{X \times_S Y/S}
$$
是同构。
\end{lemma}

\begin{proof}
由引理 \ref{morphisms-lemma-base-change-differentials}，复合
$p^*\Omega_{X/S} \to \Omega_{X \times_S Y/S} \to \Omega_{X \times_S Y/Y}$
为同构；对 $q$ 亦然。此外，由引理
\ref{morphisms-lemma-triangle-differentials}，同态
$p^*\Omega_{X/S} \to \Omega_{X \times_S Y/S}$ 的余核是
$\Omega_{X \times_S Y/X}$。结论随即成立。
\end{proof}

\begin{lemma}
\label{morphisms-lemma-finite-type-differentials}
令 $f : X \to S$ 为概形态射。若 $f$ 局部有限型，则
$\Omega_{X/S}$ 是有限型 $\mathcal{O}_X$-模。
\end{lemma}

\begin{proof}
立即由《代数》引理
\ref{algebra-lemma-differentials-finitely-generated}、引理
\ref{morphisms-lemma-differentials-affine}、引理
\ref{morphisms-lemma-locally-finite-type-characterize} 以及《性质》引理
\ref{properties-lemma-finite-type-module} 得出。
\end{proof}

\begin{lemma}
\label{morphisms-lemma-finite-presentation-differentials}
令 $f : X \to S$ 为概形态射。若 $f$ 局部有限表示，则
$\Omega_{X/S}$ 是有限表示 $\mathcal{O}_X$-模。
\end{lemma}

\begin{proof}
立即由《代数》引理
\ref{algebra-lemma-differentials-finitely-presented}、引理
\ref{morphisms-lemma-differentials-affine}、引理
\ref{morphisms-lemma-locally-finite-presentation-characterize} 以及《性质》引理
\ref{properties-lemma-finite-presentation-module} 得出。
\end{proof}

\begin{lemma}
\label{morphisms-lemma-immersion-differentials}
若 $X \to S$ 是浸入，或更一般地是单态射，则 $\Omega_{X/S}$ 为零。
\end{lemma}

\begin{proof}
这是因为此时 $\Delta_{X/S}$ 是同构，因而其余法层平凡。所以由引理
\ref{morphisms-lemma-differentials-diagonal}，有 $\Omega_{X/S} = 0$。其代数版本是
《代数》引理 \ref{algebra-lemma-trivial-differential-surjective}。
\end{proof}

\begin{lemma}
\label{morphisms-lemma-differentials-relative-immersion}
令 $i : Z \to X$ 为 $S$ 上概形的浸入。存在典范正合列
$$
\mathcal{C}_{Z/X} \to i^*\Omega_{X/S} \to \Omega_{Z/S} \to 0
$$
，其中第一个箭头由 $\text{d}_{X/S}$ 诱导，第二个箭头来自引理
\ref{morphisms-lemma-functoriality-differentials}。
\end{lemma}

\begin{proof}
这是《代数》引理 \ref{algebra-lemma-differential-seq} 的层化版本。不过，
我们应当确认第一个箭头可以全局定义。因此在此说明“由
$\text{d}_{X/S}$ 诱导”的含义。具体地，可以假设 $i$ 是闭浸入
（通过缩小 $X$）。令 $\mathcal{I} \subset \mathcal{O}_X$ 为对应于
$Z \subset X$ 的理想层。于是
$\text{d}_{X/S} : \mathcal{I} \to \Omega_{X/S}$ 把子层
$\mathcal{I}^2 \subset \mathcal{I}$ 映入
$\mathcal{I}\Omega_{X/S}$。因此它诱导同态
$\mathcal{I}/\mathcal{I}^2 \to \Omega_{X/S}/\mathcal{I}\Omega_{X/S}$，
该同态是 $\mathcal{O}_X/\mathcal{I}$-线性的。由引理
\ref{morphisms-lemma-i-star-equivalence}，这对应于所需同态
$\mathcal{C}_{Z/X} \to i^*\Omega_{X/S}$。
\end{proof}

\begin{lemma}
\label{morphisms-lemma-differentials-relative-immersion-section}
令 $i : Z \to X$ 为 $S$ 上概形的浸入，并假设 $i$（局部）有左逆。
则引理 \ref{morphisms-lemma-differentials-relative-immersion} 的典范序列
$$
0 \to \mathcal{C}_{Z/X} \to i^*\Omega_{X/S} \to \Omega_{Z/S} \to 0
$$
（局部）分裂正合。特别地，若 $s : S \to X$ 是结构态射
$X \to S$ 的截面，则同态
$\mathcal{C}_{S/X} \to s^*\Omega_{X/S}$（由 $\text{d}_{X/S}$ 诱导）是同构。
\end{lemma}

\begin{proof}
由《代数》引理 \ref{algebra-lemma-differential-seq-split}，当 $X, Z, S$
均仿射时结论成立。一般情形下，由此得出该序列局部分裂正合。若
$g : X \to Z$ 是 $i$ 的左逆，则由《模》引理
\ref{modules-lemma-check-functoriality-differentials}，$i^*c_g$ 是同态
$i^*\Omega_{X/S} \to \Omega_{Z/S}$ 的右逆。于是由《同调》引理
\ref{homology-lemma-ses-split}，该序列分裂。最后，若 $s$ 是截面，
则它是浸入 $s : Z = S \to X$，并且这是 $S$ 上的浸入（见《概形》引理
\ref{schemes-lemma-section-immersion}），此时 $\Omega_{Z/S} = 0$。
\end{proof}

\begin{remark}
\label{morphisms-remark-differentials-diagonal}
令 $X \to S$ 为概形态射。由引理
\ref{morphisms-lemma-differential-product}，有
$$
\Omega_{X \times_S X/S} =
\text{pr}_1^*\Omega_{X/S} \oplus \text{pr}_2^*\Omega_{X/S}
$$
。另一方面，对角态射 $\Delta : X \to X \times_S X$ 是浸入，且局部有左逆。
因此由引理 \ref{morphisms-lemma-differentials-relative-immersion-section}，得到典范短正合列
$$
0 \to \mathcal{C}_{X/X \times_S X} \to \Omega_{X/S} \oplus \Omega_{X/S}
\to \Omega_{X/S} \to 0
$$
。注意右侧箭头是 $(1, 1)$，它确实是分裂满射。另一方面，由引理
\ref{morphisms-lemma-differentials-diagonal}，有恒等式
$\Omega_{X/S} = \mathcal{C}_{X/X \times_S X}$。由于在这一恒等式中选择了
$\text{d}_{X/S}(f) = s_2(f) - s_1(f)$，左侧箭头遂为同态
$(-1, 1)$\footnote{具体地，局部截面
$\text{d}_{X/S}(f) = 1 \otimes f - f \otimes 1$（属于对角态射 $\Delta$ 的理想层）经由
$\text{d}_{X \times_S X/X}$ 映到局部截面
$1 \otimes 1 \otimes 1 \otimes f - 1 \otimes f \otimes 1 \otimes 1
-1 \otimes 1 \otimes f \otimes 1 + f \otimes 1 \otimes 1 \otimes 1 =
\text{pr}_2^*\text{d}_{X/S}(f) - \text{pr}_1^*\text{d}_{X/S}(f)$。}。
\end{remark}

\begin{lemma}
\label{morphisms-lemma-two-immersions}
令
$$
\xymatrix{
Z \ar[r]_i \ar[rd]_j & X \ar[d] \\
& Y
}
$$
为概形的交换图，其中 $i$ 和 $j$ 为浸入。则有典范正合列
$$
\mathcal{C}_{Z/Y} \to
\mathcal{C}_{Z/X} \to
i^*\Omega_{X/Y} \to 0
$$
，其中第一个箭头来自引理
\ref{morphisms-lemma-conormal-functorial}，第二个箭头来自引理
\ref{morphisms-lemma-differentials-relative-immersion}。
\end{lemma}

\begin{proof}
其代数版本是《代数》引理 \ref{algebra-lemma-application-NL}。
\end{proof}










\section{有限阶微分算子}
\label{morphisms-section-differential-operators}

\noindent
我们建议读者参阅交换代数章中的相应一节
（《代数》第 \ref{algebra-section-differential-operators} 节），以及模层章中的
相应一节（《模》第 \ref{modules-section-differential-operators} 节）。

\begin{lemma}
\label{morphisms-lemma-affine-case-differential-operators}
令 $R \to A$ 为环同态。以 $f : X \to S$ 表示相应的仿射概形态射。
令 $\mathcal{F}$ 和 $\mathcal{G}$ 为 $\mathcal{O}_X$-模。若
$\mathcal{F}$ 拟凝聚，则把微分算子送到其在全局截面上的作用所得的映射
$$
\text{Diff}^k_{X/S}(\mathcal{F}, \mathcal{G}) \to
\text{Diff}^k_{A/R}(\Gamma(X, \mathcal{F}), \Gamma(X, \mathcal{G}))
$$
是双射。
\end{lemma}

\begin{proof}
把 $\mathcal{F} = \widetilde{M}$ 写成某个 $A$-模 $M$ 的形式。置
$N = \Gamma(X, \mathcal{G})$。令 $D : M \to N$ 为阶数 $k$ 的微分算子。
我们须证明，存在唯一的微分算子
$\mathcal{F} \to \mathcal{G}$，其阶数为 $k$，并且它在全局截面上给出 $D$。令
$U = D(f) \subset X$ 为标准仿射开集。则 $\mathcal{F}(U) = M_f$ 是局部化。
由《代数》引理 \ref{algebra-lemma-invert-system-differential-operators}，
微分算子 $D$ 唯一延拓为微分算子
$$
D_f : \mathcal{F}(U) = \widetilde{M}(U) = M_f \to N_f = \widetilde{N}(U)
$$
。唯一性表明这些同态 $D_f$ 粘合，给出层同态
$\widetilde{M} \to \widetilde{N}$；这一粘合发生在 $X$ 的所有标准开集构成的基上。因此，由《层》第
\ref{sheaves-section-bases} 节的内容，得到唯一的层同态
$\widetilde{D} : \widetilde{M} \to \widetilde{N}$，它与这些同态相合。
由于 $\widetilde{D}$ 在标准开集上由阶数 $k$ 的微分算子给出，故
$\widetilde{D}$ 是阶数 $k$ 的微分算子（略去一个小细节）。最后，与典范的
$\mathcal{O}_X$-模同态 $c : \widetilde{N} \to \mathcal{G}$ 后复合
（《概形》引理 \ref{schemes-lemma-compare-constructions}），得到
$c \circ \widetilde{D} : \mathcal{F} \to \mathcal{G}$；由《模》引理
\ref{modules-lemma-composition-differential-operators}，这是阶数 $k$
的微分算子。存在性得证。唯一性的证明从略。
\end{proof}

\begin{lemma}
\label{morphisms-lemma-base-change-differential-operators}
令 $a : X \to S$ 与 $b : Y \to S$ 为概形态射。令 $\mathcal{F}$ 和
$\mathcal{F}'$ 为拟凝聚 $\mathcal{O}_X$-模。令
$D : \mathcal{F} \to \mathcal{F}'$ 为阶数 $k$ 的微分算子，定义在 $X/S$ 上。
令 $\mathcal{G}$ 为拟凝聚 $\mathcal{O}_Y$-模。则存在唯一的微分算子
$$
D' :
\text{pr}_1^*\mathcal{F} \otimes_{\mathcal{O}_{X \times_S Y}}
\text{pr}_2^*\mathcal{G}
\longrightarrow
\text{pr}_1^*\mathcal{F}' \otimes_{\mathcal{O}_{X \times_S Y}}
\text{pr}_2^*\mathcal{G}
$$
，它的阶数为 $k$，定义在 $X \times_S Y / Y$ 上，并且
$
D'(s \otimes t) = D(s) \otimes t
$
，其中 $s$ 是 $\mathcal{F}$ 的局部截面，$t$ 是 $\mathcal{G}$ 的局部截面。
\end{lemma}

\begin{proof}
当 $X$、$Y$、$S$ 仿射时，利用引理
\ref{morphisms-lemma-affine-case-differential-operators}，结论由相应的代数结果得出；
见《代数》引理 \ref{algebra-lemma-base-change-differential-operators}。
一般情形下，利用仿射覆盖（例如《概形》引理
\ref{schemes-lemma-affine-covering-fibre-product} 中的覆盖）全局构造 $D'$。
细节从略。
\end{proof}

\begin{remark}
\label{morphisms-remark-base-change-differential-operators}
令 $a : X \to S$ 与 $b : Y \to S$ 为概形态射。以
$p : X \times_S Y \to X$ 和 $q : X \times_S Y \to Y$ 表示投影。
在本注中，给定 $\mathcal{O}_X$-模 $\mathcal{F}$ 和 $\mathcal{O}_Y$-模
$\mathcal{G}$，置
$$
\mathcal{F} \boxtimes \mathcal{G} =
p^*\mathcal{F} \otimes_{\mathcal{O}_{X \times_S Y}} q^*\mathcal{G}
$$
。以 $\mathcal{A}_{X/S}$ 表示加性范畴，其对象为拟凝聚
$\mathcal{O}_X$-模，态射为 $X/S$ 上的有限阶微分算子。
$\mathcal{A}_{Y/S}$ 和 $\mathcal{A}_{X \times_S Y/S}$ 同理。
引理 \ref{morphisms-lemma-base-change-differential-operators} 的构造确定了函子
$$
\boxtimes : 
\mathcal{A}_{X/S} \times \mathcal{A}_{Y/S} \longrightarrow
\mathcal{A}_{X \times_S Y/S},
\quad
(\mathcal{F}, \mathcal{G}) \longmapsto \mathcal{F} \boxtimes \mathcal{G}
$$
，它在态射上双线性。若 $X = \Spec(A)$、$Y = \Spec(B)$ 且
$S = \Spec(R)$，则通过拟凝聚层与模的恒等，此函子在对象上由
$(M, N) \mapsto M \otimes_R N$ 给出，并把态射
$(D, D') : (M, N) \to (M', N')$ 送到
$D \otimes D' : M \otimes_R N \to M' \otimes_R N'$。
\end{remark}










\section{光滑态射}
\label{morphisms-section-smooth}

\noindent
令 $f : X \to Y$ 为拓扑空间的连续映射。考虑下列条件：对每个
$x \in X$，存在开邻域 $x \in U \subset X$、$f(x) \in V \subset Y$，
以及整数 $d$，使得 $f(U) \subset V$，并得到交换图
$$
\xymatrix{
X \ar[d] &
U \ar[l] \ar[d] \ar[r]_-\pi &
V \times \mathbf{R}^d \ar[ld] \\
Y & V \ar[l]
}
$$
，其中 $\pi$ 是到某个开子集上的同胚。概形的光滑态射是概形范畴中
这类映射的类似物。见引理 \ref{morphisms-lemma-smooth-locally-standard-smooth} 和引理
\ref{morphisms-lemma-smooth-etale-over-affine-space}。

\medskip\noindent
与（或许的）预期相反，光滑环同态的概念并非只用微分模定义。回忆：
$R \to A$ 是{\it 光滑环同态}，是指 $A$ 在 $R$ 上有限表示，并且 $A$
在 $R$ 上的朴素余切复形拟同构于置于次数 $0$ 的投射模；见《代数》定义
\ref{algebra-definition-smooth}。

\begin{definition}
\label{morphisms-definition-smooth}
令 $f : X \to S$ 为概形态射。
\begin{enumerate}
\item 称 $f$ 在 $x \in X$ 处{\it 光滑}，若存在仿射开邻域
$\Spec(A) = U \subset X$（它是 $x$ 的邻域）以及仿射开集
$\Spec(R) = V \subset S$，满足 $f(U) \subset V$，且诱导环同态
$R \to A$ 光滑。
\item 称 $f$ {\it 光滑}，若它在 $X$ 的每一点处都光滑。
\item 若仿射概形态射 $f : X \to S$ 同构于
某个标准光滑环同态
\begin{center}
$R \to R[x_1, \ldots, x_n]/(f_1, \ldots, f_c)$
\end{center}
（见《代数》定义
\ref{algebra-definition-standard-smooth}）所给出的态射 $X \to S$，即
$$
\Spec(R[x_1, \ldots, x_n]/(f_1, \ldots, f_c)) \to \Spec(R).
$$
则称它{\it 标准光滑}。
\end{enumerate}
\end{definition}

\noindent
这个定义的一个良好特征是：态射光滑的点所成之集自动为开集。

\medskip\noindent
注意定义中没有分离性或拟紧性的假设。因此，光滑性在源上是局部性质。
精确结果如下。

\begin{lemma}
\label{morphisms-lemma-smooth-characterize}
令 $f : X \to S$ 为概形态射。下列条件等价：
\begin{enumerate}
\item 态射 $f$ 光滑。
\item 对每一对仿射开集 $U \subset X$、$V \subset S$，若
$f(U) \subset V$，则环同态
$\mathcal{O}_S(V) \to \mathcal{O}_X(U)$ 光滑。
\item 存在开覆盖 $S = \bigcup_{j \in J} V_j$，以及开覆盖
$f^{-1}(V_j) = \bigcup_{i \in I_j} U_i$，使得每个态射
$U_i \to V_j$（$j\in J, i\in I_j$）都光滑。
\item 存在仿射开覆盖 $S = \bigcup_{j \in J} V_j$，以及仿射开覆盖
$f^{-1}(V_j) = \bigcup_{i \in I_j} U_i$，使得环同态
$\mathcal{O}_S(V_j) \to \mathcal{O}_X(U_i)$ 对所有
$j\in J, i\in I_j$ 都光滑。
\end{enumerate}
此外，若 $f$ 光滑，则对任意开子概形 $U \subset X$、$V \subset S$，
若 $f(U) \subset V$，限制 $f|_U : U \to V$ 光滑。
\end{lemma}

\begin{proof}
若能证明“$R \to A$ 光滑”是局部性质，则结论由引理
\ref{morphisms-lemma-locally-P} 得出。我们验证定义
\ref{morphisms-definition-property-local} 的条件 (a)、(b)、(c)。由《代数》引理
\ref{algebra-lemma-base-change-smooth}，光滑性在基变换下稳定，故 (a) 成立。
由《代数》引理 \ref{algebra-lemma-compose-smooth}，光滑性在复合下稳定；
并且对任意环 $R$，环同态 $R \to R_f$（标准）光滑，故 (b) 成立。
最后，由《代数》引理 \ref{algebra-lemma-locally-smooth}，性质 (c) 成立。
\end{proof}

\noindent
下面的引理把光滑态射刻画为纤维光滑的平坦有限表示态射。注意，
域上的光滑概形会在《簇》第 \ref{varieties-section-smooth} 节中更详细地讨论。

\begin{lemma}
\label{morphisms-lemma-smooth-flat-smooth-fibres}
令 $f : X \to S$ 为概形态射。若 $f$ 平坦、局部有限表示，且所有纤维
$X_s$ 都光滑，则 $f$ 光滑。
\end{lemma}

\begin{proof}
由《代数》引理 \ref{algebra-lemma-flat-fibre-smooth} 得出。
\end{proof}

\begin{lemma}
\label{morphisms-lemma-composition-smooth}
两个光滑态射的复合光滑。
\end{lemma}

\begin{proof}
在引理 \ref{morphisms-lemma-smooth-characterize} 的证明中，我们已看到光滑性是
环同态的局部性质。因此，本引理的第一项断言由引理
\ref{morphisms-lemma-composition-type-P} 以及光滑环同态性质在复合下稳定这一事实得出；
见《代数》引理 \ref{algebra-lemma-compose-smooth}。
\end{proof}

\begin{lemma}
\label{morphisms-lemma-base-change-smooth}
光滑态射的基变换光滑。
\end{lemma}

\begin{proof}
在引理 \ref{morphisms-lemma-smooth-characterize} 的证明中，我们已看到光滑性是
环同态的局部性质。因此，本引理由引理
\ref{morphisms-lemma-composition-type-P} 以及光滑环同态性质在基变换下稳定这一事实得出；
见《代数》引理 \ref{algebra-lemma-base-change-smooth}。
\end{proof}

\begin{lemma}
\label{morphisms-lemma-open-immersion-smooth}
任意开浸入都光滑。
\end{lemma}

\begin{proof}
这是因为开浸入是局部同构。
\end{proof}

\begin{lemma}
\label{morphisms-lemma-smooth-syntomic}
光滑态射是 syntomic 态射。
\end{lemma}

\begin{proof}
见《代数》引理 \ref{algebra-lemma-smooth-syntomic}。
\end{proof}

\begin{lemma}
\label{morphisms-lemma-smooth-locally-finite-presentation}
光滑态射局部有限表示。
\end{lemma}

\begin{proof}
这是因为按定义，光滑环同态有限表示。
\end{proof}

\begin{lemma}
\label{morphisms-lemma-smooth-flat}
光滑态射平坦。
\end{lemma}

\begin{proof}
合并引理 \ref{morphisms-lemma-syntomic-flat} 与 \ref{morphisms-lemma-smooth-syntomic}。
\end{proof}

\begin{lemma}
\label{morphisms-lemma-smooth-open}
光滑态射泛开。
\end{lemma}

\begin{proof}
合并引理 \ref{morphisms-lemma-smooth-flat}、
\ref{morphisms-lemma-smooth-locally-finite-presentation} 和
\ref{morphisms-lemma-fppf-open}。或者，也可合并引理
\ref{morphisms-lemma-smooth-syntomic} 与
\ref{morphisms-lemma-syntomic-open}。
\end{proof}

\noindent
下面的引理说明任意光滑态射在局部上标准光滑。因此，可用标准光滑态射作为
光滑态射的{\it 局部模型}。

\begin{lemma}
\label{morphisms-lemma-smooth-locally-standard-smooth}
\begin{slogan}
光滑态射局部标准光滑。
\end{slogan}
令 $f : X  \to S$ 为概形态射。令 $x \in X$ 为一点。令
$V \subset S$ 为 $f(x)$ 的仿射开邻域。下列条件等价：
\begin{enumerate}
\item 态射 $f$ 在 $x$ 处光滑。
\item 存在仿射开集 $U \subset X$，满足 $x \in U$、$f(U) \subset V$，
并且诱导态射 $f|_U : U \to V$ 标准光滑。
\end{enumerate}
\end{lemma}

\begin{proof}
由定义以及《代数》引理 \ref{algebra-lemma-standard-smooth} 和
\ref{algebra-lemma-smooth-syntomic} 得出。
\end{proof}

\begin{lemma}
\label{morphisms-lemma-smooth-omega-finite-locally-free}
令 $f : X \to S$ 为概形态射。假设 $f$ 光滑。则微分模
$\Omega_{X/S}$（即 $X$ 在 $S$ 上的微分模）有限局部自由，并且
$$
\text{rank}_x(\Omega_{X/S}) = \dim_x(X_{f(x)})
$$
对每个 $x \in X$ 都成立。
\end{lemma}

\begin{proof}
该断言在 $X$ 和 $S$ 上是局部的。由上述引理
\ref{morphisms-lemma-smooth-locally-standard-smooth}，可以假设 $f$ 是仿射概形的
标准光滑态射。此时结论由《代数》引理
\ref{algebra-lemma-standard-smooth} 得出（以及全局相对完全交的定义；见《代数》
定义 \ref{algebra-definition-relative-global-complete-intersection}）。
\end{proof}

\noindent
引理 \ref{morphisms-lemma-smooth-omega-finite-locally-free} 表明下述定义是有意义的。

\begin{definition}
\label{morphisms-definition-smooth-relative-dimension}
令 $d \geq 0$ 为整数。称概形态射 $f : X \to S$
为{\it 相对维数 $d$ 的光滑态射}，若 $f$ 光滑，且
$\Omega_{X/S}$ 是常秩 $d$ 的有限局部自由模。
\end{definition}

\noindent
换言之，$f$ 光滑，且非空纤维等维、维数为 $d$。由下面的引理
\ref{morphisms-lemma-smooth-at-point}，这也等价于要求：(a) $f$ 局部有限表示，
(b) $f$ 平坦，(c) 所有非空纤维等维且维数为 $d$，以及 (d)
$\Omega_{X/S}$ 是秩 $d$ 的有限局部自由模。仅仅假设 $f$ 平坦、有限表示，
且 $\Omega_{X/S}$ 是秩 $d$ 的有限局部自由模，并不足够。反例为
$\Spec(\mathbf{F}_p[t]) \to \Spec(\mathbf{F}_p[t^p])$。

\medskip\noindent
下面给出一点处光滑性的微分判据。这一结果有许多变体，都可能在某处有用。
我们会按需在这里加入它们。

\begin{lemma}
\label{morphisms-lemma-smooth-at-point}
令 $f : X \to S$ 为概形态射。令 $x \in X$。置 $s = f(x)$。
假设 $f$ 局部有限表示。下列条件等价：
\begin{enumerate}
\item 态射 $f$ 在 $x$ 处光滑。
\item 局部环同态 $\mathcal{O}_{S, s} \to \mathcal{O}_{X, x}$ 平坦，且
$X_s \to \Spec(\kappa(s))$ 在 $x$ 处光滑。
\item 局部环同态 $\mathcal{O}_{S, s} \to \mathcal{O}_{X, x}$ 平坦，且
$\mathcal{O}_{X, x}$-模 $\Omega_{X/S, x}$ 可由至多
$\dim_x(X_{f(x)})$ 个元素生成。
\item 局部环同态 $\mathcal{O}_{S, s} \to \mathcal{O}_{X, x}$ 平坦，且
$\kappa(x)$-向量空间
$$
\Omega_{X_s/s, x} \otimes_{\mathcal{O}_{X_s, x}} \kappa(x) =
\Omega_{X/S, x} \otimes_{\mathcal{O}_{X, x}} \kappa(x)
$$
可由至多 $\dim_x(X_{f(x)})$ 个元素生成。
\item 存在仿射开集 $U \subset X$ 与 $V \subset S$，使得
$x \in U$、$f(U) \subset V$，并且诱导态射
$f|_U : U \to V$ 标准光滑。
\item 存在仿射开集 $\Spec(A) = U \subset X$ 和
$\Spec(R) = V \subset S$，其中 $x \in U$ 对应于
$\mathfrak q \subset A$，且 $f(U) \subset V$，并存在表示
$$
A = R[x_1, \ldots, x_n]/(f_1, \ldots, f_c)
$$
使得
$$
g =
\det
\left(
\begin{matrix}
\partial f_1/\partial x_1 &
\partial f_2/\partial x_1 &
\ldots &
\partial f_c/\partial x_1 \\
\partial f_1/\partial x_2 &
\partial f_2/\partial x_2 &
\ldots &
\partial f_c/\partial x_2 \\
\ldots & \ldots & \ldots & \ldots \\
\partial f_1/\partial x_c &
\partial f_2/\partial x_c &
\ldots &
\partial f_c/\partial x_c
\end{matrix}
\right)
$$
在 $A$ 中的像不属于 $\mathfrak q$。
\end{enumerate}
\end{lemma}

\begin{proof}
注意，若 $f$ 在 $x$ 处光滑，则由引理
\ref{morphisms-lemma-smooth-locally-standard-smooth} 可知 (5) 成立，而 (6) 是 (5)
略微减弱的版本。此外，$f$ 光滑蕴含局部环同态
$\mathcal{O}_{S, s} \to \mathcal{O}_{X, x}$ 平坦（见引理
\ref{morphisms-lemma-smooth-flat}），并且 $\Omega_{X/S}$ 有限局部自由，其秩等于
$\dim_x(X_s)$（见引理 \ref{morphisms-lemma-smooth-omega-finite-locally-free}）。
因此 (1) 蕴含 (3) 和 (4)。由引理 \ref{morphisms-lemma-base-change-smooth}，
还可知 (1) 蕴含 (2)。

\medskip\noindent
由引理 \ref{morphisms-lemma-base-change-differentials}，微分模
$\Omega_{X_s/s}$（即纤维 $X_s$ 在 $\kappa(s)$ 上的微分模）是微分模
$\Omega_{X/S}$（即 $X$ 在 $S$ 上的微分模）的拉回。这给出本引理第 (4) 项中的等式。由引理
\ref{morphisms-lemma-finite-type-differentials}，这些模有限型。因此，模
$\Omega_{X/S, x}$ 与 $\Omega_{X_s/s, x}$ 的最小生成元数相同，并且由
Nakayama 引理（《代数》引理 \ref{algebra-lemma-NAK}）等于这个
$\kappa(x)$-向量空间的维数。这尤其表明 (3) 与 (4) 等价。

\medskip\noindent
《代数》引理 \ref{algebra-lemma-flat-fibre-smooth} 表明 (2) 蕴含 (1)。
《代数》引理 \ref{algebra-lemma-characterize-smooth-over-field} 表明 (3)
和 (4) 蕴含 (2)。最后，(6) 蕴含 (5)，例如见《代数》例
\ref{algebra-example-make-standard-smooth}；而由《代数》引理
\ref{algebra-lemma-standard-smooth}，(5) 蕴含 (1)。
\end{proof}

\begin{lemma}
\label{morphisms-lemma-set-points-where-fibres-smooth}
令
$$
\xymatrix{
X' \ar[r]_{g'} \ar[d]_{f'} & X \ar[d]^f \\
S' \ar[r]^g & S
}
$$
为概形的笛卡尔图。令 $W \subset X$ 和 $W' \subset X'$ 分别为
$f$ 与 $f'$ 光滑的点所成的开子概形。若下列任一条件成立，则
$W' = (g')^{-1}(W)$：
\begin{enumerate}
\item $f$ 平坦且局部有限表示；或
\item $f$ 局部有限表示且 $g$ 平坦。
\end{enumerate}
\end{lemma}

\begin{proof}
先假设 $f$ 局部有限型。考虑集合
$$
T = \{x \in X \mid X_{f(x)}\text{ is smooth over }\kappa(f(x))\text{ 在 }x\}
$$
以及集合 $T' \subset X'$（对应于 $f'$）。我们断言
$T' = (g')^{-1}(T)$。事实上，令 $s' \in S'$ 为一点，并令
$s = g(s')$。则有
$$
X'_{s'} =
\Spec(\kappa(s')) \times_{\Spec(\kappa(s))} X_s
$$
。换言之，基变换的纤维是纤维的基变换。因此，该断言等价于《代数》引理
\ref{algebra-lemma-smooth-field-change-local}。

\medskip\noindent
这样，情形 (1) 成立，因为在情形 (1) 中，由引理
\ref{morphisms-lemma-smooth-at-point}，$T$ 正是 $f$ 光滑的点所成的（开）集。

\medskip\noindent
在情形 (2) 中，令 $x' \in W'$。则 $g'$ 在 $x'$ 处平坦（引理
\ref{morphisms-lemma-base-change-module-flat}），并且 $g \circ f'$ 在 $x'$ 处平坦
（引理 \ref{morphisms-lemma-composition-module-flat}）。因此，由引理
\ref{morphisms-lemma-flat-permanence}，$f$ 在 $x = g'(x')$ 处平坦。另一方面，
由于 $x' \in T'$（引理 \ref{morphisms-lemma-base-change-smooth}），可知
$x \in T$。所以，由引理 \ref{morphisms-lemma-smooth-at-point}，$f$ 在 $x$ 处光滑。
\end{proof}

\noindent
下面的引理确实用到了光滑环同态的朴素余切复形之 $H^{-1}$ 消失。

\begin{lemma}
\label{morphisms-lemma-triangle-differentials-smooth}
令 $f : X \to Y$、$g : Y \to S$ 为概形态射。假设 $f$ 光滑。则
$$
0 \to f^*\Omega_{Y/S} \to \Omega_{X/S} \to \Omega_{X/Y} \to 0
$$
（见引理 \ref{morphisms-lemma-triangle-differentials}）短正合。
\end{lemma}

\begin{proof}
本引理的代数版本如下：给定环同态 $A \to B \to C$，其中
$B \to C$ 光滑，则《代数》引理
\ref{algebra-lemma-exact-sequence-differentials} 的序列
$$
0 \to C \otimes_B \Omega_{B/A} \to \Omega_{C/A} \to \Omega_{C/B} \to 0
$$
正合。这就是《代数》引理
\ref{algebra-lemma-triangle-differentials-smooth}。
\end{proof}

\begin{lemma}
\label{morphisms-lemma-differentials-relative-immersion-smooth}
令 $i : Z \to X$ 为 $S$ 上概形的浸入。假设 $Z$ 在 $S$ 上光滑。
则引理 \ref{morphisms-lemma-differentials-relative-immersion} 的典范正合列
$$
0 \to \mathcal{C}_{Z/X} \to i^*\Omega_{X/S} \to \Omega_{Z/S} \to 0
$$
短正合。
\end{lemma}

\begin{proof}
本引理的代数版本如下：给定环同态 $A \to B \to C$，其中
$A \to C$ 光滑，且 $B \to C$ 满、核为 $J$，则《代数》引理
\ref{algebra-lemma-differential-seq} 的序列
$$
0 \to J/J^2 \to C \otimes_B \Omega_{B/A} \to \Omega_{C/A} \to 0
$$
正合。这就是《代数》引理
\ref{algebra-lemma-differential-seq-smooth}。
\end{proof}

\begin{lemma}
\label{morphisms-lemma-two-immersions-smooth}
令
$$
\xymatrix{
Z \ar[r]_i \ar[rd]_j & X \ar[d] \\
& Y
}
$$
为概形的交换图，其中 $i$ 与 $j$ 是浸入，并且 $X \to Y$ 光滑。
则引理 \ref{morphisms-lemma-two-immersions} 的典范正合列
$$
0 \to  \mathcal{C}_{Z/Y} \to \mathcal{C}_{Z/X} \to i^*\Omega_{X/Y} \to 0
$$
正合。
\end{lemma}

\begin{proof}
本引理的代数版本如下：给定环同态 $A \to B \to C$，其中
$A \to C$ 满，且 $A \to B$ 光滑，则《代数》引理
\ref{algebra-lemma-application-NL} 的序列
$$
0 \to I/I^2 \to J/J^2 \to C \otimes_B \Omega_{B/A} \to 0
$$
正合。这就是《代数》引理 \ref{algebra-lemma-application-NL-smooth}。
\end{proof}

\begin{lemma}
\label{morphisms-lemma-smooth-permanence}
令
$$
\xymatrix{
X \ar[rr]_f \ar[rd]_p & &
Y \ar[dl]^q \\
& S
}
$$
为概形态射的交换图。假设：
\begin{enumerate}
\item $f$ 满且光滑；
\item $p$ 光滑；并且
\item $q$ 局部有限表示\footnote{事实上，这由 (1) 和 (2) 蕴含；见《下降》
引理 \ref{descent-lemma-flat-finitely-presented-permanence}。此外，只需假设
$f$ 满、平坦且局部有限表示；见《下降》引理
\ref{descent-lemma-smooth-permanence}。}。
\end{enumerate}
则 $q$ 光滑。
\end{lemma}

\begin{proof}
由引理 \ref{morphisms-lemma-flat-permanence}，可知 $q$ 平坦。取点 $y \in Y$。
取一点 $x \in X$，它映到 $y$。设 $f$ 的相对维数为 $a$（在 $x$ 处），且
$p$ 的相对维数为 $b$（在 $x$ 处）。由引理
\ref{morphisms-lemma-smooth-omega-finite-locally-free}，这意味着
$\Omega_{X/S, x}$ 为秩 $b$ 的自由模，且 $\Omega_{X/Y, x}$ 为秩 $a$
的自由模。由引理 \ref{morphisms-lemma-triangle-differentials-smooth} 的短正合列，
这意味着 $(f^*\Omega_{Y/S})_x$ 为秩 $b - a$ 的自由模。由 Nakayama
引理，这蕴含 $\Omega_{Y/S, y}$ 可由 $b - a$ 个元素生成。此外，
由引理 \ref{morphisms-lemma-dimension-fibre-at-a-point-additive}，有
$\dim_y(Y_s) = b - a$。所以，由引理 \ref{morphisms-lemma-smooth-at-point} 第 (2) 项，
可得 $Y \to S$ 在 $y$ 处光滑。
\end{proof}

\noindent
在下面引理的情形中，$\sigma$ 的像在 $X$ 上局部由正则序列切出；见《除子》
引理 \ref{divisors-lemma-section-smooth-regular-immersion}。

\begin{lemma}
\label{morphisms-lemma-section-smooth-morphism}
令 $f : X \to S$ 为概形态射。令 $\sigma : S \to X$ 为 $f$ 的截面。
令 $s \in S$ 为一点，使得 $f$ 在 $x = \sigma(s)$ 处光滑。则存在
仿射开邻域 $\Spec(A) = U \subset S$（属于 $s$）和
$\Spec(B) = V \subset X$（属于 $x$），使得：
\begin{enumerate}
\item $f(V) \subset U$ 且 $\sigma(U) \subset V$；
\item 对于 $I = \Ker(\sigma^\# : B \to A)$，模 $I/I^2$ 是自由 $A$-模；并且
\item 有 $B^\wedge \cong A[[x_1, \ldots, x_d]]$（作为 $A$-代数），其中
$B^\wedge$ 表示 $B$ 关于 $I$ 的完备化。
\end{enumerate}
\end{lemma}

\begin{proof}
取仿射开集 $U \subset S$，它包含 $s$。取仿射开集
$V \subset f^{-1}(U)$，它包含 $x$。取仿射开集
$U' \subset \sigma^{-1}(V)$，它包含 $s$。注意，$V' = f^{-1}(U') \cap V$ 仿射，
因为它等于纤维积 $V' = U' \times_U V$。于是 $U'$ 与 $V'$ 满足 (1)。
写 $U' = \Spec(A')$ 且 $V' = \Spec(B')$。由《代数》引理
\ref{algebra-lemma-section-smooth}，模 $I'/(I')^2$ 是有限局部自由
$A'$-模。因此，把 $U'$ 替换为较小的仿射开集 $U'' \subset U'$，并把
$V'$ 替换为 $V'' = V' \cap f^{-1}(U'')$ 后，得到
$I''/(I'')^2$ 自由的情形，即 (2) 成立。在此情形下，由《代数》引理
\ref{algebra-lemma-section-smooth}，(3) 也成立。
\end{proof}

\noindent
概形 $X$ 在一点 $x$ 处的维数（《性质》定义
\ref{properties-definition-dimension}），就是把 $X$ 视为拓扑空间时在
$x$ 处的维数；见《拓扑》定义 \ref{topology-definition-Krull}。一般而言，
这不是局部环 $\mathcal{O}_{X,x}$ 的维数。

\begin{lemma}
\label{morphisms-lemma-smoothness-dimension}
令 $f : X \to Y$ 为局部 Noether 概形之间的光滑态射。对每个点
$x$（位于 $X$），若其像 $y$ 位于 $Y$，则
$$
\dim_x(X) = \dim_y(Y) + \dim_x(X_y),
$$
其中 $X_y$ 表示 $y$ 上的纤维。
\end{lemma}

\begin{proof}
把 $X$ 替换为 $x$ 的开邻域后，存在自然数 $d$，使得
$X \to Y$ 的所有纤维在每一点处的维数都是 $d$；见引理
\ref{morphisms-lemma-smooth-omega-finite-locally-free}。于是 $f$ 平坦（引理
\ref{morphisms-lemma-smooth-flat}）、局部有限型（引理
\ref{morphisms-lemma-smooth-locally-finite-presentation}），且相对维数为 $d$。
因此结论由引理 \ref{morphisms-lemma-rel-dimension-dimension} 得出。
\end{proof}

\section{非分歧态射}
\label{morphisms-section-unramified}

\noindent
在讨论（也许）更有趣的平展态射这一类之前，我们先简要讨论非分歧态射。
回顾：若环同态 $R \to A$ 是有限型的并且 $\Omega_{A/R} = 0$，
则称它是{\it 非分歧的}（这是 \cite{Henselian} 的定义）。
若环同态 $R \to A$ 是有限表示的并且 $\Omega_{A/R} = 0$，
则称它是 {\it G-非分歧的}（这是 \cite{EGA} 的定义）。见
代数，引理 \ref{algebra-definition-unramified}。

\begin{definition}
\label{morphisms-definition-unramified}
设 $f : X \to S$ 是概形的态射。
\begin{enumerate}
\item 称 $f$ {\it 在 $x \in X$ 处非分歧}，如果存在仿射开邻域
$\Spec(A) = U \subset X$ 包含 $x$，以及仿射开集
$\Spec(R) = V \subset S$，满足 $f(U) \subset V$，并且所诱导的
环同态 $R \to A$ 非分歧。
\item 称 $f$ {\it 在 $x \in X$ 处 G-非分歧}，如果存在仿射开邻域
$\Spec(A) = U \subset X$ 包含 $x$，以及仿射开集
$\Spec(R) = V \subset S$，满足 $f(U) \subset V$，并且所诱导的
环同态 $R \to A$ 是 G-非分歧的。
\item 若 $f$ 在 $X$ 的每一点处非分歧，则称其{\it 非分歧}。
\item 若 $f$ 在 $X$ 的每一点处 G-非分歧，则称其
{\it G-非分歧}。
\end{enumerate}
\end{definition}

\noindent
注意，G-非分歧态射是非分歧态射。因此，关于非分歧态射的任何结果
都会蕴含关于 G-非分歧态射的相应结果。此外，若 $S$ 局部 Noether，
则 G-非分歧态射与非分歧态射没有区别；见引理
\ref{morphisms-lemma-noetherian-unramified}。这个定义的一个可取之处是：
一个态射非分歧（相应地，G-非分歧）的点集自动是开集。

\begin{lemma}
\label{morphisms-lemma-unramified-omega-zero}
设 $f : X \to S$ 是概形的态射。那么
\begin{enumerate}
\item $f$ 非分歧，当且仅当 $f$ 局部为有限型且
$\Omega_{X/S} = 0$；并且
\item $f$ 是 G-非分歧的，当且仅当 $f$ 局部为有限表示且
$\Omega_{X/S} = 0$。
\end{enumerate}
\end{lemma}

\begin{proof}
根据定义，环同态 $R \to A$ 非分歧（相应地，G-非分歧），当且仅当
它是有限型的（相应地，有限表示的）并且 $\Omega_{A/R} = 0$。
因此，该引理直接由定义和引理 \ref{morphisms-lemma-differentials-affine} 得出。
\end{proof}

\noindent
注意，定义中没有分离性或拟紧性假设。因此，非分歧性在源上具有
局部性质。下面是精确的结果。

\begin{lemma}
\label{morphisms-lemma-unramified-characterize}
设 $f : X \to S$ 是概形的态射。下列条件等价：
\begin{enumerate}
\item 态射 $f$ 非分歧（相应地，G-非分歧）。
\item 对每个仿射开集 $U \subset X$、$V \subset S$，若
$f(U) \subset V$，则环同态
$\mathcal{O}_S(V) \to \mathcal{O}_X(U)$ 非分歧
（相应地，G-非分歧）。
\item 存在开覆盖 $S = \bigcup_{j \in J} V_j$ 以及开覆盖
$f^{-1}(V_j) = \bigcup_{i \in I_j} U_i$，使得每个态射
$U_i \to V_j$（$j\in J, i\in I_j$）都非分歧
（相应地，G-非分歧）。
\item 存在仿射开覆盖 $S = \bigcup_{j \in J} V_j$ 以及仿射开覆盖
$f^{-1}(V_j) = \bigcup_{i \in I_j} U_i$，使得环同态
$\mathcal{O}_S(V_j) \to \mathcal{O}_X(U_i)$ 对所有
$j\in J, i\in I_j$ 都非分歧（相应地，G-非分歧）。
\end{enumerate}
此外，若 $f$ 非分歧（相应地，G-非分歧），则对任意开子概形
$U \subset X$、$V \subset S$，只要 $f(U) \subset V$，其限制
$f|_U : U \to V$ 就非分歧（相应地，G-非分歧）。
\end{lemma}

\begin{proof}
只要证明性质“$R \to A$ 非分歧”是局部性质，结论便由引理
\ref{morphisms-lemma-locally-P} 得出。我们检验定义
\ref{morphisms-definition-property-local} 的条件 (a)、(b) 和 (c)。
这些性质已在代数，引理 \ref{algebra-lemma-unramified} 中证明。
\end{proof}

\begin{lemma}
\label{morphisms-lemma-composition-unramified}
两个非分歧态射的复合仍非分歧。对 G-非分歧态射也有同样的结论。
\end{lemma}

\begin{proof}
引理 \ref{morphisms-lemma-unramified-characterize} 的证明表明，
非分歧（相应地，G-非分歧）是环同态的局部性质。
因而第一项断言由引理 \ref{morphisms-lemma-composition-type-P}
以及如下事实得出：非分歧（相应地，G-非分歧）是一个在复合下
稳定的环同态性质；见代数，引理 \ref{algebra-lemma-unramified}。
\end{proof}

\begin{lemma}
\label{morphisms-lemma-base-change-unramified}
非分歧态射的基变换仍非分歧。对 G-非分歧态射也有同样的结论。
\end{lemma}

\begin{proof}
引理 \ref{morphisms-lemma-unramified-characterize} 的证明表明，
非分歧（相应地，G-非分歧）是环同态的局部性质。
因而该引理由引理 \ref{morphisms-lemma-base-change-type-P} 以及如下事实得出：
非分歧（相应地，G-非分歧）是一个在基变换下稳定的环同态性质；
见代数，引理 \ref{algebra-lemma-unramified}。
\end{proof}

\begin{lemma}
\label{morphisms-lemma-noetherian-unramified}
设 $f : X \to S$ 是概形的态射。假设 $S$ 局部 Noether。
那么 $f$ 非分歧，当且仅当 $f$ 是 G-非分歧的。
\end{lemma}

\begin{proof}
这由定义和引理
\ref{morphisms-lemma-noetherian-finite-type-finite-presentation} 得出。
\end{proof}


\begin{lemma}
\label{morphisms-lemma-open-immersion-unramified}
任何开浸入都是 G-非分歧的。
\end{lemma}

\begin{proof}
这是因为开浸入是局部同构。
\end{proof}

\begin{lemma}
\label{morphisms-lemma-closed-immersion-unramified}
闭浸入 $i : Z \to X$ 是非分歧的。它是 G-非分歧的，当且仅当
相应的拟凝聚理想层
$\mathcal{I} = \Ker(\mathcal{O}_X \to i_*\mathcal{O}_Z)$
（作为 $\mathcal{O}_X$-模）是有限型的。
\end{lemma}

\begin{proof}
这由引理 \ref{morphisms-lemma-closed-immersion-finite-presentation} 和
代数，引理 \ref{algebra-lemma-unramified} 得出。
\end{proof}

\begin{lemma}
\label{morphisms-lemma-unramified-locally-finite-type}
非分歧态射局部为有限型。G-非分歧态射局部为有限表示。
\end{lemma}

\begin{proof}
根据定义，非分歧环同态是有限型的。
根据定义，G-非分歧环同态是有限表示的。
\end{proof}

\begin{lemma}
\label{morphisms-lemma-unramified-quasi-finite}
设 $f : X \to S$ 是概形的态射。若 $f$ 在 $x$ 处非分歧，
则 $f$ 在 $x$ 处拟有限。特别地，非分歧态射局部拟有限。
\end{lemma}

\begin{proof}
见代数，引理 \ref{algebra-lemma-unramified-quasi-finite}。
\end{proof}

\begin{lemma}
\label{morphisms-lemma-unramified-over-field}
非分歧态射的纤维具有下列性质。
\begin{enumerate}
\item 设 $X$ 是域 $k$ 上的概形。结构态射
$X \to \Spec(k)$ 非分歧，当且仅当 $X$ 是 $k$ 的有限可分
域扩张的谱的不交并。
\item 若 $f : X \to S$ 是非分歧态射，则对每个 $s \in S$，
纤维 $X_s$ 都是 $\kappa(s)$ 的有限可分域扩张的谱的不交并。
\end{enumerate}
\end{lemma}

\begin{proof}
第 (2) 项由第 (1) 项和引理
\ref{morphisms-lemma-base-change-unramified} 得出。我们来证明第 (1) 项。
先使用代数，引理 \ref{algebra-lemma-characterize-unramified}。
该引理说明，若 $X$ 是 $k$ 的有限可分域扩张的谱的不交并，
则 $X \to \Spec(k)$ 非分歧。反过来，假设
$X \to \Spec(k)$ 非分歧。由代数，引理
\ref{algebra-lemma-unramified-at-prime}，对每个 $x \in X$，
剩余域扩张 $\kappa(x)/k$ 都是有限可分的。由于
$X \to \Spec(k)$ 局部拟有限（引理
\ref{morphisms-lemma-unramified-quasi-finite}），由引理
\ref{morphisms-lemma-quasi-finite-at-point-characterize} 可知，$X$ 的所有点
都是孤立闭点。因此 $X$ 是离散空间，特别地，它是其局部环的谱的
不交并。再次使用代数，引理
\ref{algebra-lemma-unramified-at-prime}，这些局部环都是域，证毕。
\end{proof}

\noindent
下面的引理将非分歧态射刻画为具有非分歧纤维的局部有限型态射。

\begin{lemma}
\label{morphisms-lemma-unramified-etale-fibres}
设 $f : X \to S$ 是概形的态射。
\begin{enumerate}
\item 若 $f$ 非分歧，则对任意 $x \in X$，域扩张
$\kappa(x)/\kappa(f(x))$ 是有限可分的。
\item 若 $f$ 局部为有限型，并且对每个 $s \in S$，纤维
$X_s$ 都是 $\kappa(s)$ 的有限可分域扩张的谱的不交并，
则 $f$ 非分歧。
\item 若 $f$ 局部为有限表示，并且对每个 $s \in S$，纤维
$X_s$ 都是 $\kappa(s)$ 的有限可分域扩张的谱的不交并，
则 $f$ 是 G-非分歧的。
\end{enumerate}
\end{lemma}

\begin{proof}
这由代数中的引理
\ref{algebra-lemma-unramified-at-prime} 和
\ref{algebra-lemma-characterize-unramified} 得出。
\end{proof}

\noindent
下面用对角态射给出非分歧态射的一种刻画。

\begin{lemma}
\label{morphisms-lemma-diagonal-unramified-morphism}
设 $f : X \to S$ 是一个态射。
\begin{enumerate}
\item 若 $f$ 非分歧，则对角态射
$\Delta : X \to X \times_S X$ 是开浸入。
\item 若 $f$ 局部为有限型且 $\Delta$ 是开浸入，则 $f$ 非分歧。
\item 若 $f$ 局部为有限表示且 $\Delta$ 是开浸入，则 $f$
是 G-非分歧的。
\end{enumerate}
\end{lemma}

\begin{proof}
第一项断言由代数，引理
\ref{algebra-lemma-diagonal-unramified} 得出。
第二项断言来自如下事实：$\Omega_{X/S}$ 是对角态射的余法层
（引理 \ref{morphisms-lemma-differentials-diagonal}），因而若
$\Delta$ 是开浸入，它显然为零。
\end{proof}

\begin{lemma}
\label{morphisms-lemma-unramified-at-point}
设 $f : X \to S$ 是概形的态射。设 $x \in X$。
令 $s = f(x)$。假设 $f$ 局部为有限型
（相应地，局部为有限表示）。下列条件等价：
\begin{enumerate}
\item 态射 $f$ 在 $x$ 处非分歧（相应地，G-非分歧）。
\item 纤维 $X_s$ 作为 $\kappa(s)$ 上的概形在 $x$ 处非分歧。
\item $\mathcal{O}_{X, x}$-模 $\Omega_{X/S, x}$ 为零。
\item $\mathcal{O}_{X_s, x}$-模 $\Omega_{X_s/s, x}$ 为零。
\item $\kappa(x)$-向量空间
$$
\Omega_{X_s/s, x} \otimes_{\mathcal{O}_{X_s, x}} \kappa(x) =
\Omega_{X/S, x} \otimes_{\mathcal{O}_{X, x}} \kappa(x)
$$
为零。
\item 有 $\mathfrak m_s\mathcal{O}_{X, x} = \mathfrak m_x$，
并且域扩张 $\kappa(x)/\kappa(s)$ 是有限可分的。
\end{enumerate}
\end{lemma}

\begin{proof}
注意，若 $f$ 在 $x$ 处非分歧，则由定义以及
第 \ref{morphisms-section-sheaf-differentials} 节关于微分模的结果可知，
$\Omega_{X/S} = 0$ 在 $x$ 的某个邻域中成立。因此 (1) 蕴含
(3)，也蕴含 (5) 中右侧向量空间的消失。它还蕴含 (2)，因为由
引理 \ref{morphisms-lemma-base-change-differentials}，微分模
$\Omega_{X_s/s}$（即纤维 $X_s$ 在 $\kappa(s)$ 上的微分模）
是微分模 $\Omega_{X/S}$（即 $X$ 在 $S$ 上的微分模）的拉回。
这个关于微分模的事实还蕴含
第 (4) 项中所展示的向量空间等式。由引理
\ref{morphisms-lemma-finite-type-differentials}，模
$\Omega_{X/S, x}$ 和 $\Omega_{X_s/s, x}$ 都是有限型的。
因此，由 Nakayama 引理（代数，引理
\ref{algebra-lemma-NAK}），模
$\Omega_{X/S, x}$ 和 $\Omega_{X_s/s, x}$ 为零，当且仅当
(4) 中相应的 $\kappa(x)$-向量空间为零。
特别地，这表明 (3)、(4) 和 (5) 等价。$\Omega_{X/S}$ 的支集
在 $X$ 中闭；见模，引理
\ref{modules-lemma-support-finite-type-closed}。假设 (3) 蕴含
$x$ 不在该支集中。因此 $\Omega_{X/S}$ 在 $x$ 的某个邻域中为零，
这蕴含 (1)。将 (1) 与 (3) 的等价性应用于 $X_s \to s$，
便得到 (2) 与 (4) 的等价性。至此，我们已经看到
(1) -- (5) 等价。

\medskip\noindent
或者，可以使用代数，引理 \ref{algebra-lemma-unramified}
更直接地看出 (1) -- (5) 的等价性。

\medskip\noindent
(1) 与 (6) 的等价性由引理
\ref{morphisms-lemma-unramified-etale-fibres} 得出。也可更直接地由
代数中的引理 \ref{algebra-lemma-unramified-at-prime} 和
\ref{algebra-lemma-characterize-unramified} 得出。
\end{proof}

\begin{lemma}
\label{morphisms-lemma-set-points-where-fibres-unramified}
设 $f : X \to S$ 是概形的态射。假设 $f$ 局部为有限型。
下列开集的形成
\begin{align*}
T
& =
\{x \in X \mid X_{f(x)}\text{ is unramified over }\kappa(f(x))\text{ 在 }x\} \\
& =
\{x \in X \mid X\text{ is unramified over }S\text{ 在 }x\}
\end{align*}
与任意基变换交换：对任意态射 $g : S' \to S$，考虑基变换
$f' : X' \to S'$（它是 $f$ 的基变换）以及投影
$g' : X' \to X$。则相应集合 $T'$（对于态射 $f'$）等于
$T' = (g')^{-1}(T)$。
若假设 $f$ 局部为有限表示，则对 $f$ 为 G-非分歧的点组成的
开集也有同样的结论。
\end{lemma}

\begin{proof}
设 $s' \in S'$ 是一点，并令 $s = g(s')$。于是
$$
X'_{s'} =
\Spec(\kappa(s')) \times_{\Spec(\kappa(s))} X_s
$$
换言之，基变换的纤维是纤维的基变换。特别地，
$$
\Omega_{X_s/s, x} \otimes_{\mathcal{O}_{X_s, x}} \kappa(x')
=
\Omega_{X'_{s'}/s', x'} \otimes_{\mathcal{O}_{X'_{s'}, x'}} \kappa(x')
$$
见引理 \ref{morphisms-lemma-base-change-differentials}。于是由引理
\ref{morphisms-lemma-unramified-at-point}，$x' \in T'$ 当且仅当
$x \in T$。第二项断言由第一项得出，因为在这种情况下，
根据引理 \ref{morphisms-lemma-unramified-at-point}，$T$ 正是 $f$
为 G-非分歧的点组成的（开）集。
\end{proof}

\begin{lemma}
\label{morphisms-lemma-unramified-permanence}
设 $f : X \to Y$ 是 $S$ 上概形的态射。
\begin{enumerate}
\item 若 $X$ 在 $S$ 上非分歧，则 $f$ 非分歧。
\item 若 $X$ 在 $S$ 上 G-非分歧，且 $Y$ 在 $S$ 上局部为有限型，
则 $f$ 是 G-非分歧的。
\end{enumerate}
\end{lemma}

\begin{proof}
假设 $X$ 在 $S$ 上非分歧。由引理
\ref{morphisms-lemma-permanence-finite-type} 可知 $f$ 局部为有限型。
根据假设有 $\Omega_{X/S} = 0$。因而由引理
\ref{morphisms-lemma-triangle-differentials}，$\Omega_{X/Y} = 0$。
所以 $f$ 非分歧。若 $X$ 在 $S$ 上 G-非分歧且 $Y$ 在 $S$
上局部为有限型，则由引理
\ref{morphisms-lemma-finite-presentation-permanence} 可知 $f$
局部为有限表示，从而 $f$ 是 G-非分歧的。
\end{proof}

\begin{lemma}
\label{morphisms-lemma-value-at-one-point}
设 $S$ 为概形。设 $X$、$Y$ 是 $S$ 上的概形。
设 $f, g : X \to Y$ 是 $S$ 上的态射。设 $x \in X$。假设
\begin{enumerate}
\item 结构态射 $Y \to S$ 非分歧；
\item 有 $f(x) = g(x)$（在 $Y$ 中），记 $y = f(x) = g(x)$；并且
\item 所诱导的映射 $f^\sharp, g^\sharp : \kappa(y) \to \kappa(x)$
相等。
\end{enumerate}
则 $x$ 在 $X$ 中有一个开邻域，使得 $f$ 与 $g$ 在其上相等。
\end{lemma}

\begin{proof}
考虑态射 $(f, g) : X \to Y \times_S Y$。由假设 (1) 和引理
\ref{morphisms-lemma-diagonal-unramified-morphism}，
$\Delta_{Y/S}(Y)$ 的逆像在 $X$ 中是开集。
而假设 (2) 和 (3) 蕴含 $x$ 位于这个开子集中。
\end{proof}

\section{平展态射}
\label{morphisms-section-etale}

\noindent
概形的 Zariski 拓扑是一种非常粗糙的拓扑。考察
$\mathbf{C}$ 上的簇时，这一点尤其明显。事实表明，把平展态射
规定为拓扑学中局部同构的类似物，会引入一种细得多的拓扑。
在 $\mathbf{C}$ 上的簇上，用有限系数计算上同调时，这种拓扑
给出“正确的”Betti 数。另一个可观察的现象是：若
$f : X \to Y$ 是 $\mathbf{C}$ 上簇的平展态射，而 $x$
是 $X$ 的闭点，则 $f$ 诱导完备局部环的同构
$\mathcal{O}^{\wedge}_{Y, f(x)} \to \mathcal{O}^{\wedge}_{X, x}$。

\medskip\noindent
本节开始研究这些问题。事实上，我们有意把讨论限制在最低限度，
因为更有趣的结果将在别处讨论。回顾：若环同态 $R \to A$
光滑且 $\Omega_{A/R} = 0$，则称其{\it 平展}；见
代数，定义 \ref{algebra-definition-etale}。

\begin{definition}
\label{morphisms-definition-etale}
设 $f : X \to S$ 是概形的态射。
\begin{enumerate}
\item 称 $f$ {\it 在 $x \in X$ 处平展}，如果存在仿射开邻域
$\Spec(A) = U \subset X$ 包含 $x$，以及仿射开集
$\Spec(R) = V \subset S$，满足 $f(U) \subset V$，并且所诱导的
环同态 $R \to A$ 平展。
\item 若 $f$ 在 $X$ 的每一点处平展，则称其{\it 平展}。
\item 仿射概形的态射 $f : X \to S$ 称为{\it 标准平展的}，
如果 $X \to S$ 同构于
$$
\Spec(R[x]_h/(g)) \to \Spec(R)
$$
其中 $R \to R[x]_h/(g)$ 是标准平展环同态；见
代数，定义 \ref{algebra-definition-standard-etale}；也就是说，
$g$ 是首一的，并且 $g'$ 在 $R[x]_h/(g)$ 中可逆。
\end{enumerate}
\end{definition}

\noindent
一个态射平展，当且仅当它是相对维数 $0$ 的光滑态射
（见定义 \ref{morphisms-definition-smooth-relative-dimension}）。
该定义的一个可取之处是：态射平展的点集自动是开集。

\medskip\noindent
注意，定义中没有分离性或拟紧性假设。因此，平展性在源上具有
局部性质。下面是精确的结果。

\begin{lemma}
\label{morphisms-lemma-etale-characterize}
设 $f : X \to S$ 是概形的态射。下列条件等价：
\begin{enumerate}
\item 态射 $f$ 平展。
\item 对每个仿射开集 $U \subset X$、$V \subset S$，若
$f(U) \subset V$，则环同态
$\mathcal{O}_S(V) \to \mathcal{O}_X(U)$ 平展。
\item 存在开覆盖 $S = \bigcup_{j \in J} V_j$ 以及开覆盖
$f^{-1}(V_j) = \bigcup_{i \in I_j} U_i$，使得每个态射
$U_i \to V_j$（$j\in J, i\in I_j$）都平展。
\item 存在仿射开覆盖 $S = \bigcup_{j \in J} V_j$ 以及仿射开覆盖
$f^{-1}(V_j) = \bigcup_{i \in I_j} U_i$，使得环同态
$\mathcal{O}_S(V_j) \to \mathcal{O}_X(U_i)$ 对所有
$j\in J, i\in I_j$ 都平展。
\end{enumerate}
此外，若 $f$ 平展，则对任意开子概形 $U \subset X$、
$V \subset S$，只要 $f(U) \subset V$，其限制
$f|_U : U \to V$ 就平展。
\end{lemma}

\begin{proof}
只要证明性质“$R \to A$ 平展”是局部性质，结论便由引理
\ref{morphisms-lemma-locally-P} 得出。我们检验定义
\ref{morphisms-definition-property-local} 的条件 (a)、(b) 和 (c)。
这些条件都由代数，引理 \ref{algebra-lemma-etale} 得出。
\end{proof}

\begin{lemma}
\label{morphisms-lemma-composition-etale}
两个平展态射的复合仍平展。
\end{lemma}

\begin{proof}
在引理 \ref{morphisms-lemma-etale-characterize} 的证明中，我们已经看到，
平展是环同态的局部性质。因此，该引理的第一项断言由引理
\ref{morphisms-lemma-composition-type-P} 以及如下事实得出：
平展是一个在复合下稳定的环同态性质；见
代数，引理 \ref{algebra-lemma-etale}。
\end{proof}

\begin{lemma}
\label{morphisms-lemma-base-change-etale}
平展态射的基变换仍平展。
\end{lemma}

\begin{proof}
在引理 \ref{morphisms-lemma-etale-characterize} 的证明中，我们已经看到，
平展是环同态的局部性质。因此，该引理由引理
\ref{morphisms-lemma-composition-type-P} 以及如下事实得出：
平展是一个在基变换下稳定的环同态性质；见
代数，引理 \ref{algebra-lemma-etale}。
\end{proof}

\begin{lemma}
\label{morphisms-lemma-etale-smooth-unramified}
设 $f : X \to S$ 是概形的态射。设 $x \in X$。那么 $f$
在 $x$ 处平展，当且仅当 $f$ 在 $x$ 处既光滑又非分歧。
\end{lemma}

\begin{proof}
这直接由定义得出。
\end{proof}

\begin{lemma}
\label{morphisms-lemma-etale-locally-quasi-finite}
平展态射局部拟有限。
\end{lemma}

\begin{proof}
由引理 \ref{morphisms-lemma-etale-smooth-unramified}，平展态射非分歧。
由引理 \ref{morphisms-lemma-unramified-quasi-finite}，非分歧态射局部拟有限。
\end{proof}

\begin{lemma}
\label{morphisms-lemma-etale-over-field}
\begin{slogan}
域上的平展概形和平展态射的纤维的描述。
\end{slogan}
平展态射的纤维具有下列性质。
\begin{enumerate}
\item 设 $X$ 是域 $k$ 上的概形。结构态射
$X \to \Spec(k)$ 平展，当且仅当 $X$ 是 $k$
的有限可分域扩张的谱的不交并。
\item 若 $f : X \to S$ 是平展态射，则对每个 $s \in S$，
纤维 $X_s$ 都是 $\kappa(s)$ 的有限可分域扩张的谱的不交并。
\end{enumerate}
\end{lemma}

\begin{proof}
可以从引理 \ref{morphisms-lemma-unramified-over-field} 推出此结论，
其中使用上面的引理 \ref{morphisms-lemma-etale-smooth-unramified}。
下面给出一个直接证明。

\medskip\noindent
我们将使用代数，引理 \ref{algebra-lemma-etale-over-field}。
因此显然，若 $X$ 是 $k$ 的有限可分域扩张的谱的不交并，
则 $X \to \Spec(k)$ 平展。反过来，假设
$X \to \Spec(k)$ 平展。那么对任意仿射开集 $U \subset X$，
$U$ 都是 $k$ 的有限可分域扩张的谱的有限不交并。
因而 $X$ 的所有点都是闭点（例如见引理
\ref{morphisms-lemma-algebraic-residue-field-extension-closed-point-fibre}）。
所以 $X$ 是离散空间，证毕。
\end{proof}

\noindent
下面的引理将平展态射刻画为具有“平展纤维”的平坦有限表示态射。

\begin{lemma}
\label{morphisms-lemma-etale-flat-etale-fibres}
设 $f : X \to S$ 是概形的态射。若 $f$ 平坦、局部为有限表示，
并且对每个 $s \in S$，纤维 $X_s$ 都是 $\kappa(s)$
的有限可分域扩张的谱的不交并，则 $f$ 平展。
\end{lemma}

\begin{proof}
可以从代数，引理 \ref{algebra-lemma-characterize-etale}
推出此结论。下面给出另一个证明。

\medskip\noindent
由引理 \ref{morphisms-lemma-etale-over-field}，纤维 $X_s$ 平展，
从而在 $s$ 上光滑。由引理
\ref{morphisms-lemma-smooth-flat-smooth-fibres} 可知 $X \to S$ 光滑。
由引理 \ref{morphisms-lemma-unramified-etale-fibres} 可知 $f$ 非分歧。
最后应用引理 \ref{morphisms-lemma-etale-smooth-unramified}。
\end{proof}

\begin{lemma}
\label{morphisms-lemma-open-immersion-etale}
任何开浸入都平展。
\end{lemma}

\begin{proof}
这是因为开浸入是局部同构。
\end{proof}

\begin{lemma}
\label{morphisms-lemma-etale-syntomic}
平展态射是 syntomic 态射。
\end{lemma}

\begin{proof}
见代数，引理 \ref{algebra-lemma-smooth-syntomic}，并使用如下事实：
平展态射恰好是相对维数 $0$ 的光滑态射。
\end{proof}

\begin{lemma}
\label{morphisms-lemma-etale-locally-finite-presentation}
平展态射局部为有限表示。
\end{lemma}

\begin{proof}
这是因为根据定义，平展环同态是有限表示的。
\end{proof}

\begin{lemma}
\label{morphisms-lemma-etale-flat}
平展态射是平坦态射。
\end{lemma}

\begin{proof}
结合引理 \ref{morphisms-lemma-syntomic-flat} 和
\ref{morphisms-lemma-etale-syntomic}。
\end{proof}

\begin{lemma}
\label{morphisms-lemma-etale-open}
平展态射是开态射。
\end{lemma}

\begin{proof}
结合引理 \ref{morphisms-lemma-etale-flat}、
\ref{morphisms-lemma-etale-locally-finite-presentation} 和
\ref{morphisms-lemma-fppf-open}。
\end{proof}

\noindent
下面的引理说明，局部地，每个平展态射都是标准平展态射。
要在完全一般的情形下证明这个结果，实际上颇为棘手。
棘手的部分隐藏在交换代数一章中。因此，标准平展态射是一般
平展态射的一个{\it 局部模型}。

\begin{lemma}
\label{morphisms-lemma-etale-locally-standard-etale}
设 $f : X  \to S$ 是概形的态射。设 $x \in X$ 是一点。
设 $V \subset S$ 是 $f(x)$ 的仿射开邻域。下列条件等价：
\begin{enumerate}
\item 态射 $f$ 在 $x$ 处平展。
\item 存在仿射开集 $U \subset X$，满足 $x \in U$ 和
$f(U) \subset V$，使得所诱导的态射 $f|_U : U \to V$
是标准平展的（见定义 \ref{morphisms-definition-etale}）。
\end{enumerate}
\end{lemma}

\begin{proof}
这由定义和代数，命题
\ref{algebra-proposition-etale-locally-standard} 得出。
\end{proof}

\noindent
下面给出一点处平展性的微分判据。这个结果有许多变体，
它们各自可能在某处有用。我们只在需要时把它们添加在这里。

\begin{lemma}
\label{morphisms-lemma-etale-at-point}
设 $f : X \to S$ 是概形的态射。设 $x \in X$。
令 $s = f(x)$。假设 $f$ 局部为有限表示。下列条件等价：
\begin{enumerate}
\item 态射 $f$ 在 $x$ 处平展。
\item 局部环同态 $\mathcal{O}_{S, s} \to \mathcal{O}_{X, x}$
平坦，并且 $X_s \to \Spec(\kappa(s))$ 在 $x$ 处平展。
\item 局部环同态 $\mathcal{O}_{S, s} \to \mathcal{O}_{X, x}$
平坦，并且 $X_s \to \Spec(\kappa(s))$ 在 $x$ 处非分歧。
\item 局部环同态 $\mathcal{O}_{S, s} \to \mathcal{O}_{X, x}$
平坦，并且 $\mathcal{O}_{X, x}$-模 $\Omega_{X/S, x}$ 为零。
\item 局部环同态 $\mathcal{O}_{S, s} \to \mathcal{O}_{X, x}$
平坦，并且 $\kappa(x)$-向量空间
$$
\Omega_{X_s/s, x} \otimes_{\mathcal{O}_{X_s, x}} \kappa(x) =
\Omega_{X/S, x} \otimes_{\mathcal{O}_{X, x}} \kappa(x)
$$
为零。
\item 局部环同态 $\mathcal{O}_{S, s} \to \mathcal{O}_{X, x}$
平坦，且有 $\mathfrak m_s\mathcal{O}_{X, x} = \mathfrak m_x$，
并且域扩张 $\kappa(x)/\kappa(s)$ 是有限可分的。
\item 存在仿射开集 $U \subset X$ 和 $V \subset S$，满足
$x \in U$、$f(U) \subset V$，且所诱导的态射
$f|_U : U \to V$ 是相对维数 $0$ 的标准光滑态射。
\item 存在仿射开集 $\Spec(A) = U \subset X$ 和
$\Spec(R) = V \subset S$，其中 $x \in U$ 对应于
$\mathfrak q \subset A$，且 $f(U) \subset V$，使得存在表示
$$
A = R[x_1, \ldots, x_n]/(f_1, \ldots, f_n)
$$
并且
$$
g =
\det
\left(
\begin{matrix}
\partial f_1/\partial x_1 &
\partial f_2/\partial x_1 &
\ldots &
\partial f_n/\partial x_1 \\
\partial f_1/\partial x_2 &
\partial f_2/\partial x_2 &
\ldots &
\partial f_n/\partial x_2 \\
\ldots & \ldots & \ldots & \ldots \\
\partial f_1/\partial x_n &
\partial f_2/\partial x_n &
\ldots &
\partial f_n/\partial x_n
\end{matrix}
\right)
$$
在 $A$ 中的像不属于 $\mathfrak q$。
\item 存在仿射开集 $U \subset X$ 和 $V \subset S$，满足
$x \in U$、$f(U) \subset V$，且所诱导的态射
$f|_U : U \to V$ 是标准平展的。
\item 存在仿射开集 $\Spec(A) = U \subset X$ 和
$\Spec(R) = V \subset S$，其中 $x \in U$ 对应于
$\mathfrak q \subset A$，且 $f(U) \subset V$，使得存在表示
$$
A = R[x]_Q/(P) = R[x, 1/Q]/(P)
$$
其中 $P, Q \in R[x]$，$P$ 是首一的，并且
$P' = \text{d}P/\text{d}x$ 在 $A$ 中的像不属于
$\mathfrak q$。
\end{enumerate}
\end{lemma}

\begin{proof}
使用引理 \ref{morphisms-lemma-etale-locally-standard-etale} 和定义，
可见 (1) 蕴含其余所有条件。对于条件 (2) -- (10) 中的每一个，
结合引理 \ref{morphisms-lemma-smooth-at-point} 和
\ref{morphisms-lemma-unramified-at-point}，通过证明 $f$ 在 $x$
处既光滑又非分歧，再应用引理
\ref{morphisms-lemma-etale-smooth-unramified}，可见 (1) 成立。
省略若干细节。
\end{proof}

\begin{lemma}
\label{morphisms-lemma-flat-unramified-etale}
一个态射在一点处平展，当且仅当它在该点处平坦且 G-非分歧。
一个态射平展，当且仅当它平坦且 G-非分歧。
\end{lemma}

\begin{proof}
这由引理 \ref{morphisms-lemma-etale-at-point} 和
\ref{morphisms-lemma-unramified-at-point} 得出。
\end{proof}

\begin{lemma}
\label{morphisms-lemma-set-points-where-fibres-etale}
设
$$
\xymatrix{
X' \ar[r]_{g'} \ar[d]_{f'} & X \ar[d]^f \\
S' \ar[r]^g & S
}
$$
是概形的笛卡尔图。设 $W \subset X$（相应地，
$W' \subset X'$）是 $f$（相应地，$f'$）平展的点组成的开子概形。
若下列任一条件成立，则 $W' = (g')^{-1}(W)$：
\begin{enumerate}
\item $f$ 平坦且局部为有限表示；或
\item $f$ 局部为有限表示且 $g$ 平坦。
\end{enumerate}
\end{lemma}

\begin{proof}
先假设 $f$ 局部为有限型。考虑集合
$$
T = \{x \in X \mid f\text{ 在下列点非分歧： }x\}
$$
以及相应集合 $T' \subset X'$（对于 $f'$）。由引理
\ref{morphisms-lemma-set-points-where-fibres-unramified}，
$T' = (g')^{-1}(T)$。

\medskip\noindent
因此情形 (1) 成立，因为在情形 (1) 中，由引理
\ref{morphisms-lemma-flat-unramified-etale}，$T$ 正是 $f$ 平展的点组成的
（开）集。

\medskip\noindent
在情形 (2) 中，设 $x' \in W'$。那么 $g'$ 在 $x'$ 处平坦
（引理 \ref{morphisms-lemma-base-change-module-flat}），并且
$g \circ f'$ 在 $x'$ 处平坦（引理
\ref{morphisms-lemma-composition-module-flat}）。由引理
\ref{morphisms-lemma-flat-permanence} 可知，$f$ 在 $x = g'(x')$ 处平坦。
另一方面，由于 $x' \in T'$（引理
\ref{morphisms-lemma-base-change-smooth}），可知 $x \in T$。因此由引理
\ref{morphisms-lemma-etale-at-point}，$f$ 在 $x$ 处平展。
\end{proof}

\begin{lemma}
\label{morphisms-lemma-etale-permanence-better}
\begin{reference}
例如见 \cite[Remark 18.30 (2).]{GWII}。
\end{reference}
设 $f : X \to Y$ 是 $S$ 上概形的态射。若 $X$ 在 $S$ 上平展，
且 $Y$ 在 $S$ 上非分歧，则 $f$ 平展。
\end{lemma}

\begin{proof}
考虑分解 $X \to X \times_S Y \to Y$，其中第一支箭头由
$\text{id}_X$ 和 $f$ 给出，第二支箭头是投影。我们断言两支箭头
都平展，因而由引理 \ref{morphisms-lemma-composition-etale}，$f$ 平展。
确实，该投影平展，因为它是 $X \to S$ 的基变换；见引理
\ref{morphisms-lemma-base-change-etale}。第一支箭头是对角态射
$Y \to Y \times_S Y$ 的基变换，因为方块
$$
\xymatrix{
X \ar[d] \ar[r] & X \times_S Y \ar[d] \\
Y \ar[r] & Y \times_S Y
}
$$
是笛卡尔的。由引理 \ref{morphisms-lemma-diagonal-unramified-morphism}，
对角态射 $Y \to Y \times_S Y$ 是开浸入。开浸入的基变换仍是
开浸入（《概形》，引理
\ref{schemes-lemma-base-change-immersion}），而开浸入平展
（引理 \ref{morphisms-lemma-open-immersion-etale}）。
\end{proof}

\begin{lemma}
\label{morphisms-lemma-etale-permanence}
\begin{slogan}
平展态射的消去律。
\end{slogan}
设 $f : X \to Y$ 是 $S$ 上概形的态射。若 $X$ 和 $Y$
都在 $S$ 上平展，则 $f$ 平展。
\end{lemma}

\begin{proof}
由引理 \ref{morphisms-lemma-etale-smooth-unramified}，$Y \to S$
非分歧，所以可以应用引理
\ref{morphisms-lemma-etale-permanence-better}。另见
代数，引理 \ref{algebra-lemma-map-between-etale}。
\end{proof}

\begin{lemma}
\label{morphisms-lemma-etale-permanence-two}
设
$$
\xymatrix{
X \ar[rr]_f \ar[rd]_p & &
Y \ar[dl]^q \\
& S
}
$$
是概形态射的交换图。假设
\begin{enumerate}
\item $f$ 满射且平展；
\item $p$ 平展；并且
\item $q$ 局部为有限表示\footnote{事实上，这由 (1) 和 (2)
蕴含；见下降，引理
\ref{descent-lemma-flat-finitely-presented-permanence}。
此外，只需假设 $f$ 满射、平坦且局部为有限表示；见
下降，引理 \ref{descent-lemma-smooth-permanence}。}。
\end{enumerate}
那么 $q$ 平展。
\end{lemma}

\begin{proof}
由引理 \ref{morphisms-lemma-smooth-permanence} 可知 $q$ 光滑。
因此只需证明 $q$ 的相对维数为 $0$。这由引理
\ref{morphisms-lemma-dimension-fibre-at-a-point-additive} 以及如下事实得出：
$f$ 和 $p$ 的相对维数都是 $0$。
\end{proof}

\noindent
光滑态射的最后一种刻画是：光滑态射 $f : X \to S$
局部地是一个平展态射与投影 $\mathbf{A}_S^d \to S$ 的复合。

\begin{lemma}
\label{morphisms-lemma-smooth-etale-over-affine-space}
\begin{slogan}
光滑概形在平展局部形如仿射空间。
\end{slogan}
设 $\varphi : X \to Y$ 是概形的态射。设 $x \in X$。
设 $V \subset Y$ 是 $\varphi(x)$ 的仿射开邻域。
若 $\varphi$ 在 $x$ 处光滑，则存在整数 $d \geq 0$
和仿射开集 $U \subset X$，满足 $x \in U$ 和
$\varphi(U) \subset V$，并且存在交换图
$$
\xymatrix{
X \ar[d] & U \ar[l] \ar[d] \ar[r]_-\pi & \mathbf{A}^d_V \ar[ld] \\
Y & V \ar[l]
}
$$
其中 $\pi$ 平展。
\end{lemma}

\begin{proof}
由引理 \ref{morphisms-lemma-smooth-locally-standard-smooth}，可以找到引理中
所述的仿射开集 $U$，使得 $\varphi|_U : U \to V$ 标准光滑。
写成 $U = \Spec(A)$ 和 $V = \Spec(R)$，从而可以写
$$
A = R[x_1, \ldots, x_n]/(f_1, \ldots, f_c)
$$
并且
$$
g =
\det
\left(
\begin{matrix}
\partial f_1/\partial x_1 &
\partial f_2/\partial x_1 &
\ldots &
\partial f_c/\partial x_1 \\
\partial f_1/\partial x_2 &
\partial f_2/\partial x_2 &
\ldots &
\partial f_c/\partial x_2 \\
\ldots & \ldots & \ldots & \ldots \\
\partial f_1/\partial x_c &
\partial f_2/\partial x_c &
\ldots &
\partial f_c/\partial x_c
\end{matrix}
\right)
$$
在 $A$ 中的像可逆。那么显然
$R[x_{c + 1}, \ldots, x_n] \to A$ 是相对维数 $0$
的标准光滑环同态。因此它是相对维数 $0$ 的光滑环同态。
换言之，环同态 $R[x_{c + 1}, \ldots, x_n] \to A$ 平展。
由于 $\mathbf{A}^{n - c}_V = \Spec(R[x_{c + 1}, \ldots, x_n])$，
取 $d = n - c$ 即得本引理。
\end{proof}

\section{相对丰沛层}
\label{morphisms-section-relatively-ample}

\noindent
设 $X$ 是概形，$\mathcal{L}$ 是 $X$ 上的可逆层。
那么 $\mathcal{L}$ 在 $X$ 上丰沛，如果 $X$ 拟紧，并且
$X$ 的每一点都包含在形如 $X_s$ 的仿射开集中，其中
$s \in \Gamma(X, \mathcal{L}^{\otimes n})$ 且 $n \geq 1$；见性质，定义
\ref{properties-definition-ample}。我们如下把它变成一个相对概念。

\begin{definition}
\label{morphisms-definition-relatively-ample}
\begin{reference}
\cite[II Definition 4.6.1]{EGA}
\end{reference}
设 $f : X \to S$ 是概形的态射。设 $\mathcal{L}$ 是可逆
$\mathcal{O}_X$-模。若 $\mathcal{L}$ {\it 相对丰沛}，
也称{\it $f$-相对丰沛}、{\it 在 $X/S$ 上丰沛}或
{\it $f$-丰沛}，其含义是：$f : X \to S$ 拟紧，并且对每个
仿射开集 $V \subset S$，$\mathcal{L}$ 在开子概形
$f^{-1}(V)$（它是 $X$ 的开子概形）上的限制是丰沛的。
\end{definition}

\noindent
注意，$X$ 上相对丰沛层的存在性并不迫使态射 $X \to S$
是有限型的。

\begin{lemma}
\label{morphisms-lemma-ample-power-ample}
设 $X \to S$ 是概形的态射。设 $\mathcal{L}$ 是可逆
$\mathcal{O}_X$-模。设 $n \geq 1$。那么
$\mathcal{L}$ 是 $f$-丰沛的，当且仅当
$\mathcal{L}^{\otimes n}$ 是 $f$-丰沛的。
\end{lemma}

\begin{proof}
这由性质，引理 \ref{properties-lemma-ample-power-ample} 得出。
\end{proof}

\begin{lemma}
\label{morphisms-lemma-relatively-ample-separated}
设 $f : X \to S$ 是概形的态射。若存在一个 $f$-丰沛可逆层，
则 $f$ 分离。
\end{lemma}

\begin{proof}
分离性在基上是局部的（例如见《概形》，引理
\ref{schemes-lemma-characterize-separated}；这也很容易直接由定义
得出）。因此可以假设 $S$ 仿射，并且 $X$ 有丰沛可逆层。
在这种情况下，结果由性质，引理
\ref{properties-lemma-ample-separated} 得出。
\end{proof}

\noindent
相对丰沛可逆层有许多种刻画方式，类似于性质，命题
\ref{properties-proposition-characterize-ample} 中的等价条件。
我们将在需要时把这些刻画添加在这里。

\begin{lemma}
\label{morphisms-lemma-characterize-relatively-ample}
\begin{reference}
\cite[II, Proposition 4.6.3]{EGA}
\end{reference}
设 $f : X \to S$ 是概形的拟紧态射。设 $\mathcal{L}$
是 $X$ 上的可逆层。下列条件等价：
\begin{enumerate}
\item 可逆层 $\mathcal{L}$ 是 $f$-丰沛的。
\item 存在开覆盖 $S = \bigcup V_i$，使得每个
$\mathcal{L}|_{f^{-1}(V_i)}$ 相对于
$f^{-1}(V_i) \to V_i$ 都是丰沛的。
\item 存在仿射开覆盖 $S = \bigcup V_i$，使得每个
$\mathcal{L}|_{f^{-1}(V_i)}$ 都丰沛。
\item 存在拟凝聚分次 $\mathcal{O}_S$-代数 $\mathcal{A}$
以及分次 $\mathcal{O}_X$-代数的映射
$\psi : f^*\mathcal{A} \to \bigoplus_{d \geq 0} \mathcal{L}^{\otimes d}$，
使得 $U(\psi) = X$，并且
$$
r_{\mathcal{L}, \psi} :
X
\longrightarrow
\underline{\text{Proj}}_S(\mathcal{A})
$$
是开浸入（记号见构造，引理
\ref{constructions-lemma-invertible-map-into-relative-proj}）。
\item 态射 $f$ 拟分离，并且上面的第 (4) 项在
$\mathcal{A} = f_*(\bigoplus_{d \geq 0} \mathcal{L}^{\otimes d})$
和 $\psi$ 为伴随映射时成立。
\item 与 (4) 相同，但只要求 $r_{\mathcal{L}, \psi}$ 是浸入。
\end{enumerate}
\end{lemma}

\begin{proof}
由定义立即可知 (1) 蕴含 (2)，而 (2) 蕴含 (3)。
显然 (5) 蕴含 (4)。

\medskip\noindent
假设 (3) 对仿射开覆盖 $S = \bigcup V_i$ 成立。
我们将证明 (5) 成立。由于每个 $f^{-1}(V_i)$ 都有丰沛可逆层，
可知 $f^{-1}(V_i)$ 分离（性质，引理
\ref{properties-lemma-ample-separated}）。因而 $f$ 分离。
由《概形》，引理 \ref{schemes-lemma-push-forward-quasi-coherent}，
可知 $\mathcal{A} = f_*(\bigoplus_{d \geq 0} \mathcal{L}^{\otimes d})$
是拟凝聚分次 $\mathcal{O}_S$-代数。用
$\psi : f^*\mathcal{A} \to \bigoplus_{d \geq 0} \mathcal{L}^{\otimes d}$
表示伴随映射。构造，第
\ref{constructions-section-invertible-relative-proj} 节中对开集
$U(\psi)$ 的描述，以及 $\mathcal{L}|_{f^{-1}(V_i)}$
的丰沛性定义表明 $U(\psi) = X$。此外，构造，引理
\ref{constructions-lemma-invertible-map-into-relative-proj} 第 (3) 项
表明，$r_{\mathcal{L}, \psi}$ 在 $f^{-1}(V_i)$ 上的限制，
与性质，引理 \ref{properties-lemma-map-into-proj} 中的态射相同；
根据性质，引理
\ref{properties-lemma-ample-immersion-into-proj}，该态射是开浸入。
所以 (5) 成立。

\medskip\noindent
下面证明 (4) 蕴含 (1)。假设 (4)。用
$\pi : \underline{\text{Proj}}_S(\mathcal{A}) \to S$
表示结构态射。选取仿射开集 $V \subset S$。由构造，定义
\ref{constructions-definition-relative-proj}，可知
$\pi^{-1}(V) \subset \underline{\text{Proj}}_S(\mathcal{A})$
等于 $\text{Proj}(A)$，其中 $A = \mathcal{A}(V)$ 视为分次环。
因此，$r_{\mathcal{L}, \psi}$ 把 $f^{-1}(V)$ 同构地映到
$\text{Proj}(A)$ 的一个拟紧开集中。此外，
$\mathcal{L}^{\otimes d}$ 同构于
$\mathcal{O}_{\text{Proj}(A)}(d)$ 的拉回，其中某个 $d \geq 1$。
（见构造，引理
\ref{constructions-lemma-invertible-map-into-relative-proj} 第 (3) 项，
以及构造，引理
\ref{constructions-lemma-invertible-map-into-proj} 的最后一项断言。）
由性质中的引理 \ref{properties-lemma-open-in-proj-ample} 和
\ref{properties-lemma-ample-power-ample}，这蕴含
$\mathcal{L}|_{f^{-1}(V)}$ 丰沛。

\medskip\noindent
假设 (6)。由上面 (1) - (5) 的等价性可知，在 $X/S$
上相对丰沛这一性质在 $S$ 上是局部的。因此可以假设 $S$
仿射，并且必须证明 $\mathcal{L}$ 在 $X$ 上丰沛。在这种情况下，
态射 $r_{\mathcal{L}, \psi}$ 与由映射
$r_{\mathcal{L}, \psi} : X \to \text{Proj}(A)$
所给出的态射相等；后者也记作
$\psi : A = \mathcal{A}(V) \to \Gamma_*(X, \mathcal{L})$。
（见上面的引用。）如上还可看出，$\mathcal{L}^{\otimes d}$
是层 $\mathcal{O}_{\text{Proj}(A)}(d)$ 的拉回，其中某个
$d \geq 1$。此外，由于 $X$ 拟紧，
可知 $X$ 被等同于某个拟紧开子概形
$Y \subset \text{Proj}(A)$ 的闭子概形。由构造，引理
\ref{constructions-lemma-ample-on-proj}（另见性质，引理
\ref{properties-lemma-open-in-proj-ample}），可知
$\mathcal{O}_Y(d')$ 是 $Y$ 上的丰沛可逆层，其中某个
$d' \geq 1$。
由于丰沛层在闭子概形上的限制仍丰沛，见性质，引理
\ref{properties-lemma-ample-on-closed}，可知
$\mathcal{O}_Y(d')$ 的拉回丰沛。结合这些结果和性质，引理
\ref{properties-lemma-ample-power-ample}，可得
$\mathcal{L}$ 如所需是丰沛的。
\end{proof}

\begin{lemma}
\label{morphisms-lemma-ample-over-affine}
\begin{reference}
\cite[II Corollary 4.6.6]{EGA}
\end{reference}
设 $f : X \to S$ 是概形的态射。设 $\mathcal{L}$ 是可逆
$\mathcal{O}_X$-模。假设 $S$ 仿射。那么 $\mathcal{L}$
是 $f$-相对丰沛的，当且仅当 $\mathcal{L}$ 在 $X$ 上丰沛。
\end{lemma}

\begin{proof}
这直接由引理 \ref{morphisms-lemma-characterize-relatively-ample} 和定义得出。
\end{proof}

\begin{lemma}
\label{morphisms-lemma-quasi-affine-O-ample}
\begin{reference}
\cite[II Proposition 5.1.6]{EGA}
\end{reference}
设 $f : X \to S$ 是概形的态射。那么 $f$ 拟仿射，当且仅当
$\mathcal{O}_X$ 是 $f$-相对丰沛的。
\end{lemma}

\begin{proof}
这由性质，引理 \ref{properties-lemma-quasi-affine-O-ample}
和定义得出。
\end{proof}

\begin{lemma}
\label{morphisms-lemma-pullback-ample-tensor-relatively-ample}
设 $f : X \to Y$ 是概形的态射，$\mathcal{M}$ 是可逆
$\mathcal{O}_Y$-模，$\mathcal{L}$ 是可逆
$\mathcal{O}_X$-模。
\begin{enumerate}
\item 若 $\mathcal{L}$ 是 $f$-丰沛的且 $\mathcal{M}$ 丰沛，
则 $\mathcal{L} \otimes f^*\mathcal{M}^{\otimes a}$ 在
$a \gg 0$ 时丰沛。
\item 若 $\mathcal{M}$ 丰沛且 $f$ 拟仿射，则
$f^*\mathcal{M}$ 丰沛。
\end{enumerate}
\end{lemma}

\begin{proof}
假设 $\mathcal{L}$ 是 $f$-丰沛的且 $\mathcal{M}$ 丰沛。
根据假设，$Y$ 和 $f$ 都拟紧（见定义
\ref{morphisms-definition-relatively-ample} 和性质，定义
\ref{properties-definition-ample}）。因此 $X$ 拟紧。
由性质，引理 \ref{properties-lemma-ample-separated}，概形
$Y$ 分离；由引理 \ref{morphisms-lemma-relatively-ample-separated}，
态射 $f$ 分离。因此，由《概形》，引理
\ref{schemes-lemma-separated-permanence}，$X$ 分离。
选取 $x \in X$。可以选择 $m \geq 1$ 和
$t \in \Gamma(Y, \mathcal{M}^{\otimes m})$，使得 $Y_t$
仿射且 $f(x) \in Y_t$。由于 $\mathcal{L}$ 在
$f^{-1}(Y_t) = X_{f^*t}$ 上的限制是丰沛可逆层，
可以选择 $n \geq 1$ 和
$s \in \Gamma(X_{f^*t}, \mathcal{L}^{\otimes n})$，使得
$x \in (X_{f^*t})_s$ 且 $(X_{f^*t})_s$ 仿射。
由性质，引理 \ref{properties-lemma-invert-s-sections} 第 (2) 项
（上述情形满足其假设），存在整数 $e \geq 1$ 和截面
$s' \in \Gamma(X, \mathcal{L}^{\otimes n} \otimes f^*\mathcal{M}^{\otimes em})$，
它限制为 $s(f^*t)^e$（在 $X_{f^*t}$ 上）。对任意 $b > 0$，
考虑截面 $s'' = s'(f^*t)^b$，它是
$\mathcal{L}^{\otimes n} \otimes f^*\mathcal{M}^{\otimes (e + b)m}$
的截面。那么
$X_{s''} = (X_{f^*t})_s$ 是 $X$ 中包含 $x$ 的仿射开集。
选取 $b$ 使得 $n$ 整除 $e + b$，可见
$\mathcal{L}^{\otimes n} \otimes f^*\mathcal{M}^{\otimes (e + b)m}$
是 $n$ 次幂，其底为
$\mathcal{L} \otimes f^*\mathcal{M}^{\otimes a}$，其中某个
$a$；而且可以得到任何
$a$，只要它能被 $m$ 整除且足够大。由于 $X$ 拟紧，有限多个这样的仿射开集
覆盖 $X$。由此可知，对于某个充分整除且足够大的 $a$，可逆层
$\mathcal{L} \otimes f^*\mathcal{M}^{\otimes a}$ 在 $X$ 上丰沛。
另一方面，可知 $\mathcal{M}^{\otimes c}$（因而它在 $X$ 上的拉回）
对所有 $c \gg 0$ 都全局生成；见性质，命题
\ref{properties-proposition-characterize-ample}。因此，
$\mathcal{L} \otimes f^*\mathcal{M}^{\otimes a + c}$ 丰沛
（性质，引理
\ref{properties-lemma-ample-tensor-globally-generated}），其中
$c \gg 0$，从而 (1) 得证。

\medskip\noindent
第 (2) 项由引理 \ref{morphisms-lemma-quasi-affine-O-ample}、性质，引理
\ref{properties-lemma-ample-power-ample} 和第 (1) 项得出。
\end{proof}

\begin{lemma}
\label{morphisms-lemma-ample-composition}
设 $g : Y \to S$ 和 $f : X \to Y$ 是概形的态射。
设 $\mathcal{M}$ 是可逆 $\mathcal{O}_Y$-模。
设 $\mathcal{L}$ 是可逆 $\mathcal{O}_X$-模。
若 $S$ 拟紧、$\mathcal{M}$ 是 $g$-丰沛的且
$\mathcal{L}$ 是 $f$-丰沛的，则
$\mathcal{L} \otimes f^*\mathcal{M}^{\otimes a}$
是 $g \circ f$-丰沛的，其中 $a \gg 0$。
\end{lemma}

\begin{proof}
设 $S = \bigcup_{i = 1, \ldots, n} V_i$ 是有限仿射开覆盖。
由引理 \ref{morphisms-lemma-characterize-relatively-ample}，只需证明
$\mathcal{L} \otimes f^*\mathcal{M}^{\otimes a}$ 在
$(g \circ f)^{-1}(V_i)$ 上对 $i = 1, \ldots, n$ 丰沛。
因此，该引理由引理
\ref{morphisms-lemma-pullback-ample-tensor-relatively-ample} 得出。
\end{proof}

\begin{lemma}
\label{morphisms-lemma-ample-base-change}
设 $f : X \to S$ 是概形的态射。设 $\mathcal{L}$ 是可逆
$\mathcal{O}_X$-模。设 $S' \to S$ 是概形的态射。
设 $f' : X' \to S'$ 是 $f$ 的基变换，并用 $\mathcal{L}'$
表示 $\mathcal{L}$ 在 $X'$ 上的拉回。若 $\mathcal{L}$
是 $f$-丰沛的，则 $\mathcal{L}'$ 是 $f'$-丰沛的。
\end{lemma}

\begin{proof}
由引理 \ref{morphisms-lemma-characterize-relatively-ample}，只需找到仿射开覆盖
$S' = \bigcup U'_i$，使得 $\mathcal{L}'$ 在
$(f')^{-1}(U_i')$ 上的限制对所有 $i$ 都是丰沛可逆层。
可以选择 $U'_i$ 映入某个仿射开集 $U_i \subset S$。
在这种情况下，态射
$(f')^{-1}(U'_i) \to f^{-1}(U_i)$ 是仿射的，因为它是仿射态射
$U'_i \to U_i$ 的基变换（引理 \ref{morphisms-lemma-base-change-affine}）。
因此，由引理
\ref{morphisms-lemma-pullback-ample-tensor-relatively-ample}，
$\mathcal{L}'|_{(f')^{-1}(U'_i)}$ 丰沛。
\end{proof}

\begin{lemma}
\label{morphisms-lemma-ample-permanence}
设 $g : Y \to S$ 和 $f : X \to Y$ 是概形的态射。
设 $\mathcal{L}$ 是可逆 $\mathcal{O}_X$-模。
若 $\mathcal{L}$ 是 $g \circ f$-丰沛的，并且 $f$
拟紧\footnote{若 $g$ 拟分离，则这由《概形》，引理
\ref{schemes-lemma-quasi-compact-permanence} 得出。}，
则 $\mathcal{L}$ 是 $f$-丰沛的。
\end{lemma}

\begin{proof}
假设 $f$ 拟紧且 $\mathcal{L}$ 是 $g \circ f$-丰沛的。
设 $U \subset S$ 是仿射开集，设 $V \subset Y$ 是满足
$g(V) \subset U$ 的仿射开集。那么根据假设，
$\mathcal{L}|_{(g \circ f)^{-1}(U)}$ 在
$(g \circ f)^{-1}(U)$ 上丰沛。由于
$f^{-1}(V) \subset (g \circ f)^{-1}(U)$，由性质，引理
\ref{properties-lemma-ample-on-locally-closed} 可知，
$\mathcal{L}|_{f^{-1}(V)}$ 在 $f^{-1}(V)$ 上丰沛。
确切地说，由《概形》，引理
\ref{schemes-lemma-quasi-compact-permanence}，
$f^{-1}(V) \to (g \circ f)^{-1}(U)$ 是拟紧开浸入，因为
$(g \circ f)^{-1}(U)$ 分离（性质，引理
\ref{properties-lemma-ample-separated}），并且
$f^{-1}(V)$ 拟紧（因为 $f$ 拟紧）。因此，由引理
\ref{morphisms-lemma-characterize-relatively-ample}，可得
$\mathcal{L}$ 是 $f$-丰沛的。
\end{proof}

\section{相对甚丰沛层}
\label{morphisms-section-very-ample}

\noindent
回顾：给定拟凝聚层 $\mathcal{E}$（在概形 $S$ 上），
与 $\mathcal{E}$ 相伴的{\it 射影丛}是态射
$\mathbf{P}(\mathcal{E}) \to S$，其中
$\mathbf{P}(\mathcal{E}) = \underline{\text{Proj}}_S(\text{Sym}(\mathcal{E}))$；
见构造，定义 \ref{constructions-definition-projective-bundle}。

\begin{definition}
\label{morphisms-definition-very-ample}
设 $f : X \to S$ 是概形的态射。设 $\mathcal{L}$ 是可逆
$\mathcal{O}_X$-模。若 $\mathcal{L}$ {\it 相对甚丰沛}，
更确切地称为{\it $f$-相对甚丰沛}，也称{\it 在 $X/S$ 上甚丰沛}
或{\it $f$-甚丰沛}，其含义是：存在拟凝聚
$\mathcal{O}_S$-模 $\mathcal{E}$ 和浸入
$i : X \to \mathbf{P}(\mathcal{E})$（在 $S$ 上），使得
$\mathcal{L} \cong i^*\mathcal{O}_{\mathbf{P}(\mathcal{E})}(1)$。
\end{definition}



\noindent
由于这个定义中没有拟紧性假设，一般而言，相对甚丰沛可逆层
未必是相对丰沛可逆层。

\begin{lemma}
\label{morphisms-lemma-ample-very-ample}
\begin{reference}
\cite[II, Proposition 4.6.2]{EGA}
\end{reference}
设 $f : X \to S$ 是概形的态射。设 $\mathcal{L}$ 是可逆
$\mathcal{O}_X$-模。若 $f$ 拟紧且 $\mathcal{L}$
是相对甚丰沛可逆层，则 $\mathcal{L}$ 是相对丰沛可逆层。
\end{lemma}

\begin{proof}
根据定义，存在拟凝聚 $\mathcal{O}_S$-模 $\mathcal{E}$
和浸入 $i : X \to \mathbf{P}(\mathcal{E})$（在 $S$ 上），使得
$\mathcal{L} \cong i^*\mathcal{O}_{\mathbf{P}(\mathcal{E})}(1)$。
令 $\mathcal{A} = \text{Sym}(\mathcal{E})$，于是根据定义，
$\mathbf{P}(\mathcal{E}) = \underline{\text{Proj}}_S(\mathcal{A})$。
分次 $\mathcal{O}_S$-代数 $\mathcal{A}$ 配备映射
$$
\psi :
\mathcal{A} \to
\bigoplus\nolimits_{n \geq 0}
\pi_*\mathcal{O}_{\mathbf{P}(\mathcal{E})}(n) \to
\bigoplus\nolimits_{n \geq 0}
f_*\mathcal{L}^{\otimes n}
$$
其中第二支箭头使用同一
$\mathcal{L} \cong i^*\mathcal{O}_{\mathbf{P}(\mathcal{E})}(1)$。
由 $f_*$ 和 $f^*$ 的伴随性，得到态射
$\psi : f^*\mathcal{A} \to \bigoplus_{n \geq 0}\mathcal{L}^{\otimes n}$。
我们省略验证：与这个映射相伴的态射
$r_{\mathcal{L}, \psi}$ 恰好就是浸入 $i$。
因此，结果由引理 \ref{morphisms-lemma-characterize-relatively-ample}
第 (6) 项得出。
\end{proof}

\noindent
为得到该引理的正确逆命题，我们要问：给定相对丰沛可逆层
$\mathcal{L}$，是否存在整数 $n \geq 1$，使得
$\mathcal{L}^{\otimes n}$ 相对甚丰沛？一般而言并非如此。
有若干因素会阻止这一结论成立：
\begin{enumerate}
\item 即使 $S$ 仿射，也可能不存在可行的有限整数 $n$，
因为 $X \to S$ 不是有限型的；见例
\ref{morphisms-example-not-finite-type-proj}。
\item 基不拟紧意味着，即使 $f$ 为有限型，该结果也可能不成立。
确切地说，给定域 $k$，存在概形 $X_d$，它在 $k$ 上有限型，
带有丰沛可逆层 $\mathcal{O}_{X_d}(1)$，使得
$\mathcal{O}_{X_d}(1)$ 的甚丰沛张量幂中最小的是第 $d$ 次幂。
见例 \ref{morphisms-example-not-bounded}。取 $f$ 为这些概形 $X_d$
映到 $\Spec(k)$ 的各个副本的不交并的态射，即得一个例子。
\end{enumerate}
逆命题的本版本见下面的引理
\ref{morphisms-lemma-finite-type-ample-very-ample}。证明它之前先做一些准备工作。

\begin{example}
\label{morphisms-example-very-ample}
设 $S$ 是概形。设 $\mathcal{A}$ 是拟凝聚分次
$\mathcal{O}_S$-代数，它由 $\mathcal{A}_1$ 生成
（在 $\mathcal{A}_0$ 上）。令
$X = \underline{\text{Proj}}_S(\mathcal{A})$。在这种情况下，
$\mathcal{O}_X(1)$ 可逆且在 $X/S$ 上甚丰沛。确切地说，
与分次 $\mathcal{O}_S$-代数映射
$$
\text{Sym}_{\mathcal{O}_X}^*(\mathcal{A}_1)
\longrightarrow
\mathcal{A}
$$
相伴的态射是闭浸入 $X \to \mathbf{P}(\mathcal{A}_1)$，
它把 $\mathcal{O}_{\mathbf{P}(\mathcal{A}_1)}(1)$ 拉回为
$\mathcal{O}_X(1)$；见构造，引理
\ref{constructions-lemma-surjective-generated-degree-1-map-relative-proj}。
\end{example}

\begin{example}
\label{morphisms-example-not-finite-type-proj}
设 $k$ 是域。考虑分次 $k$-代数
$$
A = k[U, V, Z_1, Z_2, Z_3, \ldots]/I
\quad
\text{其中}
\quad
I = (U^2 - Z_1^2, U^4 - Z_2^2, U^6 - Z_3^2, \ldots)
$$
其分次由 $\deg(U) = \deg(V) = \deg(Z_1) = 1$
和 $\deg(Z_d) = d$ 给出。注意，$X = \text{Proj}(A)$
由 $D_{+}(U)$ 和 $D_{+}(V)$ 覆盖。因此层
$\mathcal{O}_X(n)$ 全都可逆，并同构于
$\mathcal{O}_X(1)^{\otimes n}$。特别地，$\mathcal{O}_X(1)$
是丰沛的，并且是 $f$-丰沛的，其中态射为
$f : X \to \Spec(k)$。
我们断言，$\mathcal{O}_X(1)$ 的任何幂都不是 $f$-相对甚丰沛的。
确实，容易看出 $\Gamma(X, \mathcal{O}_X(n))$ 是第 $n$
个分次分量（代数 $A$ 的）。因此，若 $\mathcal{O}_X(n)$ 甚丰沛，
则 $X$ 会是 $k$ 上某个射影空间的闭子概形，从而在 $k$
上为有限型。另一方面，$D_{+}(V)$ 是
$k[t, t_1, t_2, \ldots]/(t^2 - t_1^2, t^4 - t_2^2, t^6 - t_3^2, \ldots)$
的谱，而后者在 $k$ 上不是有限型的。
\end{example}

\begin{example}
\label{morphisms-example-not-bounded}
设 $k$ 是无限域。设 $\lambda_1, \lambda_2, \lambda_3, \ldots$
是 $k^*$ 中两两不同的元素。（这并非严格必要；事实上，即使所有
$\lambda_i$ 都等于 $1$，这个例子也完全成立。）
考虑分次 $k$-代数
$$
A_d = k[U, V, Z]/I_d
\quad
\text{其中}
\quad
I_d = (Z^2 - \prod\nolimits_{i = 1}^{2d} (U - \lambda_i V)).
$$
其分次由 $\deg(U) = \deg(V) = 1$ 和 $\deg(Z) = d$ 给出。
那么 $X_d = \text{Proj}(A_d)$ 有丰沛可逆层
$\mathcal{O}_{X_d}(1)$。我们断言，若
$\mathcal{O}_{X_d}(n)$ 甚丰沛，则 $n \geq d$。
理由是 $Z$ 的次数为 $d$，从而
$\Gamma(X_d, \mathcal{O}_{X_d}(n)) =
k[U, V]_n$ 对 $n < d$ 成立。省略细节。
\end{example}

\begin{lemma}
\label{morphisms-lemma-relatively-very-ample-separated}
设 $f : X \to S$ 是概形的态射。设 $\mathcal{L}$ 是 $X$ 上的可逆层。
若 $\mathcal{L}$ 在 $X/S$ 上相对甚丰沛，则 $f$ 分离。
\end{lemma}

\begin{proof}
分离性在基底上是局部的（见概形，节
\ref{schemes-section-separation-axioms}）。浸入是分离的
（见概形，引理 \ref{schemes-lemma-immersions-monomorphisms}）。
因此引理成立，因为局部地，$X$ 可浸入一个分次环的齐次谱，
而该齐次谱是分离的；见构造，引理
\ref{constructions-lemma-proj-separated}。
\end{proof}

\begin{lemma}
\label{morphisms-lemma-relatively-very-ample}
设 $f : X \to S$ 是概形的态射。设 $\mathcal{L}$ 是 $X$ 上的可逆层。
假设 $f$ 拟紧。下列条件等价：
\begin{enumerate}
\item $\mathcal{L}$ 在 $X/S$ 上相对甚丰沛；
\item 存在开覆盖 $S = \bigcup V_j$，使得
$\mathcal{L}|_{f^{-1}(V_j)}$ 在 $f^{-1}(V_j)/V_j$ 上相对甚丰沛，
对所有 $j$ 均成立；
\item 存在拟凝聚分次 $\mathcal{O}_S$-代数层 $\mathcal{A}$，
它由次数 $1$ 的部分在 $\mathcal{O}_S$ 上生成，并且存在分次
$\mathcal{O}_X$-代数映射
$\psi : f^*\mathcal{A} \to \bigoplus_{n \geq 0} \mathcal{L}^{\otimes n}$，
使得 $f^*\mathcal{A}_1 \to \mathcal{L}$ 满射，且相伴态射
$r_{\mathcal{L}, \psi} : X \to \underline{\text{Proj}}_S(\mathcal{A})$
是浸入；
\item $f$ 拟分离，典范映射
$\psi : f^*f_*\mathcal{L} \to \mathcal{L}$ 满射，且相伴映射
$r_{\mathcal{L}, \psi} : X \to \mathbf{P}(f_*\mathcal{L})$ 是浸入。
\end{enumerate}
\end{lemma}

\begin{proof}
显然，(1) 蕴含 (2)。同样显然，(4) 蕴含 (1)；(4) 中的拟分离性假设
用于通过概形，引理 \ref{schemes-lemma-push-forward-quasi-coherent}
保证 $f_*\mathcal{L}$ 拟凝聚。

\medskip\noindent
假设 (2)。我们证明 (4)。令 $S = \bigcup V_j$ 是 (2) 中的开覆盖。
令 $X_j = f^{-1}(V_j)$，并令 $f_j : X_j \to V_j$ 为 $f$ 的限制。
由引理 \ref{morphisms-lemma-relatively-very-ample-separated} 可知 $f$ 分离
（因为分离性在基底上是局部的）。依假设，存在拟凝聚
$\mathcal{O}_{V_j}$-模 $\mathcal{E}_j$ 和浸入
$i_j : X_j \to \mathbf{P}(\mathcal{E}_j)$，满足
$\mathcal{L}|_{X_j} \cong i_j^*\mathcal{O}_{\mathbf{P}(\mathcal{E}_j)}(1)$。
态射 $i_j$ 对应于满射
$f_j^*\mathcal{E}_j \to \mathcal{L}|_{X_j}$；见构造，节
\ref{constructions-section-projective-bundle}。这个映射伴随于一个映射
$\mathcal{E}_j \to f_*\mathcal{L}|_{V_j}$，使得复合
$$
f_j^*\mathcal{E}_j \to (f^*f_*\mathcal{L})|_{X_j} \to \mathcal{L}|_{X_j}
$$
满射。于是
$\psi : f^*f_*\mathcal{L} \to \mathcal{L}$ 满射。令
$r_{\mathcal{L}, \psi} : X \to \mathbf{P}(f_*\mathcal{L})$
为相伴态射。还须证明 $r_{\mathcal{L}, \psi}$ 是浸入；我们建议读者自行证明。
$\mathcal{O}_{V_j}$-模映射 $\mathcal{E}_j \to f_*\mathcal{L}|_{V_j}$
诱导对称代数之间的同态，后者又定义态射
$$
\mathbf{P}(f_*\mathcal{L}|_{V_j}) \supset U_j \longrightarrow
\mathbf{P}(\mathcal{E}_j)
$$
其中 $U_j$ 是构造，引理
\ref{constructions-lemma-morphism-relative-proj} 中的开子概形。
$\psi$ 与 $\mathcal{E}_j \to f_*\mathcal{L}|_{V_j}$ 的相容性说明
$r_{\mathcal{L}, \psi}(X_j) \subset U_j$，并且存在分解
$$
\xymatrix{
X_j \ar[r]^-{r_{\mathcal{L}, \psi}} & U_j \ar[r] &
\mathbf{P}(\mathcal{E}_j)
}
$$
我们省略验证。这说明 $r_{\mathcal{L}, \psi}$ 是浸入。

\medskip\noindent
至此可见，(1)、(2) 与 (4) 等价。显然 (4) 蕴含 (3)。
假设 (3)。我们证明 (1)。令 $\mathcal{A}$ 是拟凝聚分次
$\mathcal{O}_S$-代数层，它由次数 $1$ 的部分在 $\mathcal{O}_S$ 上生成。
考虑分次 $\mathcal{O}_S$-代数映射
$\text{Sym}(\mathcal{A}_1) \to \mathcal{A}$。依假设，此映射满射，
因而诱导闭浸入
$$
\underline{\text{Proj}}_S(\mathcal{A})
\longrightarrow
\mathbf{P}(\mathcal{A}_1)
$$
它把 $\mathcal{O}(1)$ 拉回为 $\mathcal{O}(1)$；见构造，引理
\ref{constructions-lemma-surjective-generated-degree-1-map-relative-proj}。
因此，(3) 蕴含 (1) 是显然的。
\end{proof}

\begin{lemma}
\label{morphisms-lemma-very-ample-base-change}
设 $f : X \to S$ 是概形的态射。设 $\mathcal{L}$ 是可逆
$\mathcal{O}_X$-模。设 $S' \to S$ 是概形的态射。令
$f' : X' \to S'$ 为 $f$ 的基变换，并用 $\mathcal{L}'$ 表示
$\mathcal{L}$ 到 $X'$ 的拉回。若 $\mathcal{L}$ 是 $f$-甚丰沛的，
则 $\mathcal{L}'$ 是 $f'$-甚丰沛的。
\end{lemma}

\begin{proof}
由定义 \ref{morphisms-definition-very-ample}，存在拟凝聚
$\mathcal{O}_S$-模 $\mathcal{E}$ 和浸入
$i : X \to \mathbf{P}(\mathcal{E})$（在 $S$ 上），使得
$\mathcal{L} \cong i^*\mathcal{O}_{\mathbf{P}(\mathcal{E})}(1)$。
$\mathbf{P}(\mathcal{E})$ 到 $S'$ 的基变换是与模 $\mathcal{E}'$
相伴的射影丛，其中该模是 $\mathcal{E}$ 的拉回；而
$\mathcal{O}_{\mathbf{P}(\mathcal{E})}(1)$ 的拉回是
$\mathcal{O}_{\mathbf{P}(\mathcal{E}')}(1)$；见构造，引理
\ref{constructions-lemma-relative-proj-base-change}。最后，浸入的基变换仍是浸入
（概形，引理 \ref{schemes-lemma-base-change-immersion}）。
\end{proof}







\section{相对于有限型态射的丰沛层与甚丰沛层}
\label{morphisms-section-ample-finite-type}

\noindent
事实上，本节大部分内容涉及尚未定义的（拟）射影态射概念。

\begin{lemma}
\label{morphisms-lemma-very-ample-finite-type-over-affine}
设 $f : X \to S$ 是概形的态射。设 $\mathcal{L}$ 是 $X$ 上的可逆层。
假设：
\begin{enumerate}
\item 可逆层 $\mathcal{L}$ 在 $X/S$ 上甚丰沛；
\item 态射 $X \to S$ 为有限型；并且
\item $S$ 仿射。
\end{enumerate}
则存在 $n \geq 0$ 以及浸入
$i : X \to \mathbf{P}^n_S$（在 $S$ 上），使得
$\mathcal{L} \cong i^*\mathcal{O}_{\mathbf{P}^n_S}(1)$。
\end{lemma}

\begin{proof}
假设 (1)、(2) 和 (3)。条件 (3) 意味着
$S = \Spec(R)$，其中 $R$ 为某个环。根据定义，条件 (1) 意味着存在拟凝聚
$\mathcal{O}_S$-模 $\mathcal{E}$ 和浸入
$\alpha : X \to \mathbf{P}(\mathcal{E})$，使得
$\mathcal{L} = \alpha^*\mathcal{O}_{\mathbf{P}(\mathcal{E})}(1)$。
写成 $\mathcal{E} = \widetilde{M}$，其中某个 $R$-模为 $M$。于是
$$
\mathbf{P}(\mathcal{E}) = \text{Proj}(\text{Sym}_R(M)).
$$
由于 $\alpha$ 是浸入，而
$\text{Proj}(\text{Sym}_R(M))$ 的拓扑由标准开集
$D_{+}(f)$ 生成，其中 $f \in \text{Sym}_R^d(M)$、$d \geq 1$，
故对每个 $x \in X$，可找到
$f \in \text{Sym}_R^d(M)$、$d \geq 1$，使得
$\alpha(x) \in D_{+}(f)$，并使
$$
\alpha|_{\alpha^{-1}(D_{+}(f))} : \alpha^{-1}(D_{+}(f)) \to D_{+}(f)
$$
为闭浸入。条件 (2) 蕴含 $X$ 拟紧。因此可找到有限多个元素
$f_j \in \text{Sym}_R^{d_j}(M)$、$d_j \geq 1$，使得对每个
$f = f_j$，上面显示的映射为闭浸入，并且
$\alpha(X) \subset \bigcup D_{+}(f_j)$。写
$U_j = \alpha^{-1}(D_{+}(f_j))$。注意，$U_j$ 是仿射概形
$D_{+}(f_j)$ 的闭子概形，因而仿射。写 $U_j = \Spec(A_j)$。
条件 (2) 还蕴含 $A_j$ 在 $R$ 上为有限型；见引理
\ref{morphisms-lemma-locally-finite-type-characterize}。选择有限多个
$x_{j, k} \in A_j$，使它们生成 $A_j$（作为 $R$-代数）。
由于 $\alpha|_{U_j}$ 是闭浸入，可见 $x_{j, k}$ 是某个元素
$$
f_{j, k}/f_j^{e_{j, k}} \in \text{Sym}_R(M)_{(f_j)}
=
\Gamma(D_{+}(f_j), \mathcal{O}_{\text{Proj}(\text{Sym}_R(M))}).
$$
的像。最后，选择 $n \geq 1$ 以及元素 $y_0, \ldots, y_n \in M$，
使得每个多项式 $f_j, f_{j, k} \in \text{Sym}_R(M)$ 都是元素
$y_t$ 的多项式，其系数在 $R$ 中。考虑分次环映射
$$
\psi : R[Y_0, \ldots, Y_n] \longrightarrow \text{Sym}_R(M),
\quad Y_i \longmapsto y_i.
$$
用 $F_j$、$F_{j, k}$ 表示 $R[Y_0, \ldots, Y_n]$ 中满足
$\psi(F_j) = f_j$ 和 $\psi(F_{j, k}) = f_{j, k}$ 的元素。
由构造，引理 \ref{constructions-lemma-morphism-proj}，得到开子概形
$$
U(\psi) \subset \text{Proj}(\text{Sym}_R(M))
$$
和态射 $r_\psi : U(\psi) \to \mathbf{P}^n_R$。此态射满足
$r_\psi^{-1}(D_{+}(F_j)) = D_{+}(f_j)$，因而可见
$\alpha(X) \subset U(\psi)$。此外，显然
$$
i = r_\psi \circ \alpha : X \longrightarrow \mathbf{P}^n_R
$$
仍是浸入，因为依构造
$i^\sharp(F_{j, k}/F_j^{e_{j, k}}) = x_{j, k} \in
A_j = \Gamma(U_j, \mathcal{O}_X)$。此外，态射 $r_\psi$ 配备映射
$\theta : r_\psi^*\mathcal{O}_{\mathbf{P}^n_R}(1)
\to \mathcal{O}_{\text{Proj}(\text{Sym}_R(M))}(1)|_{U(\psi)}$，
在此情形中该映射为同构（$\theta$ 的构造见上引引理；略去一些细节）。
由于原映射 $\alpha$ 按假设满足
$\mathcal{L} = \alpha^*\mathcal{O}_{\text{Proj}(\text{Sym}_R(M))}(1)$，
结论成立。
\end{proof}

\begin{lemma}
\label{morphisms-lemma-quasi-affine-finite-type-over-S}
设 $\pi : X \to S$ 是概形的态射。假设 $X$ 拟仿射，并且
$\pi$ 局部为有限型。则存在 $n \geq 0$ 以及浸入
$i : X \to \mathbf{A}^n_S$（在 $S$ 上）。
\end{lemma}

\begin{proof}
令 $A = \Gamma(X, \mathcal{O}_X)$。依假设，$X$ 拟紧，并被视为
$\Spec(A)$ 的开子概形；见性质，引理
\ref{properties-lemma-quasi-affine}。此外，诸开集 $X_f$ 中满足
$f \in A$ 且 $X_f$ 仿射的那些，构成 $X$ 的拓扑的一组基；
见性质，引理 \ref{properties-lemma-quasi-affine} 的证明。
因此可找到有限多个 $f_j \in A$（$j = 1, \ldots, m$），使得
$X = \bigcup X_{f_j}$，并且有 $\pi(X_{f_j}) \subset V_j$，其中
$V_j \subset S$ 是某个仿射开集。由引理
\ref{morphisms-lemma-locally-finite-type-characterize}，环映射
$\mathcal{O}(V_j) \to \mathcal{O}(X_{f_j}) = A_{f_j}$ 为有限型。
因此可选择 $a_1, \ldots, a_N \in A$，使得元素
$a_1, \ldots, a_N, 1/f_j$ 生成 $A_{f_j}$（作为
$\mathcal{O}(V_j)$-代数），对每个 $j$ 均如此。取 $n = m + N$，令
$$
i : X \longrightarrow \mathbf{A}^n_S
$$
为由全局截面 $f_1, \ldots, f_m, a_1, \ldots, a_N$ 给出的态射；
这些是 $X$ 的结构层的全局截面。
令 $D(x_j) \subset \mathbf{A}^n_S$ 为第 $j$ 个坐标函数非零处的开子概形。
那么，对 $1 \leq j \leq m$，有
$i^{-1}(D(x_j)) = X_{f_j}$，并且诱导态射
$X_{f_j} \to D(x_j)$ 经由仿射开集
\begin{center}
$\Spec(\mathcal{O}(V_j)[x_1, \ldots, x_n, 1/x_j])$
\end{center}
分解；
该开集位于 $D(x_j)$ 中。
由于环映射
$\mathcal{O}(V_j)[x_1, \ldots, x_n, 1/x_j] \to A_{f_j}$
依构造满射，故可得 $i^{-1}(D(x_j)) \to D(x_j)$
如所需是浸入。
\end{proof}

\begin{lemma}
\label{morphisms-lemma-quasi-projective-finite-type-over-S}
设 $f : X \to S$ 是概形的态射。设 $\mathcal{L}$ 是 $X$ 上的可逆层。
假设：
\begin{enumerate}
\item 可逆层 $\mathcal{L}$ 在 $X$ 上丰沛；并且
\item 态射 $X \to S$ 局部为有限型。
\end{enumerate}
则存在 $d_0 \geq 1$，使得对每个 $d \geq d_0$，存在
$n \geq 0$ 以及浸入 $i : X \to \mathbf{P}^n_S$（在 $S$ 上），满足
$\mathcal{L}^{\otimes d} \cong i^*\mathcal{O}_{\mathbf{P}^n_S}(1)$。
\end{lemma}

\begin{proof}
令
$A = \Gamma_*(X, \mathcal{L}) =
\bigoplus_{d \geq 0} \Gamma(X, \mathcal{L}^{\otimes d})$。
由性质，命题 \ref{properties-proposition-characterize-ample}，
仿射开集 $X_a$ 中满足 $a \in A_{+}$ 为齐次元的那些，
构成 $X$ 的拓扑的一组基。
因此可找到有限多个这样的元素 $a_0, \ldots, a_n \in A_{+}$，使得：
\begin{enumerate}
\item $X = \bigcup_{i = 0, \ldots, n} X_{a_i}$；
\item 每个 $X_{a_i}$ 都仿射；并且
\item 每个 $X_{a_i}$ 都映入某个仿射开集 $V_i \subset S$。
\end{enumerate}
由引理 \ref{morphisms-lemma-locally-finite-type-characterize}，可见环映射
$\mathcal{O}_S(V_i) \to \mathcal{O}_X(X_{a_i})$ 为有限型。
因此可找到有限多个元素
$f_{ij} \in \mathcal{O}_X(X_{a_i})$（$j = 1, \ldots, n_i$），
它们生成 $\mathcal{O}_X(X_{a_i})$（作为
$\mathcal{O}_S(V_i)$-代数）。由性质，引理
\ref{properties-lemma-invert-s-sections}，可将每个
$f_{ij}$ 写为 $a_{ij}/a_i^{e_{ij}}$，其中某个
$a_{ij} \in A_{+}$ 为齐次元。令 $N$ 是所有元素
$a_i$、$a_{ij}$ 的次数的一个正公倍数。考虑元素
$$
a_i^{N/\deg(a_i)}, \ a_{ij}a_i^{(N/\deg(a_i)) - e_{ij}} \in A_N.
$$
依构造，它们生成可逆层 $\mathcal{L}^{\otimes N}$（在 $X$ 上）。
因此它们给出态射
$$
j : X \longrightarrow
\mathbf{P}_S^{m}
\quad
\text{其中 } m = n + \sum n_i
$$
（在 $S$ 上）；见构造，引理
\ref{constructions-lemma-projective-space} 及定义
\ref{constructions-definition-projective-space}。此外，
$j^*\mathcal{O}_{\mathbf{P}_S}(1) = \mathcal{L}^{\otimes N}$。
我们将齐次坐标命名为 $T_0, \ldots, T_n, T_{ij}$，而不是
$T_0, \ldots, T_m$。对 $i = 0, \ldots, n$，有
$j^{-1}(D_{+}(T_i)) = X_{a_i}$。此外，拉回元素
$T_{ij}/T_i$（通过 $j^\sharp$）时，得到元素
$f_{ij} \in \mathcal{O}_X(X_{a_i})$。
因此，态射 $j$ 在 $X_{a_i}$ 上的限制给出闭浸入：将
$X_{a_i}$ 浸入仿射开集
$D_{+}(T_i) \cap \mathbf{P}^m_{V_i}$；该开集位于
$\mathbf{P}^N_S$ 中。故可得态射 $j$ 为浸入。
这说明引理对某些 $d$ 和 $n$ 成立，这在几乎所有应用中已经足够。

\medskip\noindent
这证明：对某个 $d_2 \geq 1$（即上面的 $d_2 = N$）和某个
$m \geq 0$，存在由全局截面给出的浸入
$j : X \to \mathbf{P}^m_S$；这些截面为
$s'_0, \ldots, s'_m \in \Gamma(X, \mathcal{L}^{\otimes d_2})$。
由性质，命题 \ref{properties-proposition-characterize-ample}，
存在整数 $d_1$，使得 $\mathcal{L}^{\otimes d}$ 全局生成，
对所有 $d \geq d_1$ 都如此。令 $d_0 = d_1 + d_2$。
我们断言引理对这个 $d_0$ 成立。确实，给定整数
$d \geq d_0$，可选择
$s''_1, \ldots, s''_t
\in \Gamma(X, \mathcal{L}^{\otimes d - d_2})$，使其生成
$\mathcal{L}^{\otimes d - d_2}$（在 $X$ 上）。令 $k = (m + 1)t$，并用
$s_0, \ldots, s_k$ 表示截面的集合
$s'_\alpha s''_\beta$，其中 $\alpha = 0, \ldots, m$、
$\beta = 1, \ldots, t$。这些截面生成
$\mathcal{L}^{\otimes d}$（在 $X$ 上），因而定义态射
$$
i : X \longrightarrow \mathbf{P}^{k - 1}_S
$$
使得 $i^*\mathcal{O}_{\mathbf{P}^n_S}(1) \cong \mathcal{L}^{\otimes d}$。
为看出 $i$ 是浸入，注意 $i$ 是复合
$$
X \longrightarrow \mathbf{P}^m_S \times_S \mathbf{P}^{t - 1}_S
\longrightarrow \mathbf{P}^{k - 1}_S
$$
其中第一个态射是 $(j, j')$，$j'$ 由
$s''_1, \ldots, s''_t$ 给出，而第二个态射是 Segre 嵌入
（构造，引理 \ref{constructions-lemma-segre-embedding}）。
由于 $j$ 是浸入，$(j, j')$ 也是浸入（将引理
\ref{morphisms-lemma-immersion-permanence} 应用于
$X \to \mathbf{P}^m_S \times_S \mathbf{P}^{t - 1}_S
\to \mathbf{P}^m_S$）。因此 $i$ 是浸入的复合，从而仍为浸入
（概形，引理 \ref{schemes-lemma-composition-immersion}）。
\end{proof}

\begin{lemma}
\label{morphisms-lemma-finite-type-over-affine-ample-very-ample}
设 $f : X \to S$ 是概形的态射。设 $\mathcal{L}$ 是可逆
$\mathcal{O}_X$-模。假设 $S$ 仿射且 $f$ 为有限型。下列条件等价：
\begin{enumerate}
\item $\mathcal{L}$ 在 $X$ 上丰沛；
\item $\mathcal{L}$ 是 $f$-丰沛的；
\item $\mathcal{L}^{\otimes d}$ 是 $f$-甚丰沛的，对某个 $d \geq 1$ 成立；
\item $\mathcal{L}^{\otimes d}$ 是 $f$-甚丰沛的，对所有 $d \gg 1$ 成立；
\item 对某个 $d \geq 1$，存在 $n \geq 1$ 以及浸入
$i : X \to \mathbf{P}^n_S$，使得
$\mathcal{L}^{\otimes d} \cong i^*\mathcal{O}_{\mathbf{P}^n_S}(1)$；并且
\item 对所有 $d \gg 1$，存在 $n \geq 1$ 以及浸入
$i : X \to \mathbf{P}^n_S$，使得
$\mathcal{L}^{\otimes d} \cong i^*\mathcal{O}_{\mathbf{P}^n_S}(1)$。
\end{enumerate}
\end{lemma}

\begin{proof}
(1) 与 (2) 的等价性就是引理 \ref{morphisms-lemma-ample-over-affine}。
(2) $\Rightarrow$ (6) 是引理
\ref{morphisms-lemma-quasi-projective-finite-type-over-S}。显然，(6) 蕴含 (5)。
由于 $\mathbf{P}^n_S$ 是 $S$ 上的射影丛（见构造，引理
\ref{constructions-lemma-projective-space-bundle}），由相对甚丰沛层的定义，
可见 (5) 蕴含 (3)，(6) 蕴含 (4)。显然，(4) 蕴含 (3)。
最后还须证明 (3) 蕴含 (2)，这由引理
\ref{morphisms-lemma-ample-very-ample} 及引理
\ref{morphisms-lemma-ample-power-ample} 得出。
\end{proof}

\begin{lemma}
\label{morphisms-lemma-finite-type-ample-very-ample}
设 $f : X \to S$ 是概形的态射。设 $\mathcal{L}$ 是可逆
$\mathcal{O}_X$-模。假设 $S$ 拟紧且 $f$ 为有限型。下列条件等价：
\begin{enumerate}
\item $\mathcal{L}$ 是 $f$-丰沛的；
\item $\mathcal{L}^{\otimes d}$ 是 $f$-甚丰沛的，对某个 $d \geq 1$ 成立；
\item $\mathcal{L}^{\otimes d}$ 是 $f$-甚丰沛的，对所有 $d \gg 1$ 成立。
\end{enumerate}
\end{lemma}

\begin{proof}
显然，(3) 蕴含 (2)。引理 \ref{morphisms-lemma-ample-very-ample} 保证 (2)
蕴含 (1)，因为根据定义，有限型态射拟紧。假设 $\mathcal{L}$ 是
$f$-丰沛的。选取有限仿射开覆盖
$S = V_1 \cup \ldots \cup V_m$。写 $X_i = f^{-1}(V_i)$。
由上面的引理 \ref{morphisms-lemma-finite-type-over-affine-ample-very-ample} 可见，
存在 $d_0$，使得 $\mathcal{L}^{\otimes d}$ 在 $X_i/V_i$ 上相对甚丰沛，
对所有 $d \geq d_0$ 都如此。因此，由引理
\ref{morphisms-lemma-relatively-very-ample}，可得 (1) 蕴含 (3)。
\end{proof}

\noindent
下面两个引理给出有限型态射情形中，相对甚丰沛可逆层与
相对丰沛可逆层最常用也最有用的刻画。

\begin{lemma}
\label{morphisms-lemma-characterize-very-ample-on-finite-type}
设 $f : X \to S$ 是概形的态射。设 $\mathcal{L}$ 是 $X$ 上的可逆层。
假设 $f$ 为有限型。下列条件等价：
\begin{enumerate}
\item $\mathcal{L}$ 是 $f$-相对甚丰沛的；并且
\item 存在开覆盖 $S = \bigcup V_j$；对每个 $j$，存在整数 $n_j$
以及浸入
$$
i_j :
X_j = f^{-1}(V_j) = V_j \times_S X
\longrightarrow
\mathbf{P}^{n_j}_{V_j}
$$
（在 $V_j$ 上），使得
$\mathcal{L}|_{X_j} \cong i_j^*\mathcal{O}_{\mathbf{P}^{n_j}_{V_j}}(1)$。
\end{enumerate}
\end{lemma}

\begin{proof}
取 $S$ 的一个仿射开覆盖，并将引理
\ref{morphisms-lemma-very-ample-finite-type-over-affine} 应用于 $f$ 和
$\mathcal{L}$ 的每个限制，便可见 (1) 蕴含 (2)。由引理
\ref{morphisms-lemma-relatively-very-ample} 可见 (2) 蕴含 (1)。
\end{proof}

\begin{lemma}
\label{morphisms-lemma-characterize-ample-on-finite-type}
设 $f : X \to S$ 是概形的态射。设 $\mathcal{L}$ 是 $X$ 上的可逆层。
假设 $f$ 为有限型。下列条件等价：
\begin{enumerate}
\item $\mathcal{L}$ 是 $f$-相对丰沛的；并且
\item 存在开覆盖 $S = \bigcup V_j$；对每个 $j$，存在整数
$d_j \geq 1$、$n_j \geq 0$ 以及浸入
$$
i_j :
X_j = f^{-1}(V_j) = V_j \times_S X
\longrightarrow
\mathbf{P}^{n_j}_{V_j}
$$
（在 $V_j$ 上），使得
$\mathcal{L}^{\otimes d_j}|_{X_j} \cong
i_j^*\mathcal{O}_{\mathbf{P}^{n_j}_{V_j}}(1)$。
\end{enumerate}
\end{lemma}

\begin{proof}
取 $S$ 的一个仿射开覆盖，并将引理
\ref{morphisms-lemma-finite-type-over-affine-ample-very-ample} 应用于 $f$ 和
$\mathcal{L}$ 的每个限制，便可见 (1) 蕴含 (2)。由引理
\ref{morphisms-lemma-characterize-relatively-ample} 可见 (2) 蕴含 (1)。
\end{proof}

\begin{lemma}
\label{morphisms-lemma-invertible-add-enough-ample-very-ample}
设 $f : X \to S$ 是概形的态射。设 $\mathcal{N}$、$\mathcal{L}$
为可逆 $\mathcal{O}_X$-模。假设 $S$ 拟紧、$f$ 为有限型，
且 $\mathcal{L}$ 是 $f$-丰沛的。则
$\mathcal{N} \otimes_{\mathcal{O}_X} \mathcal{L}^{\otimes d}$
是 $f$-甚丰沛的，对所有 $d \gg 1$ 都如此。
\end{lemma}

\begin{proof}
由引理 \ref{morphisms-lemma-characterize-very-ample-on-finite-type}，可归约到
$S$ 仿射的情形。结合引理
\ref{morphisms-lemma-finite-type-over-affine-ample-very-ample} 与性质，命题
\ref{properties-proposition-characterize-ample}，可找到整数 $d_0$，使得
$\mathcal{N} \otimes \mathcal{L}^{\otimes d_0}$ 全局生成。选择它的全局截面
$s_0, \ldots, s_n$，即 $\mathcal{N} \otimes \mathcal{L}^{\otimes d_0}$
的生成元。这确定一个态射 $j : X \to \mathbf{P}^n_S$（在 $S$ 上）。
由引理 \ref{morphisms-lemma-finite-type-over-affine-ample-very-ample}，还可选取整数
$d_1$，使得对所有 $d \geq d_1$，存在截面
$t_{d, 0}, \ldots, t_{d, n(d)}$（属于 $\mathcal{L}^{\otimes d}$），
它们生成该层并定义浸入
$$
j_d = \varphi_{\mathcal{L}^{\otimes d}, t_{d, 0}, \ldots, t_{d, n(d)}} :
X
\longrightarrow
\mathbf{P}^{n(d)}_S
$$
（在 $S$ 上）。于是，对 $d \geq d_0 + d_1$，可考虑态射
$$
\varphi_{\mathcal{N} \otimes \mathcal{L}^{\otimes d},
s_j \otimes t_{d - d_0, i}} :
X
\longrightarrow
\mathbf{P}^{(n + 1)(n(d - d_0) + 1) - 1}_S
$$
此态射是浸入，因为它是复合
$$
X \to \mathbf{P}^n_S \times_S \mathbf{P}^{n(d - d_0)}_S
\to \mathbf{P}^{(n + 1)(n(d - d_0) + 1) - 1}_S
$$
其中第一个态射是 $(j, j_{d - d_0})$，第二个态射是 Segre 嵌入
（构造，引理 \ref{constructions-lemma-segre-embedding}）。
由于 $j$ 是浸入，$(j, j_{d - d_0})$ 也是浸入（应用引理
\ref{morphisms-lemma-immersion-permanence}）。这是浸入的复合，因而仍为浸入
（概形，引理 \ref{schemes-lemma-composition-immersion}）。
\end{proof}







\section{拟射影态射}
\label{morphisms-section-quasi-projective}

\noindent
上一节的讨论提示了下列定义。我们采用 \cite{EGA} 中的拟射影定义；
带字母 ``H'' 的版本是 \cite{H} 中的定义。

\begin{definition}
\label{morphisms-definition-quasi-projective}
\begin{reference}
\cite[II, Definition 5.3.1]{EGA} and \cite[page 103]{H}
\end{reference}
设 $f : X \to S$ 是概形的态射。
\begin{enumerate}
\item 称 $f$ {\it 拟射影}，如果 $f$ 为有限型，且存在一个
$f$-相对丰沛的可逆 $\mathcal{O}_X$-模。
\item 称 $f$ {\it H-拟射影}，如果存在拟紧浸入
$X \to \mathbf{P}^n_S$（在 $S$ 上），其中 $n$ 为某个整数。
\footnote{这与 Hartshorne 中的定义不完全相同。确切地说，
Hartshorne 中的定义（1997 年第八次修订印刷）要求 $f$
是开浸入后接 H-射影态射的复合（见定义
\ref{morphisms-definition-projective}），而这并不蕴含 $f$ 拟紧。反方向的蕴含见引理
\ref{morphisms-lemma-H-quasi-projective-open-H-projective}。}
\item 称 $f$ {\it 局部拟射影}，如果存在开覆盖
$S = \bigcup V_j$，使得每个 $f^{-1}(V_j) \to V_j$ 都拟射影。
\end{enumerate}
\end{definition}

\noindent
正如这个定义所表明的，拟射影性质在 $S$ 上不是局部的。
稍后我们将能进一步说明拟射影概形的范畴；见态射进阶，节
\ref{more-morphisms-section-quasi-projective}。

\begin{lemma}
\label{morphisms-lemma-base-change-quasi-projective}
拟射影态射的基变换仍为拟射影态射。
\end{lemma}

\begin{proof}
这由引理 \ref{morphisms-lemma-base-change-finite-type} 和
\ref{morphisms-lemma-ample-base-change} 得出。
\end{proof}

\begin{lemma}
\label{morphisms-lemma-composition-quasi-projective}
设 $f : X \to Y$ 和 $g : Y \to S$ 是概形的态射。若 $S$ 拟紧，
且 $f$ 和 $g$ 都拟射影，则 $g \circ f$ 拟射影。
\end{lemma}

\begin{proof}
这由引理 \ref{morphisms-lemma-composition-finite-type} 和
\ref{morphisms-lemma-ample-composition} 得出。
\end{proof}

\begin{lemma}
\label{morphisms-lemma-quasi-projective-properties}
设 $f : X \to S$ 是概形的态射。若 $f$ 拟射影、H-拟射影或局部拟射影，
则 $f$ 是有限型的分离态射。
\end{lemma}

\begin{proof}
省略。
\end{proof}

\begin{lemma}
\label{morphisms-lemma-H-quasi-projective-quasi-projective}
H-拟射影态射是拟射影态射。
\end{lemma}

\begin{proof}
省略。
\end{proof}

\begin{lemma}
\label{morphisms-lemma-characterize-locally-quasi-projective}
设 $f : X \to S$ 是概形的态射。下列条件等价：
\begin{enumerate}
\item 态射 $f$ 局部拟射影。
\item 存在开覆盖 $S = \bigcup V_j$，使得每个
$f^{-1}(V_j) \to V_j$ 都 H-拟射影。
\end{enumerate}
\end{lemma}

\begin{proof}
由引理 \ref{morphisms-lemma-H-quasi-projective-quasi-projective} 可见 (2) 蕴含 (1)。
假设 (1)。这个问题在 $S$ 上是局部的，故可假设 $S$ 仿射、
$X$ 在 $S$ 上为有限型，并且 $\mathcal{L}$ 是 $X/S$ 上的相对丰沛可逆层。
由引理 \ref{morphisms-lemma-finite-type-over-affine-ample-very-ample}，可假设
$\mathcal{L}$ 在 $X$ 上丰沛。由引理
\ref{morphisms-lemma-quasi-projective-finite-type-over-S}，可见存在 $X$
到 $S$ 上某个射影空间的浸入，即 $X$ 在 $S$ 上 H-拟射影，
如所需。
\end{proof}

\begin{lemma}
\label{morphisms-lemma-quasi-affine-finite-type-quasi-projective}
\begin{reference}
\cite[II, Proposition 5.3.4 (i)]{EGA}
\end{reference}
有限型拟仿射态射是拟射影态射。
\end{lemma}

\begin{proof}
这由引理 \ref{morphisms-lemma-quasi-affine-O-ample} 得出。
\end{proof}

\begin{lemma}
\label{morphisms-lemma-quasi-projective-permanence}
设 $g : Y \to S$ 和 $f : X \to Y$ 是概形的态射。若
$g \circ f$ 拟射影且 $f$ 拟紧\footnote{若 $g$ 拟分离，则这由概形，引理
\ref{schemes-lemma-quasi-compact-permanence} 得出。}，
则 $f$ 拟射影。
\end{lemma}

\begin{proof}
由引理 \ref{morphisms-lemma-permanence-finite-type}，注意到 $f$ 为有限型。
因而引理由引理 \ref{morphisms-lemma-ample-permanence} 及定义得出。
\end{proof}





\section{固有态射}
\label{morphisms-section-proper}

\noindent
固有态射的概念在代数几何中起着重要作用。固有态射的一个重要例子
将是结构态射 $\mathbf{P}^n_S \to S$，即射影 $n$-空间的结构态射；事实上，
正是这个启发性例子引出了该定义。

\begin{definition}
\label{morphisms-definition-proper}
设 $f : X \to S$ 是概形的态射。若 $f$ 分离、为有限型且泛闭，
则称 $f$ {\it 固有}。
\end{definition}

\noindent
从零点加倍的仿射直线到仿射直线的态射为有限型且泛闭，
因此上述定义中的分离条件是必要的。本节其余部分将证明固有态射
和泛闭态射的一些基本性质。

\begin{lemma}
\label{morphisms-lemma-universally-closed-local-on-the-base}
设 $f : X \to S$ 是概形的态射。下列条件等价：
\begin{enumerate}
\item 态射 $f$ 泛闭。
\item 存在开覆盖 $S = \bigcup V_j$，使得
$f^{-1}(V_j) \to V_j$ 对所有指标 $j$ 均泛闭。
\end{enumerate}
\end{lemma}

\begin{proof}
这由定义显然可知。
\end{proof}

\begin{lemma}
\label{morphisms-lemma-proper-local-on-the-base}
设 $f : X \to S$ 是概形的态射。下列条件等价：
\begin{enumerate}
\item 态射 $f$ 固有。
\item 存在开覆盖 $S = \bigcup V_j$，使得
$f^{-1}(V_j) \to V_j$ 对所有指标 $j$ 均固有。
\end{enumerate}
\end{lemma}

\begin{proof}
省略。
\end{proof}

\begin{lemma}
\label{morphisms-lemma-composition-proper}
固有态射的复合是固有态射。泛闭态射也同样如此。
\end{lemma}

\begin{proof}
闭态射的复合是闭态射。若
$X \to Y \to Z$ 是泛闭态射，且 $Z' \to Z$ 是任意态射，
则可见
$Z' \times_Z X = (Z' \times_Z Y) \times_Y X  \to Z' \times_Z Y$
是闭的，并且 $Z' \times_Z Y \to Z'$ 是闭的。
这就证明了泛闭态射的结论。我们已经看到，“分离”和“有限型”
在复合下保持（概形，引理
\ref{schemes-lemma-separated-permanence} 及引理
\ref{morphisms-lemma-composition-finite-type}）。因此，固有态射的结论也成立。
\end{proof}

\begin{lemma}
\label{morphisms-lemma-base-change-proper}
固有态射的基变换是固有态射。泛闭态射也同样如此。
\end{lemma}

\begin{proof}
对泛闭态射而言，这由定义成立。对分离态射而言，这成立
（概形，引理 \ref{schemes-lemma-separated-permanence}）。
对有限型态射而言，这成立（引理 \ref{morphisms-lemma-base-change-finite-type}）。
因此，对固有态射也成立。
\end{proof}

\begin{lemma}
\label{morphisms-lemma-closed-immersion-proper}
闭浸入是固有的，因而更是泛闭的。
\end{lemma}

\begin{proof}
闭浸入的基变换是闭浸入（概形，引理
\ref{schemes-lemma-base-change-immersion}），因而泛闭。
闭浸入分离（概形，引理
\ref{schemes-lemma-immersions-monomorphisms}）。
闭浸入为有限型（引理 \ref{morphisms-lemma-immersion-locally-finite-type}）。
因此闭浸入固有。
\end{proof}

\begin{lemma}
\label{morphisms-lemma-image-proper-scheme-closed}
设给定概形的交换图
$$
\xymatrix{
X \ar[rr] \ar[rd] & &
Y \ar[ld] \\
& S &
}
$$
其中 $Y$ 在 $S$ 上分离。
\begin{enumerate}
\item 若 $X \to S$ 泛闭，则态射 $X \to Y$ 泛闭。
\item 若 $X$ 在 $S$ 上固有，则态射 $X \to Y$ 固有。
\end{enumerate}
特别地，在两种情形中，$X$ 在 $Y$ 中的像都是闭的。
\end{lemma}

\begin{proof}
假设 $X \to S$ 泛闭（相应地，固有）。将该态射分解为
$X \to X \times_S Y \to Y$。第一个态射是闭浸入；见概形，引理
\ref{schemes-lemma-semi-diagonal}。因此，第一个态射固有
（引理 \ref{morphisms-lemma-closed-immersion-proper}）。投影
$X \times_S Y \to Y$ 是一个泛闭（相应地，固有）态射的基变换，
因而泛闭（相应地，固有）；见引理 \ref{morphisms-lemma-base-change-proper}。
于是，作为泛闭（相应地，固有）态射的复合，$X \to Y$
泛闭（相应地，固有）（引理 \ref{morphisms-lemma-composition-proper}）。
\end{proof}

\noindent
下列引理的证明归功于 Bjorn Poonen；见
\href{https://mathoverflow.net/questions/23337/is-a-universally-closed-morphism-of-schemes-quasi-compact/23528#23528}{此处}。

\begin{lemma}
\label{morphisms-lemma-universally-closed-quasi-compact}
\begin{reference}
Due to Bjorn Poonen.
\end{reference}
概形的泛闭态射是拟紧的。
\end{lemma}

\begin{proof}
设 $f : X \to S$ 是态射。假设 $f$ 不拟紧。目标是证明 $f$
不泛闭。由概形，引理 \ref{schemes-lemma-quasi-compact-affine}，
存在仿射开集 $V \subset S$，使得 $f^{-1}(V)$ 不拟紧。
为达到目标，只须证明 $f^{-1}(V) \to V$ 不泛闭；因此可假设
$S = \Spec(A)$，其中 $A$ 表示某个环。

\medskip\noindent
写 $X = \bigcup_{i \in I} X_i$，其中诸 $X_i$ 是 $X$ 的仿射开子概形。
令 $T = \Spec(A[y_i ; i \in I])$。令 $T_i = D(y_i) \subset T$。
令 $Z$ 为闭集
$(X \times_S T) - \bigcup_{i \in I} (X_i \times_S T_i)$。
只须证明像 $f_T(Z)$ 不是闭的；这里，$Z$ 在
$f_T : X \times_S T \to T$ 下取像。

\medskip\noindent
存在点 $s \in S$，使得不存在邻域 $U$（它是 $s$ 在 $S$ 中的邻域），
满足 $X_U$ 拟紧。否则可覆盖 $S$，只需有限多个这样的 $U$，而概形，引理
\ref{schemes-lemma-quasi-compact-affine} 将蕴含 $f$ 拟紧。
固定这样的 $s \in S$。

\medskip\noindent
首先验证 $f_T(Z_s) \ne T_s$。令 $t \in T$ 为位于 $s$ 上方的点，
满足 $\kappa(t) = \kappa(s)$，并且 $y_i = 1$
在 $\kappa(t)$ 中成立，对所有 $i$ 都如此。于是有 $t \in T_i$，
对所有 $i$ 都如此；而
$Z_s \to T_s$ 在 $t$ 上方的纤维同构于
$(X - \bigcup_{i \in I} X_i)_s$，后者为空。因此
$t \in T_s - f_T(Z_s)$。

\medskip\noindent
假设 $f_T(Z)$ 在 $T$ 中闭。则存在元素
$g \in A[y_i; i \in I]$，使得 $f_T(Z) \subset V(g)$，但
$t \not \in V(g)$。因此 $g$ 在 $\kappa(t)$ 中的像非零。
特别地，$g$ 的某个系数在 $\kappa(s)$ 中的像非零，因而这个系数
在某个邻域 $U$ 上可逆，其中该邻域包含 $s$。令 $J$ 为所有满足
$j \in I$ 且 $y_j$ 出现在 $g$ 中的指标所成的有限集。由于 $X_U$ 不拟紧，
可选择点 $x \in X - \bigcup_{j \in J} X_j$，它位于某个
$u \in U$ 上方。由于 $g$ 有一个在 $U$ 上可逆的系数，
可找到点 $t' \in T$，它位于 $u$ 上方，使得
$t' \not \in V(g)$，并且 $t' \in V(y_i)$ 对所有
$i \notin J$ 都成立。这是因为
$V(y_i; i \in I, i \not\in J) = \Spec(A[y_j; j\in J])$，
而这个概形中位于 $u$ 上方的点集与
$\Spec(\kappa(u)[y_j; j \in J])$ 双射。换言之，
$t' \notin T_i$ 对每个 $i \notin J$ 都成立。由概形，引理
\ref{schemes-lemma-points-fibre-product}，可找到
点 $z$，它属于 $X \times_S T$，并映到 $x \in X$ 和 $t' \in T$。
由于 $x \not \in X_j$ 对 $j \in J$ 成立，而
$t' \not \in T_i$ 对 $i \in I \setminus J$ 成立，故 $z \in Z$。
另一方面，$f_T(z) = t' \not \in V(g)$，这与
$f_T(Z) \subset V(g)$ 矛盾。因此，“$f_T(Z)$ 闭”的假设是错误的，
确实可得 $f_T$ 不闭，如所需。
\end{proof}

\noindent
下列引理说明，固有概形在基上有限型的分离概形中的像仍然固有。

\begin{lemma}
\label{morphisms-lemma-image-proper-is-proper}
设 $S$ 是概形。设 $f : X \to Y$ 是 $S$ 上概形的态射。
若 $X$ 在 $S$ 上泛闭且 $f$ 满射，则 $Y$ 在 $S$ 上泛闭。
特别地，若 $Y$ 在 $S$ 上还分离且局部为有限型，则 $Y$ 在 $S$ 上固有。
\end{lemma}

\begin{proof}
假设 $X$ 泛闭且 $f$ 满射。记 $p : X \to S$、$q : Y \to S$
为结构态射。设 $S' \to S$ 是概形的态射。基变换
$f' : X_{S'} \to Y_{S'}$ 是满射
（引理 \ref{morphisms-lemma-base-change-surjective}），而基变换
$p' : X_{S'} \to S'$ 是闭态射。若 $T \subset Y_{S'}$ 是闭集，
则 $(f')^{-1}(T) \subset X_{S'}$ 是闭集，因而
$p'((f')^{-1}(T)) = q'(T)$ 是闭集。所以 $q'$ 是闭态射。
这证明了第一个断言。于是由引理
\ref{morphisms-lemma-universally-closed-quasi-compact}，$Y \to S$ 拟紧；
因此，若另外 $Y \to S$ 局部为有限型且分离，则根据定义，
$Y \to S$ 固有。
\end{proof}

\begin{lemma}
\label{morphisms-lemma-scheme-theoretic-image-is-proper}
设给定概形的交换图
$$
\xymatrix{
X \ar[rr]_h \ar[rd]_f & & Y \ar[ld]^g \\
& S
}
$$
假设
\begin{enumerate}
\item $X \to S$ 是泛闭态射（例如固有态射），并且
\item $Y \to S$ 分离且局部为有限型。
\end{enumerate}
则概形论像 $Z \subset Y$（即 $h$ 的概形论像）在 $S$ 上固有，且
$X \to Z$ 是满射。
\end{lemma}

\begin{proof}
$h$ 的概形论像在第 \ref{morphisms-section-scheme-theoretic-image} 节中构造。
由于 $f$ 拟紧（引理 \ref{morphisms-lemma-universally-closed-quasi-compact}），
可知 $h$ 拟紧（概形，引理
\ref{schemes-lemma-quasi-compact-permanence}）。因此
$h(X) \subset Z$ 稠密（引理
\ref{morphisms-lemma-quasi-compact-scheme-theoretic-image}）。另一方面，
$h(X)$ 在 $Y$ 中闭（引理
\ref{morphisms-lemma-image-proper-scheme-closed}），故 $X \to Z$ 是满射。
于是 $Z \to S$ 固有（引理 \ref{morphisms-lemma-image-proper-is-proper}）。
\end{proof}

\noindent
分离概形在满射泛闭态射下的目标是分离的。

\begin{lemma}
\label{morphisms-lemma-image-universally-closed-separated}
设 $S$ 是概形。设 $f : X \to Y$ 是 $S$ 上概形的满射泛闭态射。
\begin{enumerate}
\item 若 $X$ 拟分离，则 $Y$ 拟分离。
\item 若 $X$ 分离，则 $Y$ 分离。
\item 若 $X$ 在 $S$ 上拟分离，则 $Y$ 在 $S$ 上拟分离。
\item 若 $X$ 在 $S$ 上分离，则 $Y$ 在 $S$ 上分离。
\end{enumerate}
\end{lemma}

\begin{proof}
取 $S = \Spec(\mathbf{Z})$ 时，(1) 和 (2) 是 (3) 和 (4) 的推论
（见概形，定义 \ref{schemes-definition-separated}）。考虑交换图
$$
\xymatrix{
X \ar[d] \ar[rr]_{\Delta_{X/S}} & & X \times_S X \ar[d] \\
Y \ar[rr]^{\Delta_{Y/S}} & & Y \times_S Y
}
$$
左侧竖直箭头是满射（即泛满射）。右侧竖直箭头是泛闭的，因为它是
泛闭态射 $X \times_S X \to X \times_S Y \to Y \times_S Y$
的复合。因此它也拟紧；见引理
\ref{morphisms-lemma-universally-closed-quasi-compact}。

\medskip\noindent
假设 $X$ 在 $S$ 上拟分离，即 $\Delta_{X/S}$ 拟紧。
若 $V \subset Y \times_S Y$ 是拟紧开集，则
$V \times_{Y \times_S Y} X \to \Delta_{Y/S}^{-1}(V)$ 是满射，
且由上面的说明，$V \times_{Y \times_S Y} X$ 拟紧。
于是 $\Delta_{Y/S}$ 拟紧，即 $Y$ 在 $S$ 上拟分离。

\medskip\noindent
假设 $X$ 在 $S$ 上分离，即 $\Delta_{X/S}$ 是闭浸入。
那么 $X \to Y \times_S Y$ 作为闭态射的复合是闭态射。
由于 $X \to Y$ 满射，可知 $\Delta_{Y/S}(Y)$ 在
$Y \times_S Y$ 中闭。因此，由概形，定义
\ref{schemes-definition-separated} 后面的讨论，$Y$ 在 $S$ 上分离。
\end{proof}








\section{赋值判据}
\label{morphisms-section-valuative-criteria}

\noindent
我们已经在概形，第
\ref{schemes-section-valuative-criterion-universal-closedness} 节和第
\ref{schemes-section-valuative-separatedness} 节中讨论了泛闭性和分离性的
赋值判据。本节将讨论若干推论和变体。在极限，第
\ref{limits-section-Noetherian-valuative-criterion} 节中，我们将证明，
对局部 Noether 概形及有限型态射，只须考虑离散赋值环。

\begin{lemma}[固有性的赋值判据]
\label{morphisms-lemma-characterize-proper}
\begin{reference}
\cite[II Theorem 7.3.8]{EGA}
\end{reference}
设 $S$ 是概形。设 $f : X \to Y$ 是 $S$ 上概形的态射。
假设 $f$ 为有限型且拟分离。则下列条件等价：
\begin{enumerate}
\item $f$ 固有；
\item $f$ 满足赋值判据
（概形，定义 \ref{schemes-definition-valuative-criterion}）；
\item 对任意交换实线图
$$
\xymatrix{
\Spec(K) \ar[r] \ar[d] & X \ar[d] \\
\Spec(A) \ar[r] \ar@{-->}[ru] & Y
}
$$
其中 $A$ 是以 $K$ 为分式域的赋值环，都存在唯一的虚线箭头使图交换。
\end{enumerate}
\end{lemma}

\begin{proof}
(3) 是 (2) 的改述。因此，该引理是概形，命题
\ref{schemes-proposition-characterize-universally-closed}、引理
\ref{schemes-lemma-separated-implies-valuative}、引理
\ref{schemes-lemma-valuative-criterion-separatedness} 以及诸定义的形式推论。
\end{proof}

\noindent
检验赋值判据时，通常不必考虑所有可能的图。我们把下一个引理中的
赋值判据称为“精细赋值判据”。

\begin{lemma}
\label{morphisms-lemma-refined-valuative-criterion-universally-closed}
设 $f : X \to S$ 和 $h : U \to X$ 是概形的态射。
假设 $f$ 和 $h$ 拟紧，且 $h(U)$ 在 $X$ 中稠密。若对任意交换实线图
$$
\xymatrix{
\Spec(K) \ar[r] \ar[d] & U \ar[r]^h & X \ar[d]^f \\
\Spec(A) \ar[rr] \ar@{-->}[rru] & & S
}
$$
其中 $A$ 是以 $K$ 为分式域的赋值环，都存在唯一的虚线箭头使图交换，
则 $f$ 泛闭。若此外 $f$ 拟分离，则 $f$ 分离。
\end{lemma}

\begin{proof}
为证明 $f$ 泛闭，我们验证 $f$ 的赋值判据的存在性部分；由概形，命题
\ref{schemes-proposition-characterize-universally-closed}，这已经足够。
为此，考虑交换图
$$
\xymatrix{
\Spec(K) \ar[r] \ar[d] & X \ar[d] \\
\Spec(A) \ar[r] & S
}
$$
其中 $A$ 是赋值环，而 $K$ 是 $A$ 的分式域。注意，由于赋值环和域
都是既约的，根据概形，引理 \ref{schemes-lemma-map-into-reduction}，
可分别用各自的既约概形替换 $U$、$X$ 和 $S$。在这种情形下，
$h(U)$ 稠密这一假设意味着 $h : U \to X$ 的概形论像就是 $X$；
见引理 \ref{morphisms-lemma-scheme-theoretic-image-reduced}。还可以把 $S$ 替换为
$\Spec(A) \to S$ 所经过的一个仿射开集。因此可假设
$S = \Spec(R)$。

\medskip\noindent
设 $\Spec(B) \subset X$ 是 $\Spec(K) \to X$ 所经过的仿射开集。
选择一个多项式代数 $P$（定义在 $B$ 上）及一个 $B$-代数满射 $P \to K$。
于是 $\Spec(P) \to X$ 平坦。因此，由引理
\ref{morphisms-lemma-flat-base-change-scheme-theoretic-image}，态射
$U \times_X \Spec(P) \to \Spec(P)$ 的概形论像是 $\Spec(P)$。
由引理 \ref{morphisms-lemma-reach-points-scheme-theoretic-image}，可找到交换图
$$
\xymatrix{
\Spec(K') \ar[r] \ar[d] & U \times_X \Spec(P) \ar[d] \\
\Spec(A') \ar[r] & \Spec(P)
}
$$
其中 $A'$ 是赋值环，$K'$ 是 $A'$ 的分式域，并且
$\Spec(A')$ 的闭点映到 $\Spec(K) \subset \Spec(P)$。换言之，
存在 $B$-代数同态
$\varphi : K \to A'/\mathfrak m_{A'}$。选择一个赋值环
$A'' \subset A'/\mathfrak m_{A'}$，它支配 $\varphi(A)$，并以
$K'' = A'/\mathfrak m_{A'}$ 为分式域
（代数，引理 \ref{algebra-lemma-dominate}）。令
$$
C = \{\lambda \in A' \mid \lambda \bmod \mathfrak m_{A'} \in A''\}.
$$
由代数，引理 \ref{algebra-lemma-stack-valuation-rings}，这是一个赋值环。
由于 $C$ 是 $R$-代数，并以 $K'$ 为分式域，得到交换图
$$
\xymatrix{
\Spec(K') \ar[r] \ar[d] & U \ar[r] & X \ar[d] \\
\Spec(C) \ar[rr] \ar@{-->}[rru] & & S
}
$$
它具有引理陈述中的形式。因此存在按图所示的虚线箭头。由引理中的
唯一性假设，复合 $\Spec(A') \to \Spec(C) \to X$ 与给定态射
$\Spec(A') \to \Spec(P) \to \Spec(B) \subset X$ 一致。
因此，把该态射限制到
$C/\mathfrak m_{A'} = A''$ 的谱上，就得到给定态射
$\Spec(K'') = \Spec(A'/\mathfrak m_{A'}) \to \Spec(K) \to X$。
令 $x \in X$ 是 $\Spec(A'') \to X$ 下闭点的像。所诱导的环同态
$\mathcal{O}_{X, x} \to A''$ 的像是一个局部子环，它包含在
$K \subset K''$ 中。由于 $A$ 在 $K$ 中关于支配关系是极大的，
并且 $A \subset A''$，故 $A = K \cap A''$。因此
$\mathcal{O}_{X, x} \to A''$ 通过 $A \subset A''$ 分解。
这样便得到所需箭头 $\Spec(A) \to X$。

\medskip\noindent
最后，假设 $f$ 拟分离。那么 $\Delta : X \to X \times_S X$ 拟紧。
给定实线图
$$
\xymatrix{
\Spec(K) \ar[r] \ar[d] & U \ar[r]^h & X \ar[d]^\Delta \\
\Spec(A) \ar[rr] \ar@{-->}[rru] & & X \times_S X
}
$$
其中 $A$ 是以 $K$ 为分式域的赋值环，存在唯一的虚线箭头使图交换。
事实上，下方的水平箭头等价于一对态射 $\Spec(A) \to X$；
它们都可在引理的图中充当虚线箭头。因此，所要求的唯一性表明
下方的水平箭头通过 $\Delta$ 分解。于是可把刚证明的结果应用于
$\Delta : X \to X \times_S X$ 和 $h : U \to X$，得到
$\Delta$ 泛闭。显然，这意味着 $f$ 分离。
\end{proof}

\begin{remark}
\label{morphisms-remark-check-val-on-open}
引理 \ref{morphisms-lemma-refined-valuative-criterion-universally-closed} 中关于虚线箭头
唯一性的假设是必要的（细节省略）。当然，若 $f$ 分离，则唯一性有保证
（概形，引理 \ref{schemes-lemma-separated-implies-valuative}）。
\end{remark}

\begin{lemma}
\label{morphisms-lemma-morphism-defined-local-ring}
设 $S$ 是概形。设 $X$、$Y$ 是 $S$ 上的概形。设 $s \in S$，
而 $x \in X$、$y \in Y$ 是位于 $s$ 上方的点。
\begin{enumerate}
\item 设 $f, g : X \to Y$ 是 $S$ 上的态射，满足
$f(x) = g(x) = y$，并且
$f^\sharp_x = g^\sharp_x : \mathcal{O}_{Y, y} \to \mathcal{O}_{X, x}$。
在下列各情形中，存在开邻域 $U \subset X$，使得
$f|_U = g|_U$：
\begin{enumerate}
\item $Y$ 在 $S$ 上局部为有限型；
\item $X$ 整；
\item $X$ 局部 Noether；或
\item $X$ 既约且只有有限多个不可约分支。
\end{enumerate}
\item 设 $\varphi : \mathcal{O}_{Y, y} \to \mathcal{O}_{X, x}$
是局部 $\mathcal{O}_{S, s}$-代数同态。在下列各情形中，存在
开邻域 $U \subset X$，它包含 $x$，以及态射 $f : U \to Y$，
它把 $x$ 映到 $y$，并满足 $f^\sharp_x = \varphi$：
\begin{enumerate}
\item $Y$ 在 $S$ 上局部为有限表示；
\item $Y$ 局部为有限型且 $X$ 整；
\item $Y$ 局部为有限型且 $X$ 局部 Noether；或
\item $Y$ 局部为有限型且 $X$ 既约并只有有限多个不可约分支。
\end{enumerate}
\end{enumerate}
\end{lemma}

\begin{proof}
证明 (1)。可以把 $X$、$Y$、$S$ 分别替换为 $x$、$y$、$s$
的适当仿射开邻域，并归结为下列代数问题：给定环 $R$ 和两个
$R$-代数同态 $\varphi, \psi : B \to A$，满足
\begin{enumerate}
\item $R \to B$ 为有限型，或 $A$ 是整环，或 $A$ 是 Noether 环，
或 $A$ 既约且只有有限多个极小素理想；
\item 两个同态 $B \to A_\mathfrak p$ 对某个素理想
$\mathfrak p \subset A$ 相同；
\end{enumerate}
要证明，$\varphi, \psi$ 定义相同的同态 $B \to A_g$，其中可适当选择
$g \in A$，且 $g \not \in \mathfrak p$。若 $R \to B$ 为有限型，
设 $t_1, \ldots, t_m \in B$ 是 $B$ 作为 $R$-代数的一组生成元。
对每个 $j$，可找到 $g_j \in A$、$g_j \not \in \mathfrak p$，
使得 $\varphi(t_j)$ 和 $\psi(t_j)$ 在 $A_{g_j}$ 中的像相同。
然后令 $g = \prod g_j$。在其他情形中（即 $A$ 是整环、Noether 环，
或既约且只有有限多个极小素理想），可找到 $g \in A$、
$g \not \in \mathfrak p$，使得 $A_g \subset A_\mathfrak p$。
见代数，引理 \ref{algebra-lemma-subring-of-local-ring}。
因此同态 $B \to A_g$ 相同，如所需。

\medskip\noindent
证明 (2)。为此，可以把 $X$、$Y$ 和 $S$ 替换为适当的仿射开集。
设 $X = \Spec(A)$、$Y = \Spec(B)$ 且 $S = \Spec(R)$。
设 $\mathfrak p \subset A$ 是对应于 $x$ 的素理想。
设 $\mathfrak q \subset B$ 是对应于 $y$ 的素理想。于是
$\varphi$ 是局部 $R$-代数同态
$\varphi : B_\mathfrak q \to A_\mathfrak p$。
若 $R \to B$ 是有限表示的环同态，则存在
$g \in A \setminus \mathfrak p$ 及 $R$-代数同态 $B \to A_g$，使图
$$
\xymatrix{
B_\mathfrak q \ar[r]_\varphi & A_\mathfrak p \\
B \ar[u] \ar[r] & A_g \ar[u]
}
$$
交换；见代数，引理 \ref{algebra-lemma-characterize-finite-presentation} 和
\ref{algebra-lemma-localization-colimit}。所诱导的态射
$\Spec(A_g) \to \Spec(B)$ 即满足要求。若 $B$ 在 $R$ 上为有限型，
设 $t_1, \ldots, t_m \in B$ 是 $B$ 作为 $R$-代数的一组生成元。
于是可选择 $g_j \in A$、$g_j \not \in \mathfrak p$，使得
$\varphi(t_j) \in \Im(A_{g_j} \to A_\mathfrak p)$。
因此，把 $A$ 替换为 $A[1/\prod g_j]$ 后，可假设 $B$ 映入
$A \to A_\mathfrak p$ 的像中。若能找到 $g \in A$、
$g \not \in \mathfrak p$，使得 $A_g \to A_\mathfrak p$ 是单射，
就会得到所需的 $R$-代数同态 $B \to A_g$。于是，再次应用代数，引理
\ref{algebra-lemma-subring-of-local-ring} 即完成证明。
\end{proof}

\begin{lemma}
\label{morphisms-lemma-extend-across}
设 $S$ 是概形。设 $X$、$Y$ 是 $S$ 上的概形。设 $x \in X$。
设 $U \subset X$ 是开集，并设 $f : U \to Y$ 是 $S$ 上的态射。
假设
\begin{enumerate}
\item $x$ 在 $U$ 的闭包中；
\item $X$ 既约且只有有限多个不可约分支，或 $X$ 是 Noether 概形；
\item $\mathcal{O}_{X, x}$ 是赋值环；
\item $Y \to S$ 固有。
\end{enumerate}
则存在开集 $U \subset U' \subset X$，它包含 $x$，以及一个
$S$-态射 $f' : U' \to Y$，它延拓 $f$。
\end{lemma}

\begin{proof}
把 $X$ 替换为一个开邻域是无妨的：该邻域包含 $x$，并位于 $X$ 中
（略去一个小细节）。由性质，引理
\ref{properties-lemma-ring-affine-open-injective-local-ring}，
可假设 $X$ 仿射，且
$\Gamma(X, \mathcal{O}_X) \subset \mathcal{O}_{X, x}$。
特别地，$X$ 是整概形，具有唯一的泛点 $\xi$；其剩余域是分式域
$K$，即赋值环 $\mathcal{O}_{X, x}$ 的分式域。由于 $x$ 位于 $U$ 的闭包中，
可知 $U$ 非空，因而 $U$ 包含 $\xi$。于是，由固有性的赋值判据
（引理 \ref{morphisms-lemma-characterize-proper}），存在态射
$t : \Spec(\mathcal{O}_{X, x}) \to Y$，使其嵌入交换图
$$
\xymatrix{
\Spec(K) \ar[d]_\xi \ar[r] & \Spec(\mathcal{O}_{X, x}) \ar[d]_t \\
U \ar[r]^f & Y
}
$$
这个图由 $S$ 上概形的态射组成。应用引理
\ref{morphisms-lemma-morphism-defined-local-ring}，其中 $y = t(x)$ 且
$\varphi = t^\sharp_x$，可得到开邻域 $V \subset X$，它包含 $x$，
以及态射 $g : V \to Y$；该态射定义在 $S$ 上，并把 $x$ 映到
$y$，并满足 $g^\sharp_x = t^\sharp_x$。由于 $Y \to S$ 分离，
等化子 $E$（对应于 $f|_{U \cap V}$ 与 $g|_{U \cap V}$）是
$U \cap V$ 的闭子概形；见概形，引理
\ref{schemes-lemma-where-are-they-equal}。根据构造，$f$ 和 $g$
决定相同的态射 $\Spec(K) \to Y$，故 $E$ 包含整概形
$U \cap V$ 的泛点。因此 $E = U \cap V$，从而 $f$ 和 $g$
粘合为态射 $U' = U \cup V \to Y$，如所需。
\end{proof}





\section{射影态射}
\label{morphisms-section-projective}

\noindent
我们采用 \cite{EGA} 中射影态射的定义。带字母“H”的定义版本来自
\cite{H}。由此得到的定义并不相同，但两者都有用。

\begin{definition}
\label{morphisms-definition-projective}
设 $f : X \to S$ 是概形的态射。
\begin{enumerate}
\item 称 $f$ {\it 射影}，如果 $X$ 作为 $S$-概形同构于某个射影丛
$\mathbf{P}(\mathcal{E})$ 的闭子概形，其中所用的有限型
$\mathcal{O}_S$-模 $\mathcal{E}$ 是拟凝聚的。
\item 称 $f$ {\it H-射影}，如果存在整数 $n$ 以及闭浸入
$X \to \mathbf{P}^n_S$，后者定义在 $S$ 上。
\item 称 $f$ {\it 局部射影}，如果存在开覆盖 $S = \bigcup U_i$，使得每个
$f^{-1}(U_i) \to U_i$ 都射影。
\end{enumerate}
\end{definition}

\noindent
正如预期，射影态射拟射影；见引理
\ref{morphisms-lemma-projective-quasi-projective}。反之，拟射影态射常常是开浸入
与射影态射的复合；见引理
\ref{morphisms-lemma-quasi-projective-open-projective}。关于拟射影基上射影态射的
性质概览，见态射进阶，第 \ref{more-morphisms-section-projective} 节。

\begin{example}
\label{morphisms-example-projective}
设 $S$ 是概形。设 $\mathcal{A}$ 是拟凝聚分次
$\mathcal{O}_S$-代数，并由 $\mathcal{A}_1$ 在 $\mathcal{A}_0$ 上生成。
再假设 $\mathcal{A}_1$ 在 $\mathcal{O}_S$ 上为有限型。令
$X = \underline{\text{Proj}}_S(\mathcal{A})$。此时 $X \to S$ 射影。
事实上，由分次 $\mathcal{O}_S$-代数同态
$$
\text{Sym}_{\mathcal{O}_X}^*(\mathcal{A}_1)
\longrightarrow
\mathcal{A}
$$
所关联的态射是闭浸入；见构造，引理
\ref{constructions-lemma-surjective-generated-degree-1-map-relative-proj}。
\end{example}
 
\begin{lemma}
\label{morphisms-lemma-H-projective}
H-射影态射是 H-拟射影态射。H-射影态射是射影态射。
\end{lemma}

\begin{proof}
第一个断言由定义立即成立。第二个断言成立，因为
$\mathbf{P}^n_S$ 是 $S$ 上的射影丛；见构造，引理
\ref{constructions-lemma-projective-space-bundle}。
\end{proof}

\begin{lemma}
\label{morphisms-lemma-characterize-locally-projective}
设 $f : X \to S$ 是概形的态射。下列条件等价：
\begin{enumerate}
\item 态射 $f$ 局部射影。
\item 存在开覆盖 $S = \bigcup U_i$，使得每个
$f^{-1}(U_i) \to U_i$ 都 H-射影。
\end{enumerate}
\end{lemma}

\begin{proof}
由引理 \ref{morphisms-lemma-H-projective} 可见 (2) 蕴含 (1)。假设 (1)。
对每个点 $s \in S$，可找到仿射开邻域
$\Spec(R) = U \subset S$，它包含 $s$，使得 $X_U$ 同构于
$\mathbf{P}(\mathcal{E})$ 的某个闭子概形，其中所用的 $\mathcal{O}_U$-模
$\mathcal{E}$ 是有限型且拟凝聚的。写 $\mathcal{E} = \widetilde{M}$，其中
$R$-模 $M$ 是有限型的（见性质，引理
\ref{properties-lemma-finite-type-module}）。选择
$x_0, \ldots, x_n \in M$ 作为 $M$ 的一组 $R$-模生成元。
考虑分次 $R$-代数满射
$$
R[X_0, \ldots, X_n] \longrightarrow \text{Sym}_R(M).
$$
根据构造，引理
\ref{constructions-lemma-surjective-graded-rings-map-proj}，相应的态射
$$
\mathbf{P}(\mathcal{E}) \to \mathbf{P}^n_R
$$
是闭浸入。因此可知 $f^{-1}(U)$ 同构于
$\mathbf{P}^n_U$ 的闭子概形（作为 $U$ 上概形）。换言之，(2) 成立。
\end{proof}

\begin{lemma}
\label{morphisms-lemma-locally-projective-proper}
局部射影态射是固有态射。
\end{lemma}

\begin{proof}
设 $f : X \to S$ 局部射影。为证明 $f$ 固有，可以在基上局部地处理；
见引理 \ref{morphisms-lemma-proper-local-on-the-base}。因此，由上面的引理
\ref{morphisms-lemma-characterize-locally-projective}，可假设存在闭浸入
$X \to \mathbf{P}^n_S$。由引理 \ref{morphisms-lemma-composition-proper} 和
\ref{morphisms-lemma-closed-immersion-proper}，只须证明
$\mathbf{P}^n_S \to S$ 固有。由于 $\mathbf{P}^n_S \to S$ 是
$\mathbf{P}^n_{\mathbf{Z}} \to \Spec(\mathbf{Z})$ 的基变换，
由引理 \ref{morphisms-lemma-base-change-proper}，只须证明
$\mathbf{P}^n_{\mathbf{Z}} \to \Spec(\mathbf{Z})$ 固有。
由构造，引理 \ref{constructions-lemma-proj-separated}，概形
$\mathbf{P}^n_{\mathbf{Z}}$ 分离。由构造，引理
\ref{constructions-lemma-proj-quasi-compact}，概形
$\mathbf{P}^n_{\mathbf{Z}}$ 拟紧。显然
$\mathbf{P}^n_{\mathbf{Z}} \to \Spec(\mathbf{Z})$ 局部为有限型，因为
$\mathbf{P}^n_{\mathbf{Z}}$ 被仿射开集 $D_{+}(X_i)$ 覆盖，每个开集都是
有限型 $\mathbf{Z}$-代数
$$
\mathbf{Z}[X_0/X_i, \ldots, X_n/X_i].
$$
的谱。最后，还须证明
$\mathbf{P}^n_{\mathbf{Z}} \to \Spec(\mathbf{Z})$ 泛闭。这由构造，引理
\ref{constructions-lemma-proj-valuative-criterion} 及赋值判据得出
（见概形，命题
\ref{schemes-proposition-characterize-universally-closed}）。
\end{proof}

\begin{lemma}
\label{morphisms-lemma-proper-ample-locally-projective}
设 $f : X \to S$ 是概形的固有态射。若存在 $f$-丰沛可逆层，
且它定义在 $X$ 上，
则 $f$ 局部射影。
\end{lemma}

\begin{proof}
若存在 $f$-丰沛可逆层，则在 $S$ 上局部地可找到浸入
$i : X \to \mathbf{P}^n_S$；见引理
\ref{morphisms-lemma-finite-type-over-affine-ample-very-ample}。由于 $X \to S$ 固有，
态射 $i$ 是闭浸入；见引理 \ref{morphisms-lemma-image-proper-scheme-closed}。
\end{proof}

\begin{lemma}
\label{morphisms-lemma-H-projective-composition}
H-射影态射的复合是 H-射影态射。
\end{lemma}

\begin{proof}
假设 $X \to Y$ 和 $Y \to Z$ 都 H-射影。于是存在闭浸入
$X \to \mathbf{P}^n_Y$（定义在 $Y$ 上），以及闭浸入
$Y \to \mathbf{P}^m_Z$（定义在 $Z$ 上）。考虑下图
$$
\xymatrix{
X \ar[r] \ar[d] &
\mathbf{P}^n_Y \ar[r] \ar[dl] &
\mathbf{P}^n_{\mathbf{P}^m_Z} \ar[dl] \ar@{=}[r] &
\mathbf{P}^n_Z \times_Z \mathbf{P}^m_Z \ar[r] &
\mathbf{P}^{nm + n + m}_Z \ar[ddllll] \\
Y \ar[r] \ar[d] & \mathbf{P}^m_Z \ar[dl] & \\
Z & &
}
$$
这里，最右侧的顶部水平箭头是 Segre 嵌入；见构造，引理
\ref{constructions-lemma-segre-embedding}。该图把 $X$ 识别为
$\mathbf{P}^{nm + n + m}_Z$ 的闭子概形，如所需。
\end{proof}

\begin{lemma}
\label{morphisms-lemma-H-projective-base-change}
H-射影态射的基变换是 H-射影态射。
\end{lemma}

\begin{proof}
这是因为概形上射影空间的基变换仍是射影空间，并且闭浸入的基变换
仍是闭浸入；见概形，引理 \ref{schemes-lemma-base-change-immersion}。
\end{proof}

\begin{lemma}
\label{morphisms-lemma-base-change-projective}
（局部）射影态射的基变换是（局部）射影态射。
\end{lemma}

\begin{proof}
这是因为概形上射影丛的基变换仍是射影丛，有限型
$\mathcal{O}$-模的拉回仍为有限型
（模，引理 \ref{modules-lemma-pullback-finite-type}），并且闭浸入的
基变换仍是闭浸入；见概形，引理
\ref{schemes-lemma-base-change-immersion}。省略一些细节。
\end{proof}

\begin{lemma}
\label{morphisms-lemma-projective-quasi-projective}
射影态射是拟射影态射。
\end{lemma}

\begin{proof}
设 $f : X \to S$ 是射影态射。选择闭浸入
$i : X \to \mathbf{P}(\mathcal{E})$，其中 $\mathcal{E}$ 是拟凝聚的
有限型 $\mathcal{O}_S$-模。于是
$\mathcal{L} = i^*\mathcal{O}_{\mathbf{P}(\mathcal{E})}(1)$ 是
$f$-甚丰沛的。由于 $f$ 固有（引理
\ref{morphisms-lemma-locally-projective-proper}），它拟紧。因此，引理
\ref{morphisms-lemma-ample-very-ample} 蕴含 $\mathcal{L}$ 是 $f$-丰沛的。
由于 $f$ 固有，它为有限型。这样就验证了拟射影的全部定义性质，
结论成立。
\end{proof}

\begin{lemma}
\label{morphisms-lemma-H-quasi-projective-open-H-projective}
设 $f : X \to S$ 是 H-拟射影态射。则 $f$ 分解为
$X \to X' \to S$，其中 $X \to X'$ 是开浸入，而
$X' \to S$ 是 H-射影态射。
\end{lemma}

\begin{proof}
根据定义，可把 $f$ 分解为拟紧浸入
$i : X \to \mathbf{P}^n_S$，再复合投影 $\mathbf{P}^n_S \to S$。
由引理 \ref{morphisms-lemma-quasi-compact-immersion}，存在闭子概形
$X' \subset \mathbf{P}^n_S$，使得 $i$ 通过开浸入
$X \to X'$ 分解。引理得证。
\end{proof}

\begin{lemma}
\label{morphisms-lemma-quasi-projective-open-projective}
设 $f : X \to S$ 是拟射影态射，且 $S$ 拟紧并拟分离。
则 $f$ 分解为 $X \to X' \to S$，其中 $X \to X'$ 是开浸入，
而 $X' \to S$ 是射影态射。
\end{lemma}

\begin{proof}
设 $\mathcal{L}$ 为 $f$-丰沛层。由于 $f$ 为有限型且 $S$ 拟紧，
$\mathcal{L}^{\otimes n}$ 为 $f$-甚丰沛，其中 $n > 0$；见引理
\ref{morphisms-lemma-finite-type-ample-very-ample}。把
$\mathcal{L}$ 替换为 $\mathcal{L}^{\otimes n}$。写
$\mathcal{F} = f_*\mathcal{L}$。这是拟凝聚 $\mathcal{O}_S$-模；见概形，
引理 \ref{schemes-lemma-push-forward-quasi-coherent}
（拟射影态射拟紧且分离；见引理
\ref{morphisms-lemma-quasi-projective-properties}）。由性质，引理
\ref{properties-lemma-directed-colimit-finite-presentation}，可找到有向集
$I$ 及其上的有限型拟凝聚 $\mathcal{O}_S$-模系统
$\mathcal{E}_i$，使得该系统以 $I$ 为指标且
$\mathcal{F} = \colim \mathcal{E}_i$。考虑复合
$\psi_i : f^*\mathcal{E}_i \to f^*\mathcal{F} \to \mathcal{L}$。
选择有限仿射开覆盖 $S = \bigcup_{j = 1, \ldots, m} V_j$。
对每个 $j$，可选择截面
$$
s_{j, 0}, \ldots, s_{j, n_j} \in
\Gamma(f^{-1}(V_j), \mathcal{L}) = f_*\mathcal{L}(V_j) = \mathcal{F}(V_j)
$$
它们生成 $\mathcal{L}$ 在 $f^{-1}V_j$ 上的限制，并定义浸入
$$
f^{-1}V_j \longrightarrow \mathbf{P}^{n_j}_{V_j},
$$
见引理 \ref{morphisms-lemma-very-ample-finite-type-over-affine}。选择 $i$，使得存在
截面 $e_{j, t} \in \mathcal{E}_i(V_j)$，它们映到 $s_{j, t}$，
所取像位于 $\mathcal{F}$ 中；这一点对所有 $j = 1, \ldots, m$ 和
$t = 1, \ldots, n_j$ 都成立。于是同态 $\psi_i$ 满射，因为截面
$f^*e_{j, t}$ 与截面 $s_{j, t}$ 有相同的像，而后者生成
$\mathcal{L}|_{f^{-1}V_j}$。因此得到 $S$ 上的态射
$$
r_{\mathcal{L}, \psi_i} : X \longrightarrow \mathbf{P}(\mathcal{E}_i)
$$
使得在 $V_j$ 上，上述浸入分解为
$$
f^{-1}V_j \to \mathbf{P}(\mathcal{E}_i)|_{V_j} \to \mathbf{P}^{n_j}_{V_j}
$$
由引理 \ref{morphisms-lemma-immersion-permanence}，可知
$r_{\mathcal{L}, \psi_i}|_{V_j}$ 是浸入。由于 $S = \bigcup V_j$，
可知 $r_{\mathcal{L}, \psi_i}$ 是浸入。注意，
$r_{\mathcal{L}, \psi_i}$ 拟紧，因为 $X \to S$ 拟紧，且
$\mathbf{P}(\mathcal{E}_i) \to S$ 分离
（见概形，引理 \ref{schemes-lemma-quasi-compact-permanence}）。
由引理 \ref{morphisms-lemma-quasi-compact-immersion}，存在闭子概形
$X' \subset \mathbf{P}(\mathcal{E}_i)$，使得 $i$ 通过开浸入
$X \to X'$ 分解。于是 $X' \to S$ 根据定义射影，结论成立。
\end{proof}

\begin{lemma}
\label{morphisms-lemma-projective-is-quasi-projective-proper}
设 $S$ 是拟紧且拟分离的概形。设 $f : X \to S$ 是概形的态射。则
\begin{enumerate}
\item $f$ 射影，当且仅当 $f$ 拟射影且固有；并且
\item $f$ H-射影，当且仅当 $f$ H-拟射影且固有。
\end{enumerate}
\end{lemma}

\begin{proof}
若 $f$ 射影，则由引理 \ref{morphisms-lemma-projective-quasi-projective}，$f$ 拟射影，
并由引理 \ref{morphisms-lemma-locally-projective-proper}，它固有。反之，若
$X \to S$ 拟射影且固有，则由引理
\ref{morphisms-lemma-quasi-projective-open-projective}，可选择开浸入
$X \to X'$，其中 $X' \to S$ 射影。由于 $X \to S$ 固有，
可知 $X$ 在 $X'$ 中闭（引理
\ref{morphisms-lemma-image-proper-scheme-closed}），即 $X \to X'$ 是（开且）闭浸入。
由于 $X'$ 同构于 $S$ 上某个射影丛的闭子概形
（定义 \ref{morphisms-definition-projective}），可知 $X$ 也如此，即
$X \to S$ 是射影态射。这证明了 (1)。对 (2) 的证明相同，只须使用引理
\ref{morphisms-lemma-H-projective} 和
\ref{morphisms-lemma-H-quasi-projective-open-H-projective}。
\end{proof}

\begin{lemma}
\label{morphisms-lemma-composition-projective}
设 $f : X \to Y$ 和 $g : Y \to S$ 是概形的态射。若 $S$ 拟紧且拟分离，
并且 $f$ 和 $g$ 都射影，则 $g \circ f$ 射影。
\end{lemma}

\begin{proof}
由引理 \ref{morphisms-lemma-projective-quasi-projective} 和
\ref{morphisms-lemma-locally-projective-proper}，可见 $f$ 和 $g$ 拟射影且固有。
由引理 \ref{morphisms-lemma-composition-proper} 和
\ref{morphisms-lemma-composition-quasi-projective}，可见 $g \circ f$ 固有且拟射影。
因此，由引理 \ref{morphisms-lemma-projective-is-quasi-projective-proper}，
$g \circ f$ 射影。
\end{proof}

\begin{lemma}
\label{morphisms-lemma-projective-permanence}
设 $g : Y \to S$ 和 $f : X \to Y$ 是概形的态射。若 $g \circ f$ 射影，
且 $g$ 分离，则 $f$ 射影。
\end{lemma}

\begin{proof}
选择闭浸入 $X \to \mathbf{P}(\mathcal{E})$，其中 $\mathcal{E}$ 是拟凝聚的
有限型 $\mathcal{O}_S$-模。于是得到态射
$X \to \mathbf{P}(\mathcal{E}) \times_S Y$。该态射是闭浸入，因为它是复合
$$
X \to X \times_S Y \to \mathbf{P}(\mathcal{E}) \times_S Y
$$
其中第一个态射由概形，引理 \ref{schemes-lemma-semi-diagonal}
（以及 $g$ 分离这一事实）是闭浸入，第二个态射则是闭浸入的基变换。
最后，纤维积 $\mathbf{P}(\mathcal{E}) \times_S Y$ 同构于
$\mathbf{P}(g^*\mathcal{E})$，而拉回保持拟凝聚有限型模。
\end{proof}

\begin{lemma}
\label{morphisms-lemma-projective-over-quasi-projective-is-H-projective}
设 $S$ 是具有丰沛可逆层的概形。则
\begin{enumerate}
\item 任意射影态射 $X \to S$ 都 H-射影；并且
\item 任意拟射影态射 $X \to S$ 都 H-拟射影。
\end{enumerate}
\end{lemma}

\begin{proof}
关于 $S$ 的假设蕴含 $S$ 拟紧且分离；见性质，定义
\ref{properties-definition-ample}、引理
\ref{properties-lemma-ample-immersion-into-proj} 以及构造，引理
\ref{constructions-lemma-proj-separated}。因此，引理
\ref{morphisms-lemma-quasi-projective-open-projective} 适用，并可见 (1) 蕴含 (2)。
设 $\mathcal{E}$ 是有限型拟凝聚 $\mathcal{O}_S$-模。根据射影态射的定义，
只须证明 $\mathbf{P}(\mathcal{E}) \to S$ H-射影。若 $\mathcal{E}$
由有限多个全局截面生成，则相应的满射
$\mathcal{O}_S^{\oplus n} \to \mathcal{E}$ 诱导闭浸入
$$
\mathbf{P}(\mathcal{E}) \longrightarrow
\mathbf{P}(\mathcal{O}_S^{\oplus n}) = \mathbf{P}^{n - 1}_S
$$
如所需。一般地，设 $\mathcal{L}$ 是 $S$ 上的丰沛可逆层。由性质，命题
\ref{properties-proposition-characterize-ample}，存在整数 $n$，使得
$\mathcal{E} \otimes_{\mathcal{O}_S} \mathcal{L}^{\otimes n}$
由有限多个全局截面生成。由于
$\mathbf{P}(\mathcal{E}) =
\mathbf{P}(\mathcal{E} \otimes_{\mathcal{O}_S} \mathcal{L}^{\otimes n})$；见
构造，引理 \ref{constructions-lemma-twisting-and-proj}，证明完成。
\end{proof}

\begin{lemma}
\label{morphisms-lemma-proper-ample-is-proj}
设 $f : X \to S$ 是泛闭态射。设 $\mathcal{L}$ 是 $f$-丰沛可逆
$\mathcal{O}_X$-模。则引理 \ref{morphisms-lemma-characterize-relatively-ample} 中的
典范态射
$$
r : X
\longrightarrow
\underline{\text{Proj}}_S
\left(
\bigoplus\nolimits_{d \geq 0} f_*\mathcal{L}^{\otimes d}
\right)
$$
是同构。
\end{lemma}

\begin{proof}
注意 $f$ 拟紧，因为存在 $f$-丰沛可逆模这一事实迫使 $f$ 拟紧。
由所引引理，态射 $r$ 是开浸入。另一方面，由引理
\ref{morphisms-lemma-image-proper-scheme-closed}，$r$ 的像是闭的
（由构造，引理 \ref{constructions-lemma-relative-proj-separated}，
$r$ 的目标在 $S$ 上分离）。最后，由性质，引理
\ref{properties-lemma-ample-immersion-into-proj}，$r$ 的像稠密
（这里还使用了引理 \ref{morphisms-lemma-characterize-relatively-ample} 的证明中
已经表明的事实：态射 $r$ 在 $S$ 的仿射开集上由性质，引理
\ref{properties-lemma-map-into-proj} 的典范态射给出）。因此可知
$r$ 是满射开浸入，即同构。
\end{proof}

\begin{lemma}
\label{morphisms-lemma-proper-ample-delete-affine}
设 $f : X \to S$ 是泛闭态射。设 $\mathcal{L}$ 是 $f$-丰沛可逆
$\mathcal{O}_X$-模。设 $s \in \Gamma(X, \mathcal{L})$。则
$X_s \to S$ 是仿射态射。
\end{lemma}

\begin{proof}
该问题在 $S$ 上是局部的（引理 \ref{morphisms-lemma-characterize-affine}），
故可假设 $S$ 仿射。由引理 \ref{morphisms-lemma-proper-ample-is-proj}，可写
$X = \text{Proj}(A)$，其中 $A$ 是分次环，而 $s$ 对应于
$f \in A_1$，并且 $X_s = D_+(f)$（性质，引理
\ref{properties-lemma-map-into-proj}）。根据 $\text{Proj}(A)$ 的构造，
这证明了引理；见构造，第 \ref{constructions-section-proj} 节。
\end{proof}

\section{整态射与有限态射}
\label{morphisms-section-integral}

\noindent
回顾：若环同态 $R \to A$ 的每个 $A$ 中元素都满足一个系数在 $R$ 中的
首一方程，则称该同态{\it 整}。再回顾：若对环同态 $R \to A$，
$A$ 作为 $R$-模有限，则称该同态{\it 有限}。见代数，定义
\ref{algebra-definition-integral-ring-map}。

\begin{definition}
\label{morphisms-definition-integral}
设 $f : X \to S$ 是概形的态射。
\begin{enumerate}
\item 称 $f$ {\it 整}，如果 $f$ 仿射，并且对每个仿射开集
$\Spec(R) = V \subset S$，其逆像为
$\Spec(A) = f^{-1}(V) \subset X$ 时，相应的环同态 $R \to A$ 整。
\item 称 $f$ {\it 有限}，如果 $f$ 仿射，并且对每个仿射开集
$\Spec(R) = V \subset S$，其逆像为
$\Spec(A) = f^{-1}(V) \subset X$ 时，相应的环同态 $R \to A$ 有限。
\end{enumerate}
\end{definition}

\noindent
显然，整态射和有限态射分离且拟紧。有限态射显然也是有限型态射。
本节的大多数引理都是完全标准的。不过，请注意本节末尾有趣的引理
\ref{morphisms-lemma-integral-universally-closed}。

\begin{lemma}
\label{morphisms-lemma-integral-local}
设 $f : X \to S$ 是概形的态射。下列条件等价：
\begin{enumerate}
\item 态射 $f$ 整。
\item 存在仿射开覆盖 $S = \bigcup U_i$，使得每个
$f^{-1}(U_i)$ 仿射，并且
$\mathcal{O}_S(U_i) \to \mathcal{O}_X(f^{-1}(U_i))$ 整。
\item 存在开覆盖 $S = \bigcup U_i$，使得每个
$f^{-1}(U_i) \to U_i$ 都整。
\end{enumerate}
此外，若 $f$ 整，则对每个开子概形 $U \subset S$，态射
$f : f^{-1}(U) \to U$ 都整。
\end{lemma}

\begin{proof}
见代数，引理 \ref{algebra-lemma-integral-local}。省略一些细节。
\end{proof}

\begin{lemma}
\label{morphisms-lemma-finite-local}
设 $f : X \to S$ 是概形的态射。下列条件等价：
\begin{enumerate}
\item 态射 $f$ 有限。
\item 存在仿射开覆盖 $S = \bigcup U_i$，使得每个
$f^{-1}(U_i)$ 仿射，并且
$\mathcal{O}_S(U_i) \to \mathcal{O}_X(f^{-1}(U_i))$ 有限。
\item 存在开覆盖 $S = \bigcup U_i$，使得每个
$f^{-1}(U_i) \to U_i$ 都有限。
\end{enumerate}
此外，若 $f$ 有限，则对每个开子概形 $U \subset S$，态射
$f : f^{-1}(U) \to U$ 都有限。
\end{lemma}

\begin{proof}
见代数，引理 \ref{algebra-lemma-integral-local}。省略一些细节。
\end{proof}

\begin{lemma}
\label{morphisms-lemma-finite-integral}
有限态射是整态射。局部为有限型的整态射是有限态射。
\end{lemma}

\begin{proof}
见代数，引理 \ref{algebra-lemma-finite-is-integral} 和引理
\ref{algebra-lemma-characterize-finite-in-terms-of-integral}。
\end{proof}

\begin{lemma}
\label{morphisms-lemma-composition-finite}
有限态射的复合是有限态射。整态射也同样如此。
\end{lemma}

\begin{proof}
见代数，引理 \ref{algebra-lemma-finite-transitive} 和
\ref{algebra-lemma-integral-transitive}。
\end{proof}

\begin{lemma}
\label{morphisms-lemma-base-change-finite}
有限态射的基变换是有限态射。整态射也同样如此。
\end{lemma}

\begin{proof}
见代数，引理 \ref{algebra-lemma-base-change-integral}。
\end{proof}

\begin{lemma}
\label{morphisms-lemma-integral-universally-closed}
\begin{slogan}
整 $=$ 仿射 $+$ 泛闭
\end{slogan}
设 $f : X \to S$ 是概形的态射。下列条件等价：
\begin{enumerate}
\item $f$ 整；
\item $f$ 仿射且泛闭。
\end{enumerate}
\end{lemma}

\begin{proof}
设 (1) 成立。整态射按定义是仿射的。整态射的基变换仍整，故为证明 (2)，
只须证明整态射闭。这由代数，引理
\ref{algebra-lemma-integral-going-up} 和
\ref{algebra-lemma-going-up-closed} 得出。

\medskip\noindent
设 (2) 成立。可假设 $f$ 是态射 $f : \Spec(A) \to \Spec(R)$，
它由环同态 $R \to A$ 给出。取 $a$ 为 $A$ 的元素。我们须证明
$a$ 在 $R$ 上整，即在核 $I$ 中存在首一多项式；这里的同态是
$R[x] \to A$，它把 $x$ 映到 $a$。考虑环
$B = A[x]/(ax -1)$，并令 $J$ 为复合
$R[x]\to A[x] \to B$ 的核。若 $f\in J$，则存在 $q\in A[x]$，使得
$f = (ax-1)q$ 在 $A[x]$ 中成立。因此，若
$f = \sum_i f_ix^i$ 且 $q = \sum_iq_ix^i$，则对所有 $i \geq 0$，有
$f_i = aq_{i-1} - q_i$。对 $n \geq \deg q + 1$，多项式
$$
\sum\nolimits_{i \geq 0} f_i x^{n - i} =
\sum\nolimits_{i \geq 0} (a q_{i - 1} - q_i) x^{n - i} =
(a - x) \sum\nolimits_{i \geq 0} q_i x^{n - i - 1}
$$
显然属于 $I$；若 $f_0 = 1$，该多项式还首一。因此问题归结为证明
$J$ 含有常数项为 $1$ 的多项式。为此，我们证明
$\Spec(R[x]/(J + (x))$ 为空。


\medskip\noindent
由于 $f$ 泛闭，基变换 $\Spec(A[x]) \to \Spec(R[x])$ 闭。因此闭子集
$\Spec(B) \subset \Spec(A[x])$ 的像是闭子集
$\Spec(R[x]/J) \subset \Spec(R[x])$；见例
\ref{morphisms-example-scheme-theoretic-image} 和引理
\ref{morphisms-lemma-quasi-compact-scheme-theoretic-image}。特别地，
$\Spec(B) \to \Spec(R[x]/J)$ 满射。考虑下图，其中每个方块都是拉回方块：
$$
\xymatrix{
\Spec(B) \ar@{->>}[r]^g &
\Spec(R[x]/J)  \ar[r] &
\Spec(R[x])\\
\emptyset \ar[u] \ar[r] &
\Spec(R[x]/(J + (x)))\ar[u] \ar[r] &
\Spec(R) \ar[u]^0
}
$$
左下角为空，因为它是 $R\otimes_{R[x]} B$ 的谱，其中同态
$R[x]\to B$ 把 $x$ 映到一个可逆元，而同态 $R[x]\to R$ 把 $x$ 映到
$0$。由于 $g$ 满射，这蕴含 $\Spec(R[x]/(J + (x)))$ 为空，正如所需。
\end{proof}

\begin{lemma}
\label{morphisms-lemma-integral-fibres}
设 $f : X \to S$ 是整态射。则 $X$ 的每个点在其纤维中都是闭点。
\end{lemma}

\begin{proof}
见代数，引理 \ref{algebra-lemma-integral-no-inclusion}。
\end{proof}

\begin{lemma}
\label{morphisms-lemma-integral-dimension}
设 $f : X \to Y$ 是整态射。则 $\dim(X) \leq \dim(Y)$。
若 $f$ 满射，则 $\dim(X) = \dim(Y)$。
\end{lemma}

\begin{proof}
由于 $X$ 和 $Y$ 的维数是仿射开覆盖中各成员维数的上确界，可假设
$Y$ 和 $X$ 仿射。不等式由代数，引理
\ref{algebra-lemma-integral-dim-up} 得出。等式继而由代数，引理
\ref{algebra-lemma-dimension-going-up} 和
\ref{algebra-lemma-integral-going-up} 得出。
\end{proof}

\begin{lemma}
\label{morphisms-lemma-finite-quasi-finite}
有限态射是拟有限态射。
\end{lemma}

\begin{proof}
这由代数，引理 \ref{algebra-lemma-quasi-finite} 和引理
\ref{morphisms-lemma-quasi-finite-locally-quasi-compact} 推出。另一种证明是：由引理
\ref{morphisms-lemma-integral-fibres}（以及有限态射整这一事实），纤维中的所有点
都是闭点；再用引理 \ref{morphisms-lemma-quasi-finite-at-point-characterize} (3)，
可见 $f$ 在 $x$ 处拟有限，对所有 $x \in X$ 均如此。
\end{proof}

\begin{lemma}
\label{morphisms-lemma-finite-proper}
设 $f : X \to S$ 是概形的态射。下列条件等价：
\begin{enumerate}
\item $f$ 有限；
\item $f$ 仿射且固有。
\end{enumerate}
\end{lemma}

\begin{proof}
这形式地由引理 \ref{morphisms-lemma-integral-universally-closed}、有限态射整且分离这一
事实、固有态射等价于有限型、分离且泛闭的态射这一事实，以及有限型整态射
有限这一事实（引理 \ref{morphisms-lemma-finite-integral}）得出。
\end{proof}

\begin{lemma}
\label{morphisms-lemma-closed-immersion-finite}
闭浸入有限（从而尤其整）。
\end{lemma}

\begin{proof}
这是因为闭浸入仿射（引理 \ref{morphisms-lemma-closed-immersion-affine}），
而满射环同态有限且整。
\end{proof}

\begin{lemma}
\label{morphisms-lemma-finite-union-finite}
设 $X_i \to Y$，$i = 1, \ldots, n$ 是概形的有限态射。则
$X_1 \amalg \ldots \amalg X_n \to Y$ 也有限。
\end{lemma}

\begin{proof}
这由如下代数事实得出：若 $R \to A_i$，$i = 1, \ldots, n$ 是有限环同态，
则 $R \to A_1 \times \ldots \times A_n$ 也有限。
\end{proof}

\begin{lemma}
\label{morphisms-lemma-finite-permanence}
设 $f : X \to Y$ 和 $g : Y \to Z$ 是态射。
\begin{enumerate}
\item 若 $g \circ f$ 有限且 $g$ 分离，则 $f$ 有限。
\item 若 $g \circ f$ 整且 $g$ 分离，则 $f$ 整。
\end{enumerate}
\end{lemma}

\begin{proof}
设 $g \circ f$ 有限（相应地，整），且 $g$ 分离。由引理
\ref{morphisms-lemma-base-change-finite}，基变换 $X \times_Z Y \to Y$
有限（相应地，整）。态射 $X \to X \times_Z Y$ 是闭浸入，
因为 $Y \to Z$ 分离；见概形，引理
\ref{schemes-lemma-section-immersion}。闭浸入有限（相应地，整）；见引理
\ref{morphisms-lemma-closed-immersion-finite}。有限（相应地，整）态射的复合仍有限
（相应地，整）；见引理 \ref{morphisms-lemma-composition-finite}。证毕。
\end{proof}

\begin{lemma}
\label{morphisms-lemma-finite-monomorphism-closed}
设 $f : X \to Y$ 是概形的态射。若 $f$ 有限且为单态射，则 $f$ 是闭浸入。
\end{lemma}

\begin{proof}
这归结为代数，引理
\ref{algebra-lemma-finite-epimorphism-surjective}。
\end{proof}

\begin{lemma}
\label{morphisms-lemma-finite-projective}
有限态射是射影态射。
\end{lemma}

\begin{proof}
设 $f : X \to S$ 是有限态射。则 $f_*\mathcal{O}_X$ 是拟凝聚
$\mathcal{O}_S$-模（引理 \ref{morphisms-lemma-affine-equivalence-algebras}），
并且是有限型的（由有限态射的定义和性质，引理
\ref{properties-lemma-finite-type-module}）。我们断言存在 $S$ 上的闭浸入
$$
\sigma :
X
\longrightarrow
\mathbf{P}(f_*\mathcal{O}_X) =
\underline{\text{Proj}}_S(\text{Sym}^*_{\mathcal{O}_S}(f_*\mathcal{O}_X))
$$
这就完成证明。具体地，令 $\sigma$ 为（通过构造，引理
\ref{constructions-lemma-apply-relative}）对应于满射
$$
f^*f_*\mathcal{O}_X \longrightarrow \mathcal{O}_X
$$
的态射；该满射来自伴随映射 $f^*f_* \to \text{id}$。那么 $\sigma$
是闭浸入。事实上，在 $S$ 上仿射局部地，可写
$X = \Spec(A)$ 和 $S = \Spec(R)$。由于 $X$ 在 $S$ 上有限，
可选取 $a_1, \ldots, a_n \in A$，使其生成 $A$ 作为 $R$-模。于是
$R$-代数同态 $R[T_0, T_1, \ldots, T_n] \to \text{Sym}_R^*(A)$，
它把 $T_0$ 映到 $1$，并把 $T_i$ 映到 $a_i$（$1 \leq i \leq n$），
是满射。因此，由构造，引理
\ref{constructions-lemma-surjective-graded-rings-map-proj}，
$\mathbf{P}(f_*\mathcal{O}_X)$ 是 $\mathbf{P}^n_S$ 的闭子概形。
所以只须证明诱导态射 $X \to \mathbf{P}^n_S$ 是闭浸入
（例如见引理 \ref{morphisms-lemma-closed-immersion-ideals}）。读者可检验
$X \to \mathbf{P}^n_S$ 的像包含在开集
$D_+(T_0) \cong \mathbf{A}^n_S$ 中，并且 $X \to \mathbf{A}^n_S$
对应于满射 $R$-代数同态
$R[x_1, \ldots, x_n] \to \text{Sym}_R^*(A)$，它把 $x_i$ 映到 $a_i$。
因此 $X \to \mathbf{P}^n_S$ 是浸入（作为闭浸入与开浸入的复合）。由于
$X$ 在 $S$ 上有限，该浸入的像闭（这里使用引理
\ref{morphisms-lemma-finite-proper}、引理
\ref{morphisms-lemma-image-proper-scheme-closed} 和构造，引理
\ref{constructions-lemma-projective-space-separated}）。最后由概形，引理
\ref{schemes-lemma-immersion-when-closed} 得出结论。
\end{proof}

\section{泛同胚}
\label{morphisms-section-universal-homeomorphisms}

\noindent
下面的定义其实是多余的，因为泛同胚恰好就是整、泛单射且满射的态射；
见引理 \ref{morphisms-lemma-universal-homeomorphism}。

\begin{definition}
\label{morphisms-definition-universal-homeomorphism}
若概形态射 $f : X \to Y$ 的基变换
$f' : Y' \times_Y X \to Y'$ 对每个态射 $Y' \to Y$ 都是同胚，
则称它为{\it 泛同胚}。
\end{definition}

\noindent
先陈述几个必备引理。

\begin{lemma}
\label{morphisms-lemma-base-change-universal-homeomorphism}
概形的泛同胚沿任意概形态射作基变换后仍为泛同胚。
\end{lemma}

\begin{proof}
这由定义立即得出。
\end{proof}

\begin{lemma}
\label{morphisms-lemma-composition-universal-homeomorphism}
两个概形泛同胚的复合仍为泛同胚。
\end{lemma}

\begin{proof}
略。
\end{proof}

\noindent
下面的简单引理是刻画泛同胚的关键。

\begin{lemma}
\label{morphisms-lemma-homeomorphism-affine}
设 $f : X \to Y$ 是概形的态射。若 $f$ 是到 $Y$ 的某个闭子集上的同胚，
则 $f$ 仿射。
\end{lemma}

\begin{proof}
取点 $y \in Y$。若 $y \not \in f(X)$，则存在 $y$ 的仿射邻域与
$f(X)$ 不交。若 $y \in f(X)$，令 $x \in X$ 为 $X$ 中唯一映到 $y$
的点。取仿射开邻域 $y \in V$。取 $U \subset X$ 为 $x$ 的仿射开邻域，
使其映入 $V$。$f(U) \subset V \cap f(X)$ 由关于 $f$ 的假设
在诱导拓扑中开。因此可选取 $h \in \Gamma(V, \mathcal{O}_Y)$，使得
$y \in D(h)$ 且 $D(h) \cap f(X) \subset f(U)$。以
$h' \in \Gamma(U, \mathcal{O}_X)$ 表示 $f^\sharp(h)$ 在 $U$ 上的限制。
于是可见 $D(h') \subset U$ 等于 $f^{-1}(D(h))$。换言之，$Y$ 的每个点
都有逆像为仿射概形的开邻域。因此 $f$ 仿射；见引理
\ref{morphisms-lemma-characterize-affine}。
\end{proof}

\begin{lemma}
\label{morphisms-lemma-universal-homeomorphism}
设 $f : X \to Y$ 是概形的态射。下列条件等价：
\begin{enumerate}
\item $f$ 是泛同胚；
\item $f$ 整、泛单射且满射。
\end{enumerate}
\end{lemma}

\begin{proof}
设 $f$ 是泛同胚。由引理 \ref{morphisms-lemma-homeomorphism-affine}，可见 $f$
仿射。由于 $f$ 显然泛闭，引理
\ref{morphisms-lemma-integral-universally-closed} 表明 $f$ 整。还显然，$f$
泛单射且满射。

\medskip\noindent
设 $f$ 整、泛单射且满射。由引理
\ref{morphisms-lemma-integral-universally-closed}，$f$ 泛闭。它还泛双射（见引理
\ref{morphisms-lemma-base-change-surjective}），故它是泛同胚。
\end{proof}

\begin{lemma}
\label{morphisms-lemma-reduction-universal-homeomorphism}
设 $X$ 是概形。典范闭浸入 $X_{red} \to X$（见概形，定义
\ref{schemes-definition-reduced-induced-scheme}）是泛同胚。
\end{lemma}

\begin{proof}
略。
\end{proof}

\begin{lemma}
\label{morphisms-lemma-check-closed-infinitesimally}
设 $f : X \to S$ 和 $S' \to S$ 是概形的态射。假设
\begin{enumerate}
\item $S' \to S$ 是闭浸入；
\item $S' \to S$ 在点上双射；
\item $X \times_S S' \to S'$ 是闭浸入；
\item $X \to S$ 是有限型态射，或者 $S' \to S$ 是有限表示态射。
\end{enumerate}
则 $f : X \to S$ 是闭浸入。
\end{lemma}

\begin{proof}
假设 (1) 和 (2) 蕴含 $S' \to S$ 是泛同胚（例如，因为
$S_{red} = S'_{red}$，并使用引理
\ref{morphisms-lemma-reduction-universal-homeomorphism}）。因此 (3) 蕴含
$X \to S$ 是到 $S$ 的某个闭子集上的同胚。于是由引理
\ref{morphisms-lemma-homeomorphism-affine}，$X \to S$ 仿射。取仿射开集
$U \subset S$，记作 $U = \Spec(A)$。则由 (1)，
$S' = \Spec(A/I)$；并且由 (2)，$I$ 是局部幂零理想。由于 $f$ 仿射，
可写 $f^{-1}(U) = \Spec(B)$。假设 (4) 告诉我们，$B$ 是有限型
$A$-代数（引理 \ref{morphisms-lemma-locally-finite-type-characterize}），或者
$I$ 有限生成（引理 \ref{morphisms-lemma-closed-immersion-finite-presentation}）。
假设 (3) 即 $A/I \to B/IB$ 满射。若 $A \to B$ 为有限型，则由代数，
引理 \ref{algebra-lemma-surjective-mod-locally-nilpotent}；若 $I$ 有限生成
从而幂零，则由代数，引理 \ref{algebra-lemma-NAK}；均可推出
$A \to B$ 满射。这意味着 $f$ 是闭浸入；见引理
\ref{morphisms-lemma-closed-immersion}。
\end{proof}

\begin{lemma}
\label{morphisms-lemma-permanence-universal-homeomorphism}
设 $f : X \to Z$ 是两个态射 $g : X \to Y$ 和 $h : Y \to Z$ 的复合。
若态射 $\{f, g, h\}$ 中任意两个是泛同胚，则第三个也是泛同胚。
\end{lemma}

\begin{proof}
若 $g$ 和 $h$ 都是泛同胚，则由引理
\ref{morphisms-lemma-composition-universal-homeomorphism}，$f$ 也是泛同胚。

\medskip\noindent
设 $f$ 和 $g$ 都是泛同胚。我们要证明 $h$ 也是。沿概形的任意态射
$\alpha : Z' \to Z$ 对该图作基变换，得到下图，其中所有方块均笛卡尔：
$$
\xymatrix{
X' \ar[r]^{g'} \ar[d] & Y' \ar[r]^{h'} \ar[d] & Z' \ar[d] \\
X \ar[r]^{g} & Y \ar[r]^{h} & Z.
}
$$
我们的假设蕴含复合 $f'= h' \circ g' : X' \to Z'$ 和
$g' : X' \to Y'$ 都是同胚，因此 $h'$ 也是同胚。这就证明了 $h$
是泛同胚。

\medskip\noindent
最后，设 $f$ 和 $h$ 是泛同胚。我们要证明 $g$ 是泛同胚。
取概形的任意态射 $\beta : Y' \to Y$。得到下图，其中所有方块均笛卡尔：
$$
\xymatrix{
X' \ar[r]^{g'} \ar[d] & Y' \ar[d]^{\gamma} & \\
X'' \ar[r]^{g''} \ar[d] & Y'' \ar[r]^{h''} \ar[d] & Y' \ar[d]^{h \circ \beta} \\
X \ar[r]^{g} & Y \ar[r]^{h} & Z.
}
$$
这里，态射 $\gamma : Y' \to Y''$ 由纤维积的泛性质以及两个态射
$id_{Y'} : Y' \to Y'$ 和 $\beta : Y' \to Y$ 定义。我们将证明 $g'$
是同胚。由于成为同胚这一性质满足三取二性质，可见 $g''$ 是同胚。
观察上方方块，只须证明 $\gamma$ 是泛同胚。由于 $h''$ 是同胚，引理
\ref{morphisms-lemma-homeomorphism-affine} 表明它是仿射态射，从而尤其分离
（引理 \ref{morphisms-lemma-affine-separated}）。由于 $h'' \circ \gamma$ 是恒等态射，
概形，引理 \ref{schemes-lemma-section-immersion} 表明 $\gamma$ 是闭浸入。
由于 $h''$ 双射，可见 $\gamma$ 是双射闭浸入，因而是泛同胚
（例如由引理 \ref{morphisms-lemma-universal-homeomorphism} 的刻画），如所需。
\end{proof}

\section{仿射概形的泛同胚}
\label{morphisms-section-universal-homeomorphisms-affine}

\noindent
本节刻画仿射概形的泛同胚。

\begin{lemma}
\label{morphisms-lemma-subalgebra-inherits}
设 $A \to B$ 是环同态，使得诱导的概形态射
$f : \Spec(B) \to \Spec(A)$ 是泛同胚；相应地，是诱导剩余域同构的
泛同胚；相应地，泛闭；相应地，泛闭且泛单射。则对任意
$A$-子代数 $B' \subset B$，$f' : \Spec(B') \to \Spec(A)$ 也具有
相应性质。
\end{lemma}

\begin{proof}
若 $f$ 泛闭，则由引理 \ref{morphisms-lemma-integral-universally-closed}，
$B$ 在 $A$ 上整。因此 $B'$ 在 $A$ 上整，且 $f'$ 泛闭
（由同一引理）。这证明了 $f$ 泛闭的情形。

\medskip\noindent
继续可见 $B$ 在 $B'$ 上整（代数，引理
\ref{algebra-lemma-integral-permanence}），这蕴含
$\Spec(B) \to \Spec(B')$ 满射（代数，引理
\ref{algebra-lemma-integral-overring-surjective}）。因此，若
$A \to B$ 诱导剩余域的纯不可分扩张，则 $A \to B'$ 也如此。
这证明了 $f$ 泛闭且泛单射的情形；见引理
\ref{morphisms-lemma-universally-injective}。

\medskip\noindent
$f$ 为泛同胚的情形由上述论述、引理
\ref{morphisms-lemma-universal-homeomorphism} 以及如下显然观察得出：若 $f$
满射，则 $f'$ 也满射。

\medskip\noindent
若 $A \to B$ 诱导剩余域同构，则 $A \to B'$ 也如此
（见第二段的论证）。这样便证明了剩余情形。
\end{proof}

\begin{lemma}
\label{morphisms-lemma-colimit-inherits}
设 $A$ 是环。设 $B = \colim B_\lambda$ 是 $A$-代数的滤过余极限。
若每个 $f_\lambda : \Spec(B_\lambda) \to \Spec(A)$ 都是泛同胚；
相应地，是诱导剩余域同构的泛同胚；相应地，泛闭；相应地，泛闭且泛单射，
则 $f : \Spec(B) \to \Spec(A)$ 也具有相应性质。
\end{lemma}

\begin{proof}
若 $f_\lambda$ 泛闭，则由引理
\ref{morphisms-lemma-integral-universally-closed}，$B_\lambda$ 在 $A$ 上整。
因此 $B$ 在 $A$ 上整，且 $f$ 泛闭（由同一引理）。
这证明了每个 $f_\lambda$ 都泛闭的情形。

\medskip\noindent
对素理想 $\mathfrak q \subset B$，若它位于 $\mathfrak p \subset A$ 上方，
以 $\mathfrak q_\lambda \subset B_\lambda$ 表示其逆像。则
$\kappa(\mathfrak q) = \colim \kappa(\mathfrak q_\lambda)$。因此，若
$A \to B_\lambda$ 诱导剩余域的纯不可分扩张，则 $A \to B$ 也如此。
这证明了 $f_\lambda$ 泛闭且泛单射的情形；见引理
\ref{morphisms-lemma-universally-injective}。

\medskip\noindent
$f$ 为泛同胚的情形由上述论述和引理
\ref{morphisms-lemma-universal-homeomorphism} 得出，并使用如下事实：$B$ 中的素理想
等同于所有 $B_\lambda$ 中素理想的相容序列。

\medskip\noindent
若 $A \to B_\lambda$ 诱导剩余域同构，则 $A \to B$ 也如此
（见第二段的论证）。这样便证明了剩余情形。
\end{proof}

\begin{lemma}
\label{morphisms-lemma-special-elements-and-localization}
设 $A \subset B$ 是环扩张，$S \subset A$ 是乘法子集。设 $n \geq 1$，
并取 $b_i \in B$，其中 $1 \leq i \leq n$。任意
$x \in S^{-1}B$，若满足
$$
x \not \in S^{-1}A\text{ 且 } b_i x^i \in S^{-1}A\text{ 对 }i = 1, \ldots, n
$$
则等于 $s^{-1}y$，其中 $s \in S$ 且 $y \in B$，并满足
$$
y \not \in A\text{ 且 } b_i y^i \in A\text{ 对 }i = 1, \ldots, n
$$
\end{lemma}

\begin{proof}
略。提示：通分。
\end{proof}

\begin{lemma}
\label{morphisms-lemma-nth-and-nplusone-implies-square-and-cube}
设 $A \subset B$ 是环扩张。若存在 $b \in B$、$b \not \in A$ 和整数
$n \geq 2$，使得 $b^n \in A$ 且 $b^{n + 1} \in A$，则存在
$b' \in B$、$b' \not \in A$，使得 $(b')^2 \in A$ 且 $(b')^3 \in A$。
\end{lemma}

\begin{proof}
取引理中的 $b$ 和 $n$。则 $b$ 的所有充分高次幂均属于 $A$。
事实上，$(b^n)^k(b^{n + 1})^i = b^{(k + i)n + i}$，这蕴含任意幂
$b^m$ 在 $m \geq n^2$ 时都属于 $A$。因此，若 $i \geq 1$ 是使得
$b^i \not \in A$ 的最大整数，则 $(b^i)^2 \in A$ 且 $(b^i)^3 \in A$。
\end{proof}

\begin{lemma}
\label{morphisms-lemma-square-and-cube}
设 $A \subset B$ 是环扩张，使得 $\Spec(B) \to \Spec(A)$ 是诱导剩余域
同构的泛同胚。若 $A \not = B$，则存在 $b \in B$、$b \not \in A$，
使得 $b^2 \in A$ 且 $b^3 \in A$。
\end{lemma}

\begin{proof}
回顾：$A \subset B$ 整（引理
\ref{morphisms-lemma-integral-universally-closed}）。由引理
\ref{morphisms-lemma-subalgebra-inherits}，可假设 $B$ 在 $A$ 上由单个元素生成。
因此 $B$ 在 $A$ 上有限（代数，引理
\ref{algebra-lemma-characterize-finite-in-terms-of-integral}）。故
$B/A$ 作为 $A$-模的支撑闭且非空（代数，引理
\ref{algebra-lemma-support-closed} 和
\ref{algebra-lemma-support-zero}）。取 $\mathfrak p \subset A$ 为该支撑的
一个极小素理想。把 $A \subset B$ 替换为
$A_\mathfrak p \subset B_\mathfrak p$ 后（引理
\ref{morphisms-lemma-special-elements-and-localization} 允许这样做），可假设
$(A, \mathfrak m)$ 是局部环，$B$ 在 $A$ 上有限，并且 $B/A$ 的支撑是
$\{\mathfrak m\}$，这里视其为 $A$-模。由于 $B/A$ 是有限模，可见
$I = \text{Ann}_A(B/A)$ 满足 $\mathfrak m = \sqrt{I}$（代数，引理
\ref{algebra-lemma-support-closed}）。令 $\mathfrak m' \subset B$ 为位于
$\mathfrak m$ 上方的唯一素理想。由于
$\Spec(B) \to \Spec(A)$ 是同胚，可得 $\mathfrak m' = \sqrt{IB}$。
对 $f \in \mathfrak m'$，取 $n \geq 1$，使得 $f^n \in IB$。
则还有 $f^{n + 1} \in IB$。$IB \subset A$ 由我们对 $I$ 的选取成立，故
$f^n, f^{n + 1} \in A$。使用引理
\ref{morphisms-lemma-nth-and-nplusone-implies-square-and-cube}，若
$\mathfrak m' \not \subset A$，便可断定引理成立。然而，若
$\mathfrak m' \subset A$，则 $\mathfrak m' = \mathfrak m$；再由假设中
剩余域也同构，可得 $A = B$。这一矛盾完成证明。
\end{proof}

\begin{lemma}
\label{morphisms-lemma-pth-power-and-multiple}
设 $A \subset B$ 是环扩张，使得 $\Spec(B) \to \Spec(A)$ 是泛同胚。
若 $A \not = B$，则或者存在 $b \in B$、$b \not \in A$，使得
$b^2 \in A$ 且 $b^3 \in A$；或者存在素数 $p$ 和 $b \in B$、
$b \not \in A$，使得 $pb \in A$ 且 $b^p \in A$。
\end{lemma}

\begin{proof}
论证几乎与引理 \ref{morphisms-lemma-square-and-cube} 的证明完全相同，但为确保它
确实成立，我们写出全部细节。

\medskip\noindent
回顾：$A \subset B$ 整（引理
\ref{morphisms-lemma-integral-universally-closed}）。由引理
\ref{morphisms-lemma-subalgebra-inherits}，可假设 $B$ 在 $A$ 上由单个元素生成。
因此 $B$ 在 $A$ 上有限（代数，引理
\ref{algebra-lemma-characterize-finite-in-terms-of-integral}）。故
$B/A$ 作为 $A$-模的支撑闭且非空（代数，引理
\ref{algebra-lemma-support-closed} 和
\ref{algebra-lemma-support-zero}）。取 $\mathfrak p \subset A$ 为该支撑的
一个极小素理想。把 $A \subset B$ 替换为
$A_\mathfrak p \subset B_\mathfrak p$ 后（引理
\ref{morphisms-lemma-special-elements-and-localization} 允许这样做），可假设
$(A, \mathfrak m)$ 是局部环，$B$ 在 $A$ 上有限，并且 $B/A$ 的支撑是
$\{\mathfrak m\}$，这里视其为 $A$-模。由于 $B/A$ 是有限模，可见
$I = \text{Ann}_A(B/A)$ 满足 $\mathfrak m = \sqrt{I}$（代数，引理
\ref{algebra-lemma-support-closed}）。令 $\mathfrak m' \subset B$ 为位于
$\mathfrak m$ 上方的唯一素理想。由于
$\Spec(B) \to \Spec(A)$ 是同胚，可得 $\mathfrak m' = \sqrt{IB}$。
对 $f \in \mathfrak m'$，取 $n \geq 1$，使得 $f^n \in IB$。
则还有 $f^{n + 1} \in IB$。$IB \subset A$ 由我们对 $I$ 的选取成立，故
$f^n, f^{n + 1} \in A$。使用引理
\ref{morphisms-lemma-nth-and-nplusone-implies-square-and-cube}，若
$\mathfrak m' \not \subset A$，便可断定引理成立。若
$\mathfrak m' \subset A$，则 $\mathfrak m' = \mathfrak m$。由于
$A \not = B$，可知剩余域的同态
$\kappa = A/\mathfrak m \to B/\mathfrak m' = \kappa'$ 不可能是同构。
由引理 \ref{morphisms-lemma-universally-injective}，可知 $\kappa$ 的特征是某个素数
$p$，且扩张 $\kappa'/\kappa$ 纯不可分。取 $b \in B$，使其在
$\kappa'$ 中的像不属于 $\kappa$，但其 $p$ 次幂属于 $\kappa$。于是
$b \not \in A$、$b^p \in A$ 且 $pb \in A$（因为
$pb \in \mathfrak m' = \mathfrak m \subset A$），如所需。
\end{proof}

\begin{proposition}
\label{morphisms-proposition-universal-homeomorphism-equal-residue-fields}
设 $A \subset B$ 是环扩张。下列条件等价：
\begin{enumerate}
\item $\Spec(B) \to \Spec(A)$ 是诱导剩余域同构的泛同胚；
\item 每个有限子集 $E \subset B$ 都包含在扩张
$$
A[b_1, \ldots, b_n] \subset B
$$
中，并且有 $b_i^2, b_i^3 \in A[b_1, \ldots, b_{i - 1}]$，
其中 $i = 1, \ldots, n$。
\end{enumerate}
\end{proposition}

\begin{proof}
设 (1) 成立。用超限递归，对每个序数 $\alpha$ 如下构造一个
$A$-子代数 $B_\alpha \subset B$。置 $B_0 = A$。若 $\alpha$ 是极限序数，
则置 $B_\alpha = \colim_{\beta < \alpha} B_\beta$。若
$\alpha = \beta + 1$，则或者 $B_\beta = B$，这时置
$B_\alpha = B_\beta$；或者 $B_\beta \not = B$，这时应用引理
\ref{morphisms-lemma-square-and-cube}，选择 $b_\alpha \in B$、
$b_\alpha \not \in B_\beta$，使得 $b_\alpha^2, b_\alpha^3 \in B_\beta$，
并置 $B_\alpha = B_\beta[b_\alpha] \subset B$。显然
$B = \colim B_\alpha$（事实上，$B = B_\alpha$ 对某个序数 $\alpha$
成立，考察基数即可看出）。我们将用超限归纳证明：(2) 对
$A \to B_\alpha$ 对每个序数 $\alpha$ 成立。对 $\alpha = 0$ 显然成立。
设该陈述对每个 $\beta < \alpha$ 成立，并令 $E \subset B_\alpha$ 为
有限子集。若 $\alpha$ 是极限序数，则
$B_\alpha = \bigcup_{\beta < \alpha} B_\beta$，可见
$E \subset B_\beta$ 对某个 $\beta < \alpha$ 成立，这证明了该情形。若
$\alpha = \beta + 1$，则 $B_\alpha = B_\beta[b_\alpha]$。因此每个
$e \in E$ 都可写成多项式 $e = \sum d_{e, i}b_\alpha^i$，其中
$d_{e, i} \in B_\beta$。令 $D \subset B_\beta$ 为集合
$D = \{d_{e, i}\} \cup \{b_\alpha^2, b_\alpha^3\}$。由归纳假设，存在
命题中所述的 $A$-子代数 $A[b_1, \ldots, b_n] \subset B_\beta$，
它包含 $D$。于是 $A[b_1, \ldots, b_n, b_\alpha] \subset B_\alpha$
是命题中所述的 $A$-子代数，位于 $B_\alpha$ 中，并且包含 $E$。

\medskip\noindent
设 (2) 成立。把 $B = \colim B_\lambda$ 写成其有限 $A$-子代数的余极限。
由引理 \ref{morphisms-lemma-colimit-inherits}，只须证明
$\Spec(B_\lambda) \to \Spec(A)$ 是诱导剩余域同构的泛同胚。
泛闭态射的复合泛闭；诱导剩余域同构的态射也同样如此。因此只须证明：
若 $A \subset B$，且 $B$ 由单个元素 $b$ 生成，并满足
$b^2, b^3 \in A$，则 (1) 成立。这样的扩张整，故
$\Spec(B) \to \Spec(A)$ 泛闭且满射（引理
\ref{morphisms-lemma-integral-universally-closed} 和代数，引理
\ref{algebra-lemma-integral-overring-surjective}）。注意
$(b^2)^3 = (b^3)^2$ 在 $A$ 中成立。对任意环同态
$\varphi : A \to K$ 到域 $K$，可见存在 $\lambda \in K$，使得
$\varphi(b^2) = \lambda^2$ 且 $\varphi(b^3) = \lambda^3$。事实上，若
置 $\lambda = 0$ 若 $\varphi(b^2) = 0$；否则置
$\lambda = \varphi(b^3)/\varphi(b^2)$。因此 $B \otimes_A K$ 是
$K[x]/(x^2 - \lambda^2, x^3 - \lambda^3)$ 的商。该环恰有一个素理想，
其剩余域是 $K$。这蕴含 $\Spec(B) \to \Spec(A)$ 双射并诱导剩余域同构。
结合泛闭性可知 (1) 成立；见引理
\ref{morphisms-lemma-universal-homeomorphism} 和
\ref{morphisms-lemma-universally-injective}。
\end{proof}

\begin{proposition}
\label{morphisms-proposition-universal-homeomorphism}
设 $A \subset B$ 是环扩张。下列条件等价：
\begin{enumerate}
\item $\Spec(B) \to \Spec(A)$ 是泛同胚；
\item 每个有限子集 $E \subset B$ 都包含在扩张
$$
A[b_1, \ldots, b_n] \subset B
$$
中，并且对 $i = 1, \ldots, n$，有
\begin{enumerate}
\item $b_i^2, b_i^3 \in A[b_1, \ldots, b_{i - 1}]$；或
\item 存在素数 $p$，使得
$pb_i, b_i^p \in A[b_1, \ldots, b_{i - 1}]$。
\end{enumerate}
\end{enumerate}
\end{proposition}

\begin{proof}
证明与命题
\ref{morphisms-proposition-universal-homeomorphism-equal-residue-fields} 的证明完全相同，
但须作如下改动：
\begin{enumerate}
\item 用引理 \ref{morphisms-lemma-pth-power-and-multiple} 代替引理
\ref{morphisms-lemma-square-and-cube}。这意味着对每个后继序数
$\alpha = \beta + 1$，或者有 $b_\alpha^2, b_\alpha^3 \in B_\beta$，
或者存在素数 $p$，并有 $pb_\alpha, b_\alpha^p \in B_\beta$。
\item 若 $\alpha$ 是后继序数，则取
$D = \{d_{e, i}\} \cup \{b_\alpha^2, b_\alpha^3\}$，或取
$D = \{d_{e, i}\} \cup \{pb_\alpha, b_\alpha^p\}$，具体取决于
$\alpha$ 落入哪个情形。
\item 在 (2) $\Rightarrow$ (1) 的证明中，还须考虑如下情形：$B$ 在
$A$ 上由单个元素 $b$ 生成，且有 $pb, b^p \in B$，其中 $p$ 是某个素数。
此时 $A \subset B$ 在谱上诱导泛同胚，例如由代数，
引理 \ref{algebra-lemma-p-ring-map}。
\end{enumerate}
这完成证明。
\end{proof}

\begin{lemma}
\label{morphisms-lemma-universal-homeo-iso-if-invert-p}
设 $p$ 是素数。设 $A \to B$ 是诱导同构
$A[1/p] \to B[1/p]$ 的环同态（例如，若 $p$ 在 $A$ 中幂零）。
下列条件等价：
\begin{enumerate}
\item $\Spec(B) \to \Spec(A)$ 是泛同胚；
\item $A \to B$ 的核是局部幂零理想，并且对每个 $b \in B$，存在
$p$-幂 $q$，使得 $qb$ 和 $b^q$ 都在 $A \to B$ 的像中。
\end{enumerate}
\end{lemma}

\begin{proof}
若 (2) 成立，则由代数，引理 \ref{algebra-lemma-p-ring-map}，(1) 成立。
设 (1) 成立。由代数，引理
\ref{algebra-lemma-image-dense-generic-points}，$A \to B$ 的核由幂零元组成。
因此可把 $A$ 替换为 $A \to B$ 的像，并假设 $A \subset B$。由代数，引理
\ref{algebra-lemma-help-with-powers}，集合
$$
B' = \{b \in B \mid p^nb, b^{p^n} \in A\text{ 对某个 }n \geq 0\}
$$
是一个 $A$-子代数，包含于 $B$（对乘法封闭是显然的）。我们须证明
$B' = B$。否则，由引理 \ref{morphisms-lemma-pth-power-and-multiple}，存在
$b \in B$、$b \not \in B'$，使得或者 $b^2, b^3 \in B'$，或者存在素数
$\ell$，使得 $\ell b, b^\ell \in B'$。我们将证明两种情形都导致矛盾，
从而证明引理。

\medskip\noindent
由于 $A[1/p] = B[1/p]$，可选择 $p$-幂 $q$，使得 $qb \in A$。

\medskip\noindent
若 $b^2, b^3 \in B'$，则还有 $b^q \in B'$。由 $B'$ 的定义，
$(b^q)^{q'} \in A$ 对某个 $p$-幂 $q'$ 成立。于是
$qq'b, b^{qq'} \in A$，从而 $b \in B'$，矛盾。

\medskip\noindent
现在设存在素数 $\ell$，使得 $\ell b, b^\ell \in B'$。若
$\ell \not = p$，则 $\ell b \in B'$ 和 $qb \in A \subset B'$ 蕴含
$b \in B'$，矛盾。因此 $\ell = p$ 且 $b^p \in B'$，于是完全如前
得到矛盾。
\end{proof}

\begin{lemma}
\label{morphisms-lemma-make-universal-homeo}
设 $A$ 是环。设 $x, y \in A$。
\begin{enumerate}
\item 若 $x^3 = y^2$ 在 $A$ 中成立，则
$A \to B = A[t]/(t^2 - x, t^3 - y)$ 诱导剩余域上的双射，并诱导谱上的
泛同胚。
\item 若存在素数 $p$，使得 $p^px = y^p$ 在 $A$ 中成立，则
$A \to B = A[t]/(t^p - x, pt - y)$ 诱导谱上的泛同胚。
\end{enumerate}
\end{lemma}

\begin{proof}
我们用引理 \ref{morphisms-lemma-universal-homeomorphism} 的判据来检验。在两种情形中，
环同态都整。因此只须证明：给定域 $k$ 和环同态
$\varphi : A \to k$，$k$-代数 $B \otimes_A k$ 有唯一的素理想；
在情形 (1)，其剩余域等于 $k$；在情形 (2)，其剩余域在 $k$ 上纯不可分。
见引理 \ref{morphisms-lemma-universally-injective}。

\medskip\noindent
在情形 (1)，置 $\lambda = 0$ 若 $\varphi(x) = 0$；否则置
$\lambda = \varphi(y)/\varphi(x)$。那么
$B = k[t]/(t^2 - \lambda^2, t^3 - \lambda^2)$。结论于是显然。

\medskip\noindent
在情形 (2)，若 $k$ 的特征是 $p$，则得到 $\varphi(y) = 0$，且
$B = k[t]/(t^p - \varphi(x))$；这是局部 Artin $k$-代数，其剩余域或者是
$k$，或者是次数为 $p$ 的 $k$ 的纯不可分扩张。若 $k$ 的特征不是 $p$，
则置 $\lambda = \varphi(y)/p$，可见
$B = k[t]/(t - \lambda) = k$，同样得到结论。
\end{proof}

\begin{lemma}
\label{morphisms-lemma-universal-homeo-limit}
设 $A \to B$ 是环同态。
\begin{enumerate}
\item 若 $A \to B$ 在谱上诱导泛同胚，则
$B = \colim B_i$ 是有限表示 $A$-代数 $B_i$ 的滤过余极限，其中
$A \to B_i$ 在谱上诱导泛同胚。
\item 若 $A \to B$ 诱导剩余域同构，并在谱上诱导泛同胚，则
$B = \colim B_i$ 是有限表示 $A$-代数 $B_i$ 的滤过余极限，其中
$A \to B_i$ 诱导剩余域同构，并在谱上诱导泛同胚。
\end{enumerate}
\end{lemma}

\begin{proof}
(1) 的证明。我们将使用代数，引理
\ref{algebra-lemma-when-colimit} 的判据。设 $A \to C$ 为有限表示，
并设 $\varphi : C \to B$ 是 $A$-代数同态。令
$B' = \varphi(C) \subset B$ 为其像。由引理
\ref{morphisms-lemma-subalgebra-inherits}，$A \to B'$ 在谱上诱导泛同胚。
由代数，引理 \ref{algebra-lemma-ring-colimit-fp}，可写
$B' = \colim_{i \in I} B_i$，其中 $A \to B_i$ 为有限表示，且转移同态满射。
由代数，引理 \ref{algebra-lemma-characterize-finite-presentation}，
可选择指标 $0 \in I$，以及分解 $C \to B_0 \to B'$，它分解同态
$C \to B'$。我们断言，$\Spec(B_i) \to \Spec(A)$ 对所有充分大的
$i$ 都是泛同胚。该断言完成 (1) 的证明。

\medskip\noindent
断言的证明。由引理 \ref{morphisms-lemma-reduction-universal-homeomorphism}，环同态
$A_{red} \to B'_{red}$ 在谱上诱导泛同胚。因此，由代数，引理
\ref{algebra-lemma-image-dense-generic-points}，
$A_{red} \subset B'_{red}$。置 $A' = \Im(A \to B')$，则有诱导双射的满射
$A \to A' \to A_{red}$，且
$\Spec(A_{red}) = \Spec(A') = \Spec(A)$。因此
$A' \subset B'$ 在谱上诱导泛同胚。由命题
\ref{morphisms-proposition-universal-homeomorphism} 以及 $B'$ 在 $A'$ 上有限型这一事实，
可找到 $n$ 和 $b'_1, \ldots, b'_n \in B'$，使得
$B' = A'[b'_1, \ldots, b'_n]$，并且对 $j = 1, \ldots, n$，有
\begin{enumerate}
\item $(b'_j)^2, (b'_j)^3 \in A'[b'_1, \ldots, b'_{j - 1}]$；或
\item 存在素数 $p$，使得
$pb'_j, (b'_j)^p \in A'[b'_1, \ldots, b'_{j - 1}]$。
\end{enumerate}
选择 $b_1, \ldots, b_n \in B_0$，它们提升
$b'_1, \ldots, b'_n$。对 $i \geq 0$，以 $b_{j, i}$ 表示 $b_j$ 在
$B_i$ 中的像。对充分大的 $i$，我们将对 $j = 1, \ldots, n$ 有
\begin{enumerate}
\item $b_{j, i}^2, b_{j, i}^3 \in A_i[b_{1, i}, \ldots, b_{j - 1, i}]$；或
\item 存在素数 $p$，使得
$pb_{j, i}, b_{j, i}^p \in A_i[b_{1, i}, \ldots, b_{j - 1, i}]$。
\end{enumerate}
这里 $A_i \subset B_i$ 是 $A \to B_i$ 的像。注意 $A \to A_i$ 是满射
环同态，其核为局部幂零理想。必要时再增大 $i$，可假设 $B_i$ 由
$b_1, \ldots, b_n$ 在 $A_i$ 上生成，换言之，
$B_i = A_i[b_1, \ldots, b_n]$。由代数，引理
\ref{algebra-lemma-p-ring-map} 和
\ref{algebra-lemma-2-3-ring-map}，可知
$A \to A_i \to A_i[b_1] \to \ldots \to A_i[b_1, \ldots, b_n] = B_i$
在谱上诱导泛同胚。这完成断言的证明。

\medskip\noindent
(2) 的证明完全相同。
\end{proof}

\section{绝对弱正规化与半正规化}
\label{morphisms-section-seminormalization}

\noindent
受上一节所证结果的启发，我们给出如下定义。

\begin{definition}
\label{morphisms-definition-seminormal-ring}
设 $A$ 是环。
\begin{enumerate}
\item 若 $A$ 具有如下性质：对所有 $x, y \in A$，只要
$x^3 = y^2$，就存在唯一的 $a \in A$，使得 $x = a^2$ 且 $y = a^3$，
则称它{\it 半正规}。
\item 若 $A$ 满足：(a) $A$ 半正规；并且 (b) 对任意素数 $p$ 和
$x, y \in A$，只要 $p^px = y^p$，就存在唯一的 $a \in A$，使得
$x = a^p$ 且 $y = pa$，则称它{\it 绝对弱正规}。
\end{enumerate}
\end{definition}

\noindent
一个有趣的观察（见 \cite{Costa}）是：在半正规环的定义中，只须
假设\footnote{设 $A$ 是环，并假设对所有 $x, y \in A$，只要
$x^3 = y^2$，都存在 $a \in A$，使得
$x = a^2$ 且 $y = a^3$。那么 $A$ 既约：若 $x^2 = 0$，则
$x^2 = x^3$，因而存在 $a$，使得 $x = a^3$ 且 $x = a^2$。于是
$x = a^3 = ax = a^4 = x^2 = 0$。最后，若
$a_1^2 = a_2^2$ 且 $a_1^3 = a_2^3$，其中 $a_1, a_2$ 属于既约环，则
$(a_1 - a_2)^3 = a_1^3 - 3a_1^2a_2 +3a_1a_2^2 - a_2^3 =
(1 - 3 + 3 - 1)a_1^3 = 0$，故 $a_1 = a_2$。}$a$ 的存在性。
绝对弱正规概形定义于 \cite[Appendix B]{rydh_descent}。

\begin{lemma}
\label{morphisms-lemma-seminormal-local-property}
成为半正规或绝对弱正规是环的局部性质；见性质，定义
\ref{properties-definition-property-local}。
\end{lemma}

\begin{proof}
设 $A$ 半正规，且 $f \in A$。取 $x', y' \in A_f$，并满足
$(x')^3 = (y')^2$。写成 $x' = x/f^{2n}$ 和 $y' = y/f^{3n}$，
其中 $n \geq 0$ 且 $x, y \in A$。把 $x, y$ 分别替换为
$f^{2m}x, f^{3m}y$，并把 $n$ 替换为 $n + m$ 后，可见
$x^3 = y^2$ 在 $A$ 中成立。于是找到唯一的 $a \in A$，使得
$x = a^2$ 且 $y = a^3$。置 $a' = a/f^n$，便得到
$x' = (a')^2$ 且 $y' = (a')^3$，如所需。$a'$ 的唯一性由 $a$ 的
唯一性得出。完全同样地，读者可证明：若 $A$ 绝对弱正规，则 $A_f$
绝对弱正规。

\medskip\noindent
设 $A$ 是环，且 $f_1, \ldots, f_n \in A$ 生成单位理想。假设
$A_{f_i}$ 对每个 $i$ 都半正规。取 $x, y \in A$，并满足
$x^3 = y^2$。对每个 $i$，可找到唯一的 $a_i \in A_{f_i}$，使得
$x = a_i^2$ 且 $y = a_i^3$ 在 $A_{f_i}$ 中成立。由唯一性和第一段的结果（它告诉我们
$A_{f_if_j}$ 半正规），可见 $a_i$ 与 $a_j$ 映到 $A_{f_if_j}$ 的同一元素。
由代数，引理 \ref{algebra-lemma-standard-covering}，可找到唯一的
$a \in A$，它映到 $a_i$ 在 $A_{f_i}$ 中，对所有 $i$ 均如此。同理可得
$x = a^2$ 且 $y = a^3$。显然这个 $a$ 唯一。因此 $A$ 半正规。
若假设 $A_{f_i}$ 绝对弱正规，则完全相同的论证表明 $A$ 绝对弱正规。
\end{proof}

\noindent
接下来定义半正规概形和绝对弱正规概形。

\begin{definition}
\label{morphisms-definition-seminormal}
设 $X$ 是概形。
\begin{enumerate}
\item 若 $X$ 具有如下性质：每个 $x \in X$ 都有仿射开邻域
$\Spec(R) = U \subset X$，且环 $R$ 半正规，则称它{\it 半正规}。
\item 若 $X$ 具有如下性质：每个 $x \in X$ 都有仿射开邻域
$\Spec(R) = U \subset X$，且环 $R$ 绝对弱正规，则称它
{\it 绝对弱正规}。
\end{enumerate}
\end{definition}

\noindent
下面是必备引理。

\begin{lemma}
\label{morphisms-lemma-locally-seminormal}
设 $X$ 是概形。下列条件等价：
\begin{enumerate}
\item 概形 $X$ 半正规。
\item 对每个仿射开集 $U \subset X$，环 $\mathcal{O}_X(U)$ 半正规。
\item 存在仿射开覆盖 $X = \bigcup U_i$，使得每个
$\mathcal{O}_X(U_i)$ 都半正规。
\item 存在开覆盖 $X = \bigcup X_j$，使得每个开子概形 $X_j$ 都半正规。
\end{enumerate}
此外，若 $X$ 半正规，则每个开子概形都半正规。把“半正规”替换为
“绝对弱正规”，上述陈述仍成立。
\end{lemma}

\begin{proof}
结合性质，引理 \ref{properties-lemma-locally-P} 和引理
\ref{morphisms-lemma-seminormal-local-property}。
\end{proof}

\begin{lemma}
\label{morphisms-lemma-seminormal-reduced}
半正规概形或环既约。从而，绝对弱正规概形或环也同样既约。
\end{lemma}

\begin{proof}
设 $A$ 是环。若 $a \in A$ 非零但 $a^2 = 0$，则 $a^2 = 0^2$ 且
$a^3 = 0^3$，从而 $A$ 不半正规。
\end{proof}

\begin{lemma}
\label{morphisms-lemma-seminormalization-ring}
设 $A$ 是环。
\begin{enumerate}
\item 在谱上诱导泛同胚的环同态 $A \to B$ 所成的范畴有终对象
$A \to A^{awn}$。
\item 给定 (1) 的范畴中的 $A \to B$，所得同态
$B \to A^{awn}$ 是同构，当且仅当 $B$ 绝对弱正规。
\item 在剩余域上诱导同构且在谱上诱导泛同胚的环同态 $A \to B$
所成的范畴有终对象 $A \to A^{sn}$。
\item 给定 (3) 的范畴中的 $A \to B$，所得同态
$B \to A^{sn}$ 是同构，当且仅当 $B$ 半正规。
\end{enumerate}
对任意环同态 $\varphi : A \to A'$，存在唯一的同态
$\varphi^{awn} : A^{awn} \to (A')^{awn}$ 和
$\varphi^{sn} : A^{sn} \to (A')^{sn}$ 与 $\varphi$ 相容。
\end{lemma}

\begin{proof}
我们证明 (1) 和 (2)，略去 (3)、(4) 以及最后陈述的证明。考虑如下形式的
$A$-代数所成的范畴：
$$
B = A[x_1, \ldots, x_n]/J
$$
其中 $J$ 是有限生成理想，使得 $A \to B$ 在谱上定义泛同胚。我们断言该范畴
有向（范畴，定义 \ref{categories-definition-directed}）。事实上，给定
$$
B = A[x_1, \ldots, x_n]/J
\quad\text{且}\quad
B' = A[x_1, \ldots, x_{n'}]/J'
$$
则可考虑
$$
B'' = A[x_1, \ldots, x_{n + n'}]/J''
$$
其中 $J''$ 由 $J$ 的元素以及元素
$f(x_{n + 1}, \ldots, x_{n + n'})$ 生成，其中 $f \in J'$。于是有
$A$-代数同态 $B \to B''$ 和 $B' \to B''$，它们诱导同构
$B \otimes_A B' \to B''$。由引理
\ref{morphisms-lemma-base-change-universal-homeomorphism} 和
\ref{morphisms-lemma-composition-universal-homeomorphism} 可知
$\Spec(B'') \to \Spec(A)$ 是泛同胚，故 $A \to B''$ 属于我们的范畴。
最后，给定该范畴中的 $\varphi, \varphi' : B \to B'$，其中 $B$ 如上显示，
考虑商 $B''$；它是 $B'$ 除以由 $\varphi(x_i) - \varphi'(x_i)$
（$i = 1, \ldots, n$）生成的理想所得的商。由于
$\Spec(B') = \Spec(B)$，可见
$\Spec(B'') \to \Spec(B')$ 是双射闭浸入，故为泛同胚。因此 $B''$
属于我们的范畴，并且 $\varphi, \varphi'$ 被 $B' \to B''$ 等化。
这完成断言的证明。置
$$
A^{awn} = \colim B
$$
其中余极限取遍刚描述的范畴。注意，由引理
\ref{morphisms-lemma-colimit-inherits}，$A \to A^{awn}$ 在谱上诱导泛同胚
（此处用到该范畴有向）。

\medskip\noindent
给定在谱上诱导泛同胚的有限表示环同态 $A \to B$，得到典范同态
$B \to A^{awn}$；这是由 $A^{awn}$ 的构造本身给出的。由于 (1) 中每个 $A \to B$ 都是 (1) 中
有限表示的 $A \to B$ 的滤过余极限（引理
\ref{morphisms-lemma-universal-homeo-limit}），可见 $A \to A^{awn}$ 是 (1) 的范畴中的
终对象。

\medskip\noindent
设 $x, y \in A^{awn}$ 是满足 $x^3 = y^2$ 的元素。由引理
\ref{morphisms-lemma-make-universal-homeo}，
$A^{awn} \to A^{awn}[t]/(t^2 - x, t^3 - y)$ 在谱上诱导泛同胚。因此
$A \to A^{awn}[t]/(t^2 - x, t^3 - y)$ 属于 (1) 的范畴，于是得到唯一的
$A$-代数同态
$A^{awn}[t]/(t^2 - x, t^3 - y) \to A^{awn}$。$a \in A^{awn}$ 为
$t$ 的像，因而是满足 $a^2 = x$ 且 $a^3 = y$ 的唯一元素，
这些等式在 $A^{awn}$ 中成立。完全同样地，给定素数 $p$ 和
$x, y \in A^{awn}$，若 $p^px = y^p$，则找到唯一的 $a \in A^{awn}$，
使得 $a^p = x$ 且 $pq = y$。因此按定义，$A^{awn}$ 绝对弱正规。

\medskip\noindent
最后，设 $A \to B$ 属于 (1) 的范畴，且 $B$ 绝对弱正规。由于
$A^{awn} \to B^{awn}$ 在谱上诱导泛同胚，并且 $A^{awn}$ 既约（引理
\ref{morphisms-lemma-seminormal-reduced}），可得 $A^{awn} \subset B^{awn}$
（见代数，引理 \ref{algebra-lemma-image-dense-generic-points}）。若此包含
不是等式，则引理 \ref{morphisms-lemma-pth-power-and-multiple} 蕴含存在元素
$b \in B^{awn}$、$b \not \in A^{awn}$，使得或者
$b^2, b^3 \in A^{awn}$，或者 $pb, b^p \in A^{awn}$ 对某个素数 $p$ 成立。
然而，由定义 \ref{morphisms-definition-seminormal-ring} 中的存在唯一性，这迫使
$b \in A^{awn}$，从而得到矛盾，完成证明。
\end{proof}

\begin{lemma}
\label{morphisms-lemma-seminormalization}
设 $X$ 是概形。
\begin{enumerate}
\item 泛同胚 $Y \to X$ 所成的范畴有始对象 $X^{awn} \to X$。
\item 给定 (1) 的范畴中的 $Y \to X$，所得态射
$X^{awn} \to Y$ 是同构，当且仅当 $Y$ 绝对弱正规。
\item 诱导剩余域同构的泛同胚 $Y \to X$ 所成的范畴有始对象
$X^{sn} \to X$。
\item 给定 (3) 的范畴中的 $Y \to X$，所得态射
$X^{sn} \to Y$ 是同构，当且仅当 $Y$ 半正规。
\end{enumerate}
对概形的任意态射 $h : X' \to X$，存在唯一的态射
$h^{awn} : (X')^{awn} \to X^{awn}$ 和
$h^{sn} : (X')^{sn} \to X^{sn}$ 与 $h$ 相容。
\end{lemma}

\begin{proof}
我们证明 (1) 和 (2)，略去 (3) 和 (4) 的证明。设 $h : X' \to X$
是概形的态射。若 (1) 对 $X$ 和 $X'$ 成立，则
$X' \times_X X^{awn} \to X'$ 是泛同胚，故存在态射
$(X')^{awn} \to X' \times_X X^{awn}$，它在 $X'$ 上唯一；这是由
$(X')^{awn} \to X'$ 的泛性质得到的。与投影
$X' \times_X X^{awn} \to X^{awn}$ 复合，得到 $h^{awn}$。若进一步 (2)
对 $X$ 和 $X'$ 成立，且 $h$ 是开浸入，则
$X' \times_X X^{awn}$ 绝对弱正规（引理
\ref{morphisms-lemma-locally-seminormal}），于是可得
$(X')^{awn} \to X' \times_X X^{awn}$ 是同构。

\medskip\noindent
回顾：任意泛同胚仿射；见引理 \ref{morphisms-lemma-homeomorphism-affine}。因此若
$X$ 仿射，则 (1) 和 (2) 立即由引理
\ref{morphisms-lemma-seminormalization-ring} 得出。设 $X$ 是概形，并令
$\mathcal{B}$ 为 $X$ 的仿射开集所成的集合。对每个
$U \in \mathcal{B}$，得到 $U^{awn} \to U$；对
$V \subset U$，其中 $V, U \in \mathcal{B}$，由上一段的论述得到典范同构
$\rho_{V, U} : V^{awn} \to V \times_U U^{awn}$。因此，由相对粘合
（构造，引理 \ref{constructions-lemma-relative-glueing}），得到态射
$X^{awn} \to X$；它限制为 $U^{awn}$，位于 $U$ 上，并与
$\rho_{V, U}$ 相容。接着，设 $Y \to X$ 是泛同胚。则
$U \times_X Y \to U$ 对 $U \in \mathcal{B}$ 是泛同胚，并得到唯一态射
$g_U : U^{awn} \to U \times_X Y$，它在 $U$ 上。这些 $g_U$ 与态射
$\rho_{V, U}$ 相容；细节略。因此存在唯一态射
$g : X^{awn} \to Y$，它在 $X$ 上，并与 $g_U$ 一致，在 $U$ 上；见构造，评注
\ref{constructions-remark-relative-glueing-functorial}。这对 $X$ 证明了 (1)。
(2) 因其在仿射局部成立而得出。
\end{proof}

\begin{definition}
\label{morphisms-definition-seminormalization}
设 $X$ 是概形。
\begin{enumerate}
\item 引理 \ref{morphisms-lemma-seminormalization} 中构造的态射
$X^{sn} \to X$ 称为 $X$ 的{\it 半正规化}。
\item 引理 \ref{morphisms-lemma-seminormalization} 中构造的态射
$X^{awn} \to X$ 称为 $X$ 的{\it 绝对弱正规化}。
\end{enumerate}
\end{definition}

\noindent
确实，$X^{sn}$ 是 $X$ 的半正规化，也是半正规概形，而绝对弱正规化
$X^{awn}$ 是绝对弱正规概形。此外，对概形的任意态射
$h : Y \to X$，得到典范交换图
$$
\xymatrix{
Y^{awn} \ar[d]^{h^{awn}} \ar[r] &
Y^{sn} \ar[d]^{h^{sn}} \ar[r] &
Y \ar[d]^h \\
X^{awn} \ar[r] &
X^{sn} \ar[r] &
X
}
$$
；其中箭头 $h^{sn}$ 和 $h^{awn}$ 是与 $h$ 相容的唯一箭头。

\begin{lemma}
\label{morphisms-lemma-characterize-seminormal}
设 $X$ 是概形。下列条件等价：
\begin{enumerate}
\item $X$ 半正规；
\item $X$ 等于其自身的半正规化，即态射 $X^{sn} \to X$ 是同构；
\item 若 $\pi : Y \to X$ 是诱导剩余域同构的泛同胚，且 $Y$ 既约，
则 $\pi$ 是同构。
\end{enumerate}
\end{lemma}

\begin{proof}
(1) 与 (2) 的等价性由引理 \ref{morphisms-lemma-seminormalization} 显然。若 (3) 成立，
则 $X^{sn} \to X$ 是同构，从而 (2) 成立。

\medskip\noindent
设 (2) 成立，并令 $\pi : Y \to X$ 是诱导剩余域同构的泛同胚，且
$Y$ 既约。由引理 \ref{morphisms-lemma-seminormalization}，存在分解
$X \to Y \to X$，它分解 $\text{id}_X$。于是 $X \to Y$ 是闭浸入
（由概形，引理 \ref{schemes-lemma-section-immersion}，以及 $\pi$ 分离这一
事实，例如由引理 \ref{morphisms-lemma-universally-injective-separated}）。由于
$X \to Y$ 还在点上双射，而 $Y$ 既约，它必为同构。这完成证明。
\end{proof}

\begin{lemma}
\label{morphisms-lemma-characterize-absolutely-weakly-normal}
设 $X$ 是概形。下列条件等价：
\begin{enumerate}
\item $X$ 绝对弱正规；
\item $X$ 等于其自身的绝对弱正规化，即态射 $X^{awn} \to X$ 是同构；
\item 若 $\pi : Y \to X$ 是泛同胚，且 $Y$ 既约，则 $\pi$ 是同构。
\end{enumerate}
\end{lemma}

\begin{proof}
证明与引理 \ref{morphisms-lemma-characterize-seminormal} 完全相同。
\end{proof}





\section{有限局部自由态射}
\label{morphisms-section-finite-locally-free}

\noindent
许多论文中的作者使用有限平坦态射，而他们实际所指的是有限局部自由态射。
原因是，当基底局部诺特时，这两个概念相同。但一般而言并非如此，见
《习题》练习 \ref{exercises-exercise-flat-not-projective}。

\begin{definition}
\label{morphisms-definition-finite-locally-free}
设 $f : X \to S$ 是概形的态射。称 $f$ {\it 有限局部自由}，
如果 $f$ 是仿射态射，并且 $f_*\mathcal{O}_X$ 是有限局部自由
$\mathcal{O}_S$-模。在这种情况下，称 $f$ 具有{\it 秩}或
{\it 次数} $d$，如果层 $f_*\mathcal{O}_X$ 是次数为 $d$
的有限局部自由层。
\end{definition}

\noindent
注意，若 $f : X \to S$ 有限局部自由，则 $S$ 是开闭子概形
$S_d$ 的不交并，使得 $f^{-1}(S_d) \to S_d$ 是次数为 $d$
的有限局部自由态射。

\begin{lemma}
\label{morphisms-lemma-finite-flat}
设 $f : X \to S$ 是概形的态射。下列条件等价：
\begin{enumerate}
\item $f$ 有限局部自由；
\item $f$ 有限、平坦且局部有限表示。
\end{enumerate}
若 $S$ 局部诺特，则它们还等价于
\begin{enumerate}
\item[(3)] $f$ 有限且平坦。
\end{enumerate}
\end{lemma}

\begin{proof}
设 $V \subset S$ 是仿射开集。在上述三种情形中，态射都是仿射的，
因而 $f^{-1}(V)$ 是仿射的。因此可写
$V = \Spec(R)$ 和 $f^{-1}(V) = \Spec(A)$，其中取某个
$R$-代数 $A$。假设 (1) 成立。这意味着 $S$ 可由仿射开集
$V = \Spec(R)$ 覆盖，使得 $A$ 作为 $R$-模是有限自由的。
由《代数》引理 \ref{algebra-lemma-finite-finite-type}，
$R \to A$ 是有限表示的。因此 (2) 成立。反之，假设 (2) 成立。
对每个仿射开集 $V = \Spec(R)$（它是 $S$ 的开集），环同态
$R \to A$ 是有限且有限表示的，并且 $A$ 作为 $R$-模是平坦的。
由《代数》引理
\ref{algebra-lemma-finite-finitely-presented-extension}，
$A$ 作为 $R$-模是有限表示的。因此《代数》引理
\ref{algebra-lemma-finite-projective} 推出 $A$ 有限局部自由，
故 (1) 成立。诺特情形来自以下事实：诺特环上的有限模是有限表示模；
见《代数》引理
\ref{algebra-lemma-Noetherian-finite-type-is-finite-presentation}。
\end{proof}

\begin{lemma}
\label{morphisms-lemma-composition-finite-locally-free}
有限局部自由态射的复合仍是有限局部自由态射。
\end{lemma}

\begin{proof}
略。
\end{proof}

\begin{lemma}
\label{morphisms-lemma-base-change-finite-locally-free}
有限局部自由态射的基变换仍是有限局部自由态射。
\end{lemma}

\begin{proof}
略。
\end{proof}

\begin{lemma}
\label{morphisms-lemma-finite-locally-free}
设 $f : X \to S$ 是概形的有限局部自由态射。存在由开闭子概形给出的
不交并分解 $S = \coprod_{d \geq 0} S_d$，使得令
$X_d = f^{-1}(S_d)$ 后，各限制 $f|_{X_d}$ 都是有限局部自由态射
$X_d \to S_d$，其次数为 $d$。
\end{lemma}

\begin{proof}
这是因为有限局部自由层在局部具有良定义的秩。细节略。
\end{proof}

\begin{lemma}
\label{morphisms-lemma-massage-finite}
设 $f : Y \to X$ 是有限态射，且 $X$ 仿射。存在图
$$
\xymatrix{
Z' \ar[rd] &
Y' \ar[l]^i \ar[d] \ar[r] &
Y \ar[d] \\
 & X' \ar[r] & X
}
$$
其中
\begin{enumerate}
\item $Y' \to Y$ 和 $X' \to X$ 都是满的有限局部自由态射；
\item $Y' = X' \times_X Y$；
\item $i : Y' \to Z'$ 是闭浸入；
\item $Z' \to X'$ 有限局部自由；并且
\item $Z' = \bigcup_{j = 1, \ldots, m} Z'_j$ 是闭子概形的
（集合论意义下的）有限并，其中每一个都同构地映到 $X'$。
\end{enumerate}
\end{lemma}

\begin{proof}
写成 $X = \Spec(A)$ 和 $Y = \Spec(B)$。另见《代数续篇》章节
\ref{more-algebra-section-descent-flatness-integral}。
设 $x_1, \ldots, x_n \in B$ 是 $B$ 在 $A$ 上的一组生成元。
对每个 $i$，可以选取首一多项式 $P_i(T) \in A[T]$，
使得 $P(x_i) = 0$ 在 $B$ 中成立。由《代数》引理
\ref{algebra-lemma-adjoin-roots}（应用 $n$ 次），存在有限局部自由环扩张
$A \subset A'$，使得每个 $P_i$ 完全分裂：
$$
P_i(T) = \prod\nolimits_{k = 1, \ldots, d_i} (T - \alpha_{ik})
$$
，其中 $\alpha_{ik} \in A'$。令
$$
C = A'[T_1, \ldots, T_n]/(P_1(T_1), \ldots, P_n(T_n))
$$
以及 $B' = A' \otimes_A B$。同态
$C \to B'$，$T_i \mapsto 1 \otimes x_i$ 是 $A'$-代数满同态。
令 $X' = \Spec(A')$、$Y' = \Spec(B')$ 和 $Z' = \Spec(C)$，
便知 (1)--(4) 成立。(5) 成立是因为，在集合论意义下，
$\Spec(C)$ 是由下列理想切出的闭子概形的并：
$$
(T_1 - \alpha_{1k_1}, T_2 - \alpha_{2k_2}, \ldots, T_n - \alpha_{nk_n})
$$
，其中任意 $1 \leq k_i \leq d_i$。
\end{proof}

\noindent
下一个引理以正确的一般性表述在下文引理
\ref{morphisms-lemma-image-nowhere-dense-quasi-finite} 中。

\begin{lemma}
\label{morphisms-lemma-image-nowhere-dense-finite}
设 $f : Y \to X$ 是概形的有限态射。设 $T \subset Y$ 是 $Y$ 的
闭无处稠密子集。则 $f(T) \subset X$ 是 $X$ 的闭无处稠密子集。
\end{lemma}

\begin{proof}
由引理 \ref{morphisms-lemma-finite-proper}，已知 $f(T) \subset X$ 是闭的。
设 $X = \bigcup X_i$ 是仿射开覆盖。由于 $T$ 在 $Y$ 中无处稠密，
$T \cap f^{-1}(X_i)$ 在 $f^{-1}(X_i)$ 中也无处稠密。因此，若能在仿射
情形证明该定理，便知 $f(T) \cap X_i$ 无处稠密。于是由《拓扑》引理
\ref{topology-lemma-nowhere-dense-local} 可知 $T$ 在 $X$ 中无处稠密。

\medskip\noindent
假设 $X$ 仿射。选取引理 \ref{morphisms-lemma-massage-finite} 中的图
$$
\xymatrix{
Z' \ar[rd] &
Y' \ar[l]^i \ar[d]^{f'} \ar[r]_a &
Y \ar[d]^f \\
 & X' \ar[r]^b & X
}
$$
。态射 $a$、$b$ 都是开的，因为它们有限局部自由
（引理 \ref{morphisms-lemma-finite-flat} 和 \ref{morphisms-lemma-fppf-open}）。
因此 $T' = a^{-1}(T)$ 无处稠密，见《拓扑》引理
\ref{topology-lemma-open-inverse-image-closed-nowhere-dense}。
态射 $b$ 是满且开的。因此，若能证明
$f'(T') = b^{-1}(f(T))$ 无处稠密，则由《拓扑》引理
\ref{topology-lemma-open-inverse-image-closed-nowhere-dense}，
$f(T)$ 无处稠密。由于 $i$ 是闭浸入，由《拓扑》引理
\ref{topology-lemma-closed-image-nowhere-dense}，
$i(T') \subset Z'$ 是闭且无处稠密的。这样便把问题归结为下一段
所讨论的情形。

\medskip\noindent
假设 $Y = \bigcup_{i = 1, \ldots, n} Y_i$ 是闭子集的有限并，
每个子集都同构地映到 $X$。考虑 $T_i = Y_i \cap T$。
若每个 $T_i$ 在相应的 $Y_i$ 中都无处稠密，则每个 $f(T_i)$ 在 $X$
中都无处稠密，因为 $Y_i \to X$ 是同构。因此
$f(T) = f(T_i)$ 是 $X$ 中无处稠密闭子集的有限并，于是结论成立；
见《拓扑》引理 \ref{topology-lemma-nowhere-dense}。
若不然，不妨说 $T_1$ 包含一个非空开集
$V \subset Y_1$。我们将说明这会导致矛盾。
考虑 $Y_2 \cap V \subset V$。它或者是一个真闭子集，或者等于 $V$。
在第一种情形，将 $V$ 替换为 $V \setminus V \cap Y_2$，于是
$V \subset T_1$ 在 $Y_1$ 中是开集且不与 $Y_2$ 相交。
在第二种情形，$V \subset Y_1 \cap Y_2$ 在 $Y_1$ 和 $Y_2$
中都是开集。依次对 $i = 3, \ldots, n$ 重复此过程。结果得到不交并分解
$$
\{1, \ldots, n\} = I_1 \amalg I_2, \quad 1 \in I_1
$$
，以及一个开集 $V$（它是 $Y_1$ 的开集且包含于 $T_1$ 中），使得
$V \subset Y_i$ 对 $i \in I_1$ 成立，而
$V \cap Y_i = \emptyset$ 对 $i \in I_2$ 成立。令 $U = f(V)$。
这是 $X$ 的开集，因为 $f|_{Y_1} : Y_1 \to X$ 是同构。于是
$$
f^{-1}(U) = V\ \amalg\ \bigcup\nolimits_{i \in I_2} (Y_i \cap f^{-1}(U))
$$
。由于 $\bigcup_{i \in I_2} Y_i$ 是闭的，这说明
$V \subset f^{-1}(U)$ 是开的，从而 $V \subset Y$ 是开集。
这与 $T$ 在 $Y$ 中无处稠密的假设矛盾，证毕。
\end{proof}

\section{有理映射}
\label{morphisms-section-rational-maps}

\noindent
设 $X$ 是概形。注意，若 $U$、$V$ 是 $X$ 中的稠密开集，
则 $U \cap V$ 也是稠密开集。

\begin{definition}
\label{morphisms-definition-rational-map}
设 $X$、$Y$ 是概形。
\begin{enumerate}
\item 设 $f : U \to Y$、$g : V \to Y$ 是概形态射，定义在
$U$、$V$ 上；这些子集属于 $X$。称 $f$ {\it 等价}于 $g$，如果
$f|_W = g|_W$，其中 $W \subset U \cap V$ 是 $X$ 中的稠密开集。
\item 从 $X$ 到 $Y$ 的{\it 有理映射}是由 (1) 中等价关系定义的等价类。
\item 若 $X$、$Y$ 是在基概形 $S$ 上的概形，则称从 $X$ 到 $Y$ 的有理映射是
一个{\it $S$-有理映射}，即从 $X$ 到 $Y$ 的有理映射，如果该等价类存在一个代表
$f : U \to Y$，且它是一个 $S$-态射。
\end{enumerate}
\end{definition}

\noindent
对于定义 (1) 中的两个态射 $f$、$g$，我们称它们定义了同一个有理映射，
而不说它们是等价的。在某些情形中，有理映射由泛点处局部环之间的映射决定。

\begin{lemma}
\label{morphisms-lemma-rational-map-finite-presentation}
设 $S$ 是概形。设 $X$ 和 $Y$ 是在 $S$ 上的概形。假设 $X$ 有限个
不可约分支，其泛点为 $x_1, \ldots, x_n$。设 $s_i \in S$ 是
$x_i$ 的像。考虑映射
$$
\left\{
\begin{matrix}
S\text{-rational maps} \\
\text{from }X\text{ to }Y
\end{matrix}
\right\}
\longrightarrow
\left\{
\begin{matrix}
(y_1, \varphi_1, \ldots, y_n, \varphi_n)\text{ 其中 }
y_i \in Y\text{ lies over }s_i\text{ 且}\\
\varphi_i : \mathcal{O}_{Y, y_i} \to \mathcal{O}_{X, x_i}
\text{ is a local }\mathcal{O}_{S, s_i}\text{-algebra map}
\end{matrix}
\right\}
$$
，它将 $f : U \to Y$ 映为 $2n$-元组，其中
$y_i = f(x_i)$ 且 $\varphi_i = f^\sharp_{x_i}$。则有
\begin{enumerate}
\item 若 $Y \to S$ 局部有限型，则该映射是单射；
\item 若 $Y \to S$ 局部有限表示，则该映射是双射；
\item 若 $Y \to S$ 局部有限型且 $X$ 既约，则该映射是双射。
\end{enumerate}
\end{lemma}

\begin{proof}
注意，$X$ 的任意稠密开集都包含点 $x_i$，所以该构造有意义。
为证明 (1) 或 (2)，可以将 $X$ 替换为任意稠密开集。因此，若
$Z_1, \ldots, Z_n$ 是 $X$ 的不可约分支，则可以将 $X$ 替换为
$X \setminus \bigcup_{i \not = j} Z_i \cap Z_j$。
完成这一步后，$X$ 是其不可约分支（视为开闭子概形）的不交并。
于是箭头的右端和左端都是关于不可约分支的乘积，问题归结为
$X$ 不可约的情形。

\medskip\noindent
假设 $X$ 不可约，其泛点 $x$ 位于 $s \in S$ 之上。
(1) 来自引理 \ref{morphisms-lemma-morphism-defined-local-ring} 的 (1)。
(2) 和 (3) 来自同一引理的 (2)。
\end{proof}

\begin{definition}
\label{morphisms-definition-rational-function}
设 $X$ 是概形。$X$ 上的{\it 有理函数}是从 $X$ 到
$\mathbf{A}^1_{\mathbf{Z}}$ 的有理映射。
\end{definition}

\noindent
关于仿射直线 $\mathbf{A}^1$ 的定义，见《构造》定义
\ref{constructions-definition-affine-n-space}。设 $X$ 是在 $S$ 上的概形。
对任意开集 $U \subset X$，$U \to \mathbf{A}^1_{\mathbf{Z}}$ 的态射
等同于 $U \to \mathbf{A}^1_S$ 的 $S$-态射。因此，有理函数也等同于
一个 $S$-有理映射，即从 $X$ 到 $\mathbf{A}^1_S$ 的映射。

\medskip\noindent
回忆，有典范同一化
$\Mor(T, \mathbf{A}^1_{\mathbf{Z}}) = \Gamma(T, \mathcal{O}_T)$，其中 $T$ 是任意概形；
见《概形》例 \ref{schemes-example-global-sections}。因此
$\mathbf{A}^1_{\mathbf{Z}}$ 是概形范畴中的环对象。更具体地说，态射
\begin{eqnarray*}
+ : \mathbf{A}^1_{\mathbf{Z}} \times \mathbf{A}^1_{\mathbf{Z}}
& \longrightarrow &
\mathbf{A}^1_{\mathbf{Z}} \\
(f, g) & \longmapsto & f + g \\
* : \mathbf{A}^1_{\mathbf{Z}} \times \mathbf{A}^1_{\mathbf{Z}}
& \longrightarrow &
\mathbf{A}^1_{\mathbf{Z}} \\
(f, g) & \longmapsto & fg
\end{eqnarray*}
满足环中加法和乘法的全部公理（总是取含幺 $1$ 的交换环）。
因此，取值于 $\mathbf{A}^1_{\mathbf{Z}}$ 的有理映射集合也自然成为一个环。

\begin{definition}
\label{morphisms-definition-ring-of-rational-functions}
设 $X$ 是概形。$X$ 上的{\it 有理函数环}是环 $R(X)$，
其元素是有理函数，并配备上述加法和乘法。
\end{definition}

\noindent
对于只有有限个不可约分支的概形，我们可以计算这个环。

\begin{lemma}
\label{morphisms-lemma-integral-scheme-rational-functions}
设 $X$ 是只有有限个不可约分支 $X_1, \ldots, X_n$ 的概形。
若 $\eta_i \in X_i$ 是泛点，则
$$
R(X) = \mathcal{O}_{X, \eta_1} \times \ldots \times \mathcal{O}_{X, \eta_n}
$$
若 $X$ 既约，则这等于 $\prod \kappa(\eta_i)$。
若 $X$ 整，则 $R(X) = \mathcal{O}_{X, \eta} = \kappa(\eta)$
是一个域。
\end{lemma}

\begin{proof}
设 $U \subset X$ 是稠密开子集。则
$U_i = (U \cap X_i) \setminus (\bigcup_{j \not = i} X_j)$
是非空开集，因为它包含 $\eta_i$，包含于 $X_i$，并且
$\bigcup U_i \subset U \subset X$ 是稠密的。
因此，引理中的同一化来自下面这一串等式：
\begin{align*}
R(X)
& =
\colim_{U \subset X\text{ open dense}} \Mor(U, \mathbf{A}^1_\mathbf{Z}) \\
& =
\colim_{U \subset X\text{ open dense}} \mathcal{O}_X(U) \\
& =
\colim_{\eta_i \in U_i \subset X\text{ 开}} \prod \mathcal{O}_X(U_i) \\
& =
\prod \colim_{\eta_i \in U_i \subset X\text{ 开}} \mathcal{O}_X(U_i) \\
& =
\prod \mathcal{O}_{X, \eta_i}
\end{align*}
其中第二个等式来自《概形》例
\ref{schemes-example-global-sections}。最后一个断言来自
《代数》引理 \ref{algebra-lemma-minimal-prime-reduced-ring}。
\end{proof}

\begin{definition}
\label{morphisms-definition-function-field}
设 $X$ 是整概形。$X$ 的{\it 函数域}，或{\it 有理函数域}，
是域 $R(X)$。
\end{definition}

\noindent
我们有时用 $k(X)$ 而不是 $R(X)$ 表示这个域。
利用函数域可以阐明整概形上的分离条件。注意，根据引理
\ref{morphisms-lemma-integral-scheme-rational-functions}，在整概形上，每个局部环
$\mathcal{O}_{X, x}$ 都可视为 $R(X)$ 的局部子环。

\begin{lemma}
\label{morphisms-lemma-distinct-local-rings}
设 $X$ 是整分离概形。设 $Z_1$、$Z_2$ 是 $X$ 中不同的不可约闭子集。
设 $\eta_i$ 是 $Z_i$ 的泛点。若 $Z_1 \not\subset Z_2$，则有
$\mathcal{O}_{X, \eta_1} \not \subset \mathcal{O}_{X, \eta_2}$
作为 $R(X)$ 的子环。特别地，若 $Z_1 = \{x\}$ 由一个闭点 $x$ 组成，则存在一个在 $x$
的某邻域上正则的函数，但它不属于 $\mathcal{O}_{X, \eta_{2}}$。
\end{lemma}

\begin{proof}
首先注意，在假设 $X$ 分离时，存在唯一的概形态射
$\Spec(\mathcal{O}_{X, \eta_2}) \to X$（在 $X$ 上），使得复合
$$
\Spec(R(X)) \longrightarrow
\Spec(\mathcal{O}_{X, \eta_2}) \longrightarrow X
$$
是典范映射 $\Spec(R(X)) \to X$。
确实，有典范映射
$can : \Spec(\mathcal{O}_{X, \eta_2}) \to X$，见《概形》方程
\ref{schemes-equation-canonical-morphism}。给定第二个态射 $a$ 到 $X$，
由假设，$a$ 与 $can$ 在
$\Spec(\mathcal{O}_{X, \eta_2})$ 的泛点处相同。
现在，$X$ 分离保证这两个映射相同的
$\Spec(\mathcal{O}_{X, \eta_2})$ 子集是闭的；见《概形》引理
\ref{schemes-lemma-where-are-they-equal}。因此 $a = can$ 在整个
$\Spec(\mathcal{O}_{X, \eta_2})$ 上成立。

\medskip\noindent
假设 $Z_1 \not \subset Z_2$，并反设
$\mathcal{O}_{X, \eta_{1}} \subset \mathcal{O}_{X, \eta_{2}}$
作为子环包含于 $R(X)$。
于是得到第二个态射
$$
\Spec(\mathcal{O}_{X, \eta_{2}}) \longrightarrow
\Spec(\mathcal{O}_{X, \eta_{1}}) \longrightarrow
X.
$$
根据上面的结论，这个复合必须等于 $can$。
这意味着 $\eta_2$ 特化到 $\eta_1$（见《概形》引理
\ref{schemes-lemma-specialize-points}）。但这与假设
$Z_1 \not \subset Z_2$ 矛盾。
\end{proof}

\begin{definition}
\label{morphisms-definition-domain-of-definition}
设 $\varphi$ 是概形 $X$ 与 $Y$ 之间的有理映射。若 $\varphi$ 在点
$x \in X$ 处有定义，即存在代表 $(U, f)$ 的 $\varphi$ 使得 $x \in U$，
则称其{\it 在该点处有定义}。
$\varphi$ 的{\it 定义域}是所有 $\varphi$ 有定义的点组成的集合。
\end{definition}

\noindent
由此定义，一般而言并不成立：$\varphi$ 存在一个代表，其定义域包含
整个定义域。

\begin{lemma}
\label{morphisms-lemma-rational-map-from-reduced-to-separated}
设 $X$、$Y$ 是概形。假设 $X$ 既约且 $Y$ 分离。
设 $\varphi$ 是从 $X$ 到 $Y$ 的有理映射，其定义域为
$U \subset X$。则存在唯一态射 $f : U \to Y$ 表示 $\varphi$。
若 $X$ 和 $Y$ 是在分离概形 $S$ 上的概形，且 $\varphi$ 是
$S$-有理映射，则 $f$ 是在 $S$ 上的态射。
\end{lemma}

\begin{proof}
设 $(V, g)$ 和 $(V', g')$ 是 $\varphi$ 的代表。则 $g, g'$ 在
稠密开子概形 $W \subset V \cap V'$ 上相同。另一方面，
$E$ 是 $g|_{V \cap V'}$ 与 $g'|_{V \cap V'}$ 的等化子，且是
$V \cap V'$ 的闭子概形（《概形》引理
\ref{schemes-lemma-where-are-they-equal}）。现在
$W \subset E$ 蕴含集合论意义下 $E = V \cap V'$。
由于 $V \cap V'$ 既约，得到概形论意义下
$E = V \cap V'$，即
$g|_{V \cap V'} = g'|_{V \cap V'}$。因此可以将代表
$g : V \to Y$ 的 $\varphi$ 粘合为态射
$f : U \to Y$，见《概形》引理 \ref{schemes-lemma-glue}。
最后一个断言的证明略去。
\end{proof}

\noindent
一般而言，复合有理映射没有意义。原因是第一个有理映射的某个代表的像
可能与第二个有理映射的定义域没有交集。然而，若假设概形不可约并考察
支配有理映射，则可以复合有理映射。

\begin{definition}
\label{morphisms-definition-dominant-rational}
设 $X$ 与 $Y$ 是不可约概形。从 $X$ 到 $Y$ 的有理映射称为
{\it 支配的}，如果任一代表 $f : U \to Y$ 都是概形的支配态射。
\end{definition}

\noindent
由引理 \ref{morphisms-lemma-dominant-finite-number-irreducible-components}，以下要求等价：
$\eta \in X$ 是泛点并映到泛点 $\xi$，它属于 $Y$。对任一代表，都有
$f(\eta) = \xi$，其中代表写为 $f : U \to Y$。我们可以复合支配有理映射
$\varphi$；它介于不可约概形 $X$ 与 $Y$ 之间；以及任意有理映射 $\psi$，它从 $Y$ 到 $Z$。
具体地，取代表
$f : U \to Y$，其中 $U \subset X$ 是稠密开集，并取代表
$g : V \to Z$，其中 $V \subset Y$ 是稠密开集。于是
$W = f^{-1}(V) \subset X$ 是非空开集（因为它包含 $X$ 的泛点），并令
$\psi \circ \varphi$ 为 $g \circ f|_W : W \to Z$ 的等价类。验证良定义性留给读者。

\medskip\noindent
这样得到一个范畴，其对象是不可约概形，态射是支配有理映射。给定基概形 $S$，
同样可以定义一个范畴，其对象是 $S$ 上的不可约概形，态射是支配的 $S$-有理映射。

\begin{definition}
\label{morphisms-definition-birational-schemes}
设 $X$ 与 $Y$ 是不可约概形。
\begin{enumerate}
\item 如果 $X$ 与 $Y$ 在不可约概形及支配有理映射所成的范畴中同构，
  则称 $X$ 与 $Y$ 是{\it 双有理的}。
\item 假设 $X$ 与 $Y$ 是概形，基概形为 $S$。称 $X$ 与 $Y$ 是
  {\it $S$-双有理的}，如果 $X$ 与 $Y$ 在 $S$ 上的不可约概形及支配的
  $S$-有理映射所成的范畴中同构。
\end{enumerate}
\end{definition}

\noindent
若 $X$ 与 $Y$ 是双有理的不可约概形，则从 $X$ 到 $Z$ 的有理映射集合
与从 $Y$ 到 $Z$ 的有理映射集合对所有概形 $Z$ 都存在双射（关于 $Z$ 具有函子性）。对“通常的”
不可约概形而言，这只是可能的一个定义。另一种定义是要求 $X$ 与 $Y$ 具有同构的有理函数环。
对于簇，这些条件是等价的，见引理 \ref{morphisms-lemma-criterion-birational-finite-presentation}。

\begin{lemma}
\label{morphisms-lemma-birational-integral}
设 $X$ 与 $Y$ 是不可约概形。
\begin{enumerate}
\item 概形 $X$ 与 $Y$ 双有理，当且仅当它们具有同构的非空开子概形。
\item 假设 $X$ 与 $Y$ 是基概形 $S$ 上的概形。则 $X$ 与 $Y$ 是
  $S$-双有理的，当且仅当存在非空开集 $U \subset X$ 与 $V \subset Y$，
  使它们在 $S$ 上同构。
\end{enumerate}
\end{lemma}

\begin{proof}
假设 $X$ 与 $Y$ 双有理。设 $f : U \to Y$ 与 $g : V \to X$ 分别定义从
$X$ 到 $Y$ 以及从 $Y$ 到 $X$ 的互逆支配有理映射。我们可以假设 $V$ 是仿射的，
并可将 $U$ 替换为 $f^{-1}(V)$ 的一个仿射开子集。由于 $g \circ f$ 作为支配有理映射
是恒等映射，复合 $U \to V \to X$ 在 $U$ 的一个稠密开集上是恒等映射。因此，
再将 $U$ 换成较小的仿射开集后，可假设 $U \to V \to X$ 是 $U$ 到 $X$ 的包含映射。
于是 $U \to V$ 是一个浸入（将《概形》引理 \ref{schemes-lemma-section-immersion}
应用于 $U \to g^{-1}(U) \to U$）。然而，交换 $U$ 与 $V$ 的角色并重复上述论证，
可见存在非空仿射开集 $V' \subset V$，使包含映射分解为 $V' \to U \to V$。
于是 $V' \to U$ 必然是开浸入。事实上，$V' \to f^{-1}(V') \to V'$ 是单态射
（《概形》引理 \ref{schemes-lemma-immersions-monomorphisms}），并且复合为恒等映射，
所以它们是同构。因而 $V'$ 同时同构于 $X$ 与 $Y$ 的某个开子概形。
在 $S$-有理映射的情形，完全相同的论证成立。
\end{proof}

\begin{remark}
\label{morphisms-remark-generalize-category}
下面给出不可约概形与支配有理映射所成范畴的一个推广。对概形 $X$，记 $X^0$ 为满足
$x \in X$ 且 $\dim(\mathcal{O}_{X, x}) = 0$ 的点的集合；换言之，$X^0$ 是
$X$ 的不可约分支的泛点集合。于是可以考虑如下范畴：
\begin{enumerate}
\item 对象是这样的概形 $X$，其中每个拟紧开集都有有限多个不可约分支；
\item 从 $X$ 到 $Y$ 的态射是有理映射 $f : U \to Y$；它从 $X$ 到 $Y$，
  并满足 $f(U^0) = Y^0$。
\end{enumerate}
如果 $U \subset X$ 是概形的稠密开集，则 $U^0 \subset X^0$ 不一定是等式；
但若 $X$ 是该范畴的对象，则情况确实如此。因此，该范畴中任意两个态射的复合
是良定义的，并且仍是该范畴中的态射。
\end{remark}

\begin{remark}
\label{morphisms-remark-pseudo-morphisms}
定义 \ref{morphisms-definition-rational-map} 有一个变体：只考虑态射 $U \to Y$，其定义在
概形论意义下稠密的开子概形 $U \subset X$ 上。利用引理
\ref{morphisms-lemma-intersection-scheme-theoretically-dense}，可知由此得到一个等价关系。
这些等价类称为从 $X$ 到 $Y$ 的{\it 伪态射}。若 $X$ 既约，则这两个概念相同。
\end{remark}

\section{双有理态射}
\label{morphisms-section-birational}

\noindent
你可能熟悉簇的双有理映射这一概念，它具有在目标的某个开子集上为同构的性质。
然而，一般而言，双有理态射可能在任何非空开集上都不是同构，见例
\ref{morphisms-example-birational-not-iso-over-open}。下面给出形式定义。

\begin{definition}
\label{morphisms-definition-birational}
\begin{reference}
\cite[(2.2.9)]{EGA1}
\end{reference}
设 $X$、$Y$ 是概形。假设 $X$ 与 $Y$ 都只有有限多个不可约分支。
若态射 $f : X \to Y$ 满足以下条件，则称其为{\it 双有理的}：
\begin{enumerate}
\item $f$ 在 $X$ 的不可约分支的泛点集合与 $Y$ 的不可约分支的泛点集合之间诱导双射；
\item 对泛点 $\eta \in X$（它属于 $X$ 的一个不可约分支），局部环映射
  $\mathcal{O}_{Y, f(\eta)} \to \mathcal{O}_{X, \eta}$ 是同构。
\end{enumerate}
\end{definition}

\noindent
下面将看到，双有理态射在泛点上的纤维是单点集。此外，在实际中遇到的大多数情形下，
不可约概形 $X$ 与 $Y$ 之间存在双有理态射，就意味着 $X$ 与 $Y$ 是双有理概形。

\begin{lemma}
\label{morphisms-lemma-birational-dominant}
设 $f : X \to Y$ 是概形态射，且只有有限多个不可约分支。若 $f$ 是双有理的，
则 $f$ 是支配的。
\end{lemma}

\begin{proof}
由引理 \ref{morphisms-lemma-generic-points-in-image-dominant} 及定义立即得到。
\end{proof}

\begin{lemma}
\label{morphisms-lemma-birational-generic-fibres}
设 $f : X \to Y$ 是双有理概形态射，且只有有限多个不可约分支。若 $y \in Y$ 是
一个不可约分支的泛点，则基变换
$X \times_Y \Spec(\mathcal{O}_{Y, y}) \to \Spec(\mathcal{O}_{Y, y})$
是同构。
\end{lemma}

\begin{proof}
不妨设 $Y = \Spec(B)$ 是仿射且不可约的。于是 $X$ 也不可约。若对任意非空
仿射开集 $U \subset X$ 结论成立，则对 $X$ 结论也成立（略去一个小论证）。因此不妨设
$X$ 也是仿射的，记作 $X = \Spec(A)$。
令 $y \in Y$ 对应于极小素理想 $\mathfrak q \subset B$。根据假设，$A$ 有唯一的极小素理想
$\mathfrak p$，它位于 $\mathfrak q$ 之上，并且
$B_\mathfrak q \to A_\mathfrak p$ 是同构。
于是 $A_\mathfrak q \to \kappa(\mathfrak p)$ 是满射，因而
$\mathfrak p A_\mathfrak q$ 是极大理想。另一方面，
$\mathfrak p A_\mathfrak q$ 是 $A_\mathfrak q$ 的唯一极小素理想。
因此 $\mathfrak p A_\mathfrak q$ 是 $A_\mathfrak q$ 的唯一素理想，且
$A_\mathfrak q = A_\mathfrak p$。由于 $A_\mathfrak q = A \otimes_B B_\mathfrak q$，
引理得证。
\end{proof}

\begin{example}
\label{morphisms-example-birational-not-iso-over-open}
下面给出两个双有理态射的例子，它们在目标的任何开集上都不是同构。

\medskip\noindent
第一个例子。设 $k$ 是无限域。设 $A = k[x]$。令
$B = k[x, \{y_{\alpha}\}_{\alpha \in k}]/
((x-\alpha)y_\alpha, y_\alpha y_\beta)$。
存在包含 $A \subset B$ 以及回缩 $B \to A$，后者通过令所有 $y_\alpha$ 都为零得到。
态射 $\Spec(A) \to \Spec(B)$ 与态射 $\Spec(B) \to \Spec(A)$ 都是双有理的，
但在任何开集上都不是同构。

\medskip\noindent
第二个例子。设 $A$ 是整环。设 $S \subset A$ 是不含 $0$ 的乘法闭子集。
令 $B = S^{-1}A$，则态射 $f : \Spec(B) \to \Spec(A)$ 是双有理的。
若存在开集 $U$，它属于 $\Spec(A)$，使得 $f^{-1}(U) \to U$ 是同构，
则存在 $a \in A$，使 $S$ 的每个元素都在主局部化 $A_a$ 中变为可逆元。
取 $A = \mathbf{Z}$ 且 $S$ 为奇整数集合，即得反例。
\end{example}

\begin{lemma}
\label{morphisms-lemma-birational-birational}
设 $f : X \to Y$ 是概形的双有理态射，且这些概形在基概形 $S$ 上并只有有限多个不可约分支。
假设下列条件之一成立：
\begin{enumerate}
\item $f$ 局部有限型且 $Y$ 既约；
\item $f$ 局部有限表示。
\end{enumerate}
则存在稠密开集 $U \subset X$ 与 $V \subset Y$，使得
$f(U) \subset V$ 且 $f|_U : U \to V$ 是同构。特别地，若 $X$ 与 $Y$ 不可约，
则 $X$ 与 $Y$ 是 $S$-双有理的。
\end{lemma}

\begin{proof}
立即可将情形约化为 $X$ 与 $Y$ 不可约的情形，此处略去。进一步缩小后，可假设 $X$ 与
$Y$ 是仿射的，即 $X = \Spec(A)$ 且 $Y = \Spec(B)$。
根据假设，$A$（分别地，$B$）有唯一的极小素理想 $\mathfrak p$（分别地，
$\mathfrak q$），素理想 $\mathfrak p$ 位于 $\mathfrak q$ 之上，并且
$B_\mathfrak q = A_\mathfrak p$。
由引理 \ref{morphisms-lemma-birational-generic-fibres}，有
$B_\mathfrak q = A_\mathfrak q = A_\mathfrak p$。

\medskip\noindent
假设 $B \to A$ 是有限型的，设 $A = B[x_1, \ldots, x_n]$。
存在 $b_i \in B$ 及 $g_i \in B \setminus \mathfrak q$，使得 $b_i/g_i$ 映到
$x_i$ 在 $A_\mathfrak q$ 中的像。因此，$b_i  - g_ix_i$ 在
$A_{g_i'}$ 中映为零，其中 $g_i' \in B \setminus \mathfrak q$。
令 $g = \prod g_i g'_i$，则 $B_g \to A_g$ 是满射。若进一步 $Y$ 既约，则映射
$B_g \to B_\mathfrak q$ 是单射，从而 $B_g \to A_g$ 也是单射。这证明了情形 (1)。

\medskip\noindent
情形 (2) 的证明。由上一段论证可假设 $B \to A$ 是满射。由于 $f$ 局部有限表示，
核 $J \subset B$ 是有限生成理想。设 $J = (b_1, \ldots, b_r)$。由于
$B_\mathfrak q = A_\mathfrak q$，存在 $g_i \in B \setminus \mathfrak q$ 使得
$g_i b_i = 0$。令 $g = \prod g_i$，则 $B_g \to A_g$ 是同构。
\end{proof}

\begin{lemma}
\label{morphisms-lemma-criterion-birational-finite-presentation}
设 $S$ 是概形。设 $X$、$Y$ 是不可约概形，并且在 $S$ 上局部有限表示。
设 $x \in X$、$y \in Y$ 是泛点。以下条件等价：
\begin{enumerate}
\item $X$ 与 $Y$ 是 $S$-双有理的；
\item $X$ 与 $Y$ 有非空开子概形，并且这些开子概形在 $S$ 上同构；
\item $x$ 与 $y$ 映到同一点 $s$（它属于 $S$），且
  $\mathcal{O}_{X, x}$ 与 $\mathcal{O}_{Y, y}$ 作为
  $\mathcal{O}_{S, s}$-代数同构。
\end{enumerate}
\end{lemma}

\begin{proof}
我们已经在引理 \ref{morphisms-lemma-birational-integral} 中证明了 (1) 与 (2) 的等价性。
(2) 蕴含 (3) 是直接的。为完成证明，假设 (3) 成立，并证明 (1)。
由引理 \ref{morphisms-lemma-rational-map-finite-presentation}，存在有理映射
$f : U \to Y$，它把 $x \in U$ 映到 $y$，并诱导给定的同构
$\mathcal{O}_{Y, y} \cong \mathcal{O}_{X, x}$。
因此 $f$ 是双有理态射，从而由引理
\ref{morphisms-lemma-birational-birational} 在非空开子集上诱导同构。
证明完毕。
\end{proof}

\begin{lemma}
\label{morphisms-lemma-common-open}
设 $S$ 是概形。设 $X$、$Y$ 是在 $S$ 上局部有限型的整概形。
设 $x \in X$、$y \in Y$ 是泛点。以下条件等价：
\begin{enumerate}
\item $X$ 与 $Y$ 是 $S$-双有理的；
\item $X$ 与 $Y$ 有非空开子概形，并且这些开子概形在 $S$ 上同构；
\item $x$ 与 $y$ 映到同一点 $s \in S$，且
  $\kappa(x) \cong \kappa(y)$ 作为 $\kappa(s)$-扩张同构。
\end{enumerate}
\end{lemma}

\begin{proof}
我们已经在引理 \ref{morphisms-lemma-birational-integral} 中证明了 (1) 与 (2) 的等价性。
(2) 蕴含 (3) 是直接的。为完成证明，假设 (3) 成立，并证明 (1)。
由代数引理 \ref{algebra-lemma-minimal-prime-reduced-ring}，有
$\mathcal{O}_{X, x} = \kappa(x)$ 且 $\mathcal{O}_{Y, y} = \kappa(y)$。
由引理 \ref{morphisms-lemma-rational-map-finite-presentation}，存在有理映射
$f : U \to Y$，它把 $x \in U$ 映到 $y$，并诱导给定的同构
$\mathcal{O}_{Y, y} \cong \mathcal{O}_{X, x}$。
因此 $f$ 是双有理态射，从而由引理
\ref{morphisms-lemma-birational-birational} 在非空开子集上诱导同构。
证明完毕。
\end{proof}
\section{概形态射的泛有限性}
\label{morphisms-section-generically-finite}

\noindent
本节刻画局部有限型且在某种意义下“泛有限”的概形之间的映射。

\begin{lemma}
\label{morphisms-lemma-generically-finite}
设 $X$、$Y$ 是概形。设 $f : X \to Y$ 是局部有限型的。
设 $\eta \in Y$ 是 $Y$ 的某个不可约分支的泛点。以下条件等价：
\begin{enumerate}
\item 集合 $f^{-1}(\{\eta\})$ 是有限的；
\item 存在仿射开集 $U_i \subset X$（$i = 1, \ldots, n$）以及
  $V \subset Y$，使 $f(U_i) \subset V$、$\eta \in V$ 且
  $f^{-1}(\{\eta\}) \subset \bigcup U_i$，并且每个
  $f|_{U_i} : U_i \to V$ 都是有限的。
\end{enumerate}
若 $f$ 拟分离，则这些条件还等价于
\begin{enumerate}
\item[(3)] 存在仿射开集 $V \subset Y$ 以及 $U \subset X$，使
  $f(U) \subset V$、$\eta \in V$ 且 $f^{-1}(\{\eta\}) \subset U$，
  并且 $f|_U : U \to V$ 是有限的。
\end{enumerate}
若 $f$ 拟紧且拟分离，则这些条件还等价于
\begin{enumerate}
\item[(4)] 存在仿射开集 $V \subset Y$，$\eta \in V$，使得
  $f^{-1}(V) \to V$ 是有限的。
\end{enumerate}
\end{lemma}

\begin{proof}
问题在基上是局部的。因此可用 $Y$ 的一个包含 $\eta$ 的仿射邻域替换，以下证明始终假设
$Y$ 是仿射的，写作 $Y = \Spec(R)$。

\medskip\noindent
显然 (2) 蕴含 (1)。假设
$f^{-1}(\{\eta\}) = \{\xi_1, \ldots, \xi_n\}$ 是有限集。
取仿射开集 $U_i \subset X$，其中 $\xi_i \in U_i$。由代数引理
\ref{algebra-lemma-generically-finite}，将 $Y$ 替换为 $Y$ 中的某个标准开集后，
每个态射 $U_i \to Y$ 都是有限的。换言之，(2) 成立。

\medskip\noindent
显然 (3) 蕴含 (1)。现在假设 $f$ 拟分离且 (1) 成立。写成
$f^{-1}(\{\eta\}) = \{\xi_1, \ldots, \xi_n\}$。由引理
\ref{morphisms-lemma-finite-fibre}，这些 $\xi_i$ 之间没有特化关系。
由于 $\xi_i$ 映到不可约分支的泛点 $\eta$，它属于 $Y$；不存在
非平凡特化 $\xi \leadsto \xi_i$，该特化发生在 $X$ 中（因为 $\xi$ 也会映到 $\eta$）。
因此每个 $\xi_i$ 都是 $X$ 的某个不可约分支的泛点。
由于 $Y$ 仿射且 $f$ 拟分离，$X$ 也是拟分离的。由性质引理
\ref{properties-lemma-maximal-points-affine}，可找到仿射开集 $U \subset X$，其中包含每个
$\xi_i$。由代数引理 \ref{algebra-lemma-generically-finite}，将 $Y$ 替换为
$Y$ 中的某个标准开集后，态射 $U \to Y$ 是有限的。换言之，(3) 成立。

\medskip\noindent
显然 (4) 在不对 $f$ 作任何额外假设时蕴含 (1)--(3)。现在假设 $f$ 拟紧且拟分离。
我们需要证明等价条件 (1)--(3) 蕴含 (4)。
设 $U$、$V$ 如 (3) 中所述，并以 $Y$ 替换为 $V$。由于 $f$ 拟紧且 $Y$ 拟紧（因为它仿射），
$X$ 拟紧。因此 $Z = X \setminus U$ 拟紧，故态射 $f|_Z : Z \to Y$ 拟紧。
由 $Z$ 的构造，$\eta \not \in f(Z)$。由引理
\ref{morphisms-lemma-quasi-compact-generic-point-not-in-image}，存在仿射开邻域 $V'$，包含 $\eta$ 且位于 $Y$ 中，
使 $f^{-1}(V') \cap Z = \emptyset$。于是 $f^{-1}(V') \subset U$，这意味着
$f^{-1}(V') \to V'$ 是有限的。
\end{proof}

\begin{example}
\label{morphisms-example-counter-generically-finite}
令 $A = \prod_{n \in \mathbf{N}} \mathbf{F}_2$。$A$ 的每个元素都是幂等元，
因此每个素理想都是极大理想，剩余域为 $\mathbf{F}_2$。于是
$X = \Spec(A)$ 的拓扑完全不连通且拟紧。投影映射 $A \to \mathbf{F}_2$ 定义
$\Spec(A)$ 的开点。由于 $X$ 拟紧，不可能 $X$ 的所有点都是开的。取 $x \in X$ 为
一个不是开点的闭点。于是可以构造概形 $Y$，把 $X$ 的两个副本沿
$X \setminus \{x\}$ 粘合；换言之，这是 $X$，其中 $x$ 被加倍，比较
《概形》例 \ref{schemes-example-affine-space-zero-doubled}。
态射 $f : Y \to X$ 是拟紧、有限型且具有有限纤维，但不是拟分离的。
$x \in X$ 是 $X$ 的某个不可约分支的泛点（因为 $X$ 完全不连通）。
然而，引理 \ref{morphisms-lemma-generically-finite} 的性质 (3) 与 (4) 不成立。
原因是，对任意开邻域 $x \in U \subset X$，逆像 $f^{-1}(U)$ 都不是仿射的，
因为 $f^{-1}(U)$ 上的函数不能分离位于 $x$ 上方的两个点（证明略；这是一个很好的练习）。
因此，在引理的 (3) 与 (4) 中，要求 $f$ 拟分离是必要的。
\end{example}

\begin{remark}
\label{morphisms-remark-quasi-finite-finite-over-dense-open}
引理 \ref{morphisms-lemma-generically-finite} 的一个替代说法是：拟有限态射在目标的某个稠密开集上有限。
这将在《更多态射》引理
\ref{more-morphisms-lemma-quasi-finite-finite-over-dense-open} 中证明。
\end{remark}
\begin{lemma}
\label{morphisms-lemma-quasi-finiteness-over-generic-point}
设 $X$、$Y$ 是概形。设 $f : X \to Y$ 是局部有限型的。
令 $X^0$、分别的 $Y^0$ 表示 $X$、$Y$ 的不可约分支的泛点集合。令
$\eta \in Y^0$。以下条件等价：
\begin{enumerate}
\item $f^{-1}(\{\eta\}) \subset X^0$；
\item $f$ 在所有位于 $\eta$ 上方的点处拟有限；
\item $f$ 在所有 $\xi \in X^0$（它们位于 $\eta$ 上方）处拟有限。
\end{enumerate}
\end{lemma}

\begin{proof}
条件 (1) 蕴含纤维 $X_\eta$ 的点之间没有特化关系。因此由引理
\ref{morphisms-lemma-quasi-finite-at-point-characterize}，(2) 成立。
蕴含 (2) $\Rightarrow$ (3) 是直接的。
由于 $\eta$ 是 $Y$ 的泛点，$X_\eta$ 的泛点也是 $X$ 的泛点。因此由 (3) 及引理
\ref{morphisms-lemma-quasi-finite-at-point-characterize}，$X_\eta$ 的泛点也都是闭点。
于是 $X_\eta$ 的所有点都是泛点，从而 (1) 成立。
\end{proof}

\begin{lemma}
\label{morphisms-lemma-finite-over-dense-open}
设 $X$、$Y$ 是概形。设 $f : X \to Y$ 是局部有限型的。
令 $X^0$、分别的 $Y^0$ 表示 $X$、$Y$ 的不可约分支的泛点集合。假设
\begin{enumerate}
\item $X^0$ 与 $Y^0$ 有限，且 $f^{-1}(Y^0) = X^0$；
\item $f$ 拟紧或 $f$ 分离。
\end{enumerate}
则存在稠密开集 $V \subset Y$，使得 $f^{-1}(V) \to V$ 是有限的。
\end{lemma}

\begin{proof}
由于 $Y$ 只有有限多个不可约分支，可以找到一个稠密开集，它是不交的不可约分支之并。
因此不妨设 $Y$ 是具有泛点 $\eta$ 的不可约仿射概形。由于 $\eta$ 上的纤维有限，因为 $X^0$ 有限。

\medskip\noindent
假设 $f$ 分离且 $Y$ 不可约仿射。取 $V \subset Y$ 与 $U \subset X$，如引理
\ref{morphisms-lemma-generically-finite} 的 (3) 所述。由于 $f|_U : U \to V$ 有限，依据引理
\ref{morphisms-lemma-image-proper-scheme-closed} 与 \ref{morphisms-lemma-finite-proper}，
$U \subset f^{-1}(V)$ 既闭又开。因此有 $f^{-1}(V) = U \amalg W$，其中 $W$ 是
$X$ 的开子概形。然而，由于 $U$ 包含 $X$ 的所有泛点，得到 $W = \emptyset$，
如所需。

\medskip\noindent
假设 $f$ 拟紧且 $Y$ 不可约仿射。则 $X$ 拟紧，因而存在稠密开子概形 $U \subset X$，它是分离的
（性质引理 \ref{properties-lemma-quasi-compact-dense-open-separated}）。
由于泛点集合 $X^0$ 有限，得到 $X^0 \subset U$。于是
$\eta \not \in f(X \setminus U)$。由于 $X \setminus U \to Y$ 拟紧，得到一个非空开集
$V \subset Y$，使 $f^{-1}(V) \subset U$，见引理 \ref{morphisms-lemma-quasi-compact-dominant}。
以 $X$ 替换为 $f^{-1}(V)$，以 $Y$ 替换为 $V$，便约化为上一段已经处理的分离情形。
\end{proof}

\begin{lemma}
\label{morphisms-lemma-birational-isomorphism-over-dense-open}
设 $X$、$Y$ 是概形。设 $f : X \to Y$ 是双有理态射，且这两个概形只有有限多个不可约分支。
假设
\begin{enumerate}
\item $f$ 拟紧或 $f$ 分离；
\item $f$ 局部有限型且 $Y$ 既约，或 $f$ 局部有限表示。
\end{enumerate}
则存在稠密开集 $V \subset Y$，使得 $f^{-1}(V) \to V$ 是同构。
\end{lemma}

\begin{proof}
由引理 \ref{morphisms-lemma-finite-over-dense-open}，不妨设 $f$ 是有限的。
由于 $Y$ 只有有限多个不可约分支，可以找到一个稠密开集，它是不交的不可约分支之并。
因此不妨设 $Y$ 不可约。由引理 \ref{morphisms-lemma-birational-birational}，存在非空开集
$U \subset X$，使 $f|_U : U \to Y$ 是开浸入。去掉闭子集
$f$ 有限时去掉闭子集 $f(X \setminus U)$，它来自 $Y$，后得到 $f$ 是同构。
\end{proof}
\begin{lemma}
\label{morphisms-lemma-finite-degree}
设 $X$、$Y$ 是整概形。设 $f : X \to Y$ 是局部有限型的。
假设 $f$ 支配。以下条件等价：
\begin{enumerate}
\item 扩张 $R(Y) \subset R(X)$ 的超越次数为 $0$；
\item 扩张 $R(Y) \subset R(X)$ 是有限的；
\item 存在非空仿射开集 $U \subset X$ 与 $V \subset Y$，使 $f(U) \subset V$，
  且 $f|_U : U \to V$ 是有限的；
\item $X$ 的泛点是 $X$ 中映到 $Y$ 的泛点的唯一点。
\end{enumerate}
若 $f$ 分离或 $f$ 拟紧，则这些条件还等价于
\begin{enumerate}
\item[(5)] 存在非空仿射开集 $V \subset Y$，使得
  $f^{-1}(V) \to V$ 是有限的。
\end{enumerate}
\end{lemma}

\begin{proof}
取任意仿射开集 $\Spec(A) = U \subset X$ 与 $\Spec(R) = V \subset Y$，使 $f(U) \subset V$。
按定义，$R$ 与 $A$ 是整环。环映射 $R \to A$ 是有限型的
（见引理 \ref{morphisms-lemma-locally-finite-type-characterize}）。由引理
\ref{morphisms-lemma-dominant-finite-number-irreducible-components}，$X$ 的泛点映到 $Y$ 的泛点，
因此 $R \to A$ 是单射。令 $K = R(Y)$ 为 $R$ 的分式域，令 $L = R(X)$ 为 $A$ 的分式域。
于是 $L/K$ 是有限生成域扩张，从而 (1) 等价于 (2)。

\medskip\noindent
假设 (2) 成立。设 $x_1, \ldots, x_n \in A$ 是 $A$ 在 $R$ 上的生成元。
根据假设，存在非零多项式 $P_i(X) \in R[X]$，使 $P_i(x_i) = 0$。
令 $f_i \in R$ 为 $P_i$ 的首项系数。于是
$R_{f_1 \ldots f_n} \to A_{f_1 \ldots f_n}$ 是有限的，即 (3) 成立。
注意 (3) 蕴含 (2)。因此 (1)、(2)、(3) 全部等价。

\medskip\noindent
令 $\eta$ 是 $X$ 的泛点，并令 $\eta' \in Y$ 是 $Y$ 的泛点。假设 (4)。则
$\dim_\eta(X_{\eta'}) = 0$，由引理 \ref{morphisms-lemma-dimension-fibre-at-a-point} 可见
$R(X) = \kappa(\eta)$ 的超越次数为 $0$，相对于 $R(Y) = \kappa(\eta')$。
换言之，(1) 成立。现在假设等价条件 (1)、(2)、(3) 成立。设 $x \in X$ 是映到 $\eta'$ 的点。
由于 $x$ 是 $\eta$ 的特化点，得到包含关系
$R(Y) \subset \mathcal{O}_{X, x} \subset R(X)$，这意味着 $\mathcal{O}_{X, x}$ 是域，见
代数引理 \ref{algebra-lemma-integral-over-field}。因此 $x = \eta$。于是 (1)--(4) 全部等价。

\medskip\noindent
显然 (5) 在不对 $f$ 作任何额外假设时蕴含 (3)。剩下要证明的是：若 $f$ 分离或拟紧，
则等价条件 (1)--(4) 蕴含 (5)。这由引理 \ref{morphisms-lemma-finite-over-dense-open} 得出。
\end{proof}
\begin{definition}
\label{morphisms-definition-degree}
令 $X$ 与 $Y$ 为整概形。
令 $f : X \to Y$ 为局部有限型且支配的态射。
假设 $[R(X) : R(Y)] < \infty$，或引理 \ref{morphisms-lemma-finite-degree} 的等价条件 (1)--(4) 中的任一条件。
则正整数
$$
\deg(X/Y) = [R(X) : R(Y)]
$$
称为 $X$ 在 $Y$ 上的\emph{次数}。
\end{definition}

\noindent
可以将这个概念推广到态射
$f : X \to Y$，如果 (a) $Y$ 是带泛点 $\eta$ 的整概形，(b) $f$ 是局部有限型的，并且 (c) $f^{-1}(\{\eta\})$ 是有限的。
在此情形，我们定义
$$
\deg(X/Y)
=
\sum\nolimits_{\xi \in X, \ f(\xi) = \eta}
\dim_{R(Y)} (\mathcal{O}_{X, \xi}).
$$
也就是说，由于 $R(Y) = \kappa(\eta) = \mathcal{O}_{Y, \eta}$
（引理 \ref{morphisms-lemma-integral-scheme-rational-functions}），上述维数由上面的引理
\ref{morphisms-lemma-generically-finite} 可知都是有限的。
不过，对大多数应用而言，上面给出的定义才是适用的定义。

\begin{lemma}
\label{morphisms-lemma-degree-composition}
令 $X$、$Y$、$Z$ 为整概形。
令 $f : X \to Y$ 与 $g : Y \to Z$ 为局部有限型的支配态射。假设 $[R(X) : R(Y)] < \infty$ 且
$[R(Y) : R(Z)] < \infty$。则
$$
\deg(X/Z) = \deg(X/Y) \deg(Y/Z).
$$
\end{lemma}

\begin{proof}
这是域的有限扩张塔中次数的可乘性所致，见
《域》，引理 \ref{fields-lemma-multiplicativity-degrees}。
\end{proof}

\begin{remark}
\label{morphisms-remark-definition-generically-finite}
令 $f : X \to Y$ 为局部有限型的概形态射。
对于定义\emph{泛有限}态射，至少有两种可采用的性质。这两种性质分别对应于希望该性质在源上局部成立还是在目标上局部成立：
\begin{enumerate}
\item（目标上的局部性质；Ravi Vakil 建议。）
假设 $Y$ 的每个拟紧开子集都有有限多个不可约分支
（例如 $Y$ 局部 Noether）。
要求每个泛点的逆像都是有限的，见引理 \ref{morphisms-lemma-generically-finite}。
\item（源上的局部性质。）要求存在稠密开集 $U \subset X$，使得 $U \to Y$ 是局部拟有限的。
\end{enumerate}
在情形 (1) 中，不借助关于 $Y$ 的辅助条件也可以表述该要求，但对一般概形而言这可能不是正确的概念。按 (2) 表述的性质并不蕴含泛点上的纤维有限；但是，如果 $f$ 是拟紧的且 $Y$ 满足 (1) 中的条件，则确实蕴含这一点。
\end{remark}

\begin{definition}
\label{morphisms-definition-modification}
令 $X$ 为整概形。$X$ 的一个\emph{修改}是一个双有理的真态射 $f : X' \to X$，其中 $X'$ 为整概形。
\end{definition}

\noindent
令 $f : X' \to X$ 是定义中的修改。由引理 \ref{morphisms-lemma-finite-degree}，存在非空的 $U \subset X$，使得 $f^{-1}(U) \to U$ 是有限的。由一般平坦性
（命题 \ref{morphisms-proposition-generic-flatness}），我们可以假设 $f^{-1}(U) \to U$ 是平坦且有限表示的。
因此 $f^{-1}(U) \to U$ 是有限局部自由的（引理 \ref{morphisms-lemma-finite-flat}）。
由于 $f$ 是双有理的，$X'$ 在 $X$ 上的次数为 $1$。
所以 $f^{-1}(U) \to U$ 是次数为 $1$ 的有限局部自由态射，换言之，它是一个同构。
因此，我们可以将\emph{修改}重新定义为整概形之间的真态射 $f : X' \to X$，使得 $f^{-1}(U) \to U$ 是同构，其中 $U \subset X$ 是某个非空开集。

\begin{definition}
\label{morphisms-definition-alteration}
\begin{reference}
\cite[Definition 2.20]{alterations}
\end{reference}
令 $X$ 为整概形。$X$ 的一个\emph{变更（alteration）}是一个真支配态射 $f : Y \to X$，其中 $Y$ 为整概形，并且 $f^{-1}(U) \to U$ 是有限的，其中 $U \subset X$ 是某个非空开集。
\end{definition}

\noindent
这是 \cite{alterations} 中给出的定义，区别在于这里不要求 $X$ 与 $Y$ 是 Noether 概形。按照上面的论证，我们看到，变更是整概形之间的真支配态射 $f : Y \to X$，它诱导函数域的有限扩张，即满足引理 \ref{morphisms-lemma-finite-degree} 的等价条件。

\section{维数公式}
\label{morphisms-section-dimension-formula}

\noindent
对于 Noether 概形之间的态射，我们可以进一步说明局部环的维数。下面是一个重要（且不难证明）的结果。回忆 $R(X)$ 表示整概形 $X$ 的函数域。

\begin{lemma}
\label{morphisms-lemma-dimension-formula}
设 $S$ 为概形。
设 $f : X \to S$ 为概形态射。
令 $x \in X$，并置 $s = f(x)$。
假设
\begin{enumerate}
\item $S$ 是局部 Noether 的；
\item $f$ 是局部有限型的；
\item $X$ 与 $S$ 是整概形；并且
\item $f$ 是支配的。
\end{enumerate}
我们有
\begin{equation}
\label{morphisms-equation-dimension-formula}
\dim(\mathcal{O}_{X, x})
\leq
\dim(\mathcal{O}_{S, s}) + \text{trdeg}_{R(S)}R(X)
- \text{trdeg}_{\kappa(s)} \kappa(x).
\end{equation}
此外，如果 $S$ 是普遍链状的，则等号成立。
\end{lemma}

\begin{proof}
相应的代数表述是《代数》中的引理 \ref{algebra-lemma-dimension-formula}。
\end{proof}

\begin{lemma}
\label{morphisms-lemma-dimension-formula-general}
设 $S$ 为概形。设 $f : X \to S$ 为概形态射。
令 $x \in X$，并置 $s = f(x)$。假设 $S$ 是局部 Noether 的且 $f$ 是局部有限型的，
我们有
\begin{equation}
\label{morphisms-equation-dimension-formula-general}
\dim(\mathcal{O}_{X, x})
\leq
\dim(\mathcal{O}_{S, s}) + E - \text{trdeg}_{\kappa(s)} \kappa(x).
\end{equation}
其中 $E$ 是所有 $\text{trdeg}_{\kappa(f(\xi))}(\kappa(\xi))$ 的最大值，
其中 $\xi$ 遍历 $X$ 中包含 $x$ 的不可约分支的泛点。
\end{lemma}

\begin{proof}
令 $X_1, \ldots, X_n$ 为 $X$ 的不可约分支中包含 $x$ 的那些，并赋予它们约化的诱导概形结构。
这些分支对应于 $\mathfrak q_i$（即 $\mathcal{O}_{X, x}$ 的极小素理想），因此它们只有有限多个
（《概形》引理 \ref{schemes-lemma-specialize-points} 与《代数》引理
\ref{algebra-lemma-Noetherian-irreducible-components}）。
于是
$\dim(\mathcal{O}_{X, x}) = \max \dim(\mathcal{O}_{X, x}/\mathfrak q_i)
= \max \dim(\mathcal{O}_{X_i, x})$。
在定义中出现的 $\xi$ 恰好是 $E$ 所取的那些泛点，即不可约分支的泛点 $\xi_i \in X_i$。令
$Z_i = \overline{\{f(\xi_i)\}} \subset S$，并赋予它约化的诱导概形结构。
$X_i \to X \to S$ 的复合经过 $Z_i$ 分解
（《概形》引理 \ref{schemes-lemma-map-into-reduction}）。因此可将维数公式（引理
\ref{morphisms-lemma-dimension-formula}）应用于
$\dim(\mathcal{O}_{X_i, x}) \leq \dim(\mathcal{O}_{Z_i, x}) +
\text{trdeg}_{\kappa(f(\xi))}(\kappa(\xi)) -
\text{trdeg}_{\kappa(s)} \kappa(x)$。
汇总以上各式即得本引理。
\end{proof}
\noindent
一个应用是：在赋有维数函数的普遍链状概形上有限型的任意概形上构造维数函数。关于维数函数的定义，见
《拓扑》定义 \ref{topology-definition-dimension-function}。

\begin{lemma}
\label{morphisms-lemma-dimension-function-propagates}
令 $S$ 为局部 Noether 且普遍链状的概形。
令 $\delta : S \to \mathbf{Z}$ 为维数函数。
令 $f : X \to S$ 为概形态射。
假设 $f$ 局部有限型。
则映射
\begin{align*}
\delta = \delta_{X/S} : X & \longrightarrow \mathbf{Z} \\
x & \longmapsto \delta(f(x)) + \text{trdeg}_{\kappa(f(x))} \kappa(x)
\end{align*}
是 $X$ 上的维数函数。
\end{lemma}

\begin{proof}
令 $f : X \to S$ 为局部有限型的。
令 $x \leadsto y$、$x \not = y$ 是 $X$ 中的一个特化。
我们需要证明 $\delta_{X/S}(x) > \delta_{X/S}(y)$，并且 $\delta_{X/S}(x) = \delta_{X/S}(y) + 1$（当 $y$ 是 $x$ 的直接特化时）。

\medskip\noindent
选取仿射开集 $V \subset S$，其包含 $y$ 的像；并选取仿射开集 $U \subset X$，它映入 $V$ 且包含 $y$。
显然可以将 $X$ 换成 $U$，将 $S$ 换成 $V$。因此不妨设 $X = \Spec(A)$、
$S = \Spec(R)$，且 $f$ 由环映射 $R \to A$ 给出。环 $R$ 是普遍链状的
（引理 \ref{morphisms-lemma-universally-catenary-local}），且环映射 $R \to A$ 是有限型的
（引理 \ref{morphisms-lemma-locally-finite-type-characterize}）。

\medskip\noindent
令 $\mathfrak q \subset A$ 为对应于点 $x$ 的素理想，令 $\mathfrak p \subset R$ 为对应于 $f(x)$ 的素理想。
将 $\delta'$ 定义为 $\delta$ 在 $S' = \Spec(R/\mathfrak p) \subset S$ 上的限制；它是维数函数。
环 $R/\mathfrak p$ 是普遍链状的。
将 $\delta_{X/S}$ 限制到 $X' = \Spec(A/\mathfrak q)$，显然等于函数 $\delta_{X'/S'}$，后者由维数函数 $\delta'$ 构造。
因此，除上述假设外，我们还可以假设 $R \subset A$ 是整环，换言之，$X$ 与 $S$ 是整概形，且 $x$ 是 $X$ 的泛点、$f(x)$ 是 $S$ 的泛点。

\medskip\noindent
注意 $\mathcal{O}_{X, x} = R(X)$，并且由于 $x \leadsto y$ 且 $x \not = y$，环 $\mathcal{O}_{X, y}$ 的谱至少有两个点
（《概形》引理 \ref{schemes-lemma-specialize-points}），因此 $\dim(\mathcal{O}_{X, y}) > 0$。
如果 $y$ 是 $x$ 的直接特化，则 $\Spec(\mathcal{O}_{X, y}) = \{x, y\}$ 且
$\dim(\mathcal{O}_{X, y}) = 1$。

\medskip\noindent
写 $s = f(x)$，$t = f(y)$。我们计算
\begin{align*}
\delta_{X/S}(x) - \delta_{X/S}(y)
& =
\delta(s) + \text{trdeg}_{\kappa(s)} \kappa(x)
- \delta(t) - \text{trdeg}_{\kappa(t)} \kappa(y) \\
& =
\delta(s) - \delta(t) +
\text{trdeg}_{R(S)} R(X) - \text{trdeg}_{\kappa(t)} \kappa(y) \\
& =
\delta(s) - \delta(t) + \dim(\mathcal{O}_{X, y})
- \dim(\mathcal{O}_{S, t})
\end{align*}
其中最后一步使用了（\ref{morphisms-equation-dimension-formula}）中的等式。
由于 $\delta$ 是概形 $S$ 上的维数函数，且 $s \in S$ 是泛点，差值
$\delta(s) - \delta(t)$ 由《拓扑》引理
\ref{topology-lemma-dimension-function-catenary} 可知等于 $\text{codim}(\overline{\{t\}}, S)$。
由《性质》引理
\ref{properties-lemma-codimension-local-ring}，这等于 $\dim(\mathcal{O}_{S, t})$。
因此我们得到
$$
\delta_{X/S}(x) - \delta_{X/S}(y) = \dim(\mathcal{O}_{X, y})
$$
结合上面对 $\dim(\mathcal{O}_{X, y})$ 的说明即得本引理。
\end{proof}
\noindent
维数公式的另一个应用是：维数在“变更”（稍后定义）下不变。

\begin{lemma}
\label{morphisms-lemma-alteration-dimension}
令 $f : X \to Y$ 为概形态射。假设
\begin{enumerate}
\item $Y$ 是局部 Noether 的；
\item $X$ 与 $Y$ 是整概形；
\item $f$ 是支配的；并且
\item $f$ 是局部有限型的。
\end{enumerate}
则
$$
\dim(X) \leq \dim(Y) + \text{trdeg}_{R(Y)} R(X).
$$
如果 $f$ 是闭的\footnote{例如若 $f$ 是真态射，见《概形》定义
\ref{morphisms-definition-proper}。}，则等号成立。
\end{lemma}

\begin{proof}
令 $f : X \to Y$ 如引理中所述。
令 $\xi_0 \leadsto \xi_1 \leadsto \ldots \leadsto \xi_e$ 是 $X$ 中的一列特化。置 $x = \xi_e$，
$y = f(x)$。注意 $e \leq \dim(\mathcal{O}_{X, x})$，因为给出的特化发生在
$\mathcal{O}_{X, x}$ 的谱中，见《概形》引理
\ref{schemes-lemma-specialize-points}。
由维数公式（引理 \ref{morphisms-lemma-dimension-formula}），我们看到
\begin{align*}
e & \leq \dim(\mathcal{O}_{X, x}) \\
& \leq \dim(\mathcal{O}_{Y, y}) +
\text{trdeg}_{R(Y)} R(X) - \text{trdeg}_{\kappa(y)} \kappa(x) \\
& \leq \dim(\mathcal{O}_{Y, y}) + \text{trdeg}_{R(Y)} R(X)
\end{align*}
因此如愿得到 $e \leq \dim(Y) + \text{trdeg}_{R(Y)} R(X)$。

\medskip\noindent
现在再假设 $f$ 也是闭的。
设 $\overline{\xi}_0 \leadsto \overline{\xi}_1 \leadsto \ldots
\leadsto \overline{\xi}_d$ 是 $Y$ 中的一列特化。
我们要证明 $\dim(X) \geq d + r$。
不妨设 $\overline{\xi}_0 = \eta$ 是 $Y$ 的泛点。
$X_\eta$ 是在 $\kappa(\eta) = R(Y)$ 上局部有限型的概形。由于 $f$ 支配，它非空。因此由引理
\ref{morphisms-lemma-ubiquity-Jacobson-schemes}，它是 Jacobson 概形。
于是由引理 \ref{morphisms-lemma-jacobson-finite-type-points}，可以找到闭点
$\xi_0 \in X_\eta$，并且扩张 $\kappa(\eta) \subset \kappa(\xi_0)$ 是有限扩张。
注意 $\mathcal{O}_{X, \xi_0} = \mathcal{O}_{X_\eta, \xi_0}$，因为 $\eta$ 是 $Y$ 的泛点。
因此，由维数公式引理 \ref{morphisms-lemma-dimension-formula} 可知 $\dim(\mathcal{O}_{X, \xi_0}) = r$；该引理应用于概形 $X_\eta$（它在普遍链状概形
$\Spec(\kappa(\eta))$ 上；见引理 \ref{morphisms-lemma-ubiquity-uc}）及点 $\xi_0$。
这意味着我们可以找到
$\xi_{-r} \leadsto \ldots \leadsto \xi_{-1} \leadsto \xi_0$ 于 $X$ 中。
另一方面，由于 $f$ 是闭的，特化沿着 $f$ 提升，见
《拓扑》引理 \ref{topology-lemma-closed-open-map-specialization}。
因此，由于 $\xi_0$ 位于
$\eta = \overline{\xi}_0$ 上方，我们可以找到特化
$\xi_0 \leadsto \xi_1 \leadsto \ldots \leadsto \xi_d$，它们分别位于
$\overline{\xi}_0 \leadsto \overline{\xi}_1 \leadsto \ldots
\leadsto \overline{\xi}_d$ 上方。换言之，我们有
$$
\xi_{-r} \leadsto \ldots \leadsto \xi_{-1} \leadsto \xi_0
\leadsto \xi_1 \leadsto \ldots \leadsto \xi_d
$$
这意味着 $\dim(X) \geq d + r$，如愿以偿。
\end{proof}

\begin{lemma}
\label{morphisms-lemma-alteration-dimension-general}
令 $f : X \to Y$ 为概形态射。假设 $Y$ 是局部 Noether 的且 $f$ 是局部有限型的。
则
$$
\dim(X) \leq \dim(Y) + E
$$
其中 $E$ 是所有 $\text{trdeg}_{\kappa(f(\xi))}(\kappa(\xi))$ 的上确界，
其中 $\xi$ 遍历 $X$ 的不可约分支的泛点。
\end{lemma}

\begin{proof}
这是引理 \ref{morphisms-lemma-dimension-formula-general} 与《性质》引理
\ref{properties-lemma-dimension} 的直接推论。
\end{proof}
\section{相对正规化}
\label{morphisms-section-normalization-X-in-Y}

\noindent
本节构造“一个概形在另一个概形中的正规化”。

\begin{lemma}
\label{morphisms-lemma-integral-closure}
设 $X$ 为概形。设 $\mathcal{A}$ 为一个准凝聚的 $\mathcal{O}_X$-代数层。按下列规则定义的子层
$\mathcal{A}' \subset \mathcal{A}$：
$$
U \longmapsto \{f \in \mathcal{A}(U) \mid
f_x \in \mathcal{A}_x \text{ integral over } \mathcal{O}_{X, x}
\text{ 对所有 }x \in U\}
$$
是准凝聚的 $\mathcal{O}_X$-代数层，茎 $\mathcal{A}'_x$ 是 $\mathcal{O}_{X, x}$ 在
$\mathcal{A}_x$ 中的整闭包，并且对任意仿射开集 $U \subset X$，环
$\mathcal{A}'(U) \subset \mathcal{A}(U)$ 是 $\mathcal{O}_X(U)$ 在
$\mathcal{A}(U)$ 中的整闭包。
\end{lemma}

\begin{proof}
由于这些条件具有局部性质，这是一个子层。
由《代数》引理 \ref{algebra-lemma-integral-closure-is-ring}，它是一个 $\mathcal{O}_X$-代数层。
令 $U \subset X$ 为仿射开集。设 $U = \Spec(R)$，并设 $\mathcal{A}$ 是与
$R$-代数 $A$ 关联的准凝聚层。根据
《代数》引理 \ref{algebra-lemma-integral-closure-stalks}，$\mathcal{A}'$ 在
$U$ 上的值由 $A'$（即 $R$ 在 $A$ 中的整闭包）给出。这证明了引理的最后一个断言。
为证明 $\mathcal{A}'$ 准凝聚，只需证明 $\mathcal{A}'(D(f)) = A'_f$。
这由整闭包与局部化可交换这一事实得出，见
《代数》引理 \ref{algebra-lemma-integral-closure-localize}。
同一事实也表明茎具有所述形式。
\end{proof}

\begin{definition}
\label{morphisms-definition-integral-closure}
设 $X$ 为概形。设 $\mathcal{A}$ 为准凝聚的 $\mathcal{O}_X$-代数层。$\mathcal{O}_X$
在 $\mathcal{A}$ 中的\emph{整闭包}是上面引理 \ref{morphisms-lemma-integral-closure} 中构造的
准凝聚 $\mathcal{O}_X$-子代数 $\mathcal{A}' \subset \mathcal{A}$。
\end{definition}

\noindent
在上述定义的情形下，我们可以考虑相对谱的态射
$$
\xymatrix{
Y = \underline{\Spec}_X(\mathcal{A}) \ar[rr] \ar[rd] & &
X' = \underline{\Spec}_X(\mathcal{A}') \ar[ld] \\
& X &
}
$$
见引理 \ref{morphisms-lemma-affine-equivalence-algebras}。
态射 $X' \to X$ 将是 $X$ 在概形 $Y$ 中的正规化。
下面给出稍微一般的情形。设我们有概形的拟紧且拟分离态射 $f : Y \to X$。
在此情形，$\mathcal{O}_X$-代数层 $f_*\mathcal{O}_Y$ 是准凝聚的，见
《概形》引理 \ref{schemes-lemma-push-forward-quasi-coherent}。
取 $\mathcal{O}' \subset f_*\mathcal{O}_Y$ 的整闭包，得到准凝聚的 $\mathcal{O}_X$-代数层，
其相对谱就是 $X$ 在 $Y$ 中的正规化。下面是形式定义。

\begin{definition}
\label{morphisms-definition-normalization-X-in-Y}
设 $f : Y \to X$ 为概形的拟紧且拟分离态射。
令 $\mathcal{O}'$ 为 $\mathcal{O}_X$ 在 $f_*\mathcal{O}_Y$ 中的整闭包。$X$ 在 $Y$ 中的\emph{正规化}是概形
\footnote{概形 $X'$ 不必是正规概形；例如若 $Y = X$ 且 $f = \text{id}_X$，则 $X' = X$。}
$$
\nu : X' = \underline{\Spec}_X(\mathcal{O}') \to X
$$
它位于 $X$ 上。它带有初始态射 $f$ 的自然分解
$$
Y \xrightarrow{f'} X' \xrightarrow{\nu} X
$$
。\end{definition}

\noindent
该分解是典范态射
$Y \to \underline{\Spec}_X(f_*\mathcal{O}_Y)$（见《构造》引理
\ref{constructions-lemma-canonical-morphism}）与由包含映射
$\mathcal{O}' \to f_*\mathcal{O}_Y$ 产生的相对谱态射的复合。我们可以如下刻画正规化。

\begin{lemma}
\label{morphisms-lemma-characterize-normalization}
设 $f : Y \to X$ 为概形的拟紧且拟分离态射。
分解 $f = \nu \circ f'$（其中 $\nu : X' \to X$ 是 $X$ 在 $Y$ 中的正规化）由以下两个性质刻画：
\begin{enumerate}
\item 态射 $\nu$ 是整的；并且
\item 对任意分解 $f = \pi \circ g$（其中 $\pi : Z \to X$ 是整的），存在交换图
$$
\xymatrix{
Y \ar[d]_{f'} \ar[r]_g & Z \ar[d]^\pi \\
X' \ar[ru]^h \ar[r]^\nu & X
}
$$
以及唯一的态射 $h : X' \to Z$。
\end{enumerate}
此外，态射 $f' : Y \to X'$ 是支配的，并且在 (2) 中态射 $h : X' \to Z$ 是 $Z$ 在 $Y$ 中的正规化。
\end{lemma}
\begin{proof}
令 $\mathcal{O}' \subset f_*\mathcal{O}_Y$ 为定义 \ref{morphisms-definition-normalization-X-in-Y} 中的
$\mathcal{O}_X$ 的整闭包。
按构造，态射 $\nu$ 是整的，这证明了 (1)。
假设给定分解 $f = \pi \circ g$，其中 $\pi : Z \to X$ 如 (2) 中是整的。
由定义 \ref{morphisms-definition-integral}，$\pi$ 是仿射的，因此 $Z$ 是某个准凝聚的
$\mathcal{O}_X$-代数层 $\mathcal{B}$ 的相对谱。
态射 $g : Y \to Z$ 对应于一个 $\mathcal{O}_X$-代数映射
$\chi : \mathcal{B} \to f_*\mathcal{O}_Y$。由于 $\mathcal{B}(U)$ 都整于 $\mathcal{O}_X(U)$，其中 $U \subset X$ 为任意仿射开集
（由定义 \ref{morphisms-definition-integral}），
由引理 \ref{morphisms-lemma-integral-closure} 可知
$\chi(\mathcal{B}) \subset \mathcal{O}'$。
由相对谱的函子性及引理 \ref{morphisms-lemma-affine-equivalence-algebras}，这给出唯一的态射
$h : X' \to Z$。交换图的验证略去。

\medskip\noindent
(1) 与 (2) 显然刻画分解 $f = \nu \circ f'$，因为它将该分解刻画为范畴中的始对象。

\medskip\noindent
由 (2) 的泛性质可见，$f'$ 不经过 $X'$ 的真闭子概形。因此 $f'$ 的概形论像是 $X'$。
由于 $f'$ 是拟紧的（由《概形》引理
\ref{schemes-lemma-quasi-compact-permanence} 以及 $\nu$ 作为仿射态射而分离这一事实），
我们看到 $f'(Y)$ 在 $X'$ 中稠密。因此 $f'$ 是支配的。

\medskip\noindent
注意，由《概形》引理 \ref{schemes-lemma-compose-after-separated} 与
\ref{schemes-lemma-quasi-compact-permanence}，$g$ 是拟紧且拟分离的。因此引理的最后一个断言是有意义的。
由引理 \ref{morphisms-lemma-finite-permanence}，(2) 中的态射 $h$ 是整的。
给定分解 $g = \pi' \circ g'$，其中 $\pi' : Z' \to Z$ 是整的，我们得到分解
$f = (\pi \circ \pi') \circ g'$，并得到态射 $h' : X' \to Z'$。
由唯一性，$\pi' \circ h' = h$。因此，刻画 (1)、(2) 适用于态射
$h : X' \to Z$，从而得到引理的最后一个断言。
\end{proof}

\begin{lemma}
\label{morphisms-lemma-functoriality-normalization}
令
$$
\xymatrix{
Y_2 \ar[d]_{f_2} \ar[r] & Y_1 \ar[d]^{f_1} \\
X_2 \ar[r] & X_1
}
$$
为概形态射的交换图。
假设 $f_1$、$f_2$ 拟紧且拟分离。
令 $f_i = \nu_i \circ f_i'$，$i = 1, 2$，为典范分解，其中
$\nu_i : X_i' \to X_i$ 是 $X_i$ 在 $Y_i$ 中的正规化。则存在唯一的
箭头 $X'_2 \to X'_1$，使其适合于一个交换图
$$
\xymatrix{
Y_2 \ar[d]_{f_2'} \ar[r] & Y_1 \ar[d]^{f_1'} \\
X_2' \ar[d]_{\nu_2} \ar[r] & X_1' \ar[d]^{\nu_1} \\
X_2 \ar[r] & X_1
}
$$
\end{lemma}

\begin{proof}
由引理 \ref{morphisms-lemma-characterize-normalization} (1) 与
\ref{morphisms-lemma-base-change-finite}，基变换 $X_2 \times_{X_1} X'_1 \to X_2$
是整的。注意 $f_2$ 经过这个态射分解。
因此，由引理 \ref{morphisms-lemma-characterize-normalization} (2)，我们得到唯一的态射
$X'_2 \to X_2 \times_{X_1} X'_1$。
这给出了适合交换图的箭头 $X'_2 \to X'_1$，唯一性也随之成立。
\end{proof}

\begin{lemma}
\label{morphisms-lemma-normalization-localization}
设 $f : Y \to X$ 为概形的拟紧且拟分离态射。
令 $U \subset X$ 为开子概形，并置 $V = f^{-1}(U)$。
则 $U$ 在 $V$ 中的正规化是 $U$ 在 $X$ 中的 $Y$ 的正规化的逆像。
\end{lemma}

\begin{proof}
由构造立即可见。
\end{proof}

\begin{lemma}
\label{morphisms-lemma-normalization-is-normalization}
设 $f : Y \to X$ 为概形的拟紧且拟分离态射。
令 $X'$ 为 $X$ 在 $Y$ 中的正规化。则 $X'$ 在 $Y$ 中的正规化就是 $X'$。
\end{lemma}

\begin{proof}
若 $Y \to X'' \to X'$ 是 $X'$ 在 $Y$ 中的正规化，则可将引理
\ref{morphisms-lemma-characterize-normalization} 应用于复合 $X'' \to X$，从而得到典范态射
$h : X' \to X''$（在 $X$ 上）。我们略去验证态射 $h$ 与 $X'' \to X'$ 互为逆态射（使用引理中分解的唯一性）。
\end{proof}

\begin{lemma}
\label{morphisms-lemma-normalization-in-reduced}
设 $f : Y \to X$ 为概形的拟紧且拟分离态射。
设 $X' \to X$ 为 $X$ 在 $Y$ 中的正规化。若 $Y$ 约化，则 $X'$ 也约化。
\end{lemma}

\begin{proof}
这是约化环的子环仍为约化环这一事实的直接结果。
细节略。
\end{proof}
\begin{lemma}
\label{morphisms-lemma-normalization-generic}
设 $f : Y \to X$ 为概形的拟紧且拟分离态射。
设 $X' \to X$ 为 $X$ 在 $Y$ 中的正规化。$X'$ 的每个不可约分支的泛点都是 $Y$ 的某个不可约分支的泛点的像。
\end{lemma}

\begin{proof}
由引理 \ref{morphisms-lemma-normalization-localization}，不妨设 $X = \Spec(A)$ 是仿射的。
选取有限仿射开覆盖 $Y = \bigcup \Spec(B_i)$。
于是 $X' = \Spec(A')$，态射 $\Spec(B_i) \to Y \to X'$ 共同给出单射的 $A$-代数映射
$A' \to \prod B_i$。
因此引理由《代数》引理 \ref{algebra-lemma-injective-minimal-primes-in-image} 得出。
\end{proof}

\begin{lemma}
\label{morphisms-lemma-normalization-in-disjoint-union}
设 $f : Y \to X$ 为概形的拟紧且拟分离态射。
假设 $Y = Y_1 \amalg Y_2$ 是两个概形的不交并。
记 $f_i = f|_{Y_i}$。令 $X_i'$ 为 $X$ 在 $Y_i$ 中的正规化。
则 $X_1' \amalg X_2'$ 是 $X$ 在 $Y$ 中的正规化。
\end{lemma}

\begin{proof}
用整闭包的语言，这对应于以下事实：
令 $A \to B$ 为环映射。假设 $B = B_1 \times B_2$。
令 $A_i'$ 为 $A$ 在 $B_i$ 中的整闭包。则
$A_1' \times A_2'$ 是 $A$ 在 $B$ 中的整闭包。
之所以成立，是因为元素 $(1, 0)$ 与 $(0, 1)$ 属于 $B$，是幂等元，因而整于 $A$。
所以整闭包 $A'$ 是 $A$ 在 $B$ 中的一个乘积；不难看出，其因子正是上面所述的整闭包 $A'_i$（细节略）。
\end{proof}

\begin{lemma}
\label{morphisms-lemma-normalization-in-universally-closed}
设 $f : X \to S$ 为概形的拟紧、拟分离且泛闭态射。
则 $f_*\mathcal{O}_X$ 整于 $\mathcal{O}_S$。换言之，$S$ 在 $X$ 中的正规化等于构造中的引理
\ref{constructions-lemma-canonical-morphism} 给出的分解
$$
X \longrightarrow \underline{\Spec}_S(f_*\mathcal{O}_X)
\longrightarrow S
$$
。
\end{lemma}

\begin{proof}
问题在 $S$ 上是局部的，因此不妨设 $S = \Spec(R)$ 是仿射的。
令 $h \in \Gamma(X, \mathcal{O}_X)$。我们需要证明 $h$ 满足一个系数在 $R$ 中的首一方程。
将 $h$ 看作如下交换图中的态射：
$$
\xymatrix{
X \ar[rr]_h \ar[rd]_f & & \mathbf{A}^1_S \ar[ld] \\
& S &
}
$$
令 $Z \subset \mathbf{A}^1_S$ 为 $h$ 的概形论像，见定义 \ref{morphisms-definition-scheme-theoretic-image}。
态射 $h$ 是拟紧的，因为 $f$ 拟紧且 $\mathbf{A}^1_S \to S$ 分离，见
《概形》引理 \ref{schemes-lemma-quasi-compact-permanence}。
由引理 \ref{morphisms-lemma-quasi-compact-scheme-theoretic-image}，态射 $X \to Z$ 是支配的。
由引理 \ref{morphisms-lemma-image-proper-scheme-closed}，态射 $X \to Z$ 是闭的。因此（按集合论意义）$h(X) = Z$。
于是可使用引理 \ref{morphisms-lemma-image-proper-is-proper}，断定 $Z \to S$ 是泛闭的（甚至是固有的）。
由于 $Z \subset \mathbf{A}^1_S$，可见 $Z \to S$ 是仿射且固有的，因而由引理
\ref{morphisms-lemma-integral-universally-closed} 是整的。
写 $\mathbf{A}^1_S = \Spec(R[T])$，我们得到理想 $I \subset R[T]$（它是 $Z$ 的理想）
包含一个首一多项式 $P(T) \in R[T]$。
于是 $P(h) = 0$，证毕。
\end{proof}

\begin{lemma}
\label{morphisms-lemma-normalization-in-integral}
设 $f : Y \to X$ 为整态射。
则 $X$ 在 $Y$ 中的正规化等于 $Y$。
\end{lemma}

\begin{proof}
由引理 \ref{morphisms-lemma-integral-universally-closed}，这是引理
\ref{morphisms-lemma-normalization-in-universally-closed} 的特殊情形。
\end{proof}
\begin{lemma}
\label{morphisms-lemma-normal-normalization}
设 $f : Y \to X$ 为概形的拟紧且拟分离态射。令 $X'$ 为 $X$ 在 $Y$ 中的正规化。假设
\begin{enumerate}
\item $Y$ 是正规概形；
\item $Y$ 的拟紧开集具有有限多个不可约分支。
\end{enumerate}
则 $X'$ 是整正规概形的不交并。此外，态射 $Y \to X'$ 是支配的，并且在不可约分支之间诱导双射。
\end{lemma}

\begin{proof}
令 $U \subset X$ 为仿射开集。考虑 $U'$，它是 $U$ 在 $X'$ 中的逆像。
置 $V = f^{-1}(U)$。由引理 \ref{morphisms-lemma-normalization-localization}，
$V \to U' \to U$ 是 $U$ 在 $V$ 中的正规化。设 $U = \Spec(A)$。
则 $V$ 拟紧，因而由假设它有有限多个不可约分支。所以
$V = \coprod_{i = 1, \ldots n} V_i$ 是正规整概形的有限不交并，依据
《性质》引理 \ref{properties-lemma-normal-locally-finite-nr-irreducibles}。
由引理 \ref{morphisms-lemma-normalization-in-disjoint-union}，可见
$U' = \coprod_{i = 1, \ldots, n} U_i'$，其中 $U'_i$ 是 $U$ 在 $V_i$ 中的正规化。
由《性质》引理 \ref{properties-lemma-normal-integral-sections}，
$B_i = \Gamma(V_i, \mathcal{O}_{V_i})$ 是正规整环。
注意 $U_i' = \Spec(A_i')$，其中 $A_i' \subset B_i$
是 $A$ 在 $B_i$ 中的整闭包，见引理 \ref{morphisms-lemma-integral-closure}。由
《代数》引理 \ref{algebra-lemma-integral-closure-in-normal}，可见 $A_i' \subset B_i$ 是正规整环。
因此 $U' = \coprod U_i'$ 是正规整概形的有限不交并，因而是正规概形。

\medskip\noindent
由于 $X'$ 有由概形 $U'$ 构成的开覆盖，由《性质》引理
\ref{properties-lemma-locally-normal} 可知 $X'$ 是正规的。
另一方面，每个 $U'$ 都是不可约概形的有限不交并，因此 $X'$ 的每个拟紧开集都有有限多个不可约分支
（拓扑论证略）。所以由《性质》引理
\ref{properties-lemma-normal-locally-finite-nr-irreducibles}，$X'$ 是正规整概形的不交并。
从上面对 $X'$ 的描述显然可知，$Y \to X'$ 是支配的，并且在不可约分支之间诱导双射 $V \to U'$，对每个仿射开集 $U \subset X$。
态射 $Y \to X'$ 的不可约分支双射由此及拓扑论证得到（略）。
\end{proof}
\begin{lemma}
\label{morphisms-lemma-nagata-normalization-finite-general}
设 $f : X \to S$ 为态射。假设
\begin{enumerate}
\item $S$ 是 Nagata 概形；
\item $f$ 拟紧且拟分离；
\item $X$ 的拟紧开集具有有限多个不可约分支；
\item 若 $x \in X$ 是不可约分支的泛点，则域扩张
$\kappa(x)/\kappa(f(x))$ 是有限生成的；并且
\item $X$ 是约化的。
\end{enumerate}
则正规化 $\nu : S' \to S$（即 $S$ 在 $X$ 中的正规化）是有限的。
\end{lemma}

\begin{proof}
由假设 (1)，立即可将情形约化为 $S = \Spec(R)$，其中 $R$ 是 Nagata 环。
我们需要证明整闭包 $A$（即 $R$ 在 $\Gamma(X, \mathcal{O}_X)$ 中的整闭包）是有限的 $R$-模。
由假设 (2)，$f$ 拟紧，因此可写成 $X = \bigcup_{i = 1, \ldots, n} U_i$，其中每个 $U_i$ 都是仿射的。
设 $U_i = \Spec(B_i)$。由假设 (5)，每个 $B_i$ 都是约化的，并且由假设 (3) 以及
《代数》引理 \ref{algebra-lemma-irreducible}，它有有限多个极小素理想
$\mathfrak q_{i1}, \ldots, \mathfrak q_{im_i}$。
我们有
$$
\Gamma(X, \mathcal{O}_X) \subset B_1 \times \ldots \times B_n
\subset
\prod\nolimits_{i = 1, \ldots, n}
\prod\nolimits_{j = 1, \ldots, m_i} (B_i)_{\mathfrak q_{ij}}
$$
第二个包含关系来自《代数》引理 \ref{algebra-lemma-reduced-ring-sub-product-fields}。
由《代数》引理 \ref{algebra-lemma-minimal-prime-reduced-ring}，有
$\kappa(\mathfrak q_{ij}) = (B_i)_{\mathfrak q_{ij}}$。
因此，整闭包 $A$（即 $R$ 在 $\Gamma(X, \mathcal{O}_X)$ 中的整闭包）
包含于整闭包 $A_{ij}$（即 $R$ 在 $\kappa(\mathfrak q_{ij})$ 中的整闭包）的乘积中。
由于 $R$ 是 Noether 环，只需证明 $A_{ij}$ 是有限的 $R$-模，对每个 $i, j$ 均如此。
令 $\mathfrak p_{ij} \subset R$ 为 $\mathfrak q_{ij}$ 的像。
由假设 (4)，$\kappa(\mathfrak q_{ij})/\kappa(\mathfrak p_{ij})$
是有限生成的域扩张，因此 $R \to \kappa(\mathfrak q_{ij})$ 本质上是有限型的。
于是由《代数》引理
\ref{algebra-lemma-nagata-in-reduced-finite-type-finite-integral-closure}，$R \to A_{ij}$ 是有限的。
\end{proof}

\begin{lemma}
\label{morphisms-lemma-nagata-normalization-finite}
设 $f : X \to S$ 为态射。假设
\begin{enumerate}
\item $S$ 是 Nagata 概形；
\item $f$ 是有限型的；
\item $X$ 是约化的。
\end{enumerate}
则正规化 $\nu : S' \to S$（即 $S$ 在 $X$ 中的正规化）是有限的。
\end{lemma}

\begin{proof}
这是引理 \ref{morphisms-lemma-nagata-normalization-finite-general} 的特殊情形。
具体而言，(2) 成立，因为按定义有限型态射 $f$ 是拟紧的，并且由引理
\ref{morphisms-lemma-finite-type-Noetherian-quasi-separated} 是拟分离的。
条件 (3) 成立，因为 $X$ 局部 Noether，见引理 \ref{morphisms-lemma-finite-type-noetherian}。
最后，条件 (4) 成立，因为有限型态射诱导有限生成的剩余域扩张。
\end{proof}
\begin{lemma}
\label{morphisms-lemma-relative-normalization-normal-codim-1}
设 $f : Y \to X$ 为概形的有限型态射，其中 $Y$ 约化且 $X$ 为 Nagata 概形。令 $X'$ 为 $X$ 在 $Y$ 中的正规化。
令 $x' \in X'$ 为满足下列条件的点：
\begin{enumerate}
\item $\dim(\mathcal{O}_{X', x'}) = 1$，并且
\item $Y \to X'$ 在 $x'$ 上的纤维为空。
\end{enumerate}
则 $\mathcal{O}_{X', x'}$ 是离散赋值环。
\end{lemma}

\begin{proof}
可以将 $X$ 替换为包含 $x'$ 的像的仿射邻域。因此不妨设 $X = \Spec(A)$，其中 $A$ 是 Nagata 环。
由引理 \ref{morphisms-lemma-nagata-normalization-finite}，态射 $X' \to X$ 是有限的。因此可写成
$X' = \Spec(A')$，其中 $A$-代数 $A'$ 是有限的。
由引理 \ref{morphisms-lemma-normalization-is-normalization}，将 $X$ 替换为 $X'$ 后，
我们约化到下一段所述的情形。

\medskip\noindent
情形 $X = X' = \Spec(A)$，其中 $A$ 为 Noether 环。
令 $\mathfrak p \subset A$ 为对应于点 $x'$ 的素理想。
选取 $g \in \mathfrak p$，使其不包含于 $A$ 的任何极小素理想中（使用素理想避免以及 $A$ 只有有限多个极小素理想这一事实，见
《代数》引理 \ref{algebra-lemma-silly} 与
\ref{algebra-lemma-Noetherian-irreducible-components}）。
置 $Z = f^{-1}V(g) \subset Y$；它是 $Y$ 的闭子概形。
由于 $f(Z)$ 按 $g$ 的选择不包含任何泛点，且不包含 $x'$，因为 $x'$ 不在 $f$ 的像中。
$f(Z)$ 的闭包是 $f(Z)$ 中点的特化点集合，这是引理 \ref{morphisms-lemma-reach-points-scheme-theoretic-image} 的结论。
因此 $f(Z)$ 的闭包不包含 $x'$，因为条件 $\dim(\mathcal{O}_{X', x'}) = 1$ 蕴含只有 $X = X'$ 的泛点特化到 $x'$。
换言之，将 $X$ 替换为包含 $x'$ 的一个仿射开邻域后，不妨设 $f^{-1}V(g) = \emptyset$。
于是 $g$ 映为 $Y$ 上的可逆整体函数，并得到分解
$$
A \to A_g \to \Gamma(Y, \mathcal{O}_Y)
$$
由于 $X = X'$，这意味着 $A$ 等于整闭包（即 $A$ 在 $A_g$ 中的整闭包）。由
《代数》引理 \ref{algebra-lemma-integral-closure-localize}，我们断定
$A_\mathfrak p$ 是 $A_\mathfrak p$ 的整闭包（位于 $A_\mathfrak p[1/g]$ 中）。
由 $g$ 的选择、$\dim(A_\mathfrak p) = 1$ 且 $A$ 约化，可知 $A_\mathfrak p[1/g]$
是域的有限乘积（包含于 $\mathfrak p$ 的极小素理想的剩余域之乘积）。因此由
《代数》引理 \ref{algebra-lemma-characterize-reduced-ring-normal}，$A_\mathfrak p$ 是正规的，证明完成。细节略。
\end{proof}

\section{正规化}
\label{morphisms-section-normalization}
\noindent
接下来讨论概形 $X$ 的正规化。只有当 $X$ 的不可约分支局部有限时，我们才定义并构造它。
令 $X$ 为一个概形，使每个拟紧开集都有有限多个不可约分支。令
$X^{(0)} \subset X$ 为 $X$ 的不可约分支的泛点集合。令
\begin{equation}
\label{morphisms-equation-generic-points}
f :
Y = \coprod\nolimits_{\eta \in X^{(0)}} \Spec(\kappa(\eta))
\longrightarrow
X
\end{equation}
为利用《概形》第 \ref{schemes-section-points} 节中的典范映射，将泛点包含到 $X$ 中的态射。
注意，由假设该态射是拟紧的，并且由于 $Y$ 是分离的而拟分离（见《概形》第
\ref{schemes-section-separation-axioms} 节）。

\begin{definition}
\label{morphisms-definition-normalization}
设 $X$ 为一个概形，使每个拟紧开集都有有限多个不可约分支。我们将 $X$ 的\emph{正规化}定义为态射
$$
\nu : X^\nu \longrightarrow X
$$
它是 $X$ 在上面构造的态射 $f : Y \to X$（\ref{morphisms-equation-generic-points}）中的正规化。
\end{definition}

\noindent
任何局部 Noether 概形都有局部有限的不可约分支集合，因此该定义适用。
通常只对约化概形定义正规化。按上述定义，$X$ 的正规化与约化 $X_{red}$（它来自 $X$）的正规化相同。

\begin{lemma}
\label{morphisms-lemma-normalization-reduced}
设 $X$ 为一个概形，使每个拟紧开集都有有限多个不可约分支。正规化态射
$\nu$ 经过约化 $X_{red}$ 分解，并且 $X^\nu \to X_{red}$ 是 $X_{red}$ 的正规化。
\end{lemma}

\begin{proof}
令 $f : Y \to X$ 为态射（\ref{morphisms-equation-generic-points}）。由《概形》引理
\ref{schemes-lemma-map-into-reduction}，我们得到分解
$Y \to X_{red} \to X$，它分解 $f$。
由引理 \ref{morphisms-lemma-characterize-normalization}，得到典范态射 $X^\nu \to X_{red}$，
并且 $X^\nu$ 是 $X_{red}$ 在 $Y$ 中的正规化。
引理得证，因为 $Y \to X_{red}$ 与为 $X_{red}$ 构造的态射（\ref{morphisms-equation-generic-points}）相同。
\end{proof}

\noindent
若 $X$ 是约化的，则 $X$ 的正规化与 $\mathcal{O}_X$ 在亚纯函数层 $\mathcal{K}_X$ 中的整闭包的相对谱相同
（见《除子》第 \ref{divisors-section-meromorphic-functions} 节）。
也就是说，在此情形 $\mathcal{K}_X = f_*\mathcal{O}_Y$，见
《除子》引理 \ref{divisors-lemma-reduced-finite-irreducible} 及其证明。下面对此作明确描述。

\begin{lemma}
\label{morphisms-lemma-description-normalization}
设 $X$ 为约化概形，使每个拟紧开集都有有限多个不可约分支。令 $\Spec(A) = U \subset X$
为仿射开集。则
\begin{enumerate}
\item $A$ 有有限多个极小素理想
$\mathfrak q_1, \ldots, \mathfrak q_t$；
\item 总分式环 $Q(A)$（属于 $A$）是
$Q(A/\mathfrak q_1) \times \ldots \times Q(A/\mathfrak q_t)$；
\item 整闭包 $A'$（即 $A$ 在 $Q(A)$ 中的整闭包）是各整环 $A/\mathfrak q_i$
在域 $Q(A/\mathfrak q_i)$ 中的整闭包的乘积；并且
\item $\nu^{-1}(U)$ 可识别为 $A'$ 的谱，其中
$\nu : X^\nu \to X$ 是正规化态射。
\end{enumerate}
\end{lemma}

\begin{proof}
极小素理想对应于不可约分支（《代数》引理 \ref{algebra-lemma-irreducible}），
因此由假设得到 (1)。有
$(0) = \mathfrak q_1 \cap \ldots \cap \mathfrak q_t$，因为 $A$ 约化
（《代数》引理 \ref{algebra-lemma-Zariski-topology}）。
于是由《代数》引理 \ref{algebra-lemma-total-ring-fractions-no-embedded-points} 与
\ref{algebra-lemma-minimal-prime-reduced-ring}，有
$Q(A) = \prod A_{\mathfrak q_i} = \prod \kappa(\mathfrak q_i)$。
这证明了 (2)。(3) 由《代数》引理
\ref{algebra-lemma-characterize-reduced-ring-normal} 或引理 \ref{morphisms-lemma-normalization-in-disjoint-union} 得出。
(4) 成立是因为显然 $f^{-1}(U) \to U$ 是态射
$$
\Spec\left(\prod \kappa(\mathfrak q_i)\right)
\longrightarrow
\Spec(A)
$$
其中 $f : Y \to X$ 是态射（\ref{morphisms-equation-generic-points}）。
\end{proof}
\begin{lemma}
\label{morphisms-lemma-stalk-normalization}
设 $X$ 为一个概形，使每个拟紧开集都有有限多个不可约分支。设 $\nu : X^\nu \to X$
为 $X$ 的正规化。令 $x \in X$。则下列对象作为 $\mathcal{O}_{X, x}$-代数典范同构：
\begin{enumerate}
\item 茎 $(\nu_*\mathcal{O}_{X^\nu})_x$；
\item $\mathcal{O}_{X, x}$ 在 $(\mathcal{O}_{X, x})_{red}$ 的总分式环中的整闭包；
\item $\mathcal{O}_{X, x}$ 在由 $\mathcal{O}_{X, x}$ 的极小素理想的剩余域组成的乘积中的整闭包
（而这些极小素理想只有有限多个）。
\end{enumerate}
\end{lemma}

\begin{proof}
将 $X$ 替换为 $x$ 的一个仿射开邻域后，不妨设 $X$ 有有限多个不可约分支，并且 $x$ 包含在每一个分支中。
于是茎 $(\nu_*\mathcal{O}_{X^\nu})_x$ 是 $A = \mathcal{O}_{X, x}$ 在积 $L$ 中的整闭包，
其中该积由 $A$ 的极小素理想的剩余域组成。
这由正规化的构造和引理 \ref{morphisms-lemma-integral-closure} 得出。
或者，可以使用引理 \ref{morphisms-lemma-description-normalization} 以及正规化与局部化可交换这一事实
（《代数》引理 \ref{algebra-lemma-integral-closure-localize}）。
由于 $A_{red}$ 有有限多个极小素理想
（因为它们恰好对应于 $X$ 中经过 $x$ 的不可约分支的泛点），
由《代数》引理 \ref{algebra-lemma-total-ring-fractions-no-embedded-points} 可知，$L$ 是 $A_{red}$ 的总分式环。
因此，我们的环也是 $A$ 在 $A_{red}$ 的总分式环中的整闭包。
\end{proof}

\begin{lemma}
\label{morphisms-lemma-normalization-normal}
设 $X$ 为一个概形，使每个拟紧开集都有有限多个不可约分支。
\begin{enumerate}
\item 正规化 $X^\nu$ 是整正规概形的不交并。
\item 态射 $\nu : X^\nu \to X$ 是整的、满射的，并在不可约分支之间诱导双射。
\item 对任意整态射 $\alpha : X' \to X$，如果对于拟紧开集 $U \subset X$，逆像 $\alpha^{-1}(U)$
有有限多个不可约分支，并且
$\alpha|_{\alpha^{-1}(U)} : \alpha^{-1}(U) \to U$ 是双有理的\footnote{这种笨拙的表述是必要的，因为我们只在源和目标都有有限多个不可约分支时定义了态射为双有理的含义。若
$X'_{red} \to X_{red}$ 满足该条件，则已经足够。}，
则存在分解 $X^\nu \to X' \to X$，并且 $X^\nu \to X'$ 是 $X'$ 的正规化。
\item 对任意态射 $Z \to X$，其中 $Z$ 是正规概形，并且 $Z$ 的每个不可约分支都支配 $X$ 的一个不可约分支，
存在唯一分解 $Z \to X^\nu \to X$。
\end{enumerate}
\end{lemma}
\begin{proof}
令 $f : Y \to X$ 如（\ref{morphisms-equation-generic-points}）中所述。
$X^\nu$ 是整正规概形的不交并，因为 $Y$ 是正规的，且 $Y$ 的每个仿射开集都有有限多个不可约分支，见
引理 \ref{morphisms-lemma-normal-normalization}。这证明了 (1)。
或者，也可由引理 \ref{morphisms-lemma-normalization-reduced} 与
\ref{morphisms-lemma-description-normalization} 推出 (1)。

\medskip\noindent
由引理 \ref{morphisms-lemma-characterize-normalization}，态射 $\nu$ 是整的。
由引理 \ref{morphisms-lemma-normal-normalization}，态射 $Y \to X^\nu$ 在不可约分支之间诱导双射；
由 $Y$ 的构造，这意味着 $X^\nu \to X$ 在不可约分支之间诱导双射。
按构造 $f : Y \to X$ 是支配的，因此 $\nu$ 也是支配的。由于整态射是闭的
（引理 \ref{morphisms-lemma-integral-universally-closed}），这意味着 $\nu$ 是满射。这证明了 (2)。

\medskip\noindent
假设 $\alpha : X' \to X$ 如 (3) 中所述。显然 $X'$ 满足定义正规化所需的假设。
令 $f' : Y' \to X'$ 为从 $X'$ 出发按（\ref{morphisms-equation-generic-points}）构造的态射。
由于 $\alpha$ 局部双有理，显然 $Y' = Y$ 且 $f = \alpha \circ f'$。
因此分解 $X^\nu \to X' \to X$ 存在，并且由引理
\ref{morphisms-lemma-characterize-normalization}，$X^\nu \to X'$ 是 $X'$ 的正规化。这证明了 (3)。

\medskip\noindent
令 $g : Z \to X$ 为态射，其定义域是正规概形，并且每个不可约分支都支配 $X$ 的一个不可约分支。
由引理 \ref{morphisms-lemma-normalization-reduced}，有 $X^\nu = X_{red}^\nu$，并且由
《概形》引理 \ref{schemes-lemma-map-into-reduction}，$Z \to X$ 经过 $X_{red}$ 分解。
因此可将 $X$ 替换为 $X_{red}$，并假设 $X$ 约化。此外，由于分解唯一，只需在 $Z$ 上局部构造它。
令 $W \subset Z$ 与 $U \subset X$ 为满足 $g(W) \subset U$ 的仿射开集。
写 $U = \Spec(A)$、$W = \Spec(B)$，其中 $g|_W$ 由 $\varphi : A \to B$ 给出。
我们将自由使用引理 \ref{morphisms-lemma-description-normalization} 的结果。
令 $\mathfrak p_1, \ldots, \mathfrak p_t$ 为 $A$ 的极小素理想。
由于 $Z$ 正规，$B$ 是正规环，特别是约化的。此外，由假设，对任意极小素理想
$\mathfrak q \subset B$，$\varphi^{-1}(\mathfrak q)$ 是 $A$ 的极小素理想。
因此，若 $x \in A$ 是非零因子，即 $x \not \in \bigcup \mathfrak p_i$，则 $\varphi(x)$ 在 $B$ 中是非零因子。
于是得到典范环映射 $Q(A) \to Q(B)$。
由于 $B$ 是正规的，它等于自己在 $Q(B)$ 中的整闭包（见
《代数》引理 \ref{algebra-lemma-normal-ring-integrally-closed}）。
因此，整闭包 $A' \subset Q(A)$（即 $A$ 的整闭包）映入 $B$，通过典范映射 $Q(A) \to Q(B)$。
由于 $\nu^{-1}(U) = \Spec(A')$，这给出典范分解
$W \to \nu^{-1}(U) \to U$，它分解 $\nu|_W$。
唯一性的验证略去。
\end{proof}
\begin{lemma}
\label{morphisms-lemma-normalization-in-terms-of-components}
设 $X$ 为一个概形，使每个拟紧开集都有有限多个不可约分支。令 $Z_i \subset X$，$i \in I$，
为赋予约化诱导结构的 $X$ 的不可约分支。令 $Z_i^\nu \to Z_i$ 为正规化。
则 $\coprod_{i \in I} Z_i^\nu \to X$ 是 $X$ 的正规化。
\end{lemma}

\begin{proof}
由引理 \ref{morphisms-lemma-normalization-reduced}，不妨设 $X$ 约化。
于是引理可由引理 \ref{morphisms-lemma-description-normalization} 中的局部描述得到，
或者由引理 \ref{morphisms-lemma-normalization-normal} 的 (3) 得到，因为
$\coprod Z_i \to X$ 是整的且局部双有理的（$X$ 约化且具有局部有限多个不可约分支）。
\end{proof}

\begin{lemma}
\label{morphisms-lemma-normalization-birational}
设 $X$ 为具有有限多个不可约分支的约化概形。
则正规化态射 $X^\nu \to X$ 是双有理的。
\end{lemma}

\begin{proof}
由引理 \ref{morphisms-lemma-normalization-normal}，正规化在不可约分支之间诱导双射。
令 $\eta \in X$ 为 $X$ 的一个不可约分支的泛点，并令 $\eta^\nu \in X^\nu$
为 $X^\nu$ 中对应不可约分支的泛点。
则 $\eta^\nu \mapsto \eta$；为了完成证明，需要证明
$\mathcal{O}_{X, \eta} \to \mathcal{O}_{X^\nu, \eta^\nu}$
是同构，见定义 \ref{morphisms-definition-birational}。
由于 $X$ 与 $X^\nu$ 都约化，可见两个局部环都等于其剩余域
（《代数》引理 \ref{algebra-lemma-minimal-prime-reduced-ring}）。
另一方面，由正规化作为
$X$ 在 $Y = \coprod \Spec(\kappa(\eta))$ 中的正规化的构造，我们有
$\kappa(\eta) \subset \kappa(\eta^\nu) \subset \kappa(\eta)$，证明完成。
\end{proof}

\begin{lemma}
\label{morphisms-lemma-finite-birational-over-normal}
整概形之间的有限（甚至整）双有理态射 $f : X \to Y$，
若 $Y$ 为正规概形，则是同构。
\end{lemma}

\begin{proof}
令 $V \subset Y$ 为仿射开集，并令其逆像 $U \subset X$ 也是仿射开集。
由于 $f$ 是整概形之间的双有理态射，环同态
$\mathcal{O}_Y(V) \to \mathcal{O}_X(U)$ 是整环的单射，
并且诱导分式域的同构。由于 $Y$ 正规，环 $\mathcal{O}_Y(V)$
在其分式域中整闭。由于 $f$ 是有限的（或整的），
$\mathcal{O}_X(U)$ 的每个元素都在 $\mathcal{O}_Y(V)$ 上整。
因此 $\mathcal{O}_Y(V) = \mathcal{O}_X(U)$。这就证明了 $f$ 是同构，证毕。
\end{proof}

\begin{lemma}
\label{morphisms-lemma-normalization-trivial}
设 $X$ 为每个局部都只有有限多个不可约分支的概形。
正规化态射 $\nu : X^\nu \to X$ 是同构，当且仅当 $X$ 正规。
\end{lemma}

\begin{proof}
若 $\nu$ 是同构，则由引理 \ref{morphisms-lemma-normalization-normal}，$X$ 正规。
反之，设 $X$ 正规。由引理 \ref{morphisms-lemma-normalization-localization} 以及
《性质》引理 \ref{properties-lemma-normal-locally-finite-nr-irreducibles}，
不妨设 $X$ 是整概形。由引理 \ref{morphisms-lemma-normalization-normal}，
态射 $\nu$ 是整的，且 $X^\nu$ 只有一个不可约分支
（所以它是整概形）。由引理 \ref{morphisms-lemma-normalization-birational} 和
\ref{morphisms-lemma-finite-birational-over-normal} 得出结论。
\end{proof}

\begin{lemma}
\label{morphisms-lemma-Japanese-normalization}
设 $X$ 为整的 Japanese 概形。
正规化 $\nu : X^\nu \to X$ 是有限态射。
\end{lemma}

\begin{proof}
这由定义（《性质》定义 \ref{properties-definition-nagata}）和
引理 \ref{morphisms-lemma-description-normalization} 得出。具体而言，在这种情况下，
该引理说明 $\nu^{-1}(\Spec(A))$ 是 $A$ 在其分式域中的整闭包的谱。
\end{proof}

\begin{lemma}
\label{morphisms-lemma-nagata-normalization}
设 $X$ 为 Nagata 概形。
正规化 $\nu : X^\nu \to X$ 是有限态射。
\end{lemma}

\begin{proof}
注意 Nagata 概形是局部诺特的，因而定义 \ref{morphisms-definition-normalization}
确实适用。该引理现在是引理
\ref{morphisms-lemma-nagata-normalization-finite-general} 的特例，
但也可以直接如下证明。
将 $X^\nu \to X$ 写成复合
$X^\nu \to X_{red} \to X$。由于 $X_{red} \to X$ 是闭浸入，
它是有限的。因此（由引理 \ref{morphisms-lemma-composition-finite}）
只需对约化的 Nagata 概形证明该引理。
令 $\Spec(A) = U \subset X$ 为仿射开集。由引理
\ref{morphisms-lemma-description-normalization}，有
$\nu^{-1}(U) = \Spec(\prod A_i')$，其中 $A_i'$ 是
$A/\mathfrak q_i$ 在其分式域中的整闭包。由于 $A$ 是 Nagata 环
（见《性质》引理 \ref{properties-lemma-locally-nagata}），
每个环扩张 $A/\mathfrak q_i \subset A'_i$ 都是有限的。
因此 $A \to \prod A'_i$ 是有限环映射，结论成立。
\end{proof}

\begin{lemma}
\label{morphisms-lemma-normalization-geom-unibranch-univ-homeo}
设 $X$ 为不可约、几何单分支概形。
正规化态射 $\nu : X^\nu \to X$ 是泛同胚。
\end{lemma}

\begin{proof}
我们需要证明 $\nu$ 是整的、泛单射且满射，见引理
\ref{morphisms-lemma-universal-homeomorphism}。由引理
\ref{morphisms-lemma-normalization-normal}，态射 $\nu$ 是整的。
令 $x \in X$，并置 $A = \mathcal{O}_{X, x}$。
由于 $X$ 不可约，$A$ 只有一个极小素理想 $\mathfrak p$，且
$A_{red} = A/\mathfrak p$。由引理 \ref{morphisms-lemma-stalk-normalization}，
茎 $A' = (\nu_*\mathcal{O}_{X^\nu})_x$ 是 $A$ 在 $A_{red}$ 的分式域中的整闭包。
由《代数》定义 \ref{more-algebra-definition-unibranch}，可知 $A'$
只有一个素理想 $\mathfrak m'$，它位于 $\mathfrak m_x \subset A$ 之上，并且
$\kappa(\mathfrak m')/\kappa(x)$ 是纯不可分扩张。因此由引理
\ref{morphisms-lemma-universally-injective}，$\nu$ 是双射（因而满射）且泛单射。
\end{proof}





\section{Weak normalization}
\label{morphisms-section-weak-normalization}

\noindent
我们只在概形局部具有有限多个不可约分支时定义其弱正规化；这类似于正规化的情形。

\begin{lemma}
\label{morphisms-lemma-relative-weak-normalization-algebra}
设环映射 $A \to B$ 诱导谱的支配态射
$\Spec(B) \to \Spec(A)$。存在一个 $A$-子代数
$B' \subset B$，满足
\begin{enumerate}
\item $\Spec(B') \to \Spec(A)$ 是泛同胚；
\item 给定分解 $A \to C \to B$，其中
$\Spec(C) \to \Spec(A)$ 是泛同胚，则 $C \to B$ 的像包含于 $B'$ 中。
\end{enumerate}
\end{lemma}

\begin{proof}
以下反复使用引理 \ref{morphisms-lemma-reduction-universal-homeomorphism}，不再另行说明。
考虑交换图
$$
\xymatrix{
B \ar[r] & B_{red} \\
A \ar[u] \ar[r] & A_{red} \ar[u]
}
$$
对于 (2) 中任意一个分解 $A \to C \to B$（其所分解的映射为 $A \to B$），可见
$A_{red} \to C_{red} \to B_{red}$ 是
$A_{red} \to B_{red}$ 的满足 (2) 的分解。
因此，如果引理对 $A_{red} \to B_{red}$ 成立并产生
$A_{red}$-子代数 $B'_{red} \subset B_{red}$，那么令 $B' \subset B$
等于 $B'_{red}$ 的逆像，即可为 $A \to B$ 解决该引理。
于是问题归约到下一段所讨论的情形。

\medskip\noindent
设 $A$ 与 $B$ 都约化。在此情形，由《代数》引理
\ref{algebra-lemma-image-dense-generic-points}，有 $A \subset B$。
令 $A \to C \to B$ 为 (2) 中的一个分解。
于是可以将命题 \ref{morphisms-proposition-universal-homeomorphism}
应用于 $A \subset C$，从而每个 $C$ 的元素都包含于一个扩张
$A[c_1, \ldots, c_n] \subset C$ 中，并且对 $i = 1, \ldots, n$，有
\begin{enumerate}
\item $c_i^2, c_i^3 \in A[c_1, \ldots, c_{i - 1}]$，或者
\item 存在素数 $p$，使得
$pc_i, c_i^p \in A[c_1, \ldots, c_{i - 1}]$。
\end{enumerate}
因此，若将 $B' \subset B$ 定义为满足下述条件的元素 $b \in B$ 的集合，
则性质 (2) 成立：包含于一个扩张
$
A[b_1, \ldots, b_n] \subset B
$
中，且满足 (*)：对 $i = 1, \ldots, n$，有
\begin{enumerate}
\item $b_i^2, b_i^3 \in A[b_1, \ldots, b_{i - 1}]$，或者
\item 存在素数 $p$，使得
$pb_i, b_i^p \in A[b_1, \ldots, b_{i - 1}]$。
\end{enumerate}
只需检查两件事：(a) $B'$ 是 $A$-子代数；(b)
$\Spec(B') \to \Spec(A)$ 是泛同胚。
对 (a)，若 $n \geq 0$ 且 $b_1, \ldots, b_n \in B$ 满足 (*)，
并且 $m \geq 0$ 且 $b'_1, \ldots, b'_m \in B$ 满足 (*)，
则整数 $n + m$ 以及
$b_1, \ldots, b_n, b'_1, \ldots, b'_m \in B$ 也满足 (*)，所以 (a) 成立。
最后，(b) 由命题 \ref{morphisms-proposition-universal-homeomorphism}
以及 $B'$ 的构造得到。
\end{proof}

\begin{lemma}
\label{morphisms-lemma-relative-weak-normalization-localization}
设环映射 $A \to B$ 诱导谱的支配态射
$\Spec(B) \to \Spec(A)$。引理
\ref{morphisms-lemma-relative-weak-normalization-algebra} 中 $A$-子代数
$B' \subset B$ 的构造与局部化相容（解释见证明）。
\end{lemma}

\begin{proof}
令 $S \subset A$ 为乘法子集。于是 $S^{-1}A \to S^{-1}B$ 是一个
环映射，并且同样诱导支配态射
$\Spec(S^{-1}B) \to \Spec(S^{-1}A)$
（见引理 \ref{morphisms-lemma-base-change-dominant} 和
\ref{morphisms-lemma-generalizations-lift-flat}）。因此，引理
\ref{morphisms-lemma-relative-weak-normalization-algebra} 产生一个
$S^{-1}A$-子代数 $(S^{-1}B)' \subset S^{-1}B$。
所述结论是指 $S^{-1}B' = (S^{-1}B)'$，并且它们是
$S^{-1}A$-子代数，且属于 $S^{-1}B$。

\medskip\noindent
为说明这一点，我们使用引理
\ref{morphisms-lemma-relative-weak-normalization-algebra} 的证明中对 $B'$
和 $(S^{-1}B)'$ 的构造。第一步中可见，$B'$ 是
$A_{red}$-子代数 $B'_{red} \subset B_{red}$ 的逆像，
该子代数由环映射 $A_{red} \to B_{red}$ 构造；$(S^{-1}B)'$ 同理。
注意 $S^{-1}B_{red} = (S^{-1}B)_{red}$，问题归约到下一段所讨论的情形。

\medskip\noindent
若 $A$ 与 $B$ 约化，则我们把 $B'$ 构造成下列子代数的并：
$A[b_1, \ldots, b_n]$，其中对 $i = 1, \ldots, n$ 有
\begin{enumerate}
\item $b_i^2, b_i^3 \in A[b_1, \ldots, b_{i - 1}]$，或者
\item 存在素数 $p$，使得
$pb_i, b_i^p \in A[b_1, \ldots, b_{i - 1}]$。
\end{enumerate}
对 $(S^{-1}B)' \subset S^{-1}B$ 也同理。因此显然，
$B' \to B \to S^{-1}B$ 的像包含于 $(S^{-1}B)'$ 中。
要证明相应的映射 $S^{-1}B' \to (S^{-1}B)'$ 是满射，
可使用引理 \ref{morphisms-lemma-special-elements-and-localization} 逐次清除分母；
细节略去。
\end{proof}

\begin{lemma}
\label{morphisms-lemma-relative-seminormalization-algebra}
设环映射 $A \to B$ 诱导谱的支配态射
$\Spec(B) \to \Spec(A)$。存在一个 $A$-子代数
$B' \subset B$，满足
\begin{enumerate}
\item $\Spec(B') \to \Spec(A)$ 是诱导剩余域同构的泛同胚；
\item 给定分解 $A \to C \to B$，其中 $\Spec(C) \to \Spec(A)$
是诱导剩余域同构的泛同胚，则 $C \to B$ 的像包含于 $B'$ 中。
\end{enumerate}
\end{lemma}

\begin{proof}
该证明与引理 \ref{morphisms-lemma-relative-weak-normalization-algebra} 的证明完全相同，
只是使用命题 \ref{morphisms-proposition-universal-homeomorphism-equal-residue-fields}，
而不是命题 \ref{morphisms-proposition-universal-homeomorphism}。
\end{proof}

\begin{lemma}
\label{morphisms-lemma-relative-seminormalization-localization}
设环映射 $A \to B$ 诱导谱的支配态射
$\Spec(B) \to \Spec(A)$。引理
\ref{morphisms-lemma-relative-seminormalization-algebra} 中 $A$-子代数
$B' \subset B$ 的构造与局部化相容（解释见证明）。
\end{lemma}

\begin{proof}
证明与引理 \ref{morphisms-lemma-relative-weak-normalization-localization} 的证明相同。
\end{proof}

\begin{lemma}
\label{morphisms-lemma-relative-seminormalization}
设 $f : Y \to X$ 为概形的拟紧、拟分离且支配的态射。
\begin{enumerate}
\item 由分解 $Y \to X' \to X$ 构成的范畴，其中 $X' \to X$ 为泛同胚，具有始对象
$Y \to X^{Y/wn} \to X$。
\item 由分解 $Y \to X' \to X$ 构成的范畴，其中 $X' \to X$ 为诱导剩余域同构的泛同胚，具有始对象
$Y \to X^{Y/sn} \to X$。
\end{enumerate}
此外，分解 $Y \to X^{Y/wn} \to X$ 和
$Y \to X^{Y/sn} \to X$ 的构造与到 $X$ 的开子概形的基变换相容。
\end{lemma}

\begin{proof}
我们证明 (1)，略去 (2) 的证明；最后的断言也将由分解的构造推出。
以下使用引理 \ref{morphisms-lemma-universal-homeomorphism}，不再另行说明。
首先，令 $Y \to X^{Y/n} \to X$ 为 $X$ 在 $Y$ 中的正规化，见定义
\ref{morphisms-definition-normalization-X-in-Y}。对于 (1) 中的分解
$Y \to X' \to X$，由引理 \ref{morphisms-lemma-characterize-normalization}
可得到一个与给定态射相容的唯一态射 $X^{Y/n} \to X'$。
因此，只需证明将 $f$ 替换为 $X^{Y/n} \to X$ 后的引理。
这就归约到下一段研究的情形。

\medskip\noindent
假设 $f$ 是整的（证明的其余部分在 $f$ 为仿射时也更一般地成立）。
令 $U = \Spec(A)$ 为 $X$ 的仿射开集，并令
$V = f^{-1}(U) = \Spec(B)$ 为其在 $Y$ 中的逆像。
则 $A \to B$ 是一个在谱上诱导支配态射的环映射。
由引理 \ref{morphisms-lemma-relative-weak-normalization-algebra}，得到
$A$-子代数 $B' \subset B$，使得置
$U^{V/wn} = \Spec(B')$ 后，分解 $V \to U^{V/wn} \to U$
是由 $V \to U' \to U$（其中 $U' \to U$ 为泛同胚）构成的范畴中的始对象。

\medskip\noindent
若 $U_1 \subset U_2 \subset X$ 是仿射开集，令
$V_i = f^{-1}(U_i)$，由此得到一个典范态射
$$
\rho_{U_1}^{U_2} : U_1^{V_1/wn} \to U_1 \times_{U_2} U_2^{V_2/wn}
$$
它在 $U_1$ 上，并由 $U_1^{V_1/wn}$ 的泛性质给出。
这些态射满足自然的函子性；其表述和证明留给读者。
此外，态射 $\rho_{U_1}^{U_2}$ 是同构；当
$U_1 \subset U_2$ 是标准开集时，这由引理
\ref{morphisms-lemma-relative-weak-normalization-localization} 得出，
一般情形则可由这些映射的函子性以及《概形》引理
\ref{schemes-lemma-standard-open-two-affines} 归约到标准开集情形（细节略去）。
因此，由相对粘合（《构造》引理
\ref{constructions-lemma-relative-glueing}）得到一个态射
$X^{Y/wn} \to X$，它限制为 $U^{V/wn} \to U$，并且在 $U$ 上与
$\rho_{U_1}^{U_2}$ 相容。
当然，态射 $V \to U^{V/wn}$ 粘合为态射
$Y \to X^{Y/wn}$（见《构造》注
\ref{constructions-remark-relative-glueing-functorial}），
从而得到分解 $Y \to X^{Y/wn} \to X$，其中第二个态射是泛同胚。

\medskip\noindent
最后，令 $Y \to X' \to X$ 为 (1) 中的分解。
如上，考虑 $V \to U^{V/wn} \to U \subset X$，
我们得到分解 $V \to U \times_X X' \to U$，其中第二个箭头是泛同胚，
并得到一个唯一态射
$g_U : U^{V/wn} \to U \times_X X'$，它在 $U$ 上由
$U^{V/wn}$ 的泛性质给出。
这些 $g_U$ 与态射 $\rho_{U_1}^{U_2}$ 相容；细节略去。
因此存在唯一态射 $g : X^{Y/wn} \to X'$，它在 $X$ 上与
$g_U$ 在 $U$ 上一致，见《构造》注
\ref{constructions-remark-relative-glueing-functorial}。
这证明 $Y \to X^{Y/wn} \to X$ 是我们范畴中的始对象，证明完成。
\end{proof}

\begin{definition}
\label{morphisms-definition-relative-seminormalization}
设 $f : Y \to X$ 为概形的拟紧、拟分离且支配的态射。
\begin{enumerate}
\item 引理 \ref{morphisms-lemma-relative-seminormalization} 的 (2) 中构造的分解
$Y \to X^{Y/sn} \to X$ 称为 $X$ 在 $Y$ 中的\textit{半正规化}。
\item 引理 \ref{morphisms-lemma-relative-seminormalization} 的 (1) 中构造的分解
$Y \to X^{Y/wn} \to X$ 称为 $X$ 在 $Y$ 中的\textit{弱正规化}。
\end{enumerate}
\end{definition}
\noindent
以下给出对局部具有有限多个不可约分支的概形的半正规化的一种重新解释。

\begin{lemma}
\label{morphisms-lemma-seminormal-and-normalization}
设 $X$ 为每个拟紧开集都有有限多个不可约分支的概形。令
$\nu : X^\nu \to X$ 为 $X$ 的正规化。则 $X$ 在 $X^\nu$ 中的半正规化
就是 $X$ 的半正规化。公式为：
$X^{sn} = X^{X^\nu/sn}$。
\end{lemma}

\begin{proof}
令 $f : Y \to X$ 如（\ref{morphisms-equation-generic-points}）中所述，
使得 $X^\nu$ 是 $X$ 在 $Y$ 中的正规化。
半正规化 $X^{sn} \to X$ 是概形 $X$ 的由诱导剩余域同构的泛同胚
$X' \to X$ 构成的范畴中的始对象。由于 $Y$ 是
$X$ 的不可约分支的泛点处的剩余域的谱的不交并，
可见对于该范畴中的任意 $X' \to X$，我们都得到一个典范提升
$f' : Y \to X'$，它是 $f$ 的提升。于是由引理
\ref{morphisms-lemma-characterize-normalization}，得到一个典范态射 $X^\nu \to X'$。
从而再得到一个典范态射 $X^{X^\nu/sn} \to X'$，
它由 $X^{X^\nu/sn}$ 的泛性质给出。
这样，我们看到 $X^{X^\nu/sn}$ 满足与 $X^{sn}$ 相同的泛性质，结论成立。
\end{proof}

\noindent
引理 \ref{morphisms-lemma-seminormal-and-normalization} 导出以下定义。
由于我们只在 $X$ 局部具有有限多个不可约分支时构造正规化，
对于弱正规化也限制在这种情形。

\begin{definition}
\label{morphisms-definition-weak-normalization}
设 $X$ 为每个拟紧开集都有有限多个不可约分支的概形。我们将 $X$ 的
\textit{弱正规化}定义为其弱正规化
$$
X^\nu \longrightarrow X^{wn} \longrightarrow X
$$
即 $X$ 在正规化 $X^\nu$（即 $X$ 的正规化）中的弱正规化（定义 \ref{morphisms-definition-normalization}）。公式为：
$X^{wn} = X^{X^\nu/wn}$。
\end{definition}

\noindent
结合引理 \ref{morphisms-lemma-seminormal-and-normalization}，可见对于一个
局部具有有限多个不可约分支的概形 $X$，存在典范态射
$$
X^\nu \to X^{wn} \to X^{sn} \to X
$$
有了这个定义，我们就可以说明概形何时是弱正规的
（前提是它局部具有有限多个不可约分支）。

\begin{definition}
\label{morphisms-definition-weakly-normal}
设 $X$ 为每个拟紧开集都有有限多个不可约分支的概形。
称 $X$ 为\textit{弱正规}，如果弱正规化 $X^{wn} \to X$ 是同构（定义
\ref{morphisms-definition-weak-normalization}）。
\end{definition}

\noindent
由定义立即得到，对于每个拟紧开集都有有限多个不可约分支的概形 $X$，有
$$
X\text{ normal} \Rightarrow
X\text{ weakly normal} \Rightarrow
X\text{ seminormal}
$$
下面在仿射情形中说明弱正规的含义。

\begin{lemma}
\label{morphisms-lemma-affine-weakly-normal}
设 $X = \Spec(A)$ 为具有有限多个不可约分支的仿射概形。
则 $X$ 弱正规，当且仅当
\begin{enumerate}
\item $A$ 是半正规环（定义 \ref{morphisms-definition-seminormal-ring}）；
\item 对于素数 $p$ 以及满足下列条件的 $z, w \in A$：
(a) $z$ 是非零因子；(b) $w^p$ 可被 $z^p$ 整除；(c) $pw$ 可被 $z$ 整除，
都有 $w$ 可被 $z$ 整除。
\end{enumerate}
\end{lemma}

\begin{proof}
假设 $X$ 弱正规。由于弱正规概形是半正规概形，
由弱正规概形的定义可知 (1) 成立。
特别地，$A$ 是约化的。令 $p, z, w$ 如 (2) 中所述。
取 $x, y \in A$，使得 $z^p x = w^p$ 且 $zy = pw$。
于是 $p^p x = y^p$。环映射
$A \to C = A[t]/(t^p - x, pt - y)$ 在谱上诱导泛同胚。
$X^\nu$ 是 $X$ 的整闭包 $A'$ 的谱，其中它取自 $A$ 在 $A$ 的总分式环中的整闭包，见引理
\ref{morphisms-lemma-description-normalization}。
注意 $a = w/z \in A'$，因为 $a^p = x$。
因此有一个 $A$-代数同态 $A \to C \to A'$，将 $t$ 映为 $a$。
此时，弱正规的定义性质
$X = X^{wn} = X^{X^\nu/wn}$ 告诉我们，$C \to A'$ 的像包含于 $A$ 中。
于是得到所需的 $a \in A$。

\medskip\noindent
反之，假设 (1) 和 (2) 成立。令 $A'$ 如上一段中所述。
我们需要证明 $X^{X^\nu/wn} = X$。根据引理
\ref{morphisms-lemma-relative-weak-normalization-algebra} 的证明中的构造，
$X^{X^\nu/wn}$ 是 $A'$ 的子环的谱，该子环是下列子环
$A[a_1, \ldots, a_n] \subset A'$ 的并，其中对 $i = 1, \ldots, n$ 有
\begin{enumerate}
\item[(a)] $a_i^2, a_i^3 \in A[a_1, \ldots, a_{i - 1}]$，或者
\item[(b)] 存在素数 $p$，使得
$pa_i, a_i^p \in A[a_1, \ldots, a_{i - 1}]$。
\end{enumerate}
于是可以使用 (1) 和 (2) 归纳地看到
$a_1, \ldots, a_n \in A$；细节略去。
因此 $X = X^{X^\nu/wn}$，从而 $X$ 弱正规。
\end{proof}

\noindent
下面是一个必需的引理。

\begin{lemma}
\label{morphisms-lemma-locally-weakly-normal}
设 $X$ 为每个拟紧开集都有有限多个不可约分支的概形。以下条件等价：
\begin{enumerate}
\item 概形 $X$ 弱正规。
\item 对每个仿射开集 $U \subset X$，环 $\mathcal{O}_X(U)$
满足引理 \ref{morphisms-lemma-affine-weakly-normal} 的条件 (1) 和 (2)。
\item 存在仿射开覆盖 $X = \bigcup U_i$，使得每个环 $\mathcal{O}_X(U_i)$
满足引理 \ref{morphisms-lemma-affine-weakly-normal} 的条件 (1) 和 (2)。
\item 存在开覆盖 $X = \bigcup X_j$，使得每个开子概形 $X_j$ 都弱正规。
\end{enumerate}
此外，若 $X$ 弱正规，则每个开子概形也弱正规。
\end{lemma}

\begin{proof}
$X$ 弱正规的条件是态射
$X^{wn} = X^{X^\nu/wn} \to X$ 为同构。
由于 $X^\nu \to X$ 的构造与到开子概形的基变换相容，
并且 $X^{X^\nu/wn}$ 的构造与到 $X$ 的开子概形的基变换相容
（引理 \ref{morphisms-lemma-relative-seminormalization}），引理显然成立。
\end{proof}

\section{Zariski 主定理（代数版本）}
\label{morphisms-section-Zariski}

\noindent
这是可以使用纯代数方法证明的版本。
在我们能够证明更强的版本（对于非仿射态射）之前，
还需要发展更多工具。见《概形的上同调》一节
\ref{coherent-section-applications-formal-functions}
以及《更多关于态射》一节
\ref{more-morphisms-section-application-etale-neighbourhoods}。

\begin{theorem}[Zariski 主定理的代数版本]
\label{morphisms-theorem-main-theorem}
设 $f : Y \to X$ 为概形的仿射态射。
假设 $f$ 是有限型的。
令 $X'$ 为 $X$ 在 $Y$ 中的正规化。图示为：
$$
\xymatrix{
Y \ar[rd]_f \ar[rr]_{f'} & & X' \ar[ld]^\nu \\
& X &
}
$$
则存在开子概形 $U' \subset X'$，使得
\begin{enumerate}
\item $(f')^{-1}(U') \to U'$ 是同构，并且
\item $(f')^{-1}(U') \subset Y$ 正是 $f$ 为拟有限的点集。
\end{enumerate}
\end{theorem}

\begin{proof}
可以立即归约到 $X$、从而 $Y$ 都是仿射的情形。
设 $X = \Spec(R)$ 且 $Y = \Spec(A)$。
于是 $X' = \Spec(A')$，其中 $A'$ 是 $R$ 在 $A$ 中的整闭包，见定义
\ref{morphisms-definition-integral-closure} 和 \ref{morphisms-definition-normalization-X-in-Y}。
由《代数》定理 \ref{algebra-theorem-main-theorem}，
对于每个拟有限点 $y \in Y$（即 $f$ 在该点拟有限），存在开集
$U'_y \subset X'$，使得 $(f')^{-1}(U'_y) \to U'_y$ 是同构。
令 $U' = \bigcup U'_y$，其中 $y \in Y$ 遍历 $f$ 为拟有限的所有点。
还需证明 $f$ 在 $(f')^{-1}(U')$ 的所有点处都是拟有限的。
若 $y \in (f')^{-1}(U')$ 的像为 $x \in X$，则可见
$Y_x \to X'_x$ 在 $y$ 的某个邻域中是同构。因此 $Y_x$ 中没有点专化到 $y$，
因为 $f'(y)$ 在 $X'_x$ 中具有同样性质，见引理
\ref{morphisms-lemma-integral-fibres}。由引理
\ref{morphisms-lemma-quasi-finite-at-point-characterize} 的 (3)，
这意味着 $f$ 在 $y$ 处拟有限。
\end{proof}

\noindent
我们可以使用 Zariski 主定理的代数版本证明，态射拟有限的点集是开集。

\begin{lemma}
\label{morphisms-lemma-quasi-finite-points-open}
\begin{slogan}
态射的局部拟有限轨迹是开集
\end{slogan}
设 $f : X \to S$ 为概形的态射。
$X$ 中 $f$ 拟有限的点集是一个开集 $U \subset X$。
诱导的态射 $U \to S$ 局部拟有限。
\end{lemma}

\begin{proof}
假设 $f$ 在 $x$ 处拟有限。
令 $x \in U = \Spec(A) \subset X$、$V = \Spec(R) \subset S$
为定义 \ref{morphisms-definition-quasi-finite} 中的仿射开集。
由上面的定理 \ref{morphisms-theorem-main-theorem} 或《代数》引理
\ref{algebra-lemma-quasi-finite-open}，素理想 $\mathfrak q$ 中满足
$R \to A$ 拟有限者的集合是 $\Spec(A)$ 中的开集。
由于这些素理想都对应于 $X$ 中 $f$ 拟有限的点，得到第一个断言。
第二个断言显然。
\end{proof}

\noindent
稍后我们将把以下引理推广到一般的拟有限分离态射，见《更多关于态射》引理
\ref{more-morphisms-lemma-quasi-finite-separated-pass-through-finite}。

\begin{lemma}
\label{morphisms-lemma-quasi-finite-affine}
设 $f : Y \to X$ 为概形的态射。
假设
\begin{enumerate}
\item $X$ 与 $Y$ 都是仿射的，并且
\item $f$ 是拟有限的。
\end{enumerate}
则存在图
$$
\xymatrix{
Y \ar[rd]_f \ar[rr]_j & & Z \ar[ld]^\pi \\
& X &
}
$$
其中 $Z$ 是仿射的，$\pi$ 有限，且 $j$ 是开浸入。
\end{lemma}

\begin{proof}
这就是《代数》引理
\ref{algebra-lemma-quasi-finite-open-integral-closure} 在概形语言中的改写。
\end{proof}

\begin{lemma}
\label{morphisms-lemma-image-nowhere-dense-quasi-finite}
设 $f : Y \to X$ 为概形的拟有限态射。
设 $T \subset Y$ 为 $Y$ 的闭无处稠密子集。
则 $f(T) \subset X$ 是 $X$ 的无处稠密子集。
\end{lemma}

\begin{proof}
与引理 \ref{morphisms-lemma-image-nowhere-dense-finite} 的证明一样，
这立即归约到基概形 $X$ 为仿射的情形。
在这种情形下，$Y = \bigcup_{i = 1, \ldots, n} Y_i$ 是仿射开集的有限并
（因为 $f$ 拟紧）。由于每个 $T \cap Y_i$ 都无处稠密，
而且无处稠密集的有限并仍无处稠密（见《拓扑》引理
\ref{topology-lemma-nowhere-dense}），只需证明像
$f(T \cap Y_i)$ 在 $X$ 中无处稠密。
这又把我们归约到 $X$ 与 $Y$ 都为仿射的情形。
此时应用上面的引理 \ref{morphisms-lemma-quasi-finite-affine}，得到图
$$
\xymatrix{
Y \ar[rd]_f \ar[rr]_j & & Z \ar[ld]^\pi \\
& X &
}
$$
其中 $Z$ 仿射，$\pi$ 有限，且 $j$ 是开浸入。
置 $\overline{T} = \overline{j(T)} \subset Z$。
由《拓扑》引理 \ref{topology-lemma-image-nowhere-dense-open}，
$\overline{T}$ 在 $Z$ 中无处稠密。
由于 $f(T) \subset \pi(\overline{T})$，引理由有限情形中的相应结果得出，
见引理 \ref{morphisms-lemma-image-nowhere-dense-finite}。
\end{proof}

\section{泛一致有界的纤维}
\label{morphisms-section-universally-bounded}

\noindent
设 $X$ 为域 $k$ 上的概形。若 $X$ 在 $k$ 上有限，
则 $X = \Spec(A)$，其中 $A$ 是有限 $k$-代数。
另一种说法是 $X$ 在 $\Spec(k)$ 上有限局部自由，见定义
\ref{morphisms-definition-finite-locally-free}。因此
$X \to \Spec(k)$ 有一个\textit{次数}，它是整数 $d \geq 0$，
即 $d = \dim_k(A)$。我们有时称它为（有限）概形 $X$ 在 $k$ 上的\textit{次数}。

\begin{definition}
\label{morphisms-definition-universally-bounded}
设 $f : X \to Y$ 为概形的态射。
\begin{enumerate}
\item 称整数 $n$\textit{界住了 $f$ 的纤维的次数}，如果对于所有
$y \in Y$，纤维 $X_y$ 是 $\kappa(y)$ 上的有限概形，
且它在 $\kappa(y)$ 上的次数为 $\leq n$。
\item 称\textit{$f$ 的纤维泛一致有界}\footnote{这很可能不是标准术语。}，
如果存在一个整数 $n$ 界住 $f$ 的纤维的次数。
\end{enumerate}
\end{definition}

\noindent
特别地，注意每个纤维中的点数也被 $n$ 界住。
（反之不成立，即使所有纤维都是有限约化概形。）

\begin{lemma}
\label{morphisms-lemma-characterize-universally-bounded}
设 $f : X \to Y$ 为概形的态射，并设 $n \geq 0$。以下条件等价：
\begin{enumerate}
\item 整数 $n$ 界住了 $f$ 的纤维的次数；
\item 对每个态射 $\Spec(k) \to Y$，其中 $k$ 是域，
纤维积 $X_k = \Spec(k) \times_Y X$ 在 $k$ 上有限，且次数为 $\leq n$。
\end{enumerate}
在这种情况下，$f$ 的纤维泛一致有界，并且概形 $X_k$ 至多有 $n$ 个点。
更精确地，若 $X_k = \{x_1, \ldots, x_t\}$，则有
$$
n \geq \sum\nolimits_{i = 1, \ldots, t} [\kappa(x_i) : k]
$$
\end{lemma}

\begin{proof}
(2) $\Rightarrow$ (1) 显然。反向蕴含成立，因为若
$\Spec(k) \to Y$ 的像为 $y$，则
$X_k = \Spec(k) \times_{\Spec(\kappa(y))} X_y$。
按照定义，$f$ 的纤维泛一致有界是指存在某个 $n$。
最后，设 $X_k = \Spec(A)$。则 $\dim_k A = n$。
于是 $A$ 是 Artin 环，所有素理想都是极大理想 $\mathfrak m_i$，
而且 $A$ 是它在这些极大理想处的局部化之积。
见《代数》引理 \ref{algebra-lemma-finite-dimensional-algebra}
和 \ref{algebra-lemma-artinian-finite-length}。
于是 $\mathfrak m_i$ 对应于 $x_i$，有
$A_{\mathfrak m_i} = \mathcal{O}_{X_k, x_i}$，
从而存在满射
$A \to  \bigoplus \kappa(\mathfrak m_i) = \bigoplus \kappa(x_i)$。
线性代数于是给出引理断言中的不等式。
\end{proof}

\begin{lemma}
\label{morphisms-lemma-finite-locally-free-universally-bounded}
若 $f$ 是次数为 $d$ 的有限局部自由态射，
则 $d$ 界住 $f$ 的纤维的次数。
\end{lemma}

\begin{proof}
这是因为 $f$ 的任意基变换仍为次数 $d$ 的有限局部自由态射
（引理 \ref{morphisms-lemma-base-change-finite-locally-free}），
因而 $f$ 的所有纤维的次数都是 $d$。
\end{proof}

\begin{lemma}
\label{morphisms-lemma-composition-universally-bounded}
纤维泛一致有界的态射的复合，其纤维仍泛一致有界。
更精确地，假设 $n$ 界住 $f : X \to Y$ 的纤维的次数，
且 $m$ 界住 $g : Y \to Z$ 的纤维的次数。
则 $nm$ 界住 $g \circ f : X \to Z$ 的纤维的次数。
\end{lemma}

\begin{proof}
设 $f : X \to Y$ 和 $g : Y \to Z$ 的纤维泛一致有界。
假设 $\deg(X_y/\kappa(y)) \leq n$ 对所有 $y \in Y$ 成立，
并且 $\deg(Y_z/\kappa(z)) \leq m$ 对所有 $z \in Z$ 成立。
令 $z \in Z$ 为一点。由假设，概形 $Y_z$ 在
$\Spec(\kappa(z))$ 上有限。特别地，$Y_z$ 的底拓扑空间是有限离散集。
态射 $f_z : X_z \to Y_z$ 的纤维是 $f$ 在 $Y$ 中相应点处的纤维；
由上面的论证，它们是有限离散集。
因此 $X_z$ 的底拓扑空间也是有限离散集。
于是 $X_z$ 是仿射概形（这是一个很好的练习；例如也可将《性质》引理
\ref{properties-lemma-maximal-points-affine} 应用于 $X_z$ 的所有点而得出）。
写 $X_z = \Spec(A)$、$Y_z = \Spec(B)$，并置 $k = \kappa(z)$。
于是有 $k \to B \to A$，并且我们知道：(a) $\dim_k(B) \leq m$；
(b) 对每个极大理想 $\mathfrak m \subset B$，有
$\dim_{\kappa(\mathfrak m)}(A/\mathfrak mA) \leq n$。
我们断言这意味着 $\dim_k(A) \leq nm$。
注意 $B$ 是它在各个极大理想 $B_{\mathfrak m}$ 处的局部化之积；
例如，这是因为 $Y_z$ 是 $1$-点概形的不交并，或者由《代数》引理
\ref{algebra-lemma-finite-dimensional-algebra} 和
\ref{algebra-lemma-artinian-finite-length} 得出。
因此
$\dim_k(B) = \sum_{\mathfrak m} \dim_k(B_{\mathfrak m})$，并且
$\dim_k(A) = \sum_{\mathfrak m} \dim_k(A_{\mathfrak m})$，
其中两式中的 $\mathfrak m$ 都遍历 $B$（而不是 $A$）的极大理想。
由上述结论和 Nakayama 引理（《代数》引理
\ref{algebra-lemma-NAK}），可见每个 $A_{\mathfrak m}$ 都是
$B_{\mathfrak m}^{\oplus n}$ 的商模（作为 $B_{\mathfrak m}$-模）。因此
$\dim_k(A_{\mathfrak m}) \leq n \dim_k(B_{\mathfrak m})$。
综合起来可得
$$
\dim_k(A) = \sum\nolimits_{\mathfrak m} \dim_ta	(A_{\mathfrak m})
\leq \sum\nolimits_{\mathfrak m} n \dim_k(B_{\mathfrak m})
= n \dim_k(B) \leq nm
$$
正是所需结论。
\end{proof}

\begin{lemma}
\label{morphisms-lemma-base-change-universally-bounded}
纤维泛一致有界的态射的基变换，其纤维仍泛一致有界。
更精确地，若 $n$ 界住 $f : X \to Y$ 的纤维的次数，且
$Y' \to Y$ 是任意态射，则基变换
$f' : Y' \times_Y X \to Y'$ 的纤维的次数也被 $n$ 界住。
\end{lemma}

\begin{proof}
这由引理 \ref{morphisms-lemma-characterize-universally-bounded} 的结论显然成立。
\end{proof}

\begin{lemma}
\label{morphisms-lemma-descent-universally-bounded}
设 $f : X \to Y$ 为概形的态射。设 $Y' \to Y$ 为概形的态射，并令
$f' : X' = X_{Y'} \to Y'$ 为 $f$ 的基变换。
若 $Y' \to Y$ 满射且 $f'$ 的纤维泛一致有界，
则 $f$ 的纤维也泛一致有界。更精确地，若 $n$ 界住 $f'$ 的纤维的次数，
则 $n$ 也界住 $f$ 的纤维的次数。
\end{lemma}

\begin{proof}
令 $n \geq 0$ 为界住 $f'$ 的纤维次数的整数。
我们断言 $n$ 对 $f$ 也适用。事实上，若 $y \in Y$ 为一点，
则取一点 $y' \in Y'$ 位于 $y$ 之上，并注意
$$
X'_{y'} = \Spec(\kappa(y')) \times_{\Spec(\kappa(y))} X_y.
$$
由于假设 $X'_{y'}$ 是次数为 $\leq n$ 的有限概形（在 $\kappa(y')$ 上），
所以 $X_y$ 也是次数为 $\leq n$ 的有限概形（在 $\kappa(y)$ 上）。
（略去一些细节。）
\end{proof}

\begin{lemma}
\label{morphisms-lemma-immersion-universally-bounded}
浸入的纤维泛一致有界。
\end{lemma}

\begin{proof}
定义中可取整数 $n = 1$。
\end{proof}

\begin{lemma}
\label{morphisms-lemma-etale-universally-bounded}
设 $f : X \to Y$ 为概形的 'etale 态射。设 $n \geq 0$。
以下条件等价：
\begin{enumerate}
\item 整数 $n$ 界住纤维的次数；
\item 对每个域 $k$ 和态射 $\Spec(k) \to Y$，基变换
$X_k = \Spec(k) \times_Y X$ 至多有 $n$ 个点；
\item 对每个 $y \in Y$ 和每个可分代数闭包
$\kappa(y) \subset \kappa(y)^{sep}$，概形
$X_{\kappa(y)^{sep}}$ 至多有 $n$ 个点。
\end{enumerate}
\end{lemma}

\begin{proof}
这由引理 \ref{morphisms-lemma-characterize-universally-bounded}
以及下述事实得出：纤维 $X_y$ 是 $\kappa(y)$ 的有限可分域扩张的谱的不交并，
见引理 \ref{morphisms-lemma-etale-over-field}。
\end{proof}

\noindent
纤维泛一致有界是绝对概念，而不是相对概念。
这就是下一个引理的条件要求 $X$ 拟紧，而不是要求 $f$ 拟紧的原因。

\begin{lemma}
\label{morphisms-lemma-locally-quasi-finite-qc-source-universally-bounded}
设 $f : X \to Y$ 为概形的态射。假设
\begin{enumerate}
\item $f$ 局部拟有限，并且
\item $X$ 拟紧。
\end{enumerate}
则 $f$ 的纤维泛一致有界。
\end{lemma}

\begin{proof}
由于 $X$ 拟紧，存在有限仿射开覆盖
$X = \bigcup_{i = 1, \ldots, n} U_i$ 以及仿射开集
$V_i \subset Y$，$i = 1, \ldots, n$，使得 $f(U_i) \subset V_i$。
由“局部拟有限”的局部性质（见引理
\ref{morphisms-lemma-quasi-finite-at-point-characterize} 的 (4)），可见态射
$f|_{U_i} : U_i \to V_i$ 是仿射概形之间的局部拟有限态射，
因而是拟有限的，见引理 \ref{morphisms-lemma-quasi-finite-locally-quasi-compact}。
对于 $y \in Y$，显然
$X_y = \bigcup_{y \in V_i} (U_i)_y$ 是开覆盖。
因此只需对仿射概形之间的拟有限态射证明该引理
（也就是说，若 $n_i$ 对态射 $f|_{U_i} : U_i \to V_i$ 适用，
则 $\sum n_i$ 对 $f$ 适用）。

\medskip\noindent
假设 $f : X \to Y$ 是仿射概形之间的拟有限态射。
由引理 \ref{morphisms-lemma-quasi-finite-affine}，可以找到图
$$
\xymatrix{
X \ar[rd]_f \ar[rr]_j & & Z \ar[ld]^\pi \\
& Y &
}
$$
其中 $Z$ 仿射，$\pi$ 有限，且 $j$ 是开浸入。
由于 $j$ 的纤维泛一致有界
（引理 \ref{morphisms-lemma-immersion-universally-bounded}），
问题归约为证明 $\pi$ 的纤维泛一致有界
（引理 \ref{morphisms-lemma-composition-universally-bounded}）。

\medskip\noindent
这又把我们归约到形如 $\Spec(B) \to \Spec(A)$ 的态射，
其中 $A \to B$ 有限。设 $B$ 由 $x_1, \ldots, x_n$
作为 $A$-模生成，于是
$$
A^{\oplus n} \longrightarrow B, \quad
(a_1, \ldots, a_n) \longmapsto \sum a_i x_i
$$
是满的 $A$-模映射。对任意素理想 $\mathfrak p \subset A$，
它诱导满的 $\kappa(\mathfrak p)$-向量空间映射
$$
\kappa(\mathfrak p)^{\oplus n} \longrightarrow B \otimes_A \kappa(\mathfrak p)
$$
换言之，定义纤维泛一致有界的态射时，可以取整数 $n$。
\end{proof}

\begin{lemma}
\label{morphisms-lemma-universally-bounded-permanence}
考虑概形态射的交换图
$$
\xymatrix{
X \ar[rd]_g \ar[rr]_f & & Y \ar[ld]^h \\
& Z &
}
$$
若 $g$ 的纤维泛一致有界，并且 $f$ 满且平坦，
则 $h$ 的纤维也泛一致有界。更精确地，若 $n$
界住 $g$ 的纤维的次数，则 $n$ 也界住 $h$ 的纤维的次数。
\end{lemma}

\begin{proof}
假设 $g$ 的纤维泛一致有界，并且 $f$ 满且平坦。
设 $g$ 的纤维的次数被 $n \in \mathbf{N}$ 界住。
我们断言 $n$ 对 $h$ 也适用。
令 $z \in Z$。考虑概形态射 $X_z \to Y_z$。
它平坦且满射。由假设，$X_z$ 是 $\kappa(z)$ 上的有限概形，
特别地，它是 Artin 环的谱（由《代数》引理
\ref{algebra-lemma-finite-dimensional-algebra}）。
由引理 \ref{morphisms-lemma-Artinian-affine}，态射 $X_z \to Y_z$ 是仿射的，
特别是拟紧的。由引理 \ref{morphisms-lemma-fpqc-quotient-topology}，
可知 $Y_z$ 也是有限离散的，因为 $X_z$ 具有这一性质。
因此 $Y_z$ 是仿射概形（这是一个很好的练习；例如也可将《性质》引理
\ref{properties-lemma-maximal-points-affine} 应用于 $Y_z$ 的所有点而得出）。
写 $Y_z = \Spec(B)$ 且 $X_z = \Spec(A)$。
则 $A$ 在 $B$ 上忠实平坦，因此 $B \subset A$。
所以 $\dim_k(B) \leq \dim_k(A) \leq n$，正是所需结论。
\end{proof}

\section{杂项}
\label{morphisms-section-misc}

\noindent
这里收录不适合放在其他地方的结果。

\begin{lemma}
\label{morphisms-lemma-factor-through-point}
设 $f : Y \to X$ 为概形的态射，并设 $x \in X$ 为一点。
假设 $Y$ 约化，并且在集合论意义下 $f(Y)$ 包含于 $\{x\}$ 中。
则 $f$ 经过典范态射 $x = \Spec(\kappa(x)) \to X$ 分解。
\end{lemma}

\begin{proof}
略。提示：在仿射局部工作，可将问题归约到一个交换代数引理。
给定环映射 $A \to B$，其中 $B$ 约化，并且存在唯一素理想
$\mathfrak p \subset A$ 位于 $\Spec(B) \to \Spec(A)$ 的像中，
则 $A \to B$ 经过 $\kappa(\mathfrak p)$ 分解。
这是一个很好的练习。
\end{proof}

\begin{lemma}
\label{morphisms-lemma-factor-through-set-of-points}
设 $f : Y \to X$ 为概形的态射，并设 $E \subset X$。
假设 $X$ 局部诺特，$E$ 的元素之间不存在非平凡专化，
$Y$ 约化，并且 $f(Y) \subset E$。
则 $f$ 经过 $\coprod_{x \in E} x \to X$ 分解。
\end{lemma}

\begin{proof}
当 $E$ 为单点集时，这由引理 \ref{morphisms-lemma-factor-through-point} 得出。
若 $E$ 有限，则由我们关于专化的假设，$E$（赋予从 $X$ 诱导的拓扑）
是有限离散空间。因此该情形归约到单点集情形。
一般地，可归约到 $X$ 与 $Y$ 都为仿射概形的情形。
设 $f : Y \to X$ 对应于环映射 $\varphi : A \to B$。
用 $A' \subset B$ 表示 $\varphi$ 的像。令
$E' \subset \Spec(A') \subset \Spec(A)$ 为 $A'$ 的极小素理想集。
由《代数》引理 \ref{algebra-lemma-injective-minimal-primes-in-image}，
集合 $E'$ 包含于
$\Spec(B) \to \Spec(A') \subset \Spec(A)$ 的像中。
因此 $E' \subset E$。由于 $A'$ 是诺特的，由《代数》引理
\ref{algebra-lemma-Noetherian-irreducible-components}，$E'$ 有限。
由于 $\Spec(B) \to \Spec(A)$ 的像中的任何其他点
都是 $E'$ 的某个元素的专化，并且属于 $E$，
可知该像包含于 $E'$ 中（由我们关于 $E$ 中各点之间专化的假设）。
于是归约到上面已经处理的 $E$ 有限的情形。
\end{proof}
